% concatenated from main_arxiv.tex, math_commands_new.tex
\documentclass[10pt]{article} % For LaTeX2e
%\usepackage{tmlr}
% If accepted, instead use the following line for the camera-ready submission:
%\usepackage[accepted]{tmlr}
% To de-anonymize and remove mentions to TMLR (for example for posting to preprint servers), instead use the following:
\usepackage[preprint]{tmlr}

\usepackage{hyperref}
\usepackage{url}
\usepackage{microtype}
\usepackage{mathtools}
\usepackage{amssymb}
\usepackage{booktabs}
\usepackage{tikz}
\usepackage{amsthm}
\usepackage{graphicx}

%%%%
% --- Theorem-like environments (TMLR-friendly) ---
\usepackage{amsthm}

% Number everything by section: 3.1, 3.2, ...
\newtheorem{theorem}{Theorem}[section]
\newtheorem{lemma}[theorem]{Lemma}
\newtheorem{corollary}[theorem]{Corollary}

% Bold heading, normal (roman) body — matches your screenshot
\theoremstyle{definition}
\newtheorem{definition}[theorem]{Definition}
\newtheorem{remark}[theorem]{Remark}

%%%%
\theoremstyle{plain}
\newtheorem{proposition}[theorem]{Proposition}
\newtheorem{assumption}[theorem]{Assumption}   
\newtheorem{example}[theorem]{Example}
% % \theoremstyle{remark}
% % \newtheorem{remark}[theorem]{Remark}
% \theoremstyle{remark}
% \newtheorem*{remark}{\textbf{Remark}}


% math_commands.tex — TMLR-safe
% (Drop this next to main.tex and \input it once, after \usepackage{tmlr})
% Key fixes:
% - No theorem environments here (they live in main.tex).
% - Guarded definitions to avoid clobbering TMLR’s macros.
% - Fixed duplicate \calM and duplicate \rvu.
% - Added \graphicspath for figures stored one level up (../Preprint/).

% ---- PACKAGES (minimal; avoid reloading what tmlr already loads) ----
\usepackage{amsmath,amsfonts,bm}
\usepackage[T1]{fontenc}
\usepackage{algorithm,algpseudocode}
\usepackage{graphicx}
\usepackage{enumitem}
\usepackage{wrapfig}
\usepackage{xfrac}
\usepackage{subcaption}
\usepackage{longtable}
\usepackage{array}
\usepackage{xcolor}

% If your figures live in ../Preprint relative to main.tex:
\graphicspath{{../Preprint/}}

% ---- REFERENCE COMMANDS ----
\providecommand{\figref}[1]{figure~\ref{#1}}
\providecommand{\Figref}[1]{Figure~\ref{#1}}
\providecommand{\secref}[1]{section~\ref{#1}}
\providecommand{\Secref}[1]{Section~\ref{#1}}
\providecommand{\eqrefnum}[1]{\ref{#1}} % raw ref if needed
\providecommand{\eqrefword}[1]{equation~\ref{#1}}
\providecommand{\Eqref}[1]{Equation~\ref{#1}}
\providecommand{\algref}[1]{algorithm~\ref{#1}}
\providecommand{\Algref}[1]{Algorithm~\ref{#1}}

% ---- DELIMITERS & UTILITIES ----
\newcommand{\paren}[1]{\left(#1\right)}
\newcommand{\cbrace}[1]{\left\{#1\right\}}
\newcommand{\brac}[1]{\left[#1\right]}
\newcommand{\abs}[1]{\left|#1\right|}
\newcommand{\norm}[1]{\left\|#1\right\|}
\newcommand{\indy}[1]{\mathbb{I}\left\{#1\right\}}
\newcommand{\inner}[2]{\left\langle#1, #2\right\rangle}
\newcommand{\deriv}[2]{\frac{\partial #1}{\partial #2}}
\newcommand{\ceil}[1]{\left\lceil #1\right\rceil}
\newcommand{\floor}[1]{\left\lfloor #1\right\rfloor}
\newcommand{\act}[2]{\hat{\sigma}_{\kappa}\paren{#1,#2}}
\renewcommand{\vec}[1]{\texttt{vec}\paren{#1}}

% ---- BLACKBOARD, CAL, OPERATORS ----
\providecommand{\E}{\mathbb{E}}
\providecommand{\R}{\mathbb{R}}
\providecommand{\N}{\mathbb{N}}
\providecommand{\Z}{\mathbb{Z}}
\providecommand{\C}{\mathbb{C}}

\def\calB{{\mathcal{B}}}
\def\calL{{\mathcal{L}}}
\def\calM{{\mathcal{M}}}
\def\calN{{\mathcal{N}}}
\def\calP{{\mathcal{P}}}
\def\calE{{\mathcal{E}}}
\def\calT{{\mathcal{T}}}
\def\calF{{\mathcal{F}}}

\newcommand{\defeq}{\vcentcolon=}
\newcommand{\eqdef}{=\vcentcolon}
\DeclareMathOperator*{\argmax}{arg\,max}
\DeclareMathOperator*{\argmin}{arg\,min}

% Optional helper (only if you actually use it; otherwise comment out)
% \newcommand{\mycite}[1]{\textit{\textcolor{blue}{\citep{#1}}}}

% ---- BOLD UPPERCASE ----
\def\bfA{{\mathbf{A}}}\def\bfB{{\mathbf{B}}}\def\bfC{{\mathbf{C}}}\def\bfD{{\mathbf{D}}}
\def\bfE{{\mathbf{E}}}\def\bfF{{\mathbf{F}}}\def\bfG{{\mathbf{G}}}\def\bfH{{\mathbf{H}}}
\def\bfI{{\mathbf{I}}}\def\bfJ{{\mathbf{J}}}\def\bfK{{\mathbf{K}}}\def\bfL{{\mathbf{L}}}
\def\bfM{{\mathbf{M}}}\def\bfN{{\mathbf{N}}}\def\bfO{{\mathbf{O}}}\def\bfP{{\mathbf{P}}}
\def\bfQ{{\mathbf{Q}}}\def\bfR{{\mathbf{R}}}\def\bfS{{\mathbf{S}}}\def\bfT{{\mathbf{T}}}
\def\bfU{{\mathbf{U}}}\def\bfV{{\mathbf{V}}}\def\bfW{{\mathbf{W}}}\def\bfX{{\mathbf{X}}}
\def\bfY{{\mathbf{Y}}}\def\bfZ{{\mathbf{Z}}}

% ---- BOLD LOWERCASE / GREEK ----
\def\bfepsilon{{\bm{\epsilon}}}\def\bftheta{{\bm{\theta}}}\def\bfphi{{\bm{\varphi}}}
\def\bfa{{\mathbf{a}}}\def\bfb{{\mathbf{b}}}\def\bfc{{\mathbf{c}}}\def\bfd{{\mathbf{d}}}
\def\bfe{{\mathbf{e}}}\def\bff{{\mathbf{f}}}\def\bfg{{\mathbf{g}}}\def\bfh{{\mathbf{h}}}
\def\bfi{{\mathbf{i}}}\def\bfj{{\mathbf{j}}}\def\bfk{{\mathbf{k}}}\def\bfl{{\mathbf{l}}}
\def\bfm{{\mathbf{m}}}\def\bfn{{\mathbf{n}}}\def\bfo{{\mathbf{o}}}\def\bfp{{\mathbf{p}}}
\def\bfq{{\mathbf{q}}}\def\bfr{{\mathbf{r}}}\def\bfs{{\mathbf{s}}}\def\bft{{\mathbf{t}}}
\def\bfu{{\mathbf{u}}}\def\bfv{{\mathbf{v}}}\def\bfw{{\mathbf{w}}}\def\bfx{{\mathbf{x}}}
\def\bfy{{\mathbf{y}}}\def\bfz{{\mathbf{z}}}

% ---- CITATION HELPER (optional; comment out if unused) ----
% \newcommand{\mycite}[1]{\textit{\textcolor{blue}{\citep{#1}}}}

% ---- THEOREM-LIKE ENVIRONMENTS ----
% (Intentionally omitted here to avoid clashes; you already define them in main.tex)

% ---- RANDOM VARIABLE / MATRIX ALPHABETS (if you rely on them) ----
% Uncomment only if needed; these don’t clash but keep your preamble lighter if unused.
% % Random variables (roman)
% \def\ra{{\textnormal{a}}}\def\rb{{\textnormal{b}}}\def\rc{{\textnormal{c}}}\def\rd{{\textnormal{d}}}
% \def\re{{\textnormal{e}}}\def\rf{{\textnormal{f}}}\def\rg{{\textnormal{g}}}\def\rh{{\textnormal{h}}}
% \def\ri{{\textnormal{i}}}\def\rj{{\textnormal{j}}}\def\rk{{\textnormal{k}}}\def\rl{{\textnormal{l}}}
% \def\rn{{\textnormal{n}}}\def\ro{{\textnormal{o}}}\def\rp{{\textnormal{p}}}\def\rq{{\textnormal{q}}}
% \def\rr{{\textnormal{r}}}\def\rs{{\textnormal{s}}}\def\rt{{\textnormal{t}}}\def\ru{{\textnormal{u}}}
% \def\rv{{\textnormal{v}}}\def\rw{{\textnormal{w}}}\def\rx{{\textnormal{x}}}\def\ry{{\textnormal{y}}}\def\rz{{\textnormal{z}}}
%
% % Random vectors (bold)
% \def\rva{{\mathbf{a}}}\def\rvb{{\mathbf{b}}}\def\rvc{{\mathbf{c}}}\def\rvd{{\mathbf{d}}}
% \def\rve{{\mathbf{e}}}\def\rvf{{\mathbf{f}}}\def\rvg{{\mathbf{g}}}\def\rvh{{\mathbf{h}}}
% \def\rvj{{\mathbf{j}}}\def\rvk{{\mathbf{k}}}\def\rvl{{\mathbf{l}}}\def\rvm{{\mathbf{m}}}
% \def\rvn{{\mathbf{n}}}\def\rvo{{\mathbf{o}}}\def\rvp{{\mathbf{p}}}\def\rvq{{\mathbf{q}}}
% \def\rvr{{\mathbf{r}}}\def\rvs{{\mathbf{s}}}\def\rvt{{\mathbf{t}}}\def\rvu{{\mathbf{u}}}
% \def\rvv{{\mathbf{v}}}\def\rvw{{\mathbf{w}}}\def\rvx{{\mathbf{x}}}\def\rvy{{\mathbf{y}}}\def\rvz{{\mathbf{z}}}
%
% % Random matrices (bold)
% \def\rmA{{\mathbf{A}}}\def\rmB{{\mathbf{B}}}\def\rmC{{\mathbf{C}}}\def\rmD{{\mathbf{D}}}
% \def\rmE{{\mathbf{E}}}\def\rmF{{\mathbf{F}}}\def\rmG{{\mathbf{G}}}\def\rmH{{\mathbf{H}}}
% \def\rmI{{\mathbf{I}}}\def\rmJ{{\mathbf{J}}}\def\rmK{{\mathbf{K}}}\def\rmL{{\mathbf{L}}}
% \def\rmM{{\mathbf{M}}}\def\rmN{{\mathbf{N}}}\def\rmO{{\mathbf{O}}}\def\rmP{{\mathbf{P}}}
% \def\rmQ{{\mathbf{Q}}}\def\rmR{{\mathbf{R}}}\def\rmS{{\mathbf{S}}}\def\rmT{{\mathbf{T}}}
% \def\rmU{{\mathbf{U}}}\def\rmV{{\mathbf{V}}}\def\rmW{{\mathbf{W}}}\def\rmX{{\mathbf{X}}}
% \def\rmY{{\mathbf{Y}}}\def\rmZ{{\mathbf{Z}}}

% ---- MISC ----
\newcommand{\laplace}{\mathrm{Laplace}}
\newcommand{\KL}{D_{\mathrm{KL}}}
\newcommand{\Var}{\mathrm{Var}}
\newcommand{\Cov}{\mathrm{Cov}}
\newcommand{\standarderror}{\mathrm{SE}}


\title{Convergence Analysis of Two-Layer Neural Networks under\\ Gaussian Input Masking}

% For anonymized submission:
%\author{\name Anonymous Authors \email tmlr@mail \\
%       \addr Anonymous Institution}
% For camera-ready, replace with:
 \author{
   \name \!\!Afroditi Kolomvaki \email ak203@rice.edu \\
   \addr Computer Science Dept., Rice University, Houston, TX, USA
   \AND
   \name Fangshuo Liao \email fangshuo.liao@rice.edu \\
   \addr Computer Science Dept., Rice University, Houston, TX, USA
   \AND
   \name Evan Dramko \email ed55@rice.edu \\
   \addr Computer Science Dept., Rice University, Houston, TX, USA
   \AND
   \name Ziyun Guang \email cg105@rice.edu \\
   \addr Computer Science Dept., Rice University, Houston, TX, USA
   \AND
   \name Anastasios Kyrillidis \email anastasios@rice.edu \\
   \addr Computer Science Dept., Rice University, Houston, TX, USA
 }

\begin{document}
\maketitle
\begin{abstract}
We investigate the convergence guarantee of two-layer neural network training with Gaussian randomly masked inputs. This scenario corresponds to Gaussian dropout at the input level, or noisy input training common in sensor networks, privacy-preserving training, and federated learning, where each user may have access to partial or corrupted features. Using a Neural Tangent Kernel (NTK) analysis, we demonstrate that training a two-layer ReLU network with Gaussian randomly masked inputs achieves linear convergence up to an error region proportional to the mask's variance. A key technical contribution is resolving the randomness within the non-linear activation, a problem of independent interest.
\end{abstract}
%%%%%%%
\section{Introduction}
\label{sec:intro}
Neural networks (NNs) have revolutionized AI applications, where their success largely stems from their ability to learn complex patterns when trained on well-curated datasets \citep{schuhmann2022laion,li2023textbooks, gunasekar2023textbooks, edwards2024data}. 
A component to the success of NNs is its ability to model a broad range of tasks and data distributions under various scenarios. Empirical evidence has suggested neural network's ability to learn even under noisy input \citep{kariotakis2024leveraging}, gradient noise \citep{ruder2017overviewgradientdescentoptimization}, as well as modifications to the internal representations during training \citep{srivastava2014dropout,yuan2022distributed}. Leveraging such ability of the neural networks, many real-world deployment adopts a modification to the data representations during training to achieve particular goals such as robustness, privacy, or efficiency. Among the methods, perturbing the representations with an additive noise has been studied by a number of prior works \citep{gao2019convergenceadversarialtrainingoverparametrized,li2025distributionallyrobustoptimizationadversarial,li2023rethinking, madry2018towards, loo2022evolution, tsilivis2022can, ilyas2019adversarial}, showcasing both the benefit of such perturbation and the stable convergence of the training under this setting. Compared with additive noise, perturbing the representations by multiplying it with a mask has rarely been studied theoretically.  

Perturbing the representations with multiplicative noise appears in many real-world settings, either by design or unintentionally. For instance, in federated learning (FL) settings \citep{mcmahan2017communication, kairouz2021advances}, particularly vertical FL \citep{cheng2020federated, liu2021fate, romanini2021pyvertical, he2020fedml, liu2022fedbcd, liu2024vertical}, different features of the input data may be available to different parties, effectively creating a form of sparsity-inducing multiplicative masking on the input space. Moreover, the drop-out family \citep{srivastava2014dropout,rey2021gaussian} is a class of methods to prevent overfitting and improve generalization ability of neural networks during training. Lastly, training models under data-parallel protocol over a wireless channel incurs the channel effect that blurs the data passed to the workers through a multiplication \cite{TseViswanath2005}. 

Theoretically analyzing the training dynamics of neural networks under these settings are difficult, especially when the introduced randomness are intertwined with the nonlinearity of the activation function. While there has been previous work that studies the convergence of neural network training under drop-out \citep{liao2022convergence,mianjy2020convergence}, they often assume that the drop-out happens after the nonlinear activations are applied. From a technical perspective, statistics of the neural network outputs are easier to handle as the randomness are not affected by the nonlinearity. 


In this paper, we take a step further into the understanding of multiplicative perturbations in neural network training by considering noise applied before the nonlinear activation. In particular, the setting we consider is the training of a two-layer MLP where the inputs bears a multiplicative Gaussian mask. This prototype provides a simplified scenario to study the noise-inside-activation difficulty, while generalizes various training scenarios ranging from input masking \citep{kariotakis2024leveraging} to Gaussian drop-out \citep{rey2021gaussian}, if one views the input in our setting as fixed embeddings from previous layers of a deep neural network. Under this setting, we aim to answer the following question:
\begin{center}
    \textit{How do multiplicative perturbations at the input level propagate through the network\\ and affect the training dynamics?} 
\end{center}
\textbf{Our Contributions.}
Analyzing the training dynamics under the Gaussian masks over the input means that we have to study the statistical properties of random variables inside a non-linear function. 
Our work takes a step towards resolving this technical difficulty. Moreover, we utilize an NTK-based analysis \citep{du2018gradient,song2020quadraticsufficesoverparametrizationmatrix,oymak2019moderateoverparameterizationglobalconvergence,liao2022convergence} to study the training convergence of the two-layer MLP under sufficient overparameterization.

To our knowledge, this work provides the first convergence analysis for neural network training under Gaussian multiplicative input masking. 
Specifically, for inputs $\bfx$ masked by $\bfx \odot \bfc$ where $\bfc \sim \mathcal{N}(\mathbf{1}, \kappa^2\mathbf{I})$, we prove that: $i)$ The expected loss decomposes into a smoothed neural network loss plus an adaptive regularization term; $ii)$ Training achieves linear convergence to an error ball of radius $O(\kappa)$. %$iii)$ The regularization effect provides implicit robustness benefits.
Empirical results showcase and support our theory.

\textbf{Our Contributions.}
To our knowledge, this work provides the first convergence analysis for neural network training under Gaussian multiplicative input masking. Our main contributions are summarized as follows:
\begin{itemize}[leftmargin=*]
    \item \textit{Theoretical Analysis of Input Masking.} We provide a rigorous characterization of the training dynamics for two-layer ReLU networks where noise is injected \textit{before} the non-linear activation. We overcome the technical challenge of resolving the expectation of non-linear functions of random variables, proving that the expected loss decomposes into a smoothed objective plus an adaptive, data-dependent regularizer.
    \item \textit{General Stochastic Training Framework.} We develop a general convergence theorem for overparameterized neural networks trained with biased stochastic gradient estimators. This result, which establishes linear convergence to a noise-dependent error ball, is of independent interest beyond the specific setting of Gaussian masking.
    \item \textit{Explicit Convergence Guarantees.} We derive constructive bounds for the convergence rate and the final error radius. We show explicitly how these quantities depend on the mask variance $\kappa^2$, the network width $m$, and the initialization scale, demonstrating that the training converges linearly up to a floor determined by the noise level.
    \item \textit{Empirical Validation and Privacy Utility.} We confirm our theoretical predictions regarding the expected gradient and loss landscape through simulations. Furthermore, we demonstrate the practical utility of this training regime as a defense against Membership Inference Attacks (MIA), highlighting a favorable trade-off between privacy and utility.
\end{itemize}


\section{Related Work}

\textbf{Neural Network Robustness.}
The study of neural network robustness has a rich history, with early work focusing primarily on additive perturbations. Results such as \citep{bartlett2017spectrally} and \citep{miyato2018virtual} established generalization bounds for neural networks under adversarial perturbations, showing that the network's Lipschitz constant plays a crucial role in determining robustness. 
Subsequent work by \citep{cohen2019certified} introduced randomized smoothing techniques for certified robustness against $\ell_2$ perturbations, while \citep{wong2018scaling} developed methods for training provably robust deep neural networks.

Regularization techniques have emerged as powerful tools for enhancing network robustness. 
Dropout \citep{srivastava2014dropout} pioneered the idea of randomly masking \textit{internal} neurons during training, effectively creating an implicit ensemble of subnetworks \citep{yuan2022distributed, hu2023federated, kariotakis2024leveraging, wolfe2023gist, liao2022convergence, dun2023efficient, dun2022resist}. 
This connection between feature masking and regularization was further explored in \citep{ghorbani2021linearized}, who showed that dropout can be interpreted as a form of data-dependent regularization. 
Note that sparsity-inducing norms, based on Laplacian continuous distribution, have a long history in sparse recovery problems \citep{bach20112, jenatton2011structured, bach2012optimization, kyrillidis2015structured}.
Empirical studies on the effect of sparsity, represented by multiplicative Bernoulli distributions, can be found in \citep{kariotakis2024leveraging}.

\textbf{Neural Tangent Kernel (NTK).} \cite{jacot2018neural} discovers that infinite-width neural network evolves as a Gaussian process with a stable kernel computed from the outer product of the tangent features of the neural network. Later works adopted finite-width correction and applied the framework to the analysis of neural network convergence \cite{du2018gradient,du2019gradientdescentfindsglobal,oymak2019moderateoverparameterizationglobalconvergence}. 
The Neural Tangent Kernel framework is one of the few theoretical tools focused on theoretical understanding neural network training. Later works extended the proof to classification tasks, where the notion of tangent features is considered as a feature mapping onto a space where the training data are separable \cite{ji2020polylogarithmicwidthsufficesgradient}. Although the NTK framework has been treated as "lazy training" that prevents useful featuers to be learned, it enables exact analysis of the neural network training dynamic under various scenarios for different architectures \cite{nguyen2021proofglobalconvergencegradient, du2019graphneuraltangentkernel,truong2025globalconvergenceratedeep,wu2023convergenceencoderonlyshallowtransformers}. Based on the NTK framework, several paper studies the convergence of training shallow neural networks under random neuron masking (e.g. dropout \cite{srivastava2014dropout}) \cite{liao2022convergence,mianjy2020convergence}. However, these works usually considers the masking applied after the nonlinearity is applied, which allows direct computation of the statistics of the output and gradient under randomness.
%%%%%%
\section{Problem Setup}
\label{sec:problem}
Given a dataset $\left\{\paren{\bfx_i, y_i}\right\}_{i=1}^n$, we are interested in training a neural network $f\paren{\bm{\theta},\cdot}$ that maps each input $\bfx_i\in\R^d$'s to an output $f\paren{\bm{\theta},\bfx_i}$ that fits the labels $y_i\in\R$. 
We consider $f\paren{\bm{\theta},\cdot}$ as a two-layer ReLU activated Multi-Layer Perceptron (MLP) under the NTK scaling:
\[
    f\paren{\bm{\theta},\bfx} = \frac{1}{\sqrt{m}}\sum_{r=1}^ma_r\sigma\paren{\bfw_r^\top\bfx},
\]
where $\bm{\theta} = \paren{\left\{\bfw_r\right\}_{r=1}^m,\left\{a_r\right\}_{r=1}^m}$ denotes the neural network parameters, and $\bm{\sigma}\paren{\cdot} = \max\{0,\cdot\}$ denotes the ReLU activation function. 
We assume that the second-layer weights $a_r\in\{\pm 1\}$ are fixed, and only the first layer weights $\bfw_r$'s are trainable. 
Thus, we will be using $f(\mathbf{W}, \mathbf{x}) \equiv f(\bm{\theta}, \mathbf{x})$ where $\mathbf{W} \in \mathbb{R}^{m \times d}$, unless otherwise stated.
This neural network set-up is studied widely in previous works \citep{du2018gradient}. 
We consider the training of the neural network by minimizing the MSE loss $\calL\paren{\mathbf{W}}$ over the dataset $\left\{\paren{\bfx_i, y_i}\right\}_{i=1}^n$:
\begin{equation*}
    \calL\paren{\mathbf{W}} = \frac{1}{2}\sum_{i=1}^n\paren{f\paren{\mathbf{W},\bfx_i}-y_i}^2.
\end{equation*}
With the influence of the ReLU activation, the loss is both non-convex and non-smooth. 
However, a line of previous works \citep{du2018gradient, song2020quadraticsufficesoverparametrizationmatrix, oymak2019moderateoverparameterizationglobalconvergence} proves a linear convergence rate of the loss function under the assumption that the number of hidden neurons is sufficiently large by adopting an NTK-based analysis \citep{jacot2018neural}. 

\textit{While there have been theoretical approaches and assumptions that go beyond the NTK assumption, our focus is on a generalized scenario where the input data may be corrupted in each iteration under a multiplicative Gaussian noise:} 
Let $\bfc\sim\mathcal{N}\paren{\bm{1}_d,\kappa^2\bfI_d}$ be an isotropic Gaussian random vector centered at the all-one vector $\bm{1}_d$, the neural network output is given by $f\paren{\bfW,\bfx\odot\bfc}$, 
where $\odot$ denotes the Hadamard (element-wise) product between two vectors. 
Under the multiplicative noise, the neural network is trained with gradient descent where each gradient is computed based on the surrogate loss $\calL_{\bfC}\paren{\mathbf{W}}$ defined over the neural network with the masked input:
\[
    \calL_{\bfC}\paren{\mathbf{W}} = \frac{1}{2}\sum_{i=1}^n\paren{f\paren{\mathbf{W},\bfx_i\odot\bfc_i}- y_i}^2.
\]
Here $\bfC = \{\bfc_i\}_{i=1}^n$ denotes the collection of the masks for all input $\bfx_i$. 
We assume that $\bfc_i$'s are independent. 
In real-world applications, this scheme can be considered as training on an imprecise hardware, where each input data point is read-in with noise. Alternatively, one could view each $\bfx_i$ as the output of a pre-trained large model, and our training scheme can be considered as fine-tuning the last two layers with the Gaussian drop-out \citep{wang2013fast,kingma2015variational,rey2021gaussian} in the intermediate layer. 

Let $\left\{\mathbf{W}_k\right\}_{k=1}^K$ be generated from the stochastic gradient descent given by:
\begin{equation}
    \label{eq:gaussin_input_mask_gd}
    \mathbf{W}_{k+1} = \mathbf{W}_k - \eta\nabla_{\mathbf{W}}\calL_{\bfC_k}\paren{\mathbf{W}_k},
\end{equation}
where $\bfC_k$ is sampled independently in every iteration of the gradient descent.
Our goal is to study the convergence of the loss sequence $\left\{\calL\paren{\mathbf{W}_k}\right\}_{k=1}^\infty$. 
Notice that the loss involved in the weight-update is the surrogate loss $\calL_{\bfC_k}\paren{\mathbf{W}_k}$, but the loss we aim to show convergence is the original loss $\calL\paren{\mathbf{W}}$. 

Our set-up marks some differences from previous works. 
First, our set-up is distinct from unbiased estimators in current literature; our setup does not have such favorable property, since the randomness is applied at the input level of the neural network. 
Second, although there is a line of work that analyzes the convergence of vanilla drop-out tranining on two-layer neural networks \citep{liao2022convergence, mianjy2020convergence}, in their analysis the mask is applied to the hidden neurons after the activation function. 
On the contrary, our mask is applied directly to the input, which is contained in the non-linear function. 
Therefore, any analysis of the mask randomness must go through the ReLU function, which brings technical difficulty. We assume the following property for the training data.
\begin{assumption}
    \label{asump:data}
    The training dataset $\{\paren{\bfx_i, y_i}\}_{i=1}^n$ satisfies $\norm{\bfx_i}_2\leq 1, \left|y_i\right| \leq O\paren{1}$, and for any pair $i\neq j$, there exists no real number $q$ such that $\bfx_i = q \cdot \bfx_j$.
\end{assumption}
This assumption guarantees the boundedness of the dataset, and that the input data are non-degenerate, which is a standard assumption in \cite{du2018gradient,song2020quadraticsufficesoverparametrizationmatrix,liao2022convergence}.
%%%%%%%
\section{Expectation of the Loss and Gradient under Gaussian Mask}
\label{sec:surrogate}
A formal mathematical characterization of the expected loss and gradient is essential not only in prior literature of neural network training convergence \citep{liao2022convergence,mianjy2020convergence} but also in the classical analysis of SGD even in the convex domain \citep{shamir2013stochastic,garrigos2023handbook,tang2013convergence}. In this section, we focus on the derivation of the explicit form of the expected surrogate loss $\E_{\bfC}\left[\calL_{\bfC}\paren{\mathbf{W}}\right]$ and the expected surrogate gradient $\E_{\bfC}\left[\nabla_{\mathbf{W}}\calL_{\bfC}\paren{\mathbf{W}}\right]$. 
Starting with gradient calculations, the surrogate gradient with respect to the $r$-th neuron can be written as:
\begin{equation}
    \label{eq:mask_grad}
    \begin{aligned}
        \nabla_{\bfw_r}\calL_{\bfC}\paren{\mathbf{W}} & = \frac{a_r}{\sqrt{m}}\sum_{i=1}^n\paren{f\paren{\mathbf{W},\bfx_i\odot\bfc_i} - y_i} \paren{\bfx_i\odot\bfc_i}\indy{\bfw_r^\top\paren{\bfx_i\odot\bfc_i}\geq 0}
    \end{aligned}
\end{equation}
Setting $\bfc_i = \bm{1}$ for all $i\in[n]$ gives the gradient of the original loss $\nabla_{\bfw_r}\calL\paren{\mathbf{W}}$. Let $\bm{\Phi}_1\paren{\cdot}$ denote the CDF of the standard (one-dimensional) Gaussian random variable, and let $\phi,\psi:\R\rightarrow\R$ be defined as:
\begin{equation}
    \label{eq:phi_psi_defin}
    \phi\paren{x} = \exp\paren{-x^2};\quad \psi\paren{x} = |x| \cdot \phi\paren{x}.
\end{equation}
Observe that $\phi\paren{x}\in (0, 1]$ and $\psi\paren{x} \in\left(0, \sfrac{1}{\sqrt{2e}}\right]$. 
Before we state the results in this section, we need to define the following quantities.
\begin{definition}
    \label{defin:thm_quantities}
    Fix a first-layer weight $\bfW \in \mathbb{R}^{m \times d}$ and training data $\{\paren{\bfx_i, y_i}\}_{i=1}^n$. We define the:
    \begin{itemize}[leftmargin=*]
        \item Data-related quantity: $B_{\bfx} = \max_{i\in[n]}\norm{\bfx_i}_{\infty}, B_y : = \max_{i\in[n]}\left|y_i\right|$.
        \item Weight-related quantity (row-wise): $R_{\bfw} = \max_{r\in[m]}\norm{\bfw_r}_2$.
        \item Mixed quantity: $R_{\bfu} := \max_{r\in[m],i\in[n]}\norm{\bfw_r\odot\bfx_i}_2$ and 
        \begin{align*}
            \psi_{\max} = \max_{r\in[m],i\in[n]}\psi\paren{\frac{\bfw_r^\top\bfx_i}{2\kappa\norm{\bfw_r\odot\bfx_i}_2}},  \phi_{\max} = \max_{r\in[m],i\in[n]}\phi\paren{\frac{\bfw_r^\top\bfx_i}{2\kappa\norm{\bfw_r\odot\bfx_i}_2}}.
        \end{align*}
    \end{itemize}
\end{definition}

\begin{figure}[!tp] % htbp: here, top, bottom, page
    \centering
    \begin{subfigure}[b]{0.45\textwidth}
        \includegraphics[width=\textwidth]{sim2_smoothed_relu_norm1.0.png}
        \caption{}
        \label{fig:sim1_vary_kappa}
    \end{subfigure}
    \begin{subfigure}[b]{0.45\textwidth}
        \includegraphics[width=\textwidth]{./sim1_vary_kappa.png}
        \caption{}
        \label{fig:sim_smoothed_relu}
    \end{subfigure}
    \caption{(a). Effect of the noise standard deviation $\kappa$ on the shape of the smoothed activation function $\hat{\sigma}(z; \kappa) = z \cdot \Phi_1(z / (\kappa \|\mathbf{w}\odot\mathbf{x}\|_2))$, where $z = \mathbf{w}^\top\mathbf{x}$. For this visualization, $\|\mathbf{w}\odot\mathbf{x}\|_2$ is held constant at $1.0$. As $\kappa$ increases, the activation becomes progressively smoother compared to the standard ReLU (dotted black line). For small $\kappa$ (e.g., $\kappa=0.01$), $\hat{\sigma}$ closely approximates the standard ReLU. (b). Theoretical smoothed activation $\hat{\sigma}(\mathbf{w}, \mathbf{x})$ versus its empirical estimate $\mathbb{E}_{\mathbf{c}}[\sigma(\mathbf{w}^\top(\mathbf{x} \odot \mathbf{c}))]$ for a fixed pre-activation value $\mathbf{w}^\top\mathbf{x} \approx 0.77$ (actual value depends on fixed $\mathbf{w},\mathbf{x}$) as the noise standard deviation $\kappa$ varies. The close match across a range of---relatively small---$\kappa$ values validates the theoretical model for $\hat{\sigma}$. Note that this behavior consistently follows empirically for different $\mathbf{w},\mathbf{x}$ values.}
\end{figure}

\noindent \textbf{Expected Surrogate Loss.}
To start, we focus on the expected loss under the Gaussian input mask. For a fixed neural network $f\paren{\mathbf{W},\cdot}$, we have the following result.
\begin{theorem}
    \label{thm:exp_loss_approx}
    Let $\bfu_{i,r} = \bfw_r\odot\bfx_i$. Define the smoothed activation and neural network as:
    \begin{align*}
        \act{\bfw}{\bfx} &= \bfw^\top\bfx\cdot \bm{\Phi}_1\paren{\frac{\bfw^\top\bfx}{\kappa\norm{\bfw\odot\bfx}_2}},\;\;  \hat{f}\paren{\mathbf{W},\bfx}= \frac{1}{\sqrt{m}}\sum_{r=1}^ma_r\act{\bfw_r}{\bfx}.
    \end{align*}
    Let $\phi_{\max},\psi_{\max}, B_y, R_{\bfu}$, and $R_{\bfw}$ be defined in Definition~\ref{defin:thm_quantities}. If $B_y \leq 3\sqrt{m} R_{\bfw}$, then we have that:
    \begin{align*}
        & \E_{\bfC}\left[\calL_{\bfC}\paren{\bfW}\right]  = \mathcal{E} + \underbrace{\frac{1}{2}\sum_{i=1}^n\paren{\hat{f}\paren{\bfW,\bfx_i}-y_i}_2^2}_{\calT_1} + \underbrace{\frac{\kappa^2}{2m}\sum_{i=1}^n\norm{\sum_{r=1}^ma_r\bfu_{i,r}\bm{\Phi}_1\paren{\frac{\bfw_r^\top\bfx_i}{\kappa\norm{\bfu_{i,r}}_2}}}_2^2}_{\calT_2},
    \end{align*}
    with the magnitude of $\mathcal{E}$ bounded by:
    \begin{equation}
        \label{eq:mag_cap_epsilon}
        \left|\mathcal{E}\right| \leq mn\paren{\kappa^2R_{\bfu}^2\psi_{\max}^2 + \paren{\kappa^2R_{\bfu}^2 + \kappa R_{\bfw}}\phi_{\max}^2}.
    \end{equation}
\end{theorem}

\begin{remark}
The core of our analysis of the expected loss involves understanding how the ReLU activation behaves under the multiplicative Gaussian input mask. Lemma~\ref{lem:first_order_gaus_int} in the appendix provides the analytical form for the expectation of a truncated Gaussian random variable. 
This leads to the definition of a smoothed activation function, as presented in Theorem~\ref{thm:exp_loss_approx}.
Figure \ref{fig:sim_smoothed_relu} demonstrates the correspondence between the theoretical and empirical values of this smoothed activation across a range of noise levels $\kappa$ for a fixed input $\mathbf{w}^\top\mathbf{x}$: being an approximation of ReLU, as $\kappa$ increases, it is expected the two curves to deviate, yet for small enough $\kappa$ values (here, $\kappa \lessapprox 0.2$) the two curves coincide.
%Figure~\ref{fig:sim1_vary_wtx} 
Figure~\ref{fig:wtx_proxy} visually compares this theoretical smoothed activation $\hat{\sigma}$ with its empirical estimate $\mathbb{E}_{\mathbf{c}}[\sigma(\mathbf{w}^\top(\mathbf{x} \odot \mathbf{c}))]$ for a fixed $\kappa=0.2$ as the input $\mathbf{w}^\top\mathbf{x}$ varies. 
The close agreement validates our analytical derivation of $\hat{\sigma}$ and illustrates its smoothing effect compared to the standard ReLU. 

Thus, the term $\bm{\Phi}_1\left(\frac{\bfw^\top\bfx}{\kappa\norm{\bfw\odot\bfx}_2}\right)$ can be interpreted as a smoothed version of the indicator function $\mathbb{I}\{\bfw^\top\bfx \geq 0\}$. 
To visualize the impact of the noise variance $\kappa^2$ on the shape of this smoothed activation, see Figure \ref{fig:sim1_vary_kappa} for various values $\mathbf{w}, \mathbf{x}$ (For this illustration, we assume a fixed value for $\|\mathbf{w}\odot\mathbf{x}\|_2=1.0$ to isolate the effect of $z = \mathbf{w}^\top\mathbf{x}$ and $\kappa$). 
As $\kappa$ increases, the transition of $\hat{\sigma}$ around the origin becomes progressively gentler compared to the sharp kink of the standard ReLU activation.
\end{remark}

\begin{remark}
Theorem~\ref{thm:exp_loss_approx} shows that the expected loss can be approximated by the combination of terms $\mathcal{T}_1$ and $\mathcal{T}_2$, with an additive error term defined by $\mathcal{E}$. 
Notice that 
%$\bm{\Phi}_1\paren{\frac{\bfw_r^\top\bfx_i}{\kappa\norm{\bfw_r\odot\bfx_i}_2}}$ can be seen as a smoothed function of $\indy{\bfw_r^\top\bfx_i\geq 0}$, as in the case where $\bfw_r^\top\bfx_i\neq 0$
%\[
%\lim_{\kappa\rightarrow 0}\bm{\Phi}_1\paren{\frac{\bfw_r^\top\bfx_i}{\kappa\norm{\bfw_r\odot\bfx_i}_2}} = \indy{\bfw_r^\top\bfx_i\geq 0}
%\]
%and, at the same time, 
the smoothed activation $\act{\bfw}{\bfx}$ satisfies (see (\ref{eq:expect_act})):
\begin{align*}
    \hat{\sigma}(\bfw, \bfx) &= \bfw^\top\bfx \cdot \Phi_1\left(\frac{\bfw^\top\bfx}{\kappa\norm{\bfw \odot \bfx}2}\right)
    = \E_{\bfc \sim \mathcal{N}(\mathbf{1}, \kappa^2\mathbf{I})}[\sigma(\bfw^\top(\bfx \odot \bfc))]
    \pm O\paren{\kappa\norm{\bfw\odot\bfx}_2\phi\paren{\frac{\bfw^\top\bfx}{\kappa\norm{\bfw\odot\bfx}_2}}}
\end{align*}
% {
% \color{red}
% Tasos: can you clarify this part a bit more? E.g., can we say something like:
% \begin{align}
% \hat{\sigma}(\bfw, \bfx) &= \bfw^\top\bfx \cdot \Phi_1\left(\frac{\bfw^\top\bfx}{\kappa|\bfw \odot \bfx|2}\right) \\
% &= \E_{\bfc \sim \mathcal{N}(\mathbf{1}, \kappa^2\mathbf{I})}[\sigma(\bfw^\top(\bfx \odot \bfc))] + O(\kappa^2)
% \end{align}
% }, 
Therefore, here $\mathcal{T}_1$ can be seen as loss defined on the smoothed neural network $\hat{f}\paren{\mathbf{W},\cdot}$ with the same weights and dataset. 
\end{remark}
%%%%% enter new plots here!
% in preamble:
% \usepackage{subcaption}

\begin{figure}[!tp]
    \centering
    \begin{subfigure}[t]{0.5\linewidth}
        \centering
        \includegraphics[width=\linewidth]{Preprint/tilde.png}
        \caption{Exact expectation $\tilde{\sigma}(\mathbf{w},\mathbf{x})
        = z\Phi(z/\sigma)+\sigma\varphi(z/\sigma)$}
        \label{fig:wtx_exact}
    \end{subfigure}\hfill
    \begin{subfigure}[t]{0.5\linewidth}
        \centering
        \includegraphics[width=\linewidth]{Preprint/fig1.png}
        \caption{Proxy $\hat{\sigma}(\mathbf{w},\mathbf{x})= z\Phi(z/\sigma)$}
        \label{fig:wtx_proxy}
    \end{subfigure}
    \caption{Smoothed ReLU under multiplicative Gaussian input masking for fixed $\kappa=0.2$, where $z=\mathbf{w}^\top\mathbf{x}$ and $\sigma=\kappa\|\mathbf{w}\odot\mathbf{x}\|_2$.
\textbf{(a)} Exact closed-form expectation $\tilde{\sigma}(\mathbf{w},\mathbf{x})=\mathbb{E}_{\mathbf{c}}[\sigma(\mathbf{w}^\top(\mathbf{x}\odot\mathbf{c}))]=z \Phi(z/\sigma)+\sigma \varphi(z/\sigma)$ (as shown in \ref{eq:expect_act})) matches the Monte Carlo estimate.
\textbf{(b)} Proxy smoothed activation $\hat{\sigma}(\mathbf{w},\mathbf{x})=z \Phi(z/\sigma)$ (used in Theorem~\ref{thm:exp_loss_approx}) differs mainly near $z\approx 0$ due to the missing $\sigma \varphi(z/\sigma)$ term.}

    \label{fig:sim1_wtx_both}
\end{figure}

\begin{remark}
One may notice that the form of $\mathcal{T}_2$ is similar to the $\ell_2$ regularization in the ridge regression. To understand $\mathcal{T}_2$, we first notice that:
\[
\nabla_{\bfw_r}\hat{f}\paren{\mathbf{W},\bfx_i}\approx\frac{1}{\sqrt{m}}\sum_{r=1}^ma_r\bfx_i\bm{\Phi}_1\paren{\frac{\bfw_r^\top\bfx_i}{\kappa\norm{\bfu_{i,r}}_2}}.
\]
Therefore, $\mathcal{T}_2$ can approximately be written as:
\begin{align*}
    \mathcal{T}_2 \approx \frac{\kappa^2}{2}\sum_{i=1}^n\norm{\sum_{r=1}^m\nabla_{\bfw_r}{\hat f}\paren{\mathbf{W},\bfx_i}\odot\bfw_r}_2^2
    = \frac{\kappa^2}{2}\sum_{i=1}^n\sum_{j=1}^d\paren{\nabla_{\hat{\bfw}_j}f\paren{\mathbf{W},\bfx_i}^\top\hat{\bfw}_j}^2
    = \texttt{vec}\paren{\bfW}^\top\hat{\bfH}\texttt{vec}\paren{\bfW}.
\end{align*}
%%%%%%%%%%% correction: %%%%%%%%%%%%
% Here $\hat{\bfw}_j$ is the $j$th row of matrix $\bfW = \left[\bfw_1,\dots,\bfw_m\right]\in\R^{d\times m}$, $\texttt{vec}\paren{\bfW} = \text{concate}\paren{\hat{\bfw}}_1,\dots,\hat{\bfw}}_d$ is the concatenation of $\hat{\bfw}_j$'s, and $\hat{\bfH}\in\R^{md\times md}$ is the block-diagonal matrix where the $j$th diagonal block is $\hat{\bfH}_j := \sum_{i=1}^n\nabla_{\hat{\bfw}_j}f\paren{\mathbf{W},\bfx_i}\nabla_{\hat{\bfw}_j}f\paren{\mathbf{W},\bfx_i}^\top\in\R^{m\times m}$ for $j\in[d]$. 
% Intuitively, $\hat{\bfH}$ can be seen as a matrix consisting of the tangent feature's \citep{baratin2021implicit,lejeune2024adaptive} outer products. 
% As a result, $\mathcal{T}_2$ can be seen as the regularization term of $\bfW$ in terms of a norm defined by the tangent feature outer product matrix $\hat{\bfH}$.
% \end{remark}
% VS
% \\
Here $\hat{\bfw}_j$ is the $j$th row of the matrix $\bfW = \left[\bfw_1,\dots,\bfw_m\right]\in\R^{d\times m}$, 
$\texttt{vec}\paren{\bfW} = \operatorname{concat}\paren{\hat{\bfw}_1,\dots,\hat{\bfw}_d}$ 
is the concatenation of the $\hat{\bfw}_j$'s, and 
$\hat{\bfH}\in\R^{md\times md}$ is the block-diagonal matrix whose $j$th diagonal block is 
$\hat{\bfH}_j := \sum_{i=1}^n \nabla_{\hat{\bfw}_j} f\paren{\mathbf{W},\bfx_i} 
\nabla_{\hat{\bfw}_j} f\paren{\mathbf{W},\bfx_i}^\top \in\R^{m\times m}$ for $j\in[d]$. 
Intuitively, $\hat{\bfH}$ can be seen as a matrix consisting of the tangent features' \citep{baratin2021implicit,lejeune2024adaptive} outer products. 
As a result, $\mathcal{T}_2$ can be seen as the regularization term of $\bfW$ in terms of a norm defined by the tangent feature outer product matrix $\hat{\bfH}$.
\end{remark}
\begin{remark}
The magnitude of $\mathcal{E}$ is given in (\ref{eq:mag_cap_epsilon}). At a first glance, one could see that the term decreases monotonically as $\kappa$ decrease, implying a smaller error when $\kappa$ is small. As discussed in the beginning of this section, $\phi_{\max}$ and $\psi_{\max}$ are upper-bounded by some constant. Therefore, in the worst case, we have $\left|\mathcal{E}\right|\leq O\paren{mn\paren{\kappa^2R_{\bfu} + \kappa R_{\bfw}}}$, which scales linearly with $\kappa$. 
%In an ideal scenario when $\left|\bfw_r^\top\bfx_i\right|$ can be lower-bounded, $\phi_{\max}$ and $\psi_{\max}$ decreases exponentially as $\kappa$ decreases, leading to an exponential decrease in $\left|\mathcal{E}\right|$.
\end{remark}

Below, we sketch the proof of Theorem~\ref{thm:exp_loss_approx}. The full proof of Theorem~\ref{thm:exp_loss_approx} is deferred to Appendix~\ref{sec:proof_exp_loss}.

\textit{Proof sketch.} Our proof starts with the decomposition of the expected loss as:
\begin{align*}
\E_{\bfC}\left[\calL_{\bfC}\paren{\mathbf{W}}\right] & = \frac{1}{2}\sum_{i=1}^n\E_{\bfC}\left[\paren{f\paren{\mathbf{W},\bfx_i\odot\bfc_i}}^2\right]+ \frac{1}{2}\sum_{i=1}^ny_i^2 
    - \sum_{i=1}^ny_i\E_{\bfC}\left[\paren{f\paren{\mathbf{W},\bfx_i\odot\bfc_i}}\right].
\end{align*}
It boils down to analyzing the terms $\E_{\bfC}\left[\paren{f\paren{\mathbf{W},\bfx_i\odot\bfc_i}}^2\right]$ and $\E_{\bfC}\left[\paren{f\paren{\mathbf{W},\bfx_i\odot\bfc_i}}\right]$. 
Plugging in $f\paren{\mathbf{W},\bfx_i\odot\bfc_i}$, it suffices to analyze the following expectations:
\begin{align*}
    E_1 &= \E_{\bfc}\left[\sigma\paren{\bfw_r^\top\paren{\bfx\odot\bfc}}\sigma\paren{\bfw_{r'}^\top\paren{\bfx\odot\bfc}}\right];\quad E_2= \E_{\bfc}\left[\sigma\paren{\bfw_r^\top\paren{\bfx\odot\bfc}}\right].
\end{align*}
The trick of evaluating $E_1$ and $E_2$ is to notice that $\bfw_r^\top\paren{\bfx\odot\bfc} = \bfc^\top\paren{\bfw_r\odot\bfx_i}$. Since $\bfc\sim\mathcal{N}\paren{\bm{1},\kappa^2\bfI}$, we must have that $\bfc^\top\paren{\bfw_r\odot\bfx}\sim\mathcal{N}\paren{\bfw_r^\top\bfx,\kappa^2\norm{\bfw_r\odot\bfx}_2^2}$. Therefore, we can define $z_1 = \bfc^\top\paren{\bfw_r\odot\bfx}$ and $z_2 = \bfc^\top\paren{\bfw_{r'}\odot\bfx}$. Then the problem of evaluating $E_1$ and $E_2$ becomes computing:
\begin{align*}
    E_1 &= \E_{z_1,z_2}\left[z_1z_2\indy{z_1\geq 0;z_2\geq 0}\right];\quad E_2= \E_{z_1}\left[z_1\indy{z_1\geq 0}\right].
\end{align*}
Here $\text{Cov}\paren{z_1, z_2} = \paren{\bfw_r\odot\bfx_i}^\top\paren{\bfw_{r'}\odot\bfx_i}$. To complete the proof, we prove the following two lemmas.

\begin{lemma}
    \label{lem:first_order_gaus_int}
    Let $z_1\sim\mathcal{N}\paren{\mu_1,\kappa_1^2}$. Then, we have that:
    \[
        \E\left[z_1\indy{z_1\geq 0}\right]= \frac{\kappa}{\sqrt{2\pi}}\exp\paren{-\frac{\mu^2}{2\kappa^2}} + \mu\bm{\Phi}_1\paren{\frac{\mu}{\kappa}}.
    \]
\end{lemma}
%%%%

\begin{lemma}
    Let $z_1\sim\mathcal{N}\paren{\mu_1,\kappa_1^2}$ and $z_2\sim\mathcal{N}\paren{\mu_2,\kappa_2^2}$, with $\text{Cov}\paren{z_1,z_2} = \kappa_1\kappa_2\rho$. Let $\bm{\Phi}_2\paren{a,b,\rho}$ denote the joint CDF of standard Gaussian random variables $\hat{z}_1,\hat{z}_2$ with covariance $\rho$ at $z_1 = a, z_2=b$. Then, we have:
    \begin{align*}
        \E\left[z_1z_2\mathbb{I}\left\{z_1\geq 0;z_2\geq 0\right\}\right]
        &= \paren{\mu_1\mu_2+\kappa_1\kappa_2\rho}\bm{\Phi}_2\paren{\frac{\mu_1}{\kappa_1},\frac{\mu_2}{\kappa_2},\rho}+ \frac{1}{\sqrt{2\pi}}\paren{\kappa_1\mu_2T_1 +\kappa_2\mu_1T_2}\\
        & \quad\quad\quad + \frac{\kappa_1\kappa_2}{2\pi}\exp\paren{-\frac{1}{2\paren{1-\rho^2}}\paren{\frac{\mu_1^2}{\kappa_1^2} - \frac{2\rho\mu_1\mu_2}{\kappa_1\kappa_2} + \frac{\mu_2^2}{\kappa_2^2}}}
    \end{align*}
    Here, $T_1, T_2$ are defined as:
    \begin{align*}
        T_1 = \exp\paren{-\frac{\mu_1^2}{2\kappa_1^2}}\bm{\Phi}_1\paren{\frac{1}{\sqrt{1-\rho^2}}\paren{\frac{\mu_2}{\kappa_2} - \frac{\rho\mu_1}{\kappa_1}}};\;\; T_2 = \exp\paren{-\frac{\mu_2^2}{2\kappa_2^2}}\bm{\Phi}_1\paren{\frac{1}{\sqrt{1-\rho^2}}\paren{\frac{\mu_1}{\kappa_1} - \frac{\rho\mu_2}{\kappa_2}}}
    \end{align*}
\end{lemma}
Plugging $z_1 = \bfc^\top\paren{\bfw_r\odot\bfx}$ and $z_2 = \bfc^\top\paren{\bfw_{r'}^\top\bfx}$ back into $\E_{\bfC}\left[\paren{f\paren{\mathbf{W},\bfx_i\odot\bfc_i}}^2\right]$ and $\E_{\bfC}\left[\paren{f\paren{\mathbf{W},\bfx_i\odot\bfc_i}}\right]$ and bounding the emerging error terms would give the desired result.
Details are deferred into the appendix.

\textbf{Expected Surrogate Gradient.}
In the following part, we study the expectation of the surrogate gradient.

\begin{theorem}
    \label{thm:exp_grad_form}
    Assume that $\bfc_i\sim\mathcal{N}\paren{\bm{1},\kappa^2\bfI}$ for some $\kappa\leq 1$. Let  $\phi_{\max},\psi_{\max}, B_y, R_{\bfu}$, and $R_{\bfw}$ be defined in Definition~\ref{defin:thm_quantities}.
    Then, we have that:
    \begin{equation}
        \label{eq:exp_grad_form}
        \begin{aligned}
        \E_{\bfC}\left[\nabla_{\bfw_r}\calL_{\bfC}\paren{\bfW}\right] = \nabla_{\bfw_r}\calL\paren{\bfW}+ \bfg_{r}
            + \textcolor{teal}{\frac{3\kappa^2}{m}\sum_{r'=1}^ma_{r'}{\rm Diag}\paren{\bfx_i}^2\bfw_{r'}} \cdot \textcolor{teal}{\mathbb{I}\left\{\bfw_{r}^\top\bfx_i\geq 0;\bfw_{r'}^\top\bfx_i\geq 0\right\}}
        \end{aligned}
    \end{equation}
    where the magnitude of $\bfg_{r}$ can be bounded as:
    \begin{equation}
        \label{eq:grad_err_bound}
        \begin{aligned}
            \norm{\bfg_{r}}_2 \leq \paren{6n\kappa^2B_{\bfx}^2R_{\bfw} + 5n\kappa R_{\bfu}\sqrt{d}}\phi_{\max} + \frac{\sigma_{\max}\paren{\bfX}}{\sqrt{m}}\phi_{\max}\calL\paren{\mathbf{W}}^{\frac{1}{2}}+ 6n\kappa R_{\bfu}\psi_{\max}.
        \end{aligned}
    \end{equation}
\end{theorem}

\begin{remark}
Theorem~\ref{thm:exp_grad_form} shows that the expected gradient can be written as the summation of the vanilla loss gradient $\nabla_{\bfw_r}\calL\paren{\mathbf{W}}$, a term $\mathcal{T}_3$ above, and a gradient error $\bfg_r$. 
The magnitude of $\bfg_r$ is controlled in (\ref{eq:grad_err_bound}). 
As discussed previously, when $\left|\bfw_r^\top\bfx_i\right| > 0$, both $\phi_{\max}$ and $\psi_{\max}$ decreases exponentially as $\kappa$ decreases. 
Note that, although the second term scales with $\calL\paren{\mathbf{W}}$, when $\calL\paren{\mathbf{W}}$ decreases during training, that term will also contribute less to the overall gradient error.
\end{remark}

\begin{remark}
One could observe that the \textcolor{teal}{third term} on the right-hand side of (\ref{eq:exp_grad_form}) is the gradient of the function $\mathcal{R}\paren{\bfW}$ with respect to $\bfw_r$, where $\mathcal{R}\paren{\bfW}$ is given by:
\[
    \mathcal{R}\paren{\bfW} = \frac{3\kappa^2}{m}\norm{\sum_{r'=1}^ma_{r'}\bfw_{r'}\odot\bfx_i\mathbb{I}\left\{\bfw_{r'}^\top\bfx_i\geq 0\right\}}_2.
\]
Therefore, $\mathcal{R}\paren{\bfW}$ can be seen as a scaled version of the $\mathcal{T}_2$-term in Theorem~\ref{thm:exp_loss_approx}. This again verifies the regularization effect of the Gaussian random mask.
\end{remark}

The proof of Theorem~\ref{thm:exp_grad_form} is provided in Appendix~\ref{sec:proof_exp_grad}, and we provide the proof sketch below.

\textit{Proof sketch.} By the form of the surrogate gradient in (\ref{eq:mask_grad}), we focus on the following two terms:
\begin{align*}
    \mathcal{T}_1= f\paren{\mathbf{W},\bfx_i\odot\bfc_i}\paren{\bfx_i\odot\bfc_i}\mathbb{I}\left\{\bfw_r^\top\paren{\bfx_i\odot\bfc_i}\geq 0\right\};  \mathcal{T}_2= y_i\paren{\bfx_i\odot\bfc_i}\mathbb{I}\left\{\bfw_r^\top\paren{\bfx_i\odot\bfc_i}\geq 0\right\}.
\end{align*}
Plug in the form of $f(\bfW,\bfx_i\odot\bfc_i)$, we can write $\mathcal{T}_1$ as:
\begin{align*}
    \mathcal{T}_1 &= \frac{1}{\sqrt{m}}\sum_{r'=1}^ma_{r'} \sigma(\bfw_{r'}^\top\paren{\bfx_i\odot\bfc_i})\cdot\paren{\bfx_i\odot\bfc_i}\mathbb{I}\left\{\bfw_r^\top\paren{\bfx_i\odot\bfc_i}\geq 0\right\}\\
    & = \frac{1}{\sqrt{m}}\sum_{r'=1}^ma_{r'}\text{Diag}\paren{\bfx_i}\bfc_i\bfc_i^\top\mathbb{I}\left\{\bfc_i^\top\paren{\bfw_r\odot\bfx_i}\geq 0\right\}\cdot \mathbb{I}\left\{\bfc_i^\top\paren{\bfw_{r'}\odot\bfx_i\geq 0}\right\}\paren{\bfw_{r'}\odot\bfx_i},
\end{align*}
while $\mathcal{T}_2$ can be written as:
\[
    \mathcal{T}_2 = y_i\bfx_i\odot \paren{\bfc_i\mathbb{I}\left\{\bfc_i^\top\paren{\bfw_r\odot\bfx_i}\geq 0\right\}}.
\]
This allows us to focus on the following quantities instead
\begin{align*}
\E&\left[\bfc\bfc^\top\mathbb{I}\left\{\bfc^\top\bfu\geq 0;\bfc^\top\bfv\geq 0\right\}\right];  \E \left[\bfc\mathbb{I}\left\{\bfc^\top\bfu\geq 0\right\}\right]
\end{align*}
with multi-variate Gaussian random variable analysis. The rest of the proof then proceeds similarly as in the proof of Theorem~\ref{thm:exp_loss_approx}.
%%%%%
\section{Convergence Guarantee of Training with Gaussian Mask}
\label{sec:conv_guarantee}
Here, we study the convergence property of a general framework of stochastic neural network training, which gives us a theoretical result that can be of independent interest.
Recall also the setting in Section \ref{sec:problem}: Consider $f\paren{\mathbf{W},\cdot}$ as a two-layer ReLU activated MLP, as described above.
Let $\bm{\xi}$ denote the randomness in one step of (stochastic) gradient descent. 
Let $\hat{\calL}\paren{\mathbf{W},\bm{\xi}}$ and $\nabla_{\bfw_r}\hat{\calL}\paren{\mathbf{W},\bm{\xi}}$ denote the stochastic loss and the stochastic gradient induced by $\xi$, respectively. 
We consider the sequence $\left\{\mathbf{W}_k\right\}_{k=1}^K$ generated by the following updates:
\begin{equation}
    \label{eq:general_sgd}
    \bfW_{k+1} = \bfW_k - \eta\nabla_{\bfW}\hat{\calL}\paren{\mathbf{W}_k,\bm{\xi}_k}.
\end{equation}
% \textcolor{red}{
% Notice that here $\nabla_{\bfw_r}\hat{\calL}\paren{\mathbf{W},\bm{\xi}}$ may \textit{not} be the gradient of $\hat{\calL}\paren{\mathbf{W},\bm{\xi}}$ with respect to $\bfw_r$. Afroditi: Could we explain this better? I don't understand it.}


Instead, the connection between $\nabla_{\bfw_r}\hat{\calL}\paren{\mathbf{W},\bm{\xi}}$ and $\hat{\calL}\paren{\mathbf{W},\bm{\xi}}$ with respect to $\bfw_r$, along with other requirements, are stated in the assumption below.

\begin{assumption}
    \label{asump:randomness}
    For all $\bm{\xi}, \bfW$, we assume that the following properties hold:
    \begin{gather}
        \label{eq:general_exp_loss_ub}
        \E_{\bm{\xi}}\left[\hat{\calL}\paren{\mathbf{W},\bm{\xi}}\right] \leq 2\calL\paren{\mathbf{W}}+ \varepsilon_1, \\
        \label{eq:general_exp_grad_err}
        \norm{\E_{\bm{\xi}}\left[\nabla_{\bfw_r}\hat{\calL}\paren{\mathbf{W},\bm{\xi}}\right] - \nabla_{\bfw_r}\calL\paren{\mathbf{W}}}_2 \leq \varepsilon_3\calL\paren{\mathbf{W}}^{\frac{1}{2}} + \varepsilon_2, \\
        \label{eq:general_smoothness}
        \norm{\nabla_{\bfw_r}\hat{\calL}\paren{\mathbf{W},\bm{\xi}}}_2^2 \leq \gamma\hat{\calL}\paren{\mathbf{W},\bm{\xi}}.
    \end{gather}
\end{assumption}

Here, (\ref{eq:general_exp_loss_ub}) and (\ref{eq:general_exp_grad_err}) provide an upper bound on the expected loss and the error of the expected gradient. 
(\ref{eq:general_smoothness}) can be seen as a relaxed form of the smoothness. 
Our analysis is based on the standard NTK-type argument as in \citep{du2018gradient,song2020quadraticsufficesoverparametrizationmatrix,liao2022convergence}, which considers the infinite-width NTK $\bfH^\infty$ given by:
\begin{equation}
    \label{eq:infinite_NTK_def}
    \bfH^\infty_{ij} = \bfx_i^\top\bfx_j \E_{\bfw\sim\mathcal{N}\paren{\bm{0},\bfI}}\left[\indy{\bfw^\top\bfx_i\geq 0;\bfw^\top\bfx_j\geq 0}\right].
\end{equation}
It is shown in \cite{du2018gradient} that $\bfH^\infty$ is positive definite. We define $\lambda_0:=\lambda_{\min}\paren{\bfH^\infty} > 0$.

\begin{theorem}
    \label{thm:gen_sto_conv}
    Assume that the first-layer weights of a neural network are initialized according to $\bfw_{0,r}\sim\mathcal{N}\paren{\bm{0}, \tau^2\bfI}$ for some $\tau > 0$, and the second-layer weights are initialized according to $\bfa_r\sim{\rm Unif}\{\pm 1\}$. 
    Let the number of hidden neurons satisfy $m = \Omega\paren{\frac{n^4K^2}{\lambda_0^4\delta^2\tau^2}}$ and the step size satisfy $\eta = O\paren{\frac{\lambda_0}{n^2}}$. 
    Assume that Assumptions~\ref{asump:data}, \ref{asump:randomness} hold for some $\gamma = C_1\cdot\frac{n}{m}$ with some small enough $\varepsilon_1, \varepsilon_2, \varepsilon_3$ satisfying:
    \begin{equation}
        \varepsilon_1\leq O\paren{\frac{\delta m}{nK^4}}, \quad \varepsilon_2 \leq O\paren{\frac{\delta\lambda_0}{nK^2}}, \quad \varepsilon_3\leq O\paren{\frac{\lambda_0}{\sqrt{mn}}}.
    \end{equation}
    Then, with probability at least $1 - 2\delta - n^2\exp\paren{-\frac{n^3}{\delta^2\tau^2\lambda_0^3}}$, for all $k\in[K]$, the sequence $\{\mathbf{W}_k\}_{k=1}^K$ generated by (\ref{eq:general_sgd}) satisfies:
    \begin{equation}
        \label{eq:gen_sto_conv_form}
        \begin{aligned}
            \E_{\bm{\xi}_0,\dots,\bm{\xi}_{k-1}}\left[\calL\paren{\mathbf{W}_k}\right] & \leq \paren{1 - \frac{\eta\lambda_0}{2}}^k\calL\paren{\mathbf{W}_0}+ O\paren{\frac{mn}{\lambda_0^2}\cdot\varepsilon_2^2 + \varepsilon_1}.
        \end{aligned}
    \end{equation}
    Furthermore, we can guarantee that $\norm{\bfw_{k,r} - \bfw_{0,r}}_2\leq O\paren{\frac{\tau\lambda_0}{n}}$ for all $r\in[m]$ and $k\in[K]$.
\end{theorem}
In short, Theorem~\ref{thm:gen_sto_conv} shows that under a small enough $\varepsilon_1,\varepsilon_2, \varepsilon_3$ and $\gamma$, if the neural network is sufficiently overparameterized, then, with a small enough step size $\eta$, we can guarantee the convergence under the training given by (\ref{eq:general_sgd}) and that the change in each $\bfw_r$ is bounded by $O\paren{\frac{\tau\lambda_0}{n}}$. 
As shown in (\ref{eq:gen_sto_conv_form}), the expected loss converges linearly up to a ball around the global minimum with radius given by $O\paren{\frac{mn}{\lambda_0^2}\cdot\varepsilon_2^2 + \varepsilon_1}$. 
This error region monotonically decreases as the error in the expected loss and gradient, namely $\varepsilon_1$ and $\varepsilon_2$, decreases.

\noindent \textbf{Training Convergence with Gaussian Input Mask.}
We apply the general result in Theorem~\ref{thm:gen_sto_conv} to the scenarios of Gaussian input masking, as given by (\ref{eq:gaussin_input_mask_gd}). To apply Theorem~\ref{thm:gen_sto_conv}, one need to make sure that the requirements in Assumption~\ref{asump:randomness} are guaranteed. Here, we present two corollaries as extensions of Theorem~\ref{thm:exp_loss_approx} and Theorem~\ref{thm:exp_grad_form}, with the goal of showing (\ref{eq:general_exp_loss_ub}) and (\ref{eq:general_exp_grad_err}).

\begin{corollary}
    \label{cor:exp_loss_ub}
    Let $B_y, \phi_{\max},\psi_{\max},R_{\bfu}$, and $R_{\bfw}$ be defined in Definition~\ref{defin:thm_quantities}. If $B_y\leq 3\sqrt{m}R_{\bfw}$, then we have
    \begin{align*}
        \E_{\bfC}\left[\calL_{\bfC}\paren{\mathbf{W}}\right] \leq 2\calL\paren{\mathbf{W}} + 2mn\kappa^2R_{\bfu}^2
        + mn\paren{\kappa^2R_{\bfu}^2 + \kappa R_{\bfw}}\phi_{\max}^2 +
        mn\kappa^2\paren{R_{\bfu}^2+1}\psi_{\max}^2
    \end{align*}
\end{corollary}

Corollary~\ref{cor:exp_loss_ub} follows simply from Theorem~\ref{thm:exp_loss_approx} by upper-bounding the difference between the smoothed neural network function $\hat{f}\paren{\mathbf{W},\cdot}$ and the vanilla neural network function $f\paren{\mathbf{W},\cdot}$, and by upper-bounding the regularization term. In particular, Corollary~\ref{cor:exp_loss_ub} implies the bound of the error $\varepsilon_1$ as $2mn\kappa^2R_{\bfu} + mn\paren{\kappa^2R_{\bfu}^2 + \kappa R_{\bfw}}\phi_{\max}^2 +  mn\kappa^2\paren{R_{\bfu}^2+1}\psi_{\max}^2$. 

\begin{corollary}
    \label{cor:exp_grad_err_ub}
    Let $B_y, \phi_{\max},\psi_{\max},R_{\bfu}$, and $R_{\bfw}$ be defined in Definition~\ref{defin:thm_quantities}. If $B_y\leq 3\sqrt{m}R_{\bfw}$, then we have
    \begin{align*} \norm{\E_{\bfC}\left[\nabla_{\bfw_r}\calL_{C}\paren{\mathbf{W}}\right] - \nabla_{\bfw_r}\calL\paren{\mathbf{W}}}_2
        &\leq O\paren{\paren{n\kappa^2B_{\bfx}^2R_{\bfw} + n\kappa R_{\bfu}\sqrt{d}}\phi_{\max}} + O\paren{\frac{\sigma_{\max}\paren{\bfX}\phi_{\max}}{\sqrt{m}}}\calL\paren{\mathbf{W}}^{\frac{1}{2}}\\
        & \quad\quad\quad  + O\paren{n\kappa R_{\bfu}\psi_{\max}+\kappa^2\sqrt{m}B_{\bfx}^2R_{\bfw}}
    \end{align*}    
\end{corollary}

Similar to Corollary~\ref{cor:exp_loss_ub}, Corollary~\ref{cor:exp_grad_err_ub} follows from upper bounding the regularization term in Theorem~\ref{thm:exp_grad_form}. By Corollary~\ref{cor:exp_grad_err_ub}, we can write $\varepsilon_2$ and $\varepsilon_3$ in Assumption~\ref{asump:randomness} as $\varepsilon_2 = O\paren{\paren{n\kappa^2B_{\bfx}^2R_{\bfw} + n\kappa R_{\bfu}\sqrt{d}}\phi_{\max}} + O\paren{n\kappa R_{\bfu}\psi_{\max}+\kappa^2\sqrt{m}B_{\bfx}^2R_{\bfw}}$ and likely $\varepsilon_3 = O\paren{\frac{\sigma_{\max}\paren{\bfX}\phi_{\max}}{\sqrt{m}}}$. The proof of Corollary~\ref{cor:exp_loss_ub} and Corollary~\ref{cor:exp_grad_err_ub} are deferred to Appendix~\ref{sec:proof_conv_form}. To complete the requirements in Assumption~\ref{asump:randomness}, we can show the following lemma for (\ref{eq:general_smoothness}).

\begin{lemma}
    \label{lem:general_smoothness}
    Assume that Assumption~\ref{asump:data} holds. Then, we have:
    \[
\norm{\nabla_{\bfw_r}\calL_{\bfC}\paren{\mathbf{W}}}_2^2\leq C_1\frac{n}{m}\calL_{\bfC}\paren{\mathbf{W}}.
    \]
\end{lemma}
The proof can be found in the appendix \ref{lem:general_smoothness}
%%%% previous proof
% \begin{proof}
%     By the form of $\nabla_{\bfw_r}\calL_{\bfC}\paren{\mathbf{W}}$ in (\ref{eq:mask_grad}), we have:
%     \begin{align*}
%         \norm{\nabla_{\bfw_r}\calL_{\bfC}\paren{\mathbf{W}}}_2& \leq \frac{C\left|a_r\right|}{\sqrt{m}}\sum_{i=1}^n\left|f\paren{\mathbf{W},\bfx_i\odot\bfc_i}-y_i\right|\cdot\norm{\bfx}_i\\
%         & \leq \frac{1}{\sqrt{m}}\sum_{i=1}^n\left|f\paren{\mathbf{W},\bfx_i\odot\bfc_i}-y_i\right|\\
%         & \leq \frac{\sqrt{n}}{\sqrt{m}}\calL_{\bfC}\paren{\mathbf{W}},
%     \end{align*}
%     where, in the first inequality, we use the fact that the indicator function is upper-bounded by 1, and, in the second inequality, we use the fact that $a_r = \pm 1$ and $\norm{\bfx_i}_2\leq 1$.
% \end{proof}
%%%%%

With the help of Corollary~\ref{cor:exp_loss_ub},\ref{cor:exp_grad_err_ub}, and Lemma~\ref{lem:general_smoothness}, we can derive the convergence guarantee of training the two-layer ReLU neural network under Gaussian input mask.

\begin{theorem}
    \label{thm:conv_gaussian_mask}
    Assume that the first-layer weights are initialized according to $\bfw_{0,r}\sim\mathcal{N}\paren{\bm{0}, \tau^2\bfI}$ for some $\tau > 0$, and the second-layer weights are initialized according to $\bfa_r\sim{\rm Unif}\{\pm 1\}$. Let the number of hidden neurons satisfy $m = \Omega\paren{\frac{n^4K^2}{\lambda_0^4\delta^2\tau^2}}$ and the step size satisfy $\eta = O\paren{\frac{\lambda_0}{n^2}}$. Assume that for all $\bfW\in\left\{\bfW_k\right\}_{k=1}^K$, the following hold:
    \begin{align}
        \label{eq:kappa_req}
        &\kappa = O\paren{\frac{\sqrt{\delta\lambda_0}}{\tau^2K^2\paren{m^{\frac{1}{4}}\sqrt{d}+nd}\paren{\hat{\phi}_{\max}+\hat{\psi}_{\max}}}} \\
        \label{eq:mix_req}
&\sigma_{\max}\paren{\bfX}\hat{\phi}_{\max}\leq O\paren{\frac{\lambda_0}{\sqrt{n}}}.
    \end{align}
    Then, we have that, with probability at least $1 - 2\delta - n^2\exp\paren{-\frac{n^3}{\delta^2\tau^2\lambda_0^3}}$, for all $k\in[K]$, the sequence $\{\mathbf{W}_k\}_{k=1}^K$ generated by (\ref{eq:general_sgd}) satisfies:
    \begin{align}
        \label{eq:gaus_mask_conv_form}
            \E_{\bfC_0,\dots,\bfC_{k-1}}\left[\calL\paren{\mathbf{W}_k}\right] &\leq \paren{1 - \frac{\eta\lambda_0}{2}}^k\calL\paren{\mathbf{W}_0} + \textcolor{purple}{O\paren{\kappa\tau^2mn^2d^2\paren{\hat{\phi}_{\max}^2 + \hat{\psi}_{\max}^2}}} \\ &\quad \quad \quad + \textcolor{teal}{O\paren{\kappa^2\tau^2m^2nd}} + \textcolor{orange}{O\big(\kappa\tau m n \sqrt{d}  \hat\phi_{\max}^2\big)}        
    \end{align}
    where $\hat{\phi}_{\max} = \max_{k\in[K]}\phi_{\max}\paren{\bfW_k}$ and $\hat{\psi}_{\max} = \max_{k\in[K]}\psi_{\max}\paren{\bfW_k}$.
\end{theorem}

In short, Theorem~\ref{thm:conv_gaussian_mask} guarantees the convergence of training a two-layer ReLU neural network under Gaussian input mask in the form of (\ref{eq:gaus_mask_conv_form}) under the condition of sufficient overparameterization, proper step size, and the requirement in (\ref{eq:kappa_req}) and (\ref{eq:mix_req}). 
In particular, (\ref{eq:kappa_req}) requires a sufficiently small Gaussian variance $\kappa$. 
The condition in (\ref{eq:mix_req}) requires either a small maximum singular value of the input data matrix $\bfX$, or a small $\phi_{\max}$. 
Lastly, (\ref{eq:gaus_mask_conv_form}) shows a linear convergence of the expected loss up to some error region. 
Notice that the \textcolor{purple}{first part} of the error region depends both on $\kappa,\tau$ and on $\phi_{\max}$ and $\psi_{\max}$, and the \textcolor{teal}{second part} of the error region depends solely on $\kappa$ and $\tau$. 
This means that one can guarantee an arbitrarily small error region when the Gaussian noise $\kappa$ and the initialization scale $\tau$ is sufficiently small.

\begin{remark}
Both the requirement and the error region in Theorem~\ref{thm:conv_gaussian_mask} depend on the quantity $\hat{\phi}_{\max}$ and $\hat{\psi}_{\max}$. 
Recall that:
\begin{align*}
    \hat{\phi}_{\max}= \max_{k,r,i} \left\{\exp\paren{-\frac{\paren{\bfw_r^\top\bfx_i}}{4\kappa^2\norm{\bfw_r\odot\bfx_i}_2^2}}\right\}; \;\;\hat{\psi}_{\max}= \max_{k,r,i} \left\{\left|\bfw_r^\top\bfx_i\right| \cdot \exp\paren{-\frac{\paren{\bfw_r^\top\bfx_i}}{4\kappa^2\norm{\bfw_r\odot\bfx_i}_2^2}}\right\}.
\end{align*}
Both quantities decay exponentially fast as $\kappa$ decays, when $\bfw_r^\top\bfx_i \neq 0$ for all $k,r,i$. 
As the sequence $\left\{\bfW_k\right\}_{k=1}^K$ is generated under the randomness of $\bfC_k$'s, intuitively it is almost never the case where $\bfw_{k,r}^\top\bfx_i = 0$. Therefore, in most cases Theorem~\ref{thm:conv_gaussian_mask} should require only a log-dependency of $\kappa$ on other parameters in order for (\ref{eq:kappa_req}) and (\ref{eq:mix_req}) to be satisfied. However, it should be noticed that $\kappa$ still need to decay in powerlaw if one want to sufficiently decrease the \textcolor{teal}{second part} of the error region. 
\end{remark}

\section{Experiments}
\subsection{Empirical Validation of Training Convergence with Gaussian Mask.}
Theorem~\ref{thm:conv_gaussian_mask} asserts that training a sufficiently overparameterized two-layer ReLU network with Gaussian multiplicative input noise results in linear convergence of the expected loss to an error ball. The radius of this error ball is proportional to the noise variance (controlled by $\kappa$) and other network and data-dependent terms. We empirically verify this convergence behavior.

\textit{Simulation Setup.} 
We train a two-layer ReLU MLP: As a toy example, the network has $d=20$ input features and $m=100$ hidden units. 
The training dataset comprised $n=500$ synthetic samples, with input features $\bfx_i$ normalized such that $\|\bfx_i\|_2 \le 1$, and target values $y_i$ generated from a non-linear function of $\bfx_i$ with small added noise. 
First-layer weights $\mathbf{W}$ were initialized using Kaiming uniform initialization, and second-layer weights $a_r \in \{\pm 1\}$ were fixed. 
The network was trained for $2000$ iterations using full-batch gradient descent with a learning rate of $0.005$. 
We performed separate training runs for different noise levels: $\kappa \in \{0.0, 0.05, 0.2, 0.4, 0.6, 1.0, 2.0\}$. For each run, we tracked the evolution of the clean training loss $\mathcal{L}(\mathbf{W}_k)$.

\textit{Results and Discussion.} The training trajectories, plotted in Figure~\ref{fig:training_trajectories}, illustrate the theoretical predictions.
For clean training ($\kappa=0.0$), the loss exhibits an initial linear convergence phase.
When multiplicative Gaussian noise is introduced, the initial linear convergence trend is preserved. 
However, as training progresses, the loss converges not to the same minimal value but to a distinct error ball, plateauing at a value higher than the clean case.
As expected, the size of this error ball, indicated by the final converged loss value, systematically increases with the noise level $\kappa$. This direct relationship between $\kappa$ and the size of the error ball provides strong empirical support for the convergence guarantees outlined in Theorem~\ref{thm:conv_gaussian_mask}.

\begin{figure}[!tp]
    \centering
    \begin{subfigure}[b]{0.57\textwidth}
        \includegraphics[width=\linewidth]{./sim5_pytorch_trajectories_n500_d20_m100_kappas7.png} 
        \caption{}
        \label{fig:training_trajectories}
    \end{subfigure}
    \begin{subfigure}[b]{0.38\textwidth}
        \includegraphics[width=\linewidth]{Preprint/Dynamics.pdf} 
        \caption{}
        \label{fig:mask}
    \end{subfigure}
    \caption{(a). Training loss $\mathcal{L}(\mathbf{W}_k)$ (log-scale) versus training iteration for a two-layer ReLU network ($n=500, d=20, m=100$) trained with full-batch gradient descent under different levels of input multiplicative Gaussian noise standard deviation $\kappa$. (b). Distributed training with Gaussian mask for differen $\kappa$ and number of local steps.}
\end{figure}

\subsection{Impact of Multiplicative Gaussian Noise on Model Accuracy}
We trained a 1-hidden-layer MLP (4096 hidden units, GELU activation, dropout=0.2) on CIFAR-10 for 80 epochs using AdamW optimization with cosine annealing, label smoothing, and standard data augmentation. During training, we injected multiplicative Gaussian noise $x \leftarrow x \cdot (1 + \kappa Z)\quad \text{where } Z \sim \mathcal{N}(0,1).$ and evaluated the impact of noise strength $\kappa$ on clean test accuracy.
The results reveal that small amounts of MG noise $\kappa \approx 0.2$ improve generalization, acting as an effective regularizer beyond the existing random cropping and flipping. This suggests that modest input perturbations help the model learn more robust features that transfer better to the test set. However, accuracy degrades monotonically beyond this point, dropping to 49.88\% at $\kappa = 1.8.$ 
\begin{figure}[h]
    \centering
    \begin{subfigure}[b]{0.45\textwidth}
        \includegraphics[width=\linewidth]{Preprint/acc_vs_k_mlp.png} 
        \caption{MLP}
        \label{fig:mpl_auc}
    \end{subfigure}
    \begin{subfigure}[b]{0.53\textwidth}
        \includegraphics[width=\linewidth]{Preprint/cnn_last_ones1.png} 
        \caption{CNN}
        \label{fig:acc_vs_k_cnn}
    \end{subfigure}
    \caption{Test accuracy versus multiplicative Gaussian noise strength $\kappa$ for (a) a 1-hidden-layer MLP and (b) a CNN, trained on CIFAR-10. Small noise levels ($\kappa \approx 0.2)$ can improve generalization for the MLP, likely due to regularization effects. In contrast, the CNN exhibits robustness by maintaining baseline accuracy. Beyond this point, accuracy degrades monotonically for both architectures as noise corrupts the training signal.}
    \label{fig:acc_vs_k}
\end{figure}
Next, to assess the impact of multiplicative Gaussian noise on a well-regularized convolutional architecture, we trained a CNN on CIFAR-10 with varying noise strengths $\kappa$. The architecture consists of $4$ convolutional layers ($32\rightarrow 32 \rightarrow 64 \rightarrow 64$ filters with $3\times 3$ kernels), batch normalization after each convolutional layer, three dropout layers with rates $0.2, 0.3$ and $0.5$ respectively, and a fully connected layer with ReLU activation. We trained the CNN for 80 epochs using the Adam optimizer with $lr=10^{-3}$, weight decay$=10^{-5}$ and batch size$=128$. During training, we applied Multiplicative Gaussian (MG) noise $x \leftarrow x \cdot (1 + \kappa Z)\quad \text{where } Z \sim \mathcal{N}(0,1)$, while all accuracy measurements were performed on the clean test set (noiseless). As shown in Figure \ref{fig:acc_vs_k_cnn}, we observe that the model exhibits robustness to low-magnitude MG noise, maintaining its baseline accuracy ($\approx 71\%$) at $\kappa=0.2$. Beyond that point though, accuracy degrades monotonically and excessive noise corrupts the training signal.

\subsection{Application: Distributed Training over Wireless Channels}
Communicating signals over wireless channels incurs fading phenomena to the signals \citep{TseViswanath2005}. Specifically, a time-varying signal $x(t)$ transmitted over channel given by $h(t)$ and additive noise $z(t)$ result in $y(t) = x(t)h(t) + z(t)$. For data parallel distributed training over wireless channels, each input data is passed through the channel to the workers to perform local training. We consider using the Gaussian masked input training scheme studied in this paper as a simplified setup to model the channel fading in wireless communication. In particular, we let $\bfx$ be the signals transmitted $x(t)$, and $\bfc\sim\mathcal{N}\paren{\bm{1}, \kappa^2\bfI_d}$ be the channel effect $h(t)$. For simplicity, we set the additive noise to $0$. The time-depending behavior of $x(t)$ and $h(t)$ are transformed into the masking scheme that in each iteration, a new mask is applied to the sample.

Under this setup, we train a two-layer MLP with 128 hidden neurons for the MNIST dataset using FedAvg. That is, we assume that the total training process is partitioned into multiple global iterations, where in each global iteration, the central server passes the updated model parameter together with the current copy of training data through a wireless channel. Each worker receives the training data with channel fading (in our case, modeled with Gaussian multiplicative noise), and updated its local copy of the model using gradient descent starting from the parameter shared by the server for some number of local steps. After the local update, the workers sends the updated parameter to the central server to perform an aggregation by averaing the worker's weights. Under this setup, we train an aggregated model with 5 workers and the choice of \{1, 20, 40\} local steps with a batch size of 128. We also vary the variance of the Gaussian mask to study the relationship between the number of local steps and the noise scales. For each combination of local steps and Gaussian variance, we run 5 trials and record the mean and standard deviation of the resulting accuracy.

We plot the result in Figure~\ref{fig:mask}. In general, for all choices of the number of local steps, we observe a decay in the test accuracy as the input masking variance grows larger, indicating the negative influence of the noise to the overall training performance. In particular, we can also observe that, in the low noise regime ($\kappa \approx 0$), the resulting final accuracy is not influenced much by the number of local iterations. However, as the noise scale grows larger (larger $\kappa$), the final accuracy decays drastically as we increase the number of local iterations. We hypothesis that this behavior is due to the fact more local iterations allows the workers to fit more to the noise instead of the original signals in the data.

\subsection{Efficacy of Multiplicative Gaussian Noise Against Membership Inference Attacks}
\label{sec:exp_mia}

% While our theoretical results focus on the convergence properties of training with multiplicative Gaussian (MG) noise, a key motivation for studying such noise mechanisms is their potential application in privacy-preserving machine learning. 
In this section, we empirically evaluate the effectiveness of training with input multiplicative Gaussian (MG) noise as a defense against Membership Inference Attacks (MIAs). In simple words, the primary goal of a MIA is to determine if a specific data point $\mathbf{x}$ was part of the training set of a target model $f$.

\noindent \textbf{Threat Model and Attack Methodology.}
We adopt the black-box shadow model attack methodology (Adversary 1) from the ML-Leaks framework~\citep{salem2019mlleaks}. In this setup, the adversary aims to determine whether a specific data point was part of a target model's training set using only the model's output posteriors. Because the adversary lacks the target's training labels, they employ a shadow model trained on a proxy dataset to mimic the target’s behavior. By observing how the shadow model treats its own members versus non-members, the adversary generates labeled data to train an attack model. This binary classifier learns to identify the statistical "signatures" of membership—such as increased confidence or reduced entropy—enabling it to perform membership inference on the original target model. A detailed breakdown of the data partitioning and the five-stage attack pipeline is provided in Appendix~\ref{sec:mia_details}.

\noindent \textbf{Experimental Setup.} 
We evaluate the effectiveness of Multiplicative Gaussian (MG) noise as a privacy defense using the CIFAR-10 dataset. The data is partitioned into four disjoint sets to train and audit both target and shadow models independently. Our evaluation covers two architectures---a fully-connected MLP and a multi-block CNN---to ensure the defense generalizes across different model complexities. We measure privacy leakage by training a Logistic Regression attack model against target models subjected to varying noise intensities ($\kappa \in \{0.0, 0.5, 1.2, 1.8\}$) and training durations (20--120 epochs). The defense is quantified via Precision, Recall, and AUC, where an AUC of 0.5 indicates perfect privacy (random guessing). Detailed hyperparameters, partitioning sizes, and architectural specifications are provided in Appendix~\ref{sec:mia_details}.

\subsubsection{Results and Discussion}

Our experiments confirm that training with multiplicative Gaussian noise systematically enhances a model's resilience to membership inference attacks. The results for the MLP and CNN model are presented in Figure~\ref{fig:mg_mlp} and \ref{fig:mg_cnn} respectively, with the AUC values of our experiments illustrated in Figure~\ref{fig:attack_auc_vs_epochs}.


% \begin{table}[h]
% \centering
% \caption{}
% \label{tab:mia_auc_mlp}
% \begin{tabular}{l|cccc}
% \toprule
% \textbf{Epochs} & \textbf{$\kappa=0.0$} (Standard) & \textbf{$\kappa=0.5$} & \textbf{$\kappa=1.2$} & \textbf{$\kappa=1.8$} \\
% \midrule
% 20 & $0.578$ & $0.563$ & $0.538$ & $0.522$ \\
% 40 & $0.670$ & $0.611$ & $0.550$ & $0.522$ \\
% 60 & $0.720$ & $0.646$ & $0.567$ & $0.541$ \\
% 80 & $0.716$ & $0.664$ & $0.573$ & $0.533$ \\
% 100 & $0.702$ & $0.678$ & $0.581$ & $0.537$ \\
% 120 & $0.782$ & $0.692$ & $0.585$ & $0.543$ \\
% \bottomrule
% \end{tabular}
% \end{table}
%%%
\begin{figure}[h]
    \centering
    \begin{subfigure}[b]{0.45\textwidth}
        \includegraphics[width=\linewidth]{Preprint/attack_auc_vs_epochs_nn_multi_kappaFINAL.png} 
        \caption{MLP}
        \label{fig:mpl_auc}
    \end{subfigure}
    \begin{subfigure}[b]{0.45\textwidth}
        \includegraphics[width=\linewidth]{Preprint/attack_auc_vs_epochs_cnn_multi_kappaFINAL.png} 
        \caption{CNN}
        \label{fig:cnn_auc}
    \end{subfigure}
    \caption{Attack AUC on the target model when it is (a) an MLP and (b) a CNN. Higher values indicate greater privacy leakage. Training with MG noise (larger $\kappa$) consistently reduces attack success.}
    \label{fig:attack_auc_vs_epochs}
\end{figure}
%%%
\begin{figure}[h]
  \centering
\includegraphics[width=0.95\linewidth]{Preprint/nn_plots.png}
  \caption{MLP target. Multiplicative Gaussian noise provides resilience against the attacks as $\kappa$ increases.}
  \label{fig:mg_mlp}
\end{figure}
\begin{figure}[h]
  \centering
\includegraphics[width=0.95\linewidth]{Preprint/cnn_plots.png}
  \caption{CNN target. We observe a similar trend as the MLP}
  \label{fig:mg_cnn}
\end{figure}

% \textcolor{red}{First, we observe the expected behavior for standard training ($\kappa=0.0$). As the number of training epochs increases from 20 to 120, the model overfits more to its training data. Jasper: where can we observe the "overfitting"?} This increased overfitting gap is successfully exploited by the attacker, whose AUC rises from $0.578$ to a significant $0.782$. This validates that our attack implementation correctly captures privacy leakage.
Figure ~\ref{fig:mg_cnn} demonstrates the attack success for $\kappa=0.0$, in which case as the number of training epochs increases from 20 to 120, AUC rises from $0.578$ to a significant $0.782$. This validates that our attack implementation correctly captures privacy leakage.

The central finding is the consistent defensive benefit of multiplicative Gaussian noise. As shown in Figures~\ref{fig:mg_mlp} for MLP, for any given number of training epochs, applying MG noise (increasing $\kappa$) retains (relatively) constant both the precision and recall of the attack. For example, after 120 epochs of training, the standard model is highly vulnerable (Attack AUC = $0.782$). In contrast, the model trained with $\kappa=0.5$ reduces this leakage (AUC = $0.692$), and models with stronger noise achieve even better privacy (AUC = $0.585$ for $\kappa=1.2$ and AUC = $0.543$ for $\kappa=1.8$). This demonstrates a clear dose-response relationship: greater noise variance leads to stronger privacy protection against MIAs. Similar trends were observed for the CNN architecture (see Figure~\ref{fig:mg_cnn}).
%%%%% ACCURACY DROP SECTION %%%%%%%%
Figure ~\ref{fig:acc_vs_k} illustrates the classic privacy-utility trade-off. In essence, by sacrificing some model utility, training with multiplicative Gaussian noise effectively obfuscates the statistical signature of data membership, thereby mitigating privacy risks.
% \begin{figure}[h]
%     \centering
%     \begin{subfigure}[b]{0.48\textwidth}
%         \includegraphics[width=\linewidth]{} 
%         \caption{MLP}
%         \label{fig:mpl_acc}
%     \end{subfigure}
%     \begin{subfigure}[b]{0.48\textwidth}
%         \includegraphics[width=\linewidth]{} 
%         \caption{CNN}
%         \label{fig:cnn_acc}
%     \end{subfigure}
%     \caption{Injecting multiplicative gaussian noise at the input of the model results in a test accuracy drop.}
%     \label{fig:acc_vs_kappa}
% \end{figure}
%%%%%%%%%%%%%%%%%%%%%%%%%%%%%%%%%

% \textcolor{red}{This privacy benefit stems from the regularizing effect of the input noise. 
% Training on perturbed data forces the model to learn a smoother decision boundary, reducing the gap between its confidence on member and non-member points. Jasper: where could we observe this behavior?}
% This is corroborated by the lower utility \textcolor{red}{(test accuracy of $\approx 35-40\%$) Is this data presented?} of the trained target/shadow models, illustrating the classic privacy-utility trade-off. 
% In essence, by sacrificing some model utility, training with multiplicative Gaussian noise effectively obfuscates the statistical signature of data membership, thereby mitigating privacy risks.


%\begin{figure}[h]
%  \centering
%\includegraphics[width=0.95\linewidth]{Preprint/nn plots.png}
%  \caption{MLP target. Multiplicative Gaussian noise provides resilience against the attacks as $\kappa$ increases.}
%  \label{fig:mg_mlp}
%\end{figure}
%\begin{figure}[h]
%  \centering
%\includegraphics[width=0.95\linewidth]{Preprint/cnn plots.png}
%  \caption{CNN target. We observe a similar trend as the MLP}
%  \label{fig:mg_cnn}
%\end{figure}



%{
%\color{red}
%\subsection{Differentially private feature selection}

%We empirically evaluate the utility and feature selection capabilities of our input multiplicative Gaussian noise mechanism, comparing it against a standard additive Gaussian noise approach for achieving differential privacy (DP).

%\noindent\textbf{Problem Setting and Privacy Mechanisms.}
%Given a dataset $\mathcal{D} = \{(\mathbf{x}_i, y_i)\}_{i=1}^N$, where $\mathbf{x}_i \in \mathbb{R}^d$ are feature vectors and $y_i$ are labels, our objective is to train an MLP model $f(\mathbf{W}; \cdot)$ and subsequently infer feature importances from $\mathbf{W}$ under $(\epsilon, \delta)$-differential privacy.

%\textit{1) Input Multiplicative Gaussian Noise (Our Method):} Inputs are privatized as $\mathbf{x}_{i, \text{mult-priv}} = \mathbf{x}_i \odot \mathbf{c}_i$, where $\mathbf{c}_i \sim \mathcal{N}(\mathbf{1}, \kappa^2\mathbf{I}_d)$. 
%We define as the privacy budget the quantity $\epsilon_{\text{mult}} \approx \|\mathbf{x}\|_{\infty}/\kappa$. 

%\textit{2) Input Additive Gaussian Noise (Baseline DP):} Inputs are privatized as $\mathbf{x}_{i, \text{add-priv}} = \mathbf{x}_{i, \text{clipped}} + \mathbf{z}_i$, where $\mathbf{z}_i \sim \mathcal{N}(\mathbf{0}, \sigma_{\rm add}^2\mathbf{I}_d)$. 
%Inputs $\mathbf{x}_i$ are first clipped to an $\ell_2$-norm of $S_{\rm clip}=4.0$. 
%The noise scale $\sigma_{\rm add}$ is then calibrated for a target $(\epsilon_{\text{add}}, \delta=10^{-5})$-DP using the Gaussian mechanism: $\sigma_{\rm add} = (\sqrt{2 \ln(1.25/\delta)} \cdot S_{\rm clip}) / \epsilon_{\text{add}}$.

%\noindent\textbf{Experimental Setup.}
%Here, we show the Adult Income dataset case ($d=65$ processed features after grouping rare categories, $\|\mathbf{x}\|_{\infty} \approx 13.17$); more results in Appendix \ref{appendix:additional_datasets}. 
%A shallow MLP (one hidden layer with $32$ ReLU neurons, $\ell_2$-norm regularization $10^{-3}$) is trained for $100$ epochs (Adam, $L_R = 0.001$, batch $256$).

%\textit{Non-Private Baseline:} An MLP with identical architecture ($\ell_2$-norm regularization of $10^{-2}$ with $200$ iterations with early stopping) achieved $0.8537$ accuracy and $0.9100$ AUC. 
%Its top-5 features included \texttt{num\_capital-gain} and \texttt{num\_age}.

%\textit{Evaluation:} We train private MLPs using both noise mechanisms for target $\epsilon \in \{0.5, \dots, 130.0\}$. Utility (accuracy, AUC) and feature selection consistency (Spearman's $\rho$, Kendall's $\tau$, Jaccard Index for top-5/10 features against the non-private MLP) are evaluated on a test set.

%\noindent\textbf{Results and Discussion.}
%\textit{Model Utility:} Figures~\ref{fig:dpfs_comp_utility_acc} and \ref{fig:dpfs_comp_utility_auc} compare model utility. 
%For weak privacy ($\epsilon \ge 50$), both methods achieve utility close to the non-private baseline (e.g., at $\epsilon \approx 130$, Multiplicative: Acc $0.8533$/AUC $0.9107$; Additive: Acc $0.8528$/AUC $0.9114$). 
%As privacy strengthens ($\epsilon$ decreases), multiplicative noise consistently maintains superior utility. At $\epsilon=10.0$, multiplicative noise yields Acc $0.8492$/AUC $0.9063$, significantly outperforming additive noise (Acc $0.7577$/AUC $0.8810$). 
%This advantage is more pronounced for $\epsilon \le 2.0$; at $\epsilon=0.5$, multiplicative noise achieves Acc $0.7555$/AUC $0.8886$, while additive noise AUC degrades to $0.7783$.

%\begin{figure}[!htp]
%    \centering
%    \includegraphics[width=0.65\linewidth]{dpfs_comparison_accuracy_vs_epsilon.png}
%    \caption{Accuracy vs. $\epsilon$ for MLP models ($d=65, m=32$) trained with multiplicative and additive input noise. Dashed line: non-private MLP baseline ($0.8537$).}
%    \label{fig:dpfs_comp_utility_acc}
%\end{figure}

%\begin{figure}[!htp]
%    \centering
%    \includegraphics[width=0.65\linewidth]{dpfs_comparison_auc_vs_epsilon.png}
%    \caption{ROC AUC vs. $\epsilon$ for MLP models ($d=65, m=32$) trained with multiplicative and additive input noise. Dashed line: non-private MLP baseline ($0.9100$).}
%    \label{fig:dpfs_comp_utility_auc}
%\end{figure}
%\newpage
%\textit{Feature Selection Consistency:} The stability of feature importance rankings is shown in Figures~\ref{fig:dpfs_comp_rankcorr_spearman}, \ref{fig:dpfs_comp_rankcorr_kendall}, \ref{fig:dpfs_comp_jaccard_top5}, and \ref{fig:dpfs_comp_jaccard_top10}.
%Multiplicative noise demonstrates better preservation of feature ranking in lower noise regimes ($\kappa < 1$, corresponding to $\epsilon > 13.17$). 
%For instance, at $\epsilon \approx 130$ ($\kappa \approx 0.10$), Spearman's $\rho$ for multiplicative noise is $0.455$, substantially higher than $0.264$ for additive noise (Figure~\ref{fig:dpfs_comp_rankcorr_spearman}). 
%Kendall's Tau exhibits a similar trend.
%The Jaccard Index for top-5 features (Figure~\ref{fig:dpfs_comp_jaccard_top5}) for multiplicative noise reaches $0.667$ (4 out of 5 common features) for $\epsilon \ge 25$ ($\kappa \le 0.53$), significantly better than additive noise (e.g., $0.429$ at $\epsilon \approx 100$). 
%For top-10 features (Figure~\ref{fig:dpfs_comp_jaccard_top10}), multiplicative noise also shows higher overlap (e.g., $0.429$ vs $0.333$ at $\epsilon \approx 130$).
%However, as privacy becomes stricter ($\epsilon \le 10$, $\kappa \ge 1.3$), feature selection consistency degrades considerably for both methods. 
%Rank correlations drop below $0.2$, and Jaccard indices for top-k features become very low (e.g., top-5 Jaccard for multiplicative noise stabilizes at $0.111$, indicating only one common feature). 
%In this high-noise regime, additive noise often results in near-zero or negative rank correlations, performing worse than multiplicative noise.

%\begin{figure}[!htp]
%    \centering
%    \includegraphics[width=0.65\linewidth]{dpfs_comparison_spearman_vs_epsilon.png}
%    \caption{Spearman's $\rho$ rank correlation (vs. non-private MLP baseline) for models trained with multiplicative and additive noise, versus $\epsilon$.}
%    \label{fig:dpfs_comp_rankcorr_spearman}
%\end{figure}

%\begin{figure}[!htp]
%    \centering
%    \includegraphics[width=0.65\linewidth]{dpfs_comparison_kendall_vs_epsilon.png}
%    \caption{Kendall's Tau rank correlation for models trained with multiplicative and additive noise, versus $\epsilon$.}
%    \label{fig:dpfs_comp_rankcorr_kendall}
% \end{figure}


%\begin{figure}[t]
%  \centering
%  \begin{subfigure}[t]{0.48\linewidth}
%    \centering
%    \includegraphics[width=\linewidth]{dpfs_comparison_jaccard_top5_vs_epsilon.png}
%    \caption{Top-5}
%    \label{fig:dpfs_comp_jaccard_top5}
%  \end{subfigure}\hfill
%  \begin{subfigure}[t]{0.48\linewidth}
%    \centering
%    \includegraphics[width=\linewidth]{dpfs_comparison_jaccard_top10_vs_epsilon.png}
%    \caption{Top-10}
%    \label{fig:dpfs_comp_jaccard_top10}
%  \end{subfigure}
%  \caption{Jaccard Index for top-5 (\subref{fig:dpfs_comp_jaccard_top5}) and top-10 (\subref{fig:dpfs_comp_jaccard_top10}) selected features (vs. non-private MLP baseline) for models trained with multiplicative and additive noise, versus $\epsilon$.}
%  \label{fig:jaccard_pair}
%\end{figure}

%Overall, for the Adult dataset with a shallow MLP, multiplicative Gaussian input noise offers some compelling features over input additive Gaussian noise in terms of predictive utility, especially under stricter privacy guarantees ($\epsilon \le 10$).
%Furthermore, in lower noise regimes ($\kappa < 1$, corresponding to $\epsilon > 13$), multiplicative noise demonstrates superior feature selection consistency. 
%While identifying a stable set of important features remains challenging for both methods under strong privacy constraints, multiplicative noise generally appears less disruptive to feature hierarchies. 
%We repeat this process for two more datasets, Breast Cancer Wisconsin (Diagnostic) Dataset and SpamBase Datasets, where similar arguments are made; see the Appendix.
%}
\section{Conclusion}
\label{sec:conclusion}

This work investigates the fundamental question of how independent multiplicative Gaussian masking affects training dynamics. Focusing on a two-layer ReLU network in the NTK regime, we demonstrate that the masked objective admits a closed-form decomposition into a smoothed square loss plus an explicit, data-dependent regularizer. This structure allows training with the gradient from the masked objective to achieve linear convergence toward a noise-controlled error ball given small step size and large enough over-parameterization. Beyong our theory, We provided experimental result to validate the convergence to small ball, and presented applications in distributed training under channel fading and how the masking can defend against attacks. 

\textbf{Limitations and Future Work.} Our theoretical analysis is currently constrained to the small-noise regime and the extreme over-parameterization typical of NTK models. Furthermore, our proofs rely on the independence of masks across iterations and do not yet incorporate a formal privacy accounting pipeline, such as subsampling or composition. Despite these constraints, multiplicative Gaussian masking represents a provable and practically viable method for injecting input-level uncertainty. These results provide a principled foundation for future exploration of deep networks and more complex noise settings in feature-wise training.

% This work revisits the elementary (yet largely unexplored) question:
% \emph{What happens when every training example is multiplicatively
% perturbed by an independent Gaussian mask?}
% Focussing on the simplest non-trivial architecture, a two–layer ReLU
% network in the NTK regime, we show that:

% \begin{enumerate}[leftmargin=*]
%     \item The masked objective admits a \emph{closed form} decomposition into a smoothed square loss plus an explicit, data–dependent regularizer;
%     \item The gradient enjoys an analogous decomposition, so that a carefully chosen step size yields linear convergence, up to a noise--controlled error ball;
%     \item All bounds are constructive: error radii, step sizes and width requirements are given by explicit algebraic expressions;
%     \item The Gaussian mask can be interpreted as privacy pre-processing step, whenever a preceding clipping step is applied, yet we are not studying whether this holds in theory and it is the case; this is left for future work. Empirical evidence on three tabular data sets confirms a favorable privacy--utility trade-off.
% \end{enumerate}

% \paragraph{Limitations.} Our work has limitations. $i)$ \emph{Small–noise regime.}  The convergence proof requires $\kappa=O(1/\mathrm{poly}(n))$. In practice, larger variances are attractive for privacy but are \emph{not} covered theoretically. $ii)$ \emph{Width and depth.} As in most NTK analyses, we assume extreme over-parameterization and only two layers.  Whether the logarithmic dependence on $\kappa$ persists in deep networks is open. $iii)$ \emph{Independence of masks.} The model redraws a fresh mask at every iteration; on-device learning often uses a \emph{single} noisy sample.  Our bounds do not directly transfer to that setting. $iv)$ \emph{Privacy.} We demonstrate empirical privacy utility but do not provide an end-to-end privacy accounting pipeline (e.g., subsampling and composition across SGD iterations).
% The above constitute fruitful ground for future research. 

% To conclude, multiplicative Gaussian masking provides a provable and practically appealing way to inject uncertainty at the input level.  Our results constitute a first principled step that, we hope, will stimulate further theoretical and empirical work on training under feature-wise noise.

\bibliographystyle{plainnat}
\bibliography{refs}

\newpage
\title{Convergence Analysis of Two-Layer Neural Networks under\\ Gaussian Input Masking\\(Supplementary Material)}
\maketitle

\appendix
%\vspace{0.5em}
%This Appendix provides additional experiments and analysis on two additional datasets (Breast Cancer Wisconsin and Spambase), complementing the results in the main text.

\section{Additional Experimental Results and Related Details.}
\subsection{Empirical Validation of Expected Gradient Properties (Theorem~\ref{thm:exp_grad_form}).}
Theorem~\ref{thm:exp_grad_form} characterizes the expected gradient under Gaussian input masking as $\E_{\bfC}\left[\nabla_{\bfw_r}\calL_{\bfC}\paren{\bfW}\right] = \nabla_{\bfw_r}\calL\paren{\bfW}+ \mathcal{T}_{3,r} + \bfg_{r}$. Here, $\nabla_{\bfw_r}\calL\paren{\bfW}$ is the clean input gradient, $\mathcal{T}_{3,r}$ is a systematic deviation term proportional to $\kappa^2$, and $\bfg_{r}$ is a residual error bounded by Eq.~(\ref{eq:grad_err_bound}).

\textit{Simulation Setup.} We used a two-layer ReLU MLP with $d=20$ input features, $m=100$ hidden units, on $n=500$ synthetic samples ($\|\bfx_i\|_2 \le 1$, $y_i \sim \mathcal{N}(0, 0.5^2)$). 
First-layer weights $\mathbf{W}$ are from $\mathcal{N}(0, 0.1^2)$; second-layer $a_r \in \{\pm 1\}$ are fixed. 
We analyze $\nabla_{\mathbf{w}_r}\mathcal{L}_{\mathbf{C}}(\mathbf{W})$ by averaging $N=2000$ Monte Carlo samples for $\kappa \in [0.001, 1.0]$, for a representative neuron $r$.
For this setup, the clean loss $\mathcal{L}(\mathbf{W}) \approx 71.52$ and $\|\nabla_{\mathbf{w}_r}\mathcal{L}(\mathbf{W})\|_2 \approx 0.981$.

\textit{Results and Discussion.} Our simulations validate the decomposition in Theorem~\ref{thm:exp_grad_form}.
Figure~\ref{fig:grad_components} displays the $\ell_2$-norms of the gradient components versus $\kappa$. The clean gradient norm is constant. The $\mathcal{T}_{3,r}$'s norm, $\|\mathcal{T}_{3,r}\|_2$, scales with $\kappa$ (e.g., from $\approx 8.1 \times 10^{-7}$ at $\kappa=0.001$ to $\approx 0.81$ at $\kappa=1.0$), confirming its theoretical dependence. 
The norm of the empirically estimated expected masked gradient, $\|\mathbb{E}_{\mathbf{C}}[\nabla_{\mathbf{w}_r}\mathcal{L}_{\mathbf{C}}(\mathbf{W})]\|_2$, follows the clean gradient for small $\kappa$ and reflects the vector sum with the growing teal term for larger $\kappa$ in Eq.~\ref{eq:exp_grad_form}.

\begin{figure}[htbp]
    \centering
    \begin{subfigure}[t]{0.49\linewidth}
        \centering
        \includegraphics[width=\linewidth]{./sim4_grad_components.png}
        \caption{$\ell_2$-norms of key components of the expected gradient $\mathbb{E}_{\mathbf{C}}[\nabla_{\mathbf{w}_r}\mathcal{L}_{\mathbf{C}}(\mathbf{W})]$:the clean gradient norm ($\|\nabla_{\mathbf{w}_r}\calL(\mathbf{W})\|_2$), the $\mathcal{T}_{3,r}$'s norm ($\|\mathcal{T}_{3,r}\|_2$), and the total expected masked gradient norm. $\mathcal{T}_{3,r}$ scales with $\kappa^2$.}
        \label{fig:grad_components}
    \end{subfigure}\hfill
    \begin{subfigure}[t]{0.49\linewidth}
        \centering
        \includegraphics[width=\linewidth]{./sim4_gr_error_bound.png}
        \caption{Comparison of the $\ell_2$-norm of the empirically estimated gradient error term, $\|\mathbf{g}_r\|_2$, against its theoretical upper bound from Eq.~(\ref{eq:grad_err_bound}). The empirical error (solid line) remains below the derived bound (dashed line) across all tested $\kappa$. Log-log scale.}
        \label{fig:gr_error_bound}
    \end{subfigure}

    \caption{$\ell_2$-norms of the gradient components 
    (left) and residual error bound check (right).}
    \label{fig:sim4_combined}
\end{figure}
Figure~\ref{fig:gr_error_bound} examines the residual error term $\mathbf{g}_r$. 
It compares the $\ell_2$-norm of the empirically estimated $\mathbf{g}_0$ with its theoretical upper bound from Eq.~(\ref{eq:grad_err_bound}). 
The estimated error norm, $\|\mathbf{g}_r\|_{\text{est}}$, increases with $\kappa$ (from $\approx 1.65 \times 10^{-2}$ at $\kappa=0.001$ to $\approx 0.53$ at $\kappa=1.0$).
Importantly, the theoretical bound on $\|\mathbf{g}_r\|_2$ consistently upper-bounds the empirical error across the entire range of $\kappa$. 
For instance, at $\kappa=0.001$, $\|\mathbf{g}_r\|_{\text{est}} \approx 0.0165$ while its bound is $\approx 7.85$.

In summary, the simulations confirm that the expected gradient under Gaussian input masking, with sufficiently small $\kappa$ values, is well-approximated by the sum of the clean gradient and the $\kappa^2$-dependent term, with a residual error that is effectively bounded by our theoretical derivation.

\subsection{Experimental Details for Section~\ref{sec:exp_mia}}
\label{sec:mia_details}
We adopt the black-box threat model and the shadow model attack methodology (Adversary 1) proposed in the ML-Leaks framework~\citep{salem2019mlleaks}.

\noindent\textit{Threat Model.} The adversary has black-box access to a trained target model $f$. This means the adversary can query the model with any input $\mathbf{x}$ and observe its output posterior probability vector $\mathbf{p} = f(\mathbf{x})$ (i.e., the softmax output over the classes), but has no access to the model's parameters, gradients or original training data. The adversary's goal is to train an \textit{attack model} $A$ that, given the posterior $\mathbf{p}_f(\mathbf{x})$ from the target model for a point $\mathbf{x}$, predicts whether $\mathbf{x}$ was a member of the target's training set.

\noindent\textit{Shadow Model Attack Pipeline.} Since the attacker does not have access to the target model's training set, they cannot directly generate labeled data (member vs. non-member posteriors) to train their attack model. The shadow model technique circumvents this by creating a proxy environment where the attacker controls data membership. The pipeline is as follows:
\begin{enumerate}[leftmargin=*]
    \item \textit{Data Partitioning:} The attacker possesses a dataset $D_{\text{shadow}}$, disjoint from the target's training set but drawn from the same data distribution. This set is split into $D^{\text{Train}}_{\text{Shadow}}$ and $D^{\text{Out}}_{\text{Shadow}}$.
    \item \textit{Shadow Model Training:} A shadow model $S$, which mimics the target model's architecture and training process, is trained on $D^{\text{Train}}_{\text{Shadow}}$.
    \item \textit{Attack Dataset Generation:} The trained shadow model $S$ is queried on its own training data (members, $D^{\text{Train}}_{\text{Shadow}}$) and its hold-out data (non-members, $D^{\text{Out}}_{\text{Shadow}}$). The resulting posterior vectors $\mathbf{p}_S(\mathbf{x})$ are collected. Following \citep{salem2019mlleaks}, the top-3 sorted probabilities of each posterior are used as features: $\phi(\mathbf{p}_S(\mathbf{x})) = (\text{p}_{(1)}, \text{p}_{(2)}, \text{p}_{(3)})$. These feature vectors are labeled ``1'' if $\mathbf{x} \in D^{\text{Train}}_{\text{Shadow}}$ and ``0'' otherwise.
    \item \textit{Attack Model Training:} A binary classifier, the attack model $A$, is trained on this generated dataset of ``(feature, label)'' pairs. It learns to distinguish the statistical ``signature'' of a member's posterior from a non-member's. This signature often manifests as higher confidence (larger $\text{p}_{(1)}$) and lower entropy for members, a result of the target/shadow model overfitting to its training data.
    \item \textit{Inference on Target Model:} To attack the original target model $f$, the adversary queries it with a point of interest $\mathbf{x}$, extracts the features $\phi(\mathbf{p}_f(\mathbf{x}))$, and feeds them to the trained attack model $A$ to get a membership prediction.
\end{enumerate}

\noindent \textbf{Dataset and Partitioning.} 
We use the CIFAR-10 dataset, consisting of 60,000 images. The full pool is shuffled and divided into four disjoint sets of 10,520 images each: \texttt{target\_train} (training MG-protected models), \texttt{target\_test} (non-member audit data), \texttt{shadow\_train} (training shadow models), and \texttt{shadow\_test} (shadow non-member data). All data is normalized using the mean and standard deviation of their respective training sets. For protected models, inputs $\mathbf{x}$ are modified via elementwise multiplication with a random mask $\mathbf{m}$, where $m_i \sim \mathcal{N}(1, \kappa^2)$.

\noindent \textbf{Model Architectures.}
\begin{itemize}
    \item \textbf{MLP (``nn''):} A fully-connected network with one hidden layer of 100 neurons (Tanh activation). Input layer: 3,072 features.
    \item \textbf{CNN (``cnn''):} Two \texttt{Conv-ReLU-MaxPool} blocks, followed by a Tanh-activated fully-connected layer with 100 hidden units.
\end{itemize}

\noindent \textbf{Training and Hyperparameters.} 
Both models are trained using the Adam optimizer (Learning Rate: $10^{-3}$, $\ell_2$ Regularization: $10^{-7}$) for intervals between 20 and 120 epochs. The attack model $A$ is a \texttt{LogisticRegression} classifier (scikit-learn), trained on a balanced dataset of member and non-member posteriors.
\subsubsection{Multiplicative Gaussian noise vs Differential privacy}
We evaluate the privacy-utility tradeoff of Multiplicative Gaussian (MG) noise against Differential Privacy (DP-SGD) using the CIFAR-10 image classification dataset. Following the standard ML-Leaks evaluation protocol [Shokri et al., 2017], we utilize the dataset partitioning described in the previous section and employ a \texttt{Standard CNN} architecture, i.e. a shallow baseline consisting of two convolutional layers ($5 \times 5$ kernels, 32 filters) followed by max-pooling and a fully connected layer. For the multiplicative gaussian defense, we sweep the noise parameter $\kappa \in \{ 0.0, 0.2, 0.4, 0.6, 0.8, 1.0, 1.2 \}$ where $\kappa=0.0$ represents the undefended baseline. For each training batch, we apply element-wise multiplicative noise to input features
$\tilde{x} = x \odot \bigl(1 + \kappa Z\bigr)$ ›where $Z \sim \mathcal{N}(0, I)$ is the standard Gaussian noise. For the the Differential Privacy, (DP-SGD) part, we sweep the noise multiplier $\sigma \in \{ 0.3, 0.5, 0.8, 1.0, 1.5, 2.0, 3.0 \}$ using the opacus library (\cite{yousefpour2021opacus}). We set the per-sample gradient clipping norm $C=1.0$ and target $\delta=10^{-5}$. The empirical findings for this experimet are illustrated in Figure \ref{fig:my_image_1}. The plots map the privacy leakage (Attack Precision and Recall) against the model's utility (Target Accuracy) across the swept noise parameters.
\begin{figure}[htbp]
  \centering
  \includegraphics[width=0.8\linewidth]{Preprint/mg_vs_dp_cnn_isoacc_tradeoff.png}
  \caption{Privacy-utility tradeoff for the Standard CNN. The left panel shows Attack Precision vs. Accuracy, and the right panel shows Attack Recall vs. Accuracy.}
  \label{fig:my_image_1}
\end{figure}

% \textbf{Iso-Accuracy Comparison Methodology}: Comparing privacy mechanisms at arbitrary training points can be misleading, as a "better" privacy score may be correlated with lower model utility. To ensure a fair comparison, we implement an Iso-Accuracy Analysis:  We identify the shared accuracy interval where both MG and DP-SGD results coexist. For example, if MG achieves accuracies in $[44\%, 61\%]$ and DP achieves $[39\%, 55\%]$, the overlap is $[44\%, 55\%]$. Within this interval, we generate a grid of 25 equidistant accuracy points and use linear interpolation to estimate the corresponding Membership Inference Attack (MIA) precision and recall for both mechanisms at these exact accuracy levels.
% %%%
% \begin{figure}[htbp]
%   \centering
%   \includegraphics[width=0.8\linewidth]{Preprint/mg_vs_dp_cnn_isoacc_iso_accuracy.png}
%   \caption{Iso-accuracy analysis plots}
%   \label{fig:my_image_2}
% \end{figure}

% %%%% deeper arch

% \begin{figure}[t]
%   \centering
%   \includegraphics[width=0.8\linewidth]{Preprint/improved_cnn1.png}
%   \vspace{0.5em}
%   \includegraphics[width=0.8\linewidth]{Preprint/improved_cnn2.png}
%   \caption{ImprovedCNN plots.}
%   \label{fig:improved_cnn}
% \end{figure}
\textbf{Evaluation on High-Capacity Architecture}: 
We repeated the evaluation using an \texttt{Improved CNN} architecture. This model features a deeper convolutional structure (convolutional blocks with increasing filter sizes ($32 \rightarrow 64 \rightarrow 128$) using $3 \times 3$ kernels) and also incorporates Batch Normalization and Dropout to achieve higher baseline utility. Figure \ref{fig:improved_cnn_1} below illustrates the results for the improved CNN architecture. We observe that the performance gap between the two methods narrows and the MG noise curve remains much closer to the near-random guess of the attacker for a longer stretch of the accuracy spectrum. 
\begin{figure}[htbp]
  \centering
  \includegraphics[width=0.8\linewidth]{Preprint/improved_cnn1.png}
  \caption{Privacy-utility tradeoff for the Improved CNN. Unlike the Standard CNN, the MG noise curve stays much closer to the DP-SGD curve across the accuracy range}
  \label{fig:improved_cnn_1}
\end{figure}
% The results for the high-capacity model (Fig. \ref{fig:improved_cnn_1}) indicate that the MG noise curve shifts closer to the DP-SGD curve, so as the model architecture improves, the performance gap between MG and DP-SGD narrows.

\section{Proofs in Section~\ref{sec:surrogate}}
In this section, we first prove an exact form of the expected surrogate loss function.
\subsection{General Form of Expected Loss}
\begin{lemma}
    \label{lem:exp_loss}
    Let $\bfu_{i,r} = \bfw_r\odot\bfx_i$, let $\act{\bfw}{\bfx} = \bfw^\top\bfx\cdot\bm{\Phi}_1\paren{\frac{\bfw^\top\bfx}{\kappa\norm{\bfw\odot\bfx}_2}}$, and let $\hat{f}\paren{\bm{\theta},\bfx} = \frac{1}{\sqrt{m}}\sum_{r=1}^ma_r\act{\bfw_r}{\bfx}$. Then we have
    \begin{align*}
        \E_{\bfC}\left[\calL_{\bfC}\paren{\bm{\theta}}\right] & = \frac{1}{2}\sum_{i=1}^n\paren{\hat{f}\paren{\bm{\theta},\bfx_i}-y_i}_2^2 + \frac{\kappa^2}{2m}\sum_{i=1}^n\norm{\sum_{r=1}^ma_r\bfu_{i,r}\bm{\Phi}\paren{\frac{\bfw_r^\top\bfx_i}{\kappa\norm{\bfu_{i,r}}_2}}}_2^2\\
        & \quad\quad\quad + \frac{1}{m}\sum_{i=1}^n\sum_{r,r'=1}^ma_ra_{r'}\paren{C_{i,r,r'} + \frac{\kappa^2}{2\pi}E_{i,r,r'}}\\
        & \quad\quad\quad + \frac{2\kappa}{\sqrt{2\pi m}}\sum_{i=1}^n\sum_{r=1}^ma_rG_{i,r}\paren{\frac{1}{\sqrt{m}}\sum_{r'=1}^ma_r'T_{i,r,r'} - y}
    \end{align*}
    where $C_{i,r,r'}, E_{i,r,r'}, T_{i,r,r'}$ and $G_{i,r}$ are defined as
    \begin{gather*}
        C_{i,r,r'} = \paren{\paren{\bfw_r^\top\bfx_i}\paren{\bfw_{r'}^\top\bfx_i} + \kappa^2\bfu_{i,r}^\top\bfu_{i,r'}}\mathcal{C}\paren{\frac{\bfw_r^\top\bfx_i}{\kappa\norm{\bfu_{i,r}}_2},\frac{\bfw_{r'}^\top\bfx_i}{\kappa\norm{\bfu_{i,r'}}_2},\frac{\bfu_{i,r}^\top\bfu_{i,r'}}{\norm{\bfu_{i,r}}_2\norm{\bfu_{i,r'}}_2}}\\
        E_{i,r,r'} = \norm{\bfu_{i,r}}_2\norm{\bfu_{i,r'}}_2\exp\paren{-\frac{\norm{\bfu_{i,r'}}_2^2\paren{\bfw_r^\top\bfx_i}^2-2\paren{\bfu_{i,r}^\top\bfu_{i,r'}}\paren{\bfw_r^\top\bfx_i}\paren{\bfw_{r'}^\top\bfx_i} + \norm{\bfu_{i,r}}_2^2\paren{\bfw_{r'}^\top\bfx_i}^2}{2\kappa^2\paren{\norm{\bfu_{i,r}}_2^2\norm{\bfu_{i,r'}}_2^2-\paren{\bfu_{i,r}^\top\bfu_{i,r'}}^2}}}\\
        T_{i,r,r'} = \bfw_{r'}^\top\bfx_i\cdot \bm{\Phi}_1\paren{\frac{\norm{\bfu_{i,r}}_2\norm{\bfu_{i,r'}}_2\cdot\bfw_{r'}^\top\bfx-\bfu_{i,r}^\top\bfu_{i,r'}\cdot\bfw_r^\top\bfx_i}{\kappa^2\norm{\bfu_{i,r}}_2\norm{\bfu_{i,r'}}_2\paren{\norm{\bfu_{i,r}}_2^2\norm{\bfu_{i,r'}}_2^2-\paren{\bfu_{i,r}^\top\bfu_{i,r'}}^2}}}\\
        G_{i,r} = \norm{\bfu_{i,r}}_2\exp\paren{-\frac{\paren{\bfw_r^\top\bfx_i}^2}{2\kappa^2\norm{\bfu_r}_2^2}}
    \end{gather*}
\end{lemma}
\begin{proof}
    By definition, we have
    \[
        \E_{\bfC}\left[\calL_{\bfC}\paren{\bm{\theta}}\right] = \frac{1}{2}\sum_{i=1}^n\E_{\bfc_i}\left[\paren{f\paren{\bm{\theta},\bfx_i\odot\bfc_i} - y_i}^2\right]
    \]
    For simplicity, we fix $i\in[n]$, and study $\E_{\bfc}\left[\paren{f\paren{\bm{\theta},\bfx\odot\bfc} - y}^2\right]$. In the analysis below, we let $\bfu_r = \bfw_r\odot\bfx$. In particular, we have
    \begin{align*}
        \E_{\bfc}\left[\paren{f\paren{\bm{\theta},\bfx\odot\bfc} - y}^2\right] & = \E_{\bfc}\left[f\paren{\bm{\theta},\bfx\odot\bfc}^2\right] - 2y\E_{\bfc}\left[f\paren{\bm{\theta},\bfx\odot\bfc}\right] + y^2
    \end{align*}
    Here we shall evaluate the two expectations separately. To start, for the first-order term, we have
    \begin{equation}
        \label{eq:lem_exp_loss_1}
        \E_{\bfc}\left[f\paren{\bm{\theta},\bfx\odot\bfc}\right] = \frac{1}{\sqrt{m}}\sum_{r=1}^ma_r\E\left[\sigma\paren{\bfw_r^\top\paren{\bfc\odot\bfx}}\right] = \frac{1}{\sqrt{m}}\sum_{r=1}^ma_r\E\left[\sigma\paren{\bfc^\top\paren{\bfw_r\odot\bfx}}\right]
    \end{equation}
    Notice that since $\bfc\sim\mathcal{N}\paren{\bm{1}, \kappa^2\bfI}$. By Lemma~\ref{lem:inner_gaussian_rv} we have that $\bfc^\top\paren{\bfw_r\odot\bfx}\sim\mathcal{N}\paren{\bfw_r^\top\bfx,\kappa^2\norm{\bfw_r\odot\bfx}_2^2}$, since $\bm{1}^\top\paren{\bfw_r\odot\bfx} = \bfw_r^\top\bfx$. Applying Lemma~\ref{lem:exp_relu} with $z=\bfc^\top\paren{\bfw_r\odot\bfx}$, mean $\bfw_r^\top\bfx$, and standard deviation $\kappa\norm{\bfw_r\odot\bfx}$, we have that
    \begin{equation}
        \label{eq:expect_act}
        \E\left[\sigma\paren{\bfc^\top\paren{\bfw_r\odot\bfx}}\right] = \frac{\kappa\norm{\bfw_r\odot\bfx}_2}{\sqrt{2
        \pi}}\exp\paren{-\frac{\paren{\bfw_r^\top\bfx}^2}{2\kappa^2\norm{\bfw_r\odot\bfx}_2^2}} + \bfw_r^\top\bfx\bm{\Phi}_1\paren{\frac{\bfw_r^\top\bfx}{\kappa\norm{\bfw_r\odot\bfx}_2}}
    \end{equation}
    Plugging in to the form of (\ref{eq:lem_exp_loss_1}) gives
    \begin{equation}
        \E_{\bfc}\left[f\paren{\bm{\theta},\bfx\odot\bfc}\right] = \frac{1}{\sqrt{m}}\sum_{r=1}^ma_r\act{\bfw_r}{\bfx} + \frac{\kappa}{\sqrt{2\pi m}}\sum_{r=1}^ma_r\norm{\bfu_r}_2\exp\paren{-\frac{\paren{\bfw_r^\top\bfx}^2}{2\kappa^2\norm{\bfu_r}_2^2}}
    \end{equation}
    Next, we focus on the second-order term. Since $\bfw_r^\top\paren{\bfc\odot\bfx} = \bfc^\top\paren{\bfw_r\odot\bfx}$, we have
    \[
        \E_{\bfc}\left[f\paren{\bm{\theta},\bfx\odot\bfc}^2\right] = \frac{1}{m}\sum_{r,r'=1}^ma_ra_{r'}\E\left[\sigma\paren{\bfc^\top\bfu_r}\sigma\paren{\bfc^\top\bfu_{r'}}\right]
    \]
    Let $z_1 = \bfc^\top\bfu_r$ and $z_2 = \bfc^\top\bfu_{r'}$, we have that $z_1 \sim\mathcal{N}\paren{\bfw_r^\top\bfx,\kappa^2\norm{\bfu_r}_2^2}, z_2\sim\mathcal{N}\paren{\bfw_{r'}^\top\bfx,\kappa^2\norm{\bfu_{r'}}_2^2}$, and $\text{Cov}\paren{z_1,z_2} = \kappa^2\bfu_r^\top\bfu_{r'}$. Applying Lemma~\ref{lem:exp_prod_relu} with $a = b = 0$ gives
    \begin{align*}
        & \E\left[\sigma\paren{\bfc^\top\bfu_r}\sigma\paren{\bfc^\top\bfu_{r'}}\right]\\
        &\quad\quad\quad = \paren{\paren{\bfw_r^\top\bfx}\paren{\bfw_{r'}^\top\bfx} + \kappa^2\bfu_r^\top\bfu_{r'}}\bm{\Phi}_2\paren{\frac{\bfw_r^\top\bfx}{\kappa\norm{\bfu_r}_2},\frac{\bfw_{r'}^\top\bfx}{\kappa\norm{\bfu_{r'}}_2},\frac{\bfu_r^\top\bfu_{r'}}{\norm{\bfu_r}_2\norm{\bfu_{r'}}_2}}\\
        &\quad\quad\quad\quad\quad  + \frac{\kappa^2}{2\pi}\norm{\bfu_r}_2\norm{\bfu_{r'}}_2\exp\paren{-\frac{\norm{\bfu_{r'}}_2^2\paren{\bfw_r^\top\bfx}^2-2\paren{\bfu_r^\top\bfu_{r'}}\paren{\bfw_r^\top\bfx}\paren{\bfw_{r'}^\top\bfx} + \norm{\bfu_r}_2^2\paren{\bfw_{r'}^\top\bfx}^2}{2\kappa^2\paren{\norm{\bfu_r}_2^2\norm{\bfu_{r'}}_2^2-\paren{\bfu_r^\top\bfu_{r'}}^2}}}\\
        & \quad\quad\quad\quad\quad  + \frac{\kappa}{\sqrt{2\pi}}\paren{\norm{\bfu_r}_2\cdot\bfw_{r'}^\top\bfx \cdot \hat{T}_{1,r,r'} + \norm{\bfu_{r'}}_2\cdot\bfw_{r}^\top\bfx \cdot \hat{T}_{2,r,r'}}
    \end{align*}
    where $\hat{T}_{1,r,r'}, \hat{T}_{2,r,r'}$ are defined as
    \begin{align*}
        \hat{T}_{1,r,r'} = \exp\paren{-\frac{\paren{\bfw_r^\top\bfx}^2}{2\kappa^2\norm{\bfu_r}_2^2}}\bm{\Phi}_1\paren{\frac{\norm{\bfu_r}_2^2\cdot\bfw_{r'}^\top\bfx-\bfu_r^\top\bfu_{r'}\cdot\bfw_r^\top\bfx}{\kappa\norm{\bfu_r}_2\paren{\norm{\bfu_r}_2^2\norm{\bfu_{r'}}_2^2-\paren{\bfu_r^\top\bfu_{r'}}^2}^{\frac{1}{2}}}}\\
        \hat{T}_{2,r,r'} = \exp\paren{-\frac{\paren{\bfw_{r'}^\top\bfx}^2}{2\kappa^2\norm{\bfu_{r'}}_2^2}}\bm{\Phi}_1\paren{\frac{\norm{\bfu_{r'}}_2^2\cdot\bfw_{r}^\top\bfx-\bfu_r^\top\bfu_{r'}\cdot\bfw_{r'}^\top\bfx}{\kappa\norm{\bfu_{r'}}_2\paren{\norm{\bfu_r}_2^2\norm{\bfu_{r'}}_2^2-\paren{\bfu_r^\top\bfu_{r'}}^2}^{\frac{1}{2}}}}
    \end{align*}
    For the simplicity of notations, we define $T_{1,r,r'} = \norm{\bfu_r}_2\cdot\bfw_{r'}^\top\bfx \cdot \hat{T}_{1,r,r'}$ and $T_{2,r,r'} = \norm{\bfu_{r'}}_2\cdot\bfw_{r}^\top\bfx \cdot \hat{T}_{2,r,r'}$. Moreover, we define
    \[
        E_{r,r'} = \norm{\bfu_r}_2\norm{\bfu_{r'}}_2\exp\paren{-\frac{\norm{\bfu_{r'}}_2^2\paren{\bfw_r^\top\bfx}^2-2\paren{\bfu_r^\top\bfu_{r'}}\paren{\bfw_r^\top\bfx}\paren{\bfw_{r'}^\top\bfx} + \norm{\bfu_r}_2^2\paren{\bfw_{r'}^\top\bfx}^2}{2\kappa^2\paren{\norm{\bfu_r}_2^2\norm{\bfu_{r'}}_2^2-\paren{\bfu_r^\top\bfu_{r'}}^2}}}
    \]
    Lastly, we use the definition of the Gaussian Copula function $\mathcal{C}\paren{a, b, \rho} = \bm{\Phi}_2\paren{a,b,\rho} - \bm{\Phi}_1\paren{a}\bm{\Phi}_1\paren{b}$ and define
    \[
        C_{r,r'} = \paren{\paren{\bfw_r^\top\bfx}\paren{\bfw_{r'}^\top\bfx} + \kappa^2\bfu_r^\top\bfu_{r'}}\mathcal{C}\paren{\frac{\bfw_r^\top\bfx}{\kappa\norm{\bfu_r}_2},\frac{\bfw_{r'}^\top\bfx}{\kappa\norm{\bfu_{r'}}_2},\frac{\bfu_r^\top\bfu_{r'}}{\norm{\bfu_r}_2\norm{\bfu_{r'}}_2}}
    \]
    Under these definitions, we have that
    \begin{align*}
        \E\left[\sigma\paren{\bfc^\top\bfu_r}\sigma\paren{\bfc^\top\bfu_{r'}}\right] & = \paren{\paren{\bfw_r^\top\bfx}\paren{\bfw_{r'}^\top\bfx} + \kappa^2\bfu_r^\top\bfu_{r'}}\bm{\Phi}_1\paren{\frac{\bfw_r^\top\bfx}{\kappa\norm{\bfu_r}_2}}\bm{\Phi}_1\paren{\frac{\bfw_{r'}^\top\bfx}{\kappa\norm{\bfu_{r'}}_2}}\\
        & \quad\quad\quad + C_{r,r'} + \frac{\kappa^2}{2\pi}E_{r,r'} + \frac{\kappa}{\sqrt{2\pi}}\paren{T_{1,r,r'} + T_{2,r,r'}}\\
        & = \act{\bfw_r}{\bfx}\act{\bfw_{r'}}{\bfx} + \kappa^2\bfu_r^\top\bfu_{r'}\bm{\Phi}_1\paren{\frac{\bfw_r^\top\bfx}{\kappa\norm{\bfu_r}_2}}\bm{\Phi}_1\paren{\frac{\bfw_{r'}^\top\bfx}{\kappa\norm{\bfu_{r'}}_2}}\\
        & \quad\quad\quad + C_{r,r'} + \frac{\kappa^2}{2\pi}E_{r,r'} + \frac{\kappa}{\sqrt{2\pi}}\paren{T_{1,r,r'} + T_{2,r,r'}}\\
    \end{align*}
    Plugging back into the expression of $\E_{\bfc}\left[f\paren{\bm{\theta},\bfx\odot\bfc}^2\right]$ gives
    \begin{align*}
        & \E_{\bfc}\left[f\paren{\bm{\theta},\bfx\odot\bfc}^2\right]\\
        &\quad\quad = \frac{1}{m}\sum_{r,r'=1}^ma_ra_{r'}\act{\bfw_r}{\bfx}\act{\bfw_{r'}}{\bfx} + \frac{1}{m}\sum_{r,r'=1}^ma_ra_{r'}\kappa^2\bfu_r^\top\bfu_{r'}\bm{\Phi}_1\paren{\frac{\bfw_r^\top\bfx}{\kappa\norm{\bfu_r}_2}}\bm{\Phi}_1\paren{\frac{\bfw_{r'}^\top\bfx}{\kappa\norm{\bfu_{r'}}_2}}\\
        & \quad\quad\quad\quad + \frac{1}{m}\sum_{r,r'=1}^ma_ra_{r'}\paren{C_{r,r'} + \frac{\kappa^2}{2\pi}E_{r,r'} + \frac{\kappa}{\sqrt{2\pi}}\paren{T_{1,r,r'} + T_{2,r,r'}}}\\
        &\quad\quad = \paren{\frac{1}{\sqrt{m}}\sum_{r=1}^ma_r\act{\bfw_r}{\bfx}}^2 + \kappa^2\norm{\frac{1}{\sqrt{m}}\sum_{r=1}^ma_r\bfu_r\bm{\Phi}\paren{\frac{\bfw_r^\top\bfx}{\kappa\norm{\bfu_r}_2}}}_2^2\\
        & \quad\quad\quad\quad + \frac{1}{m}\sum_{r,r'=1}^ma_ra_{r'}\paren{C_{r,r'} + \frac{\kappa^2}{2\pi}E_{r,r'} + \frac{\kappa}{\sqrt{2\pi}}\paren{T_{1,r,r'} + T_{2,r,r'}}}\\
    \end{align*}
    Combining the expression of $\E_{\bfc}\left[f\paren{\bm{\theta},\bfx\odot\bfc}^2\right]$ and $\E_{\bfc}\left[f\paren{\bm{\theta},\bfx\odot\bfc}\right]$, and noticing $T_{1,r,r'} = T_{2,r',r}$, we have
    \begin{align*}
        \E_{\bfc}\left[\paren{f\paren{\bm{\theta},\bfx\odot\bfc} - y}^2\right] & = \paren{\frac{1}{\sqrt{m}}\sum_{r=1}^ma_r\act{\bfw_r}{\bfx}}^2 + \kappa^2\norm{\frac{1}{\sqrt{m}}\sum_{r=1}^ma_r\bfu_r\bm{\Phi}\paren{\frac{\bfw_r^\top\bfx}{\kappa\norm{\bfu_r}_2}}}_2^2\\
        & \quad\quad\quad + \frac{1}{m}\sum_{r,r'=1}^ma_ra_{r'}\paren{C_{r,r'} + \frac{\kappa^2}{2\pi}E_{r,r'} + \frac{\kappa}{\sqrt{2\pi}}\paren{T_{1,r,r'} + T_{2,r,r'}}}\\
        &\quad\quad\quad - \frac{2y}{\sqrt{m}}\sum_{r=1}^ma_r\act{\bfw_r}{\bfx} - \frac{2\kappa y}{\sqrt{2\pi m}}\sum_{r=1}^ma_r\norm{\bfu_r}_2\exp\paren{-\frac{\paren{\bfw_r^\top\bfx}^2}{\norm{\bfu_r}_2^2}} + y^2\\
        & = \paren{\frac{1}{\sqrt{m}}\sum_{r=1}^ma_r\act{\bfw_r}{\bfx} - y}^2 +  \frac{\kappa^2}{m}\norm{\sum_{r=1}^ma_r\bfu_r\bm{\Phi}\paren{\frac{\bfw_r^\top\bfx}{\kappa\norm{\bfu_r}_2}}}_2^2\\
        & \quad\quad\quad + \frac{1}{m}\sum_{r,r'=1}^ma_ra_{r'}\paren{C_{r,r'} + \frac{\kappa^2}{2\pi}E_{r,r'} + \frac{2\kappa}{\sqrt{2\pi}}T_{1,r,r'}}\\
        &\quad\quad\quad - \frac{2\kappa y}{\sqrt{2\pi m}}\sum_{r=1}^ma_r\norm{\bfu_r}_2\exp\paren{-\frac{\paren{\bfw_r^\top\bfx}^2}{2\kappa^2\norm{\bfu_r}_2^2}}
    \end{align*}
    To extend to the case of $\bfx_i,y_i$, we need to re-define
    \begin{gather*}
        C_{i,r,r'} = \paren{\paren{\bfw_r^\top\bfx_i}\paren{\bfw_{r'}^\top\bfx_i} + \kappa^2\bfu_{i,r}^\top\bfu_{i,r'}}\mathcal{C}\paren{\frac{\bfw_r^\top\bfx_i}{\kappa\norm{\bfu_{i,r}}_2},\frac{\bfw_{r'}^\top\bfx_i}{\kappa\norm{\bfu_{i,r'}}_2},\frac{\bfu_{i,r}^\top\bfu_{i,r'}}{\norm{\bfu_{i,r}}_2\norm{\bfu_{i,r'}}_2}}\\
        E_{i,r,r'} = \norm{\bfu_{i,r}}_2\norm{\bfu_{i,r'}}_2\exp\paren{-\frac{\norm{\bfu_{i,r'}}_2^2\paren{\bfw_r^\top\bfx_i}^2-2\paren{\bfu_{i,r}^\top\bfu_{i,r'}}\paren{\bfw_r^\top\bfx_i}\paren{\bfw_{r'}^\top\bfx_i} + \norm{\bfu_{i,r}}_2^2\paren{\bfw_{r'}^\top\bfx_i}^2}{2\kappa^2\paren{\norm{\bfu_{i,r}}_2^2\norm{\bfu_{i,r'}}_2^2-\paren{\bfu_{i,r}^\top\bfu_{i,r'}}^2}}}\\
        T_{i,r,r'} = \bfw_{r'}^\top\bfx_i\cdot \bm{\Phi}_1\paren{\frac{\norm{\bfu_{i,r}}_2^2\cdot\bfw_{r'}^\top\bfx_i-\bfu_{i,r}^\top\bfu_{i,r'}\cdot\bfw_r^\top\bfx_i}{\kappa\norm{\bfu_{i,r}}_2\paren{\norm{\bfu_{i,r}}_2^2\norm{\bfu_{i,r'}}_2^2-\paren{\bfu_{i,r}^\top\bfu_{i,r'}}^2}^{\frac{1}{2}}}}\\
        G_{i,r} = \norm{\bfu_{i,r}}_2\exp\paren{-\frac{\paren{\bfw_r^\top\bfx_i}^2}{2\kappa^2\norm{\bfu_r}_2^2}}
    \end{gather*}
    Moreover, let $\hat{f}\paren{\bm{\theta},\bfx} = \frac{1}{\sqrt{m}}\sum_{r=1}^ma_r\act{\bfw_r}{\bfx}$. Then we have that
    \begin{align*}
        \E_{\bfC}\left[\calL_{\bfC}\paren{\bm{\theta}}\right] & = \frac{1}{2}\sum_{i=1}^n\paren{\hat{f}\paren{\bm{\theta},\bfx_i}-y_i}_2^2 + \frac{\kappa^2}{2m}\sum_{i=1}^n\norm{\sum_{r=1}^ma_r\bfu_{i,r}\bm{\Phi}\paren{\frac{\bfw_r^\top\bfx_i}{\kappa\norm{\bfu_{i,r}}_2}}}_2^2\\
        & \quad\quad\quad + \frac{1}{m}\sum_{i=1}^n\sum_{r,r'=1}^ma_ra_{r'}\paren{C_{i,r,r'} + \frac{\kappa^2}{2\pi}E_{i,r,r'} + \kappa\sqrt{\frac{2}{\pi}}T_{i,r,r'}G_{i,r}}\\
        & \quad\quad\quad - \kappa y\sqrt{\frac{2}{\pi m}}\sum_{i=1}^n\sum_{r=1}^ma_rG_{i,r}\\
        & = \frac{1}{2}\sum_{i=1}^n\paren{\hat{f}\paren{\bm{\theta},\bfx_i}-y_i}_2^2 + \frac{\kappa^2}{2m}\sum_{i=1}^n\norm{\sum_{r=1}^ma_r\bfu_{i,r}\bm{\Phi}\paren{\frac{\bfw_r^\top\bfx_i}{\kappa\norm{\bfu_{i,r}}_2}}}_2^2\\
        & \quad\quad\quad + \frac{1}{m}\sum_{i=1}^n\sum_{r,r'=1}^ma_ra_{r'}\paren{C_{i,r,r'} + \frac{\kappa^2}{2\pi}E_{i,r,r'}}\\
        & \quad\quad\quad + \frac{2\kappa}{\sqrt{2\pi m}}\sum_{i=1}^n\sum_{r=1}^ma_rG_{i,r}\paren{\frac{1}{\sqrt{m}}\sum_{r'=1}^ma_r'T_{i,r,r'} - y}
    \end{align*}
\end{proof}
%%%%%%%%%%%%%%%%%%
\subsection{Proof of Theorem~\ref{thm:exp_loss_approx}}
\label{sec:proof_exp_loss}
\begin{proof}
    Let $\bfu_{i,r} = \bfw_r\odot\bfx_i$. By Lemma~\ref{lem:exp_loss}, we have that
    \begin{align*}
        \E_{\bfC}\left[\calL_{\bfC}\paren{\bm{\theta}}\right] & = \frac{1}{2}\sum_{i=1}^n\paren{\hat{f}\paren{\bm{\theta},\bfx_i}-y_i}_2^2 + \frac{\kappa^2}{2m}\sum_{i=1}^n\norm{\sum_{r=1}^ma_r\bfu_{i,r}\bm{\Phi}\paren{\frac{\bfw_r^\top\bfx_i}{\kappa\norm{\bfu_{i,r}}_2}}}_2^2\\
        & \quad\quad\quad + \frac{1}{m}\sum_{i=1}^n\sum_{r,r'=1}^ma_ra_{r'}\paren{C_{i,r,r'} + \frac{\kappa^2}{2\pi}E_{i,r,r'}}\\
        & \quad\quad\quad + \frac{2\kappa}{\sqrt{2\pi m}}\sum_{i=1}^n\sum_{r=1}^ma_rG_{i,r}\paren{\frac{1}{\sqrt{m}}\sum_{r'=1}^ma_r'T_{i,r,r'} - y}
    \end{align*}
    where $C_{i,r,r'}, E_{i,r,r'}, T_{i,r,r'}$ and $G_{i,r}$ are defined as
    \begin{gather*}
        C_{i,r,r'} = \paren{\paren{\bfw_r^\top\bfx_i}\paren{\bfw_{r'}^\top\bfx_i} + \kappa^2\bfu_{i,r}^\top\bfu_{i,r'}}\mathcal{C}\paren{\frac{\bfw_r^\top\bfx_i}{\kappa\norm{\bfu_{i,r}}_2},\frac{\bfw_{r'}^\top\bfx_i}{\kappa\norm{\bfu_{i,r'}}_2},\frac{\bfu_{i,r}^\top\bfu_{i,r'}}{\norm{\bfu_{i,r}}_2\norm{\bfu_{i,r'}}_2}}\\
        E_{i,r,r'} = \norm{\bfu_{i,r}}_2\norm{\bfu_{i,r'}}_2\exp\paren{-\frac{\norm{\bfu_{i,r'}}_2^2\paren{\bfw_r^\top\bfx_i}^2-2\paren{\bfu_{i,r}^\top\bfu_{i,r'}}\paren{\bfw_r^\top\bfx_i}\paren{\bfw_{r'}^\top\bfx_i} + \norm{\bfu_{i,r}}_2^2\paren{\bfw_{r'}^\top\bfx_i}^2}{2\kappa^2\paren{\norm{\bfu_{i,r}}_2^2\norm{\bfu_{i,r'}}_2^2-\paren{\bfu_{i,r}^\top\bfu_{i,r'}}^2}}}\\
        T_{i,r,r'} = \bfw_{r'}^\top\bfx_i\cdot \bm{\Phi}_1\paren{\frac{\norm{\bfu_{i,r}}_2^2\cdot\bfw_{r'}^\top\bfx_i-\bfu_{i,r}^\top\bfu_{i,r'}\cdot\bfw_r^\top\bfx_i}{\kappa\norm{\bfu_{i,r}}_2\paren{\norm{\bfu_{i,r}}_2^2\norm{\bfu_{i,r'}}_2^2-\paren{\bfu_{i,r}^\top\bfu_{i,r'}}^2}^{\frac{1}{2}}}}\\
        G_{i,r} = \norm{\bfu_{i,r}}_2\exp\paren{-\frac{\paren{\bfw_r^\top\bfx_i}^2}{2\kappa^2\norm{\bfu_r}_2^2}}
    \end{gather*}
    Therefore, the proof of the theorem relies on the upper bound of $C_{i,r,r'}, E_{i,r,r'}, T_{i,r,r'}$ and $G_{i,r}$. To upper-bound $C_{i,r,r'}$, we utilize the result in \cite{} that
    \[
        \left|C\paren{a, b,\rho}\right| \leq \frac{\left|\arcsin \rho\right|}{2\pi}\exp\paren{-\frac{a^2-2\rho ab + b^2}{2\paren{1-\rho^2}}} \leq \frac{|\rho|}{4}\exp\paren{-\frac{a^2+b^2}{4}}
    \]
    where we used Lemma~\ref{lem:arcsin_bound} that $|\arcsin x|\leq \frac{\pi}{2}\cdot |x|$. Plugging in $a = \frac{\bfw_r^\top\bfx_i}{\kappa\norm{\bfu_{i,r}}_2}, b = \frac{\bfw_{r'}^\top\bfx_i}{\kappa\norm{\bfu_{i,r'}}_2}$ and $\rho = \frac{\bfu_{i,r}^\top\bfu_{i,r'}}{\norm{\bfu_{i,r}}_2\norm{\bfu_{i,r'}}_2}$ gives
    \begin{align*}
        \left|C_{i,r,r'}\right| & \leq \left|\paren{\bfw_r^\top\bfx_i}\paren{\bfw_{r'}^\top\bfx_i} + \kappa^2\bfu_{i,r}^\top\bfu_{i,r'}\right|\cdot\frac{\left|\bfu_{i,r}^\top\bfu_{i,r'}\right|}{4\norm{\bfu_{i,r}}_2\norm{\bfu_{i,r'}}_2}\exp\paren{-\frac{1}{4\kappa^2}\paren{\frac{\paren{\bfw_r^\top\bfx_i}^2}{\norm{\bfu_{i,r}}_2^2} + \frac{\paren{\bfw_{r'}^\top\bfx_i}^2}{\norm{\bfu_{i,r'}}_2^2}}}\\
        & \leq \frac{1}{4}\paren{\left|\paren{\bfw_r^\top\bfx_i}\paren{\bfw_{r'}^\top\bfx_i}\right| + \kappa^2\norm{\bfu_{i,r}}_2\norm{\bfu_{i,r'}}_2}\phi\paren{\frac{\bfw_r^\top\bfx_i}{2\kappa\norm{\bfu_{i,r}}}_2}\phi\paren{\frac{\bfw_{r'}^\top\bfx_i}{2\kappa\norm{\bfu_{i,r'}}_2}}\\
        & = \frac{\kappa^2}{4}\norm{\bfu_{i,r}}_2\norm{\bfu_{i,r'}}_2\paren{\frac{\left|\bfw_r^\top\bfx_i\right|}{\kappa\norm{\bfu_{i,r}}_2}\cdot\frac{\left|\bfw_{r'}^\top\bfx_i\right|}{\kappa\norm{\bfu_{i,r}}_2} + 1}\phi\paren{\frac{\bfw_r^\top\bfx_i}{2\kappa\norm{\bfu_{i,r}}_2}}\phi\paren{\frac{\bfw_{r'}^\top\bfx_i}{2\kappa\norm{\bfu_{i,r'}}_2}}\\
        & = \frac{\kappa^2}{4}\norm{\bfu_{i,r}}_2\norm{\bfu_{i,r'}}_2\paren{\psi\paren{\frac{\bfw_r^\top\bfx_i}{2\kappa\norm{\bfu_{i,r}}_2}}\psi\paren{\frac{\bfw_{r'}^\top\bfx_i}{2\kappa\norm{\bfu_{i,r'}}_2}} + \phi\paren{\frac{\bfw_r^\top\bfx_i}{2\kappa\norm{\bfu_{i,r}}_2}}\phi\paren{\frac{\bfw_{r'}^\top\bfx_i}{2\kappa\norm{\bfu_{i,r'}}_2}}}
    \end{align*}
    where we use the definition $P_{i,r} = \left|\bfw_r^\top\bfx_i\right|\cdot \exp\paren{-\frac{\paren{\bfw_r^\top\bfx_i}^2}{4\kappa^2\norm{\bfu_{i,r}}_2^2}}$.
    For the term $E_{i,r,r'}$, we notice that by letting $a = \frac{\bfw_r^\top\bfx_i}{\kappa\norm{\bfu_{i,r}}_2}, b = \frac{\bfw_{r'}^\top\bfx_i}{\kappa\norm{\bfu_{i,r'}}_2}$ and $\rho = \frac{\bfu_{i,r}^\top\bfu_{i,r'}}{\norm{\bfu_{i,r}}_2\norm{\bfu_{i,r'}}}$, we have
    \[
        \exp\paren{-\frac{\norm{\bfu_{i,r'}}_2^2\paren{\bfw_r^\top\bfx_i}^2-2\paren{\bfu_{i,r}^\top\bfu_{i,r'}}\paren{\bfw_r^\top\bfx_i}\paren{\bfw_{r'}^\top\bfx_i} + \norm{\bfu_{i,r}}_2^2\paren{\bfw_{r'}^\top\bfx_i}^2}{2\kappa^2\paren{\norm{\bfu_{i,r}}_2^2\norm{\bfu_{i,r'}}_2^2-\paren{\bfu_{i,r}^\top\bfu_{i,r'}}^2}}} = \exp\paren{-\frac{a^2- 2\rho ab + b^2}{2\paren{1-\rho^2}}}
    \]
    Using $\exp\paren{-\frac{a^2- 2\rho ab + b^2}{2\paren{1-\rho^2}}}\leq \exp\paren{-\frac{a^2+b^2}{4}}$, we have that
    \begin{align*}
        \left|E_{i,r,r'}\right| & \leq \norm{\bfu_{i,r}}_2\norm{\bfu_{i,r'}}_2\exp\paren{-\frac{1}{4\kappa^2}\paren{\frac{\paren{\bfw_r^\top\bfx_i}^2}{\norm{\bfu_{i,r}}_2^2} + \frac{\paren{\bfw_{r'}^\top\bfx_i}^2}{\norm{\bfu_{i,r'}}_2^2}}}\\
        & = \norm{\bfu_{i,r}}_2\norm{\bfu_{i,r'}}_2\phi\paren{\frac{\bfw_r^\top\bfx_i}{2\kappa\norm{\bfu_{i,r}}_2}}\phi\paren{\frac{\bfw_{r'}^\top\bfx_i}{2\kappa\norm{\bfu_{i,r'}}_2}}
    \end{align*}
    Therefore, we have
    \begin{align*}
        & \left|C_{i,r,r'} + \frac{\kappa^2}{2\pi}E_{i,r,r'}\right|\\
        &\quad\quad\quad \leq \frac{\kappa^2}{4}\norm{\bfu_{i,r}}_2\norm{\bfu_{i,r'}}_2\paren{\psi\paren{\frac{\bfw_r^\top\bfx_i}{2\kappa\norm{\bfu_{i,r}}_2}}\psi\paren{\frac{\bfw_{r'}^\top\bfx_i}{2\kappa\norm{\bfu_{i,r'}}_2}} + \phi\paren{\frac{\bfw_r^\top\bfx_i}{2\kappa\norm{\bfu_{i,r}}_2}}\phi\paren{\frac{\bfw_{r'}^\top\bfx_i}{2\kappa\norm{\bfu_{i,r'}}_2}}}
    \end{align*}
    This gives that
    \begin{equation}
        \label{eq:C_plus_E_bound}
        \begin{aligned}
            & \left|\sum_{r,r'=1}^ma_ra_{r'}\paren{C_{i,r,r'} + \frac{\kappa^2}{2\pi}E_{i,r,r'}}\right|\\
            &\quad\quad\quad  \leq \frac{\kappa^2}{4}\paren{\paren{\sum_{r=1}^m\norm{\bfu_{i,r}}_2\psi\paren{\frac{\bfw_r^\top\bfx_i}{2\kappa\norm{\bfu_{i,r}}_2}}}^2 + \paren{\sum_{r=1}^m\norm{\bfu_{i,r}}_2\phi\paren{\frac{\bfw_r^\top\bfx_i}{2\kappa\norm{\bfu_{i,r}}_2}}}^2}
        \end{aligned}
    \end{equation}
    By definition, we have $\norm{\bfu_{i,r}}_2\leq R_{\bfu}$. Therefore
    \[
        \left|\sum_{r,r'=1}^ma_ra_{r'}\paren{C_{i,r,r'} + \frac{\kappa^2}{2\pi}E_{i,r,r'}}\right| \leq \frac{1}{4}\kappa^2R_{\bfu}^2\paren{\paren{\sum_{r=1}^m\psi\paren{\frac{\bfw_r^\top\bfx_i}{2\kappa\norm{\bfu_{i,r}}_2}}}^2 + \paren{\sum_{r=1}^m\phi\paren{\frac{\bfw_r^\top\bfx_i}{2\kappa\norm{\bfu_{i,r}}_2}}}^2}
    \]
    Next, we focus on the term $T_{i,r,r'}$ and $G_{i,r}$. By the property of CDF, we have that $\left|T_{i,r,r'}\right| \leq \left|\bfw_{r'}^\top\bfx_i\right| \leq \norm{\bfw_r}_2$. Therefore
    \[
        \left|\frac{1}{\sqrt{m}}\sum_{r'=1}^ma_{r'}T_{i,r,r'} - y_i\right| \leq \frac{1}{\sqrt{m}}\sum_{r'=1}^m\norm{\bfw_{r'}}_2 + \left|y_i\right| \leq \sqrt{m}R_{\bfw} + B_y \leq 2\sqrt{m}R_{\bfw}
    \]
    where we applied $\norm{\bfw_r}_2\leq R_{\bfw}$ and $B_y\leq 3\sqrt{m} R_{\bfw}$. Thus
    \begin{align*}
        \left|\sum_{r=1}^ma_rG_{i,r}\paren{\frac{1}{\sqrt{m}}\sum_{r'=1}^ma_{r'}T_{i,r,r'} - y_i}\right| & \leq \sum_{r=1}^mG_{i,r}\cdot 2\sqrt{m}R_{\bfw} \leq \sqrt{m}R_{\bfw}\sum_{r=1}^m\phi\paren{\frac{\bfw_r^\top\bfx_i}{2\kappa\norm{\bfu_{i,r}}_2}}
    \end{align*}
    where we used $\norm{\bfu_{i,r}}_2\leq R_{\bfu}$.
    Combining the inequality above and (\ref{eq:C_plus_E_bound}), we have
    \begin{align*}
        \left|\mathcal{E}\right| & \leq \frac{n\kappa^2R_{\bfu}^2}{4m}\paren{\paren{\sum_{r=1}^m\psi\paren{\frac{\bfw_r^\top\bfx_i}{2\kappa\norm{\bfu_{i,r}}_2}}}^2 + \paren{\sum_{r=1}^m\phi\paren{\frac{\bfw_r^\top\bfx_i}{2\kappa\norm{\bfu_{i,r}}_2}}}^2} + \frac{n\kappa R_{\bfw}}{2}\sum_{r=1}^m\phi\paren{\frac{\bfw_r^\top\bfx_i}{2\kappa\norm{\bfu_{i,r}}_2}}
    \end{align*}
    Applying the definition of $\psi_{\max}$ and $\phi_{\max}$ gives the desired results.
\end{proof}
%%%%%
\subsection{Proof of Theorem~\ref{thm:exp_grad_form}}
\label{sec:proof_exp_grad}
\begin{proof}
    By the form of the gradient, we have
    \begin{equation}
        \label{eq:exp_grad}
        \begin{aligned}
            \E_{\bfC_k}\left[\nabla_{\bfw_r}\calL_{\bfC}\paren{\bm{\theta}}\right] & = \frac{a_r}{\sqrt{m}}\sum_{i=1}^n\E_{\bfC}\left[\paren{f\paren{\bm{\theta},\bfx_i\odot\bfc_i} - y_i}\bfx_i\odot\bfc_i\mathbb{I}\left\{\inner{\bfw_{r}}{\bfx_i\odot\bfc_i}\geq 0\right\}\right]\\
            & = \frac{a_r}{\sqrt{m}}\sum_{i=1}^n\underbrace{\E_{\bfc_{i}}\left[f\paren{\bm{\theta},\bfx_i\odot\bfc_i}\bfx_i\odot\bfc_i\mathbb{I}\left\{\inner{\bfw_{r}}{\bfx_i\odot\bfc_i}\geq 0\right\}\right]}_{\mathcal{T}_{1,i}}\\
            &\quad\quad\quad - \frac{a_r}{\sqrt{m}}\sum_{i=1}^ny_i\underbrace{\E_{\bfc_{i}}\left[\bfx_i\odot\bfc_i\mathbb{I}\left\{\inner{\bfw_{r}}{\bfx_i\odot\bfc_i}\geq 0\right\}\right]}_{\mathcal{T}_{2,i}}
        \end{aligned}
    \end{equation}
    Let $\bfu_{r,i} = \bfw_{r}\odot\bfx_i$. For $\mathcal{T}_{1,i}$, we further have
    \begin{align*}
        \mathcal{T}_{1,i} & = \frac{1}{\sqrt{m}}\sum_{r'=1}^ma_{r'}\E_{\bfc_{i}}\left[\sigma\paren{\bfw_{r'}^\top\paren{\bfx_i\odot\bfc_{i}}}\bfx_i\odot\bfc_{i}\mathbb{I}\left\{\bfw_{r}^\top\paren{\bfx_i\odot\bfc_{i}}\geq 0\right\}\right]\\
        & = \frac{1}{\sqrt{m}}\sum_{r'=1}^ma_{r'}\E_{\bfc_{i}}\left[\paren{\bfx_i\odot\bfc_{i}}\paren{\bfx_i\odot\bfc_{i}}^\top\bfw_{r'}\mathbb{I}\left\{\bfw_{r}^\top\paren{\bfx_i\odot\bfc_{i}}\geq 0;\bfw_{r'}^\top\paren{\bfx_i\odot\bfc_{i}}\geq 0\right\}\right]\\
        & = \frac{1}{\sqrt{m}}\sum_{r'=1}^ma_{r'}\paren{\E_{\bfc_{i}}\left[\bfc_{i}\bfc_{i}^\top\mathbb{I}\left\{\bfw_{r}^\top\paren{\bfx_i\odot\bfc_{i}}\geq 0;\bfw_{r'}^\top\paren{\bfx_i\odot\bfc_{i}}\geq 0\right\}\right]\odot\paren{\bfx_i\bfx_i^\top}}\bfw_{r'}\\
        & = \frac{1}{\sqrt{m}}\sum_{r'=1}^ma_{r'}\paren{\E_{\bfc_{i}}\left[\bfc_{i}\bfc_{i}^\top\mathbb{I}\left\{\bfu_{r,i}^\top\bfc_{i}\geq 0;\bfu_{r',i}^\top\bfc_{i}\geq 0\right\}\right]\odot\paren{\bfx_i\bfx_i^\top}}\bfw_{r'}\\
        & = \frac{1}{\sqrt{m}}\sum_{r=1}^ma_{r'}\text{Diag}\paren{\bfx}_i\E_{\bfc_{i}}\left[\bfc_{i}\bfc_{i}^\top\mathbb{I}\left\{\bfu_{r,i}^\top\bfc_{i}\geq 0;\bfu_{r',i}^\top\bfc_{i}\geq 0\right\}\right]\bfu_{r',i}
    \end{align*}
    For $\mathcal{T}_{2,i}$, we can easily obtain
    \[
        \mathcal{T}_{2,i} = \E_{\bfc_{i}}\left[\bfc_{i}\mathbb{I}\left\{\bfu_{r,i}^\top\bfc_{i}\geq 0\right\}\right]\odot\bfx_i
    \]
    Abstractly, we are thus interested in the following quantity:
    \[
        \E_{\bfc}\left[\bfc\bfc^\top\mathbb{I}\left\{\bfc^\top\bfu\geq 0;\bfc^\top\bfv\geq 0\right\}\right]; \quad\E_{\bfc}\left[\bfc\mathbb{I}\left\{\bfc^\top\bfu\geq 0\right\}\right]
    \]
    where $\bfc\sim\mathcal{N}\paren{\bm{\mu},\kappa^2\bfI}$, and $\bfu,\bfv$ are fixed vectors. Let $z_1 = \bfc^\top\bfu$ and $z_2 =\bfc^\top\bfv$. Then we have
    \[
        z_1\sim\mathcal{N}\paren{\bfc^\top\bfu,\kappa^2\norm{\bfu}_2^2};\quad z_2\sim\mathcal{N}\paren{\bfc^\top\bfv,\kappa^2\norm{\bfv}_2^2}
    \]
    According to Lemma~\ref{lem:exp_grad_2nd} and Lemma~\ref{lem:exp_grad_1st_prod}, and by defining $\bm{\mu} = \bm{1}$, we have that $\bm{\Delta}^{(1)}_{k,r,i}\in\R^{d}$ and $\bm{\Delta}^{(2)}_{k,r,r',i}\in\R^{d\times d}$ defined below
    \begin{align*}
        \bm{\Delta}^{(1)}_{k,r,i} & : = \E_{\bfc_{i}}\left[\bfc_{i}\mathbb{I}\left\{\bfc_{i}^\top\bfu_{r,i}\geq 0\right\}\right] - \bm{1}\cdot \bm{\Phi}_1\paren{\frac{\bfw_r^\top\bfx_i}{\kappa\norm{\bfw_r\odot\bfx_i}_2}}\\
        \bm{\Delta}_{k,r,r',i} & : = \E_{\bfc}\left[\bfc\bfc^\top\mathbb{I}\left\{\bfu_{r,i}^\top\bfc_{i}\geq 0;\bfu_{r',i}^\top\bfc_{i}\geq 0\right\}\right]\bfu_{r',i}\\
        & \quad\quad\quad - \paren{\bm{1}\bm{1}^\top\bfu_{r',i} + 3\kappa^2\bfu_{r',i}}\bm{\Phi}_1\paren{\frac{\bfw_r^\top\bfx_i}{\kappa\norm{\bfw_r\odot\bfx_i}_2}}\bm{\Phi}_1\paren{\frac{\bfw_{r'}^\top\bfx_i}{\kappa\norm{\bfw_{r'}\odot\bfx_i}_2}}
    \end{align*}
    satisfies
    \begin{align*}
        \norm{\bm{\Delta}^{(1)}_{r,i}}_{\infty} & \leq \kappa R_{\bfu}\phi_{\max}\\
        \norm{\bm{\Delta}^{(2)}_{r,r',i}}_{\infty} & \leq 4\kappa\norm{\bfv}_2\paren{\sqrt{d}\phi_{\max} + \psi_{\max}}
    \end{align*}
    Here we used $\norm{\bm{\mu}}_{\infty} = 1$ and $\bm{\mu}^\top\paren{\bfw_{r}\odot \bfx_i} = \bfw_{k,r}^\top\bfx_i$ when $\bm{\mu} = 1$. 
    Therefore, for $\mathcal{T}_{1,i}$, we have
    \begin{align*}
        \mathcal{T}_{1,i} & = \frac{1}{\sqrt{m}}\sum_{r'=1}^ma_{r'}\text{Diag}\paren{\bfx_i}\paren{\paren{\bfw_{r'}^\top\bfx_i\cdot \bm{1}}\bm{\Phi}_1\paren{\frac{\bfw_{r'}^\top\bfx_i}{\kappa\norm{\bfw_{r'}s\odot\bfx_i}_2}}\bm{\Phi}_1\paren{\frac{\bfw_{r'}^\top\bfx_i}{\kappa\norm{\bfw_{r'}\odot\bfx_i}_2}} + \bm{\Delta}_{r,r',i}^{(2)}}\\
        & \quad\quad\quad + \frac{3\kappa^2}{\sqrt{m}}\sum_{r'=1}^ma_{r'}\text{Diag}\paren{\bfx_i}\paren{\bfw_r\odot\bfx_i}\bm{\Phi}_1\paren{\frac{\bfw_r^\top\bfx_i}{\kappa\norm{\bfw_r\odot\bfx_i}_2}}\bm{\Phi}_1\paren{\frac{\bfw_{r'}^\top\bfx_i}{\kappa\norm{\bfw_{r'}\odot\bfx_i}_2}}\\
        & = \frac{1}{\sqrt{m}}\sum_{r'=1}^ma_{r'}\bfw_{r'}^\top\bfx_i\cdot\bfx_i\bm{\Phi}_1\paren{\frac{\bfw_{r'}^\top\bfx_i}{\kappa\norm{\bfw_{r'}s\odot\bfx_i}_2}}\bm{\Phi}_1\paren{\frac{\bfw_{r'}^\top\bfx_i}{\kappa\norm{\bfw_{r'}\odot\bfx_i}_2}} + \frac{1}{\sqrt{m}}\sum_{r'=1}^ma_{r'}\paren{\bfx_i\odot \bm{\Delta}_{r,r',i}^{(2)}}\\
        &\quad\quad\quad + \frac{3\kappa^2}{\sqrt{m}}\sum_{r'=1}^ma_{r'}\text{Diag}\paren{\bfx_i}^2\bfw_{r'}\bm{\Phi}_1\paren{\frac{\bfw_{r'}^\top\bfx_i}{\kappa\norm{\bfw_{r'}s\odot\bfx_i}_2}}\bm{\Phi}_1\paren{\frac{\bfw_{r'}^\top\bfx_i}{\kappa\norm{\bfw_{r'}\odot\bfx_i}_2}}\\
        & = f\paren{\bm{\theta},\bfx_i}\bfx_i\indy{\bfw_{r}^\top\bfx_i\geq 0} + \frac{3\kappa^2}{\sqrt{m}}\sum_{r'=1}^ma_{r'}\text{Diag}\paren{\bfx_i}^2\bfw_{r'}\mathbb{I}\left\{\bfw_{r}^\top\bfx_i\geq 0;\bfw_{r'}^\top\bfx_i\geq 0\right\}\\
        & \quad\quad\quad + \frac{1}{\sqrt{m}}\sum_{r'=1}^ma_{r'}\paren{\bfx_i\odot \bm{\Delta}_{r,r',i}^{(2)}} + \bfg_{1,i} + \bfg_{2,i}\\
    \end{align*}
    where 
    \begin{gather*}
        \bfg_{1,i} = \frac{1}{\sqrt{m}}\sum_{r'=1}^ma_{r'}\bfw_{r'}^\top\bfx_i\cdot \bfx_i\paren{\bm{\Phi}_1\paren{\frac{\bfw_{r'}^\top\bfx_i}{\kappa\norm{\bfw_{r'}s\odot\bfx_i}_2}}\bm{\Phi}_1\paren{\frac{\bfw_{r'}^\top\bfx_i}{\kappa\norm{\bfw_{r'}\odot\bfx_i}_2}} - \indy{\bfw_r^\top\bfx_i\geq 0;\bfw_{r'}^\top\bfx_i\geq 0}}\\
        \bfg_{2,i} = \frac{3\kappa^2}{\sqrt{m}}\sum_{r'=1}^ma_{r'}\text{Diag}\paren{\bfx_i}^2\bfw_{r'}\paren{\bm{\Phi}_1\paren{\frac{\bfw_{r'}^\top\bfx_i}{\kappa\norm{\bfw_{r'}s\odot\bfx_i}_2}}\bm{\Phi}_1\paren{\frac{\bfw_{r'}^\top\bfx_i}{\kappa\norm{\bfw_{r'}\odot\bfx_i}_2}} - \indy{\bfw_r^\top\bfx_i\geq 0;\bfw_{r'}^\top\bfx_i\geq 0}}
    \end{gather*}
    Likely, for $\mathcal{T}_{2,i}$ we have
    \begin{align*}
        \mathcal{T}_{2,i} & = \paren{\bm{1}\bm{\Phi}_1\paren{\frac{\bfw_r^\top\bfx_i}{\kappa\norm{\bfw_r\odot\bfx_i}_2}}+\bm{\Delta}^{(1)}_{r,i}}\odot\bfx_i\\
        & = \bfx_i\bm{\Phi}_1\paren{\frac{\bfw_r^\top\bfx_i}{\kappa\norm{\bfw_r\odot\bfx_i}_2}} + \bm{\Delta}^{(1)}_{r,i}\odot\bfx_i\\
        & = \bfx_i\indy{\bfw_r^\top\bfx_i\geq 0} + \bm{\Delta}^{(1)}_{r,i}\odot\bfx_i + \bfg_{3,i}
    \end{align*}
    where $\bfg_{3,i} = \bfx_i\cdot \paren{\bm{\Phi}_1\paren{\frac{\bfw_r^\top\bfx_i}{\kappa\norm{\bfw_r\odot\bfx_i}_2}} - \indy{\bfw_r^\top\bfx_i\geq 0}}$.
    Therefore, the final gradient is given by
    \begin{align*}
        \E_{\bfC}\left[\nabla_{\bfw_r}\calL_{\bfC}\paren{\bm{\theta}}\right] & = \frac{a_r}{\sqrt{m}}\sum_{i=1}^n\paren{f\paren{\bm{\theta},\bfx_i} -y_i}\bfx_i\mathbb{I}\left\{\bfw_{r}^\top\bfx_i\right\}\\
        &\quad\quad\quad + \frac{a_r}{\sqrt{m}}\sum_{i=1}^n\paren{\frac{1}{\sqrt{m}}\sum_{r'=1}^ma_{r'}\paren{\bfx_i\odot \bm{\Delta}_{r,r',i}^{(2)}} + y_i\bm{\Delta}^{(1)}_{r,i}\odot\bfx_i}\\
        &\quad\quad\quad + \frac{3\kappa^2}{\sqrt{m}}\sum_{r'=1}^ma_{r'}\text{Diag}\paren{\bfx_i}^2\bfw_{r'}\mathbb{I}\left\{\bfw_{r}^\top\bfx_i\geq 0;\bfw_{r'}^\top\bfx_i\geq 0\right\}\\
        & \quad\quad\quad + \frac{a_r}{\sqrt{m}}\sum_{i=1}^n\paren{\bfg_{1,i} + \bfg_{2,i} - y_i\cdot \bfg_{3,i}}\\
        & = \nabla_{\bfw_r}\calL\paren{\bm{\theta}} + \frac{3\kappa^2}{\sqrt{m}}\sum_{r'=1}^ma_{r'}\text{Diag}\paren{\bfx_i}^2\bfw_{r'}\mathbb{I}\left\{\bfw_{r}^\top\bfx_i\geq 0;\bfw_{r'}^\top\bfx_i\geq 0\right\}\\
        &\quad\quad\quad + \underbrace{\frac{a_r}{\sqrt{m}}\sum_{i=1}^n\paren{\frac{1}{\sqrt{m}}\sum_{r'=1}^ma_{r'}\paren{\bfx_i\odot \bm{\Delta}_{r,r',i}^{(2)}} + y_i\bm{\Delta}^{(1)}_{r,i}\odot\bfx_i}}_{\bfg_{4}}\\
        & \quad\quad\quad + \frac{a_r}{\sqrt{m}}\sum_{i=1}^n\paren{\bfg_{1,i} + \bfg_{2,i} - y_i\cdot \bfg_{3,i}}
    \end{align*}
    Notice that we can re-write $\bfg_{1,i}$ as
    \begin{align*}
        \bfg_{1,i} & = \bfx_i\paren{\bm{\Phi}_1\paren{\frac{\bfw_r^\top\bfx_i}{\kappa\norm{\bfw_r\odot\bfx_i}_2}} - \indy{\bfw_r^\top\bfx_i\geq 0}}\cdot \frac{1}{\sqrt{m}}\sum_{r'=1}^ma_{r'}\bfw_{r'}^\top\bfx_i\indy{\paren{\bfw_{r'}^\top\bfx_i\geq 0}}\\
        & \quad\quad\quad + \bfx_i\bm{\Phi}_1\paren{\frac{\bfw_r^\top\bfx_i}{\kappa\norm{\bfw_r\odot\bfx_i}_2}}\cdot \frac{1}{\sqrt{m}}\sum_{r'=1}^ma_{r'}\bfw_{r'}^\top\bfx_i\paren{\bm{\Phi}_1\paren{\frac{\bfw_{r'}^\top\bfx_i}{\kappa\norm{\bfw_{r'}\odot\bfx_i}_2}} - \indy{\bfw_{r'}^\top\bfx_i\geq 0}}\\
        & = \bfx_i\bm{\Phi}_1\paren{\frac{\bfw_r^\top\bfx_i}{\kappa\norm{\bfw_r\odot\bfx_i}_2}}\cdot \frac{1}{\sqrt{m}}\sum_{r'=1}^ma_{r'}\bfw_{r'}^\top\bfx_i\paren{\bm{\Phi}_1\paren{\frac{\bfw_{r'}^\top\bfx_i}{\kappa\norm{\bfw_{r'}\odot\bfx_i}_2}} - \indy{\bfw_{r'}^\top\bfx_i\geq 0}}\\
        & \quad\quad\quad + \bfx_i\paren{\bm{\Phi}_1\paren{\frac{\bfw_r^\top\bfx_i}{\kappa\norm{\bfw_r\odot\bfx_i}_2}} - \indy{\bfw_r^\top\bfx_i\geq 0}}\cdot f\paren{\bm{\theta},\bfx_i}
    \end{align*}
    Then, by the definition of $\bfg_{3,i}$, we have that
    \begin{align*}
        \bfg_{1,i} - y_i\cdot\bfg_{3,i} & = \bfx_i\bm{\Phi}_1\paren{\frac{\bfw_r^\top\bfx_i}{\kappa\norm{\bfw_r\odot\bfx_i}_2}}\cdot \frac{1}{\sqrt{m}}\sum_{r'=1}^ma_{r'}\bfw_{r'}^\top\bfx_i\paren{\bm{\Phi}_1\paren{\frac{\bfw_{r'}^\top\bfx_i}{\kappa\norm{\bfw_{r'}\odot\bfx_i}_2}} - \indy{\bfw_{r'}^\top\bfx_i\geq 0}}\\
        & \quad\quad\quad + \paren{f\paren{\bm{\theta},\bfx_i} - y_i}\bfx_i\paren{\bm{\Phi}_1\paren{\frac{\bfw_r^\top\bfx_i}{\kappa\norm{\bfw_r\odot\bfx_i}_2}} - \indy{\bfw_r^\top\bfx_i\geq 0}}
    \end{align*}
    Using Lemma~\ref{lem:gaus_cdf_approx}, we have that 
    \[
        \left|\bm{\Phi}_1\paren{a} - \indy{a\geq 0}\right|\leq \exp\paren{-\frac{a^2}{2}} \leq \phi\paren{\frac{a}{2}}
    \]
    Therefore, we have that 
    \begin{align*}
        \norm{\sum_{i=1}^n\paren{\bfg_{1,i} - y_i\cdot\bfg_{3,i}}}_2 & \leq \frac{n}{\sqrt{m}}\norm{ \sum_{r'=1}^ma_{r'}\bfw_{r'}^\top\bfx_i\paren{\bm{\Phi}_1\paren{\frac{\bfw_{r'}^\top\bfx_i}{\kappa\norm{\bfw_{r'}\odot\bfx_i}_2}} - \indy{\bfw_{r'}^\top\bfx_i\geq 0}}}_2 \\
        & \quad\quad\quad + \norm{\sum_{i=1}^n\paren{f\paren{\bm{\theta},\bfx_i} - y_i}\cdot\bfx_i\paren{\bm{\Phi}_1\paren{\frac{\bfw_r^\top\bfx_i}{\kappa\norm{\bfw_r\odot\bfx_i}_2}} - \indy{\bfw_r^\top\bfx_i\geq 0}}}_2\\
        & \leq \frac{n}{\sqrt{m}}\sum_{r'=1}^m\left|\bfw_{r'}^\top\bfx_i\right|\phi\paren{\frac{\bfw_{r'}^\top\bfx_i}{2\kappa\norm{\bfw_{r'}\odot\bfx_i}_2}} + \norm{\text{Diag}\paren{\bm{\Delta}}\bfX\paren{f\paren{\bm{\theta} - \bfy}}}_2\\
        & = \kappa\sqrt{m}R_{\bfu}\psi_{\max} + \sigma_{\max}\paren{\bfX}\phi_{\max}\calL\paren{\bm{\theta}}^{\frac{1}{2}}
    \end{align*}
    Moreover, we can bound $\bfg_{2,i}$ as
    \begin{align*}
        \norm{\bfg_{2,i}} & \leq \frac{3\kappa^2}{\sqrt{m}}\sum_{r'=1}^m\norm{\bfx}_{\infty}^2\norm{\bfw_r}\cdot 2\phi_{\max}\leq 6\kappa^2\sqrt{m}B_{\bfx}^2R_{\bfw}\phi_{\max}
    \end{align*}
    Lastly, we can bound $\bfg_{3}$ as
    \begin{align*}
        \norm{\bfg_{3}}_2 & \leq \frac{1}{m}\sum_{i=1}^n\sum_{r'=1}^m\norm{\bfx_i\odot \bm{\Delta}_{r,r',i}^{(2)}} + \frac{1}{\sqrt{m}}\sum_{i=1}^n\left|y_i\right|\norm{\bm{\Delta}^{(1)}_{r,i}\odot\bfx_i}_2\\
        & \leq \frac{1}{m}\sum_{i=1}^n\sum_{r'=1}^m\norm{\bm{\Delta}_{r,r',i}^{(2)}}_{\infty}\norm{\bfx_i}_2 + \frac{1}{\sqrt{m}}\sum_{i=1}^n\left|y_i\right|\norm{\bm{\Delta}_{r,r',i}^{(1)}}_{\infty}\norm{\bfx_i}_2\\
        & \leq \frac{1}{m}\sum_{i=1}^n\sum_{r'=1}^m\norm{\bm{\Delta}_{r,r',i}^{(2)}}_{\infty} + \frac{B_y}{\sqrt{m}}\sum_{i=1}^n\norm{\bm{\Delta}_{r,r',i}^{(1)}}_{\infty}\\
        & \leq 4n\kappa R_{\bfu}\paren{\sqrt{d}\phi_{\max} + \psi_{\max}} + \frac{B_y}{\sqrt{m}}\cdot n\kappa R_{\bfu}\phi_{\max}\\
        & \leq 5n\kappa R_{\bfu}\paren{\sqrt{d}\phi_{\max} + \psi_{\max}}
    \end{align*}
    when $B_y \leq \sqrt{md}$. Therefore, we have that
    \begin{align*}
        & \norm{\E_{\bfC}\left[\nabla_{\bfw_r}\calL_{\bfC}\paren{\bm{\theta}}\right] - \paren{\nabla_{\bfw_r}\calL\paren{\bm{\theta}} + \frac{3\kappa^2}{\sqrt{m}}\sum_{r'=1}^ma_{r'}\text{Diag}\paren{\bfx_i}^2\bfw_{r'}\mathbb{I}\left\{\bfw_{r}^\top\bfx_i\geq 0;\bfw_{r'}^\top\bfx_i\geq 0\right\}}}_2 \\
        & \quad\quad\quad \leq n\kappa R_{\bfu}\psi_{\max} + \frac{\sigma_{\max}\paren{\bfX}\phi_{\max}}{\sqrt{m}}\calL\paren{\bm{\theta}}^{\frac{1}{2}} + 6n\kappa^2B_{\bfx}^2R_{\bfw}\phi_{\max} + 5n\kappa R_{\bfu}\paren{\sqrt{d}\phi_{\max} + \psi_{\max}}\\
        & \quad\quad\quad \leq \paren{\frac{\sigma_{\max}\paren{\bfX}}{\sqrt{m}}\calL\paren{\bm{\theta}}^{\frac{1}{2}} + 6n\kappa^2B_{\bfx}^2R_{\bfw} + 5n\kappa R_{\bfu}\sqrt{d}}\phi_{\max} + 6n\kappa R_{\bfu}\psi_{\max}
    \end{align*}
\end{proof}
%%%%%%
\section{Proofs in Section~\ref{sec:conv_guarantee}}
\subsection{Proof of Theorem~\ref{thm:gen_sto_conv}}
\begin{proof}
    To start the proof, we define the following quantity in the standard NTK-based analysis of two-layer ReLU neural network. Let $R = C_1\cdot \frac{\tau\lambda_0}{n}$ for some $C_1 > 0$, we define event $A_{i,r}$ and set $S_i, S_i^\perp$ as
    \begin{gather}
        A_{i,r} = \left\{\exists \bfw\in\mathcal{B}\paren{\bfw_{0,r},R}: \indy{\bfw_{0,r}^\top\bfx_i \geq 0}\neq \indy{\bfw^\top\bfx_i\geq 0}\right\}\\
        S_i = \left\{r\in [m]: \neg A_{i,r}\right\};\quad  S_i^\perp = [m]\setminus S_i
    \end{gather}
    Lemma 16 from \cite{} shows that with probability at least $1 - n\exp\paren{-\frac{mR}{\tau}}$, we have that $\left|S_i^\perp\right|\leq \frac{4mR}{\tau}$. In the following of the proof, we assume that such event holds. 
    Define $K' = \min\left\{k\in\mathbb{N}: \exists r\in[m] \text{ s.t. } \norm{\bfw_{k,r} - \bfw_{0,r}}_2 > R\right\}$. Then for all $k < K'$, we have that $\bfw_{k,r}\in\mathcal{B}\paren{\bfw_{0,r}, R}$. Fix any $k < K' - 1$. Consider the expansion of $\calL\paren{\bm{\theta}_{k+1}}$ as the following
    \begin{equation}
        \label{eq:thm_conv_proof_1}
        \begin{aligned}
            \calL\paren{\bm{\theta}_{k+1}} & = \frac{1}{2}\sum_{i=1}^n\paren{f\paren{\bm{\theta}_{k+1},\bfx_i} - y_i}^2\\
            & = \frac{1}{2}\sum_{i=1}^n\paren{\paren{f\paren{\bm{\theta}_{k+1},\bfx_i} - f\paren{\bm{\theta}_{k},\bfx_i}} + \paren{f\paren{\bm{\theta}_k,\bfx_i} - y_i}}^2\\
            & = \frac{1}{2}\sum_{i=1}^n\paren{f\paren{\bm{\theta}_{k+1},\bfx_i} - f\paren{\bm{\theta}_{k},\bfx_i}}^2 + \sum_{i=1}^n\paren{f\paren{\bm{\theta}_{k+1},\bfx_i} - f\paren{\bm{\theta}_{k},\bfx_i}}\paren{f\paren{\bm{\theta}_k,\bfx_i} - y_i}\\
            &\quad\quad\quad + \frac{1}{2}\sum_{i=1}^n\paren{f\paren{\bm{\theta}_k,\bfx_i} - y_i}^2
        \end{aligned}
    \end{equation}
    We will analyze the three terms separately. To start, notice that
    \begin{equation}
        \label{eq:thm_conv_proof_2}
        \frac{1}{2}\sum_{i=1}^n\paren{f\paren{\bm{\theta}_k,\bfx_i} - y_i}^2 = \calL\paren{\bm{\theta}_k}
    \end{equation}
    For the first term, by the definition of $f\paren{\bm{\theta},\bfx}$, we have
    \begin{align*}
        \left|f\paren{\bm{\theta}_{k+1},\bfx_i} - f\paren{\bm{\theta}_{k},\bfx_i}\right| & = \left|\frac{1}{\sqrt{m}}\sum_{r=1}^ma_r\paren{\sigma\paren{\bfw_{k+1,r}^\top\bfx_i} - \sigma\paren{\bfw_{k,r}^\top\bfx_i}}\right|\\
        & \leq \frac{1}{\sqrt{m}}\sum_{r=1}^m\left|\sigma\paren{\bfw_{k+1,r}^\top\bfx_i} - \sigma\paren{\bfw_{k,r}^\top\bfx_i}\right|\\
        & \leq \frac{1}{\sqrt{m}}\sum_{r=1}^m\left|\paren{\bfw_{k+1}-\bfw_k}^\top\bfx_i\right|\\
        & \leq \frac{1}{\sqrt{m}}\sum_{r=1}^m\norm{\bfw_{k+1}-\bfw_k}\\
        & = \frac{\eta}{\sqrt{m}}\sum_{r=1}^m\norm{\nabla_{\bfw_r}\hat{\calL}\paren{\bm{\theta}_k,\bm{\xi}_k}}_2
    \end{align*}
    where in the first inequality we use the fact that $a = \pm 1$, and in the second inequality we use the $1$-Lipschitzness of ReLU. Applying Assumption~\ref{asump:randomness}, we have that
    \begin{equation}
        \label{eq:thm_conv_proof_3}
        \begin{aligned}
            \sum_{i=1}^n\paren{f\paren{\bm{\theta}_{k+1},\bfx_i} - f\paren{\bm{\theta}_{k},\bfx_i}}^2 & \leq \frac{\eta^2}{m}\sum_{i=1}^n\paren{\sum_{r=1}^m\norm{\nabla_{\bfw_r}\hat{\calL}\paren{\bm{\theta}_k,\bm{\xi}_k}}_2}\\
            & \leq \frac{\eta^2n}{m}\cdot \paren{m\cdot \sqrt{\gamma\hat{\calL}\paren{\bm{\theta}_k,\bm{\xi}_k}}}^2\\
            & = \eta^2 mn\gamma\hat{\calL}\paren{\bm{\theta}_k,\bm{\xi}_k}
        \end{aligned}
    \end{equation}
    Lastly, to analyze the second term, we use the following definition of $I_{i,k}$ and $I_{i,k}^\perp$
    \begin{align*}
        I_{i,k} = \frac{1}{\sqrt{m}}\sum_{r\in S_i}a_r\sigma\paren{\bfw_{k,r}^\top\bfx_i};\quad I_{i,k}^\perp = \frac{1}{\sqrt{m}}\sum_{r\in S_i^\perp}a_r\sigma\paren{\bfw_{k,r}^\top\bfx_i}
    \end{align*}
    Then we have that $f\paren{\bm{\theta}_k,\bfx_i} = I_{i,k} + I_{i,k}^\perp$. Therefore
    \[
        f\paren{\bm{\theta}_{k+1},\bfx_i} - f\paren{\bm{\theta}_k,\bfx_i} = \paren{I_{i,k+1} - I_{i,k}} + \paren{I_{i,k+1}^\perp - I_{i,k}^\perp}
    \]
    By the $1$-Lipschitzness of ReLU, we have that
    \begin{align*}
        \left|I_{i,k+1}^\perp - I_{i,k}^\perp\right| & = \left|\frac{1}{\sqrt{m}}\sum_{r\in S_i^\perp}a_r\paren{\sigma\paren{\bfw_{k+1,r}^\top\bfx_i} - \sigma\paren{\bfw_{k,r}^\top\bfx_i}}\right|\\
        & \leq \frac{1}{\sqrt{m}}\sum_{r\in S_i^\perp}\left|\sigma\paren{\bfw_{k+1,r}^\top\bfx_i} - \sigma\paren{\bfw_{k,r}^\top\bfx_i}\right|\\
        & \leq \frac{1}{\sqrt{m}}\sum_{r\in S_i^\perp}\left|\paren{\bfw_{k+1,r}-\bfw_{k,r}}^\top\bfx_i\right|\\
        & \leq \frac{\eta}{\sqrt{m}}\sum_{r\in S_i^\perp}\norm{\nabla_{\bfw_r}\hat{\calL}\paren{\bm{\theta}_k,\bm{\xi}_k}}_2\\
        & \leq \frac{\eta\sqrt{\gamma}}{\sqrt{m}}\left|S_i^\perp\right|\hat{\calL}\paren{\bm{\theta}_k,\bm{\xi}_k}^{\frac{1}{2}}
    \end{align*}
    Applying $\left|S_i^\perp\right|\leq \frac{4mR}{\tau}$ gives $\left|I_{i,k+1}^\perp - I_{i,k}^\perp\right| \leq \frac{4\eta R}{\tau}\sqrt{\gamma m}\hat{\calL}\paren{\bm{\theta}_k,\bm{\xi}_k}^{\frac{1}{2}}$.This gives that
    \begin{equation}
        \label{eq:thm_conv_proof_4}
        \begin{aligned}
            & \sum_{i=1}^n\paren{f\paren{\bm{\theta}_{k+1},\bfx_i} - f\paren{\bm{\theta}_k,\bfx_i}}\paren{f\paren{\bm{\theta}_k,\bfx_i} - y_i}\\
            & \quad\quad\quad = \sum_{i=1}^n\paren{I_{i,k+1} - I_{i,k}}\paren{f\paren{\bm{\theta}_k,\bfx_i} - y_i} + \sum_{i=1}^n\paren{I_{i,k+1}^\perp - I_{i,k}^\perp}\paren{f\paren{\bm{\theta}_k,\bfx_i} - y_i}\\
            & \quad\quad\quad\leq \sum_{i=1}^n\paren{I_{i,k+1} - I_{i,k}}\paren{f\paren{\bm{\theta}_k,\bfx_i} - y_i} + \paren{\sum_{i=1}^n\paren{I_{i,k+1}^\perp - I_{i,k}^\perp}^2}^{\frac{1}{2}}\paren{\sum_{i=1}^n\paren{f\paren{\bm{\theta}_k,\bfx_i} - y_i}^2}^{\frac{1}{2}}\\
            & \quad\quad\quad \leq \sum_{i=1}^n\paren{I_{i,k+1} - I_{i,k}}\paren{f\paren{\bm{\theta}_k,\bfx_i} - y_i} + \frac{4\eta R}{\tau}\sqrt{\gamma mn}\hat{\calL}\paren{\bm{\theta}_k,\bm{\xi}_k}^{\frac{1}{2}}\calL\paren{\bm{\theta}_k}^{\frac{1}{2}}
        \end{aligned}
    \end{equation}
    Plugging (\ref{eq:thm_conv_proof_2}), (\ref{eq:thm_conv_proof_3}), and (\ref{eq:thm_conv_proof_4}) into (\ref{eq:thm_conv_proof_1}) gives
    \begin{align*}
        \calL\paren{\bm{\theta}_{k+1}} & \leq \calL\paren{\bm{\theta}_k} + \eta^2mn\gamma\hat{\calL}\paren{\bm{\theta}_k,\bm{\xi}_k} + \frac{4\eta R}{\tau}\sqrt{\gamma mn}\hat{\calL}\paren{\bm{\theta}_k,\bm{\xi}_k}^{\frac{1}{2}}\calL\paren{\bm{\theta}_k}^{\frac{1}{2}}\\
        & \quad\quad\quad +\sum_{i=1}^n\paren{I_{i,k+1} - I_{i,k}}\paren{f\paren{\bm{\theta}_k,\bfx_i} - y_i}
    \end{align*}
    Under Jensen's inequality, we have that $\E_{\bm{\xi}_k}\left[\hat{\calL}\paren{\bm{\theta}_k,\bm{\xi}_k}^{\frac{1}{2}}\right] \leq \E_{\bm{\xi}_k}\left[\hat{\calL}\paren{\bm{\theta}_k,\bm{\xi}_k}\right]^{\frac{1}{2}}$. Using the property that
    \[
        \E_{\bm{\xi}_k}\left[\hat{\calL}\paren{\bm{\theta}_k,\bm{\xi}_k}\right] \leq 2\calL\paren{\bm{\theta}_k} + \varepsilon_1
    \]
    from Assumption~\ref{asump:randomness}, we can also obtain that
    \[
        \E_{\bm{\xi}_k}\left[\hat{\calL}\paren{\bm{\theta}_k,\bm{\xi}_k}^{\frac{1}{2}}\right] \leq \paren{2\calL\paren{\bm{\theta}_k} + \varepsilon_1}^{\frac{1}{2}}
    \]
    Therefore, taking the expectation of $\calL\paren{\bm{\theta}_{k+1}}$ gives
    \begin{equation}
        \label{eq:thm_conv_proof_5}
        \begin{aligned}
            \E_{\bm{\xi}_k}\left[\calL\paren{\bm{\theta}_{k+1}}\right] & \leq \calL\paren{\bm{\theta}_k} + \eta^2mn\gamma\paren{2\calL\paren{\bm{\theta}_k} + \varepsilon} + \frac{4\eta R}{\tau}\sqrt{\gamma mn}\paren{2\calL\paren{\bm{\theta}_k} + \varepsilon_1}^{\frac{1}{2}}\calL\paren{\bm{\theta}_k}^{\frac{1}{2}}\\
            &\quad\quad\quad + \sum_{i=1}^n\E_{\bm{\xi}_k}\left[I_{i,k+1}-I_{i,k}\right]\paren{f\paren{\bm{\theta}_k,\bfx_i} - y_i}\\
            & \leq \calL\paren{\bm{\theta}_k} + \eta^2mn\gamma\paren{2\calL\paren{\bm{\theta}_k} + \varepsilon_1} + \frac{10\eta R}{\tau}\sqrt{\gamma mn}\calL\paren{\bm{\theta}_k} + \frac{4\eta R}{\tau}\sqrt{\gamma mn}\cdot \varepsilon_1\\
            &\quad\quad\quad + \sum_{i=1}^n\E_{\bm{\xi}_k}\left[I_{i,k+1}-I_{i,k}\right]\paren{f\paren{\bm{\theta}_k,\bfx_i} - y_i}\\
            & = \paren{1 + 2\eta^2mn\gamma + 10C\eta\lambda_0\sqrt{\frac{\gamma m}{n}}}\calL\paren{\bm{\theta}_k} + \paren{\eta^2mn\gamma + 4C\eta\lambda_0\sqrt{\frac{\gamma m}{n}}}\varepsilon_1\\
             &\quad\quad\quad + \sum_{i=1}^n\E_{\bm{\xi}_k}\left[I_{i,k+1}-I_{i,k}\right]\paren{f\paren{\bm{\theta}_k,\bfx_i} - y_i}
        \end{aligned}
    \end{equation}
    where in the last inequality we use the property that $\sqrt{a(a+b)}\leq \frac{5}{4}a + b$.
    Recall that $\bfw_{k+1,r},\bfw_{k,r}\in\mathcal{B}\paren{\bfw_{0,r}, R}$. Therefore, for $r\in S_i$, we must have that $\indy{\bfw_{k+1,r}^\top\bfx_i\geq 0} = \indy{\bfw_{0,r}^\top\bfx_i\geq 0} = \indy{\bfw_{k,r}^\top\bfx_i\geq 0}$. Thus, we have
    \begin{equation}
        \label{eq:thm_conv_proof_5}
        \begin{aligned}
            \E_{\bm{\xi}_k}\left[I_{i,k+1} - I_{i,k}\right] & = \frac{1}{\sqrt{m}}\sum_{r\in S_i}a_r\E_{\bm{\xi}_k}\left[\bfw_{k+1,r}-\bfw_{k,r}\right]^\top\bfx_i\indy{\bfw_{k,r}^\top\bfx_i\geq 0}\\
            & = -\frac{\eta}{\sqrt{m}}\sum_{r\in S_i}a_r\E_{\bm{\xi}_k}\left[\nabla_{\bfw_r}\hat{\calL}\paren{\bm{\theta}_k,\bm{\xi}_k}\right]^\top\bfx_i\indy{\bfw_{k,r}^\top\bfx_i\geq 0}
        \end{aligned}
    \end{equation}
    Let $\bfg_{k,r} = \E_{\bm{\xi}_k}\left[\nabla_{\bfw_r}\hat{\calL}\paren{\bm{\theta}_k,\bm{\xi}_k}\right] - \nabla_{\bfw_r}\calL\paren{\bm{\theta}_k}$. Then by Assumption~\ref{asump:randomness} we have that $\norm{\bfg_{k,r}}_2\leq \varepsilon_3\calL\paren{\bm{\theta}}^{\frac{1}{2}} + \varepsilon_2$. Using $\bfg_{k,r}$, we can write (\ref{eq:thm_conv_proof_5}) as
    \begin{equation}
        \label{eq:thm_conv_proof_6}
        \E_{\bm{\xi}_k}\left[I_{i,k+1} - I_{i,k}\right] =  -\eta\sum_{r\in S_i}\frac{a_r}{\sqrt{m}}\nabla_{\bfw_r}\calL\paren{\bm{\theta}_k}^\top\bfx_i\indy{\bfw_{k,r}^\top\bfx_i\geq 0}  -\frac{\eta}{\sqrt{m}}\sum_{r\in S_i}a_r\bfg_{k,r}^\top\bfx_i\indy{\bfw_{k,r}^\top\bfx_i\geq 0}
    \end{equation}
    By definition, we have
    \[
        \nabla_{\bfw_r}\calL\paren{\bm{\theta}_k} = \frac{a_r}{\sqrt{m}}\sum_{j=1}^n\paren{f\paren{\bm{\theta}_k,\bfx_i}-y_i}\bfx_j\indy{\bfw_{k,r}^\top\bfx_j\geq 0}
    \]
    Therefore, we have that
    \[
        \frac{a_r}{\sqrt{m}}\nabla_{\bfw_r}\calL\paren{\bm{\theta}_k}^\top\bfx_i\indy{\bfw_{k,r}^\top\bfx_i\geq 0} = \frac{1}{m}\sum_{j=1}^n\paren{f\paren{\bm{\theta}_k,\bfx_i}-y_i}\bfx_i^\top\bfx_j\indy{\bfw_{k,r}^\top\bfx_i\geq 0\bfw_{k,r}^\top\bfx_j\geq 0}
    \]
    Combining with (\ref{eq:thm_conv_proof_6}), we have that
    \begin{align*}
        & \sum_{i=1}^n\E_{\bm{\xi}_k}\left[I_{i,k+1} - I_{i,k}\right]\paren{f\paren{\bm{\theta}_k,\bfx_i}-y_i}\\
        &\quad\quad\quad  = -\frac{\eta}{m}\sum_{i,j=1}^n\sum_{r\in S_i}\paren{f\paren{\bm{\theta}_k,\bfx_i}-y_i}\paren{f\paren{\bm{\theta}_k,\bfx_j}-y_j}\bfx_i^\top\bfx_j\indy{\bfw_{k,r}^\top\bfx_i\geq 0;\bfw_{k,r}^\top\bfx_j\geq 0}\\
        &\quad\quad\quad\quad\quad - \frac{\eta}{\sqrt{m}}\sum_{i=1}^n\sum_{r\in S_i}\paren{f\paren{\bm{\theta}_k,\bfx_i}-y_i}a_r\bfg_{k,r}^\top\bfx_i\indy{\bfw_{k,r}^\top\bfx_i\geq0}\\
        &\quad\quad\quad  \leq -\eta\sum_{i,j=1}^n\paren{f\paren{\bm{\theta}_k,\bfx_i}-y_i}\underbrace{\paren{\frac{\bfx_i^\top\bfx_j}{m}\sum_{r\in S_i}\indy{\bfw_{k,r}^\top\bfx_i\geq 0;\bfw_{k,r}^\top\bfx_j\geq 0}}}_{\bfH_{k,ij}}\paren{f\paren{\bm{\theta}_k,\bfx_j}-y_j}\\
        &\quad\quad\quad\quad\quad + \frac{\eta}{\sqrt{m}}\sum_{i=1}^m\sum_{r=1}^n\left|f\paren{\bm{\theta}_k,\bfx_i}-y_i\right|\norm{\bfg_{k,r}}_2\\
        & \quad\quad\quad  \leq-\eta\lambda_{\min}\paren{\bfH_k}\sum_{i=1}^n\paren{f\paren{\bm{\theta}_k,\bfx_i}-y_i}^2 + \eta\varepsilon_2\sqrt{mn}\paren{\sum_{r=1}^n\paren{f\paren{\bm{\theta}_k,\bfx_i}-y_i}^2}^{\frac{1}{2}} + \eta\varepsilon_3\sqrt{mn}\calL\paren{\bm{\theta}_k}\\
        & = -\paren{2\eta\lambda_{\min}\paren{\bfH_k} + \eta\varepsilon_3\sqrt{mn}}\calL\paren{\bm{\theta}_k} + 2\eta\varepsilon_2\sqrt{mn}\calL\paren{\bm{\theta}_k}^{\frac{1}{2}}
    \end{align*}
    Using the property that $ab \leq \frac{a^2}{2} + \frac{b^2}{2}$, we have that for any $C' > 0$,
    \[
        2\eta\varepsilon_2\sqrt{mn}\calL\paren{\bm{\theta}_k}^{\frac{1}{2}} \leq C'\eta\lambda_0\sqrt{\frac{\gamma m}{n}} + \frac{\eta\varepsilon_2^2 n}{C'\lambda_0}\sqrt{\frac{mn}{\gamma}}
    \]
    Moreover, by Lemma~\ref{lem:kernel_diff}, we have that when $m = \Omega\paren{\frac{n^2}{\lambda_0^2}\log\frac{n}{\delta}}$, with probability at least $1 - \delta - n^2\exp{-\frac{mR}{\tau}}$, it holds that 
    \[
        \norm{\bfH_k - \bfH^\infty}_F \leq \frac{\lambda_0}{6} + \frac{1}{m}\paren{\sum_{i,j=1}^n\left|S_i^\perp\right|^2}^{\frac{1}{2}} + \frac{2nR}{\tau}
    \]
    Plugging in $\left|S_i^\perp\right|\leq \frac{4mR}{\tau}$ and $R\leq C_1\cdot\frac{\tau\lambda_0}{n}$, we have that $\norm{\bfH_k - \bfH^\infty}_F \leq \frac{\lambda_0}{2}$ for small enough $C_1$. Thus, we have that $\lambda_{\min}\paren{\bfH_k}\geq \frac{\lambda_0}{2}$. 
    Therefore, we have that
    \[
        \sum_{i=1}^n\E_{\bm{\xi}_k}\left[I_{i,k+1} - I_{i,k}\right]\paren{f\paren{\bm{\theta}_k,\bfx_i}-y_i} \leq \paren{C'\eta\lambda_0\sqrt{\frac{\gamma m}{n}} + \eta\varepsilon_3\sqrt{mn} - \eta\lambda_0}\calL\paren{\bm{\theta}_k} + \frac{\eta\varepsilon_2^2 n}{C'\lambda_0}\sqrt{\frac{mn}{\gamma}}
    \]
    Plugging this back into (\ref{eq:thm_conv_proof_6}) gives
    \begin{align*}
        \E_{\bm{\xi}_k}\left[\calL\paren{\bm{\theta}_{k+1}}\right] & \leq \paren{1 + 2\eta^2mn\gamma + 10C\eta\lambda_0\sqrt{\frac{\gamma m}{n}}}\calL\paren{\bm{\theta}_k} + \paren{\eta^2mn\gamma + 4C\eta\lambda_0\sqrt{\frac{\gamma m}{n}}}\varepsilon_1\\
        &\quad\quad\quad  + \paren{C'\eta\lambda_0\sqrt{\frac{\gamma m}{n}} + \eta\varepsilon_3\sqrt{mn} - \eta\lambda_0}\calL\paren{\bm{\theta}_k}+ \frac{\eta\varepsilon_2^2 n}{C'\lambda_0}\sqrt{\frac{mn}{\gamma}}\\
        & = \paren{1 - \eta\lambda_0 + \eta\varepsilon_3\sqrt{mn}+ 2\eta^2mn\gamma + \paren{10C + C'}\eta\lambda_0\sqrt{\frac{\gamma m}{n}}}\calL\paren{\bm{\theta}_k}\\
        &\quad\quad\quad + \paren{\eta^2mn\gamma + 4C\eta\lambda_0\sqrt{\frac{\gamma m}{n}}}\varepsilon_1+\frac{\eta\varepsilon_2^2 n}{C'\lambda_0}\sqrt{\frac{mn}{\gamma}}
    \end{align*}
    Apply $\varepsilon_3\leq C_{\varepsilon}\cdot\frac{\lambda_0}{\sqrt{mn}}, \gamma = C_1\cdot\frac{n}{m}$ and $\eta = C_2\cdot \frac{\lambda_0}{n^2}$ gives
    \begin{align*}
            \E_{\bm{\xi}_k}\left[\calL\paren{\bm{\theta}_{k+1}}\right] & \leq \paren{1 - \paren{1 - 2C_1C_2 - \paren{10C + C'}\sqrt{C_1}}\eta\lambda_0}\calL\paren{\bm{\theta}_k}\\
            & \quad\quad\quad + \eta\lambda_0\paren{\frac{mn\varepsilon_2^2}{C'\sqrt{C_1}\lambda_0^2} + \paren{C_1C_2 + 4C\sqrt{C_1}}\varepsilon_1}
    \end{align*}
    Choosing a small enough $C_1, C_2, C, C'$ gives
    \begin{equation}
        \E_{\bm{\xi}_k}\left[\calL\paren{\bm{\theta}_{k+1}}\right] \leq \paren{1 - \frac{\eta\lambda_0}{2}}\calL\paren{\bm{\theta}_k} + \frac{1}{2}\hat{C}\eta\lambda_0\paren{\frac{mn}{\lambda_0^2}\cdot\varepsilon_2^2 + \varepsilon_1}   
    \end{equation}
    for a large enough $\hat{C}$. Thus, unrolling the iterations gives
    \begin{equation}
        \label{eq:thm_conv_proof_7}
        \E_{\bm{\xi}_0,\dots,\bm{\xi}_{k-1}}\left[\calL\paren{\bfW_{k}}\right] \leq \paren{1 - \frac{\eta\lambda_0}{2}}^k\calL\paren{\bfW_0} + \hat{C}\paren{\frac{mn}{\lambda_0^2}\cdot\varepsilon_2^2 + \varepsilon_1}
    \end{equation}
    for all $k < K'$. Next, we shall lower bound $K'$. For all $k \leq K'$, we have that
    \begin{align*}
        \norm{\bfw_{k,r} - \bfw_{0,r}}_2 & \leq \sum_{t=0}^{k-1}\norm{\bfw_{t+1,r} -\bfw_{t,r}}_2 = \eta\sum_{t=0}^{k-1}\norm{\nabla_{\bfw_r}\hat{\calL}\paren{\bfW_t,\bm{\xi}_t}}_2 \leq \eta\sqrt{\gamma}\sum_{t=0}^{k-1}\hat{\calL}\paren{\bfW_t,\bm{\xi}_t}^{\frac{1}{2}}
    \end{align*}
    By (\ref{eq:thm_conv_proof_7}), we have
    \begin{align*}
        \E_{\bm{\xi}_0,\dots,\bm{\xi}_{t-1}}\left[\hat{\calL}\paren{\bfW_t,\bm{\xi}_t}^{\frac{1}{2}}\right] & \leq \paren{\calL\paren{\bfW_t} + \varepsilon_1}^{\frac{1}{2}}\\
        & \leq \paren{2\paren{1 - \frac{\eta\lambda_0}{2}}^t\calL\paren{\bfW_0} + \paren{\hat{C}+1}\paren{\frac{mn}{\lambda_0^2}\cdot\varepsilon_2^2 + \varepsilon_1}}^{\frac{1}{2}}\\
        & \leq 2\paren{1 - \frac{\eta\lambda_0}{4}}^t\calL\paren{\bfW_0}^{\frac{1}{2}} + \sqrt{\hat{C}+1}\paren{\frac{\varepsilon_2}{\lambda_0}\sqrt{mn} + \sqrt{\varepsilon_1}}
    \end{align*}
    Therefore, we have
    \begin{align*}
        \E_{\bm{\xi}_0,\dots,\bm{\xi}_{k-1}}\left[\norm{\bfw_{k,r} - \bfw_{0,r}}_2\right] & \leq \eta\sqrt\gamma\sum_{t=0}^{k-1}\E_{\bm{\xi}_0,\dots,\bm{\xi}_{t-1}}\left[\hat{\calL}\paren{\bfW_t,\bm{\xi}_t}^{\frac{1}{2}}\right]\\
        & \leq 2\eta\sqrt{\gamma}\calL\paren{\bfW_0}^{\frac{1}{2}}\sum_{t=0}^\infty\paren{1 - \frac{\eta\lambda_0}{4}}^t + k\sqrt{\paren{\hat{C}+1}\gamma}\paren{\frac{\varepsilon_2}{\lambda_0}\sqrt{mn} + \sqrt{\varepsilon_1}}\\
        & = \frac{8\sqrt{\gamma}}{\lambda_0}\calL\paren{\bfW_0}^{\frac{1}{2}}+k\sqrt{\paren{\hat{C}+1}\gamma}\paren{\frac{\varepsilon_2}{\lambda_0}\sqrt{mn} + \sqrt{\varepsilon_1}}
    \end{align*}
    By Lemma 26 in \cite{}, we have that $\E_{\bfW_0,\bfa}\left[\calL\paren{\bfW_0}^2\right] = O\paren{n}$. $\gamma = C_1\cdot\frac{n}{m}$, we have that
    \[
        \E_{\bfW_0,\bfa,\bm{\xi}_0,\dots,\bm{\xi}_{k-1}}\left[\norm{\bfw_{k,r} - \bfw_{0,r}}_2\right] \leq O\paren{\frac{n}{\lambda_0}\sqrt{m}} + O\paren{k\paren{\frac{\varepsilon_2n}{\lambda_0} + \sqrt{\frac{\varepsilon_1n}{m}}}}
    \]
    Thus, by Markov's inequality, we have that with probability at least $1 - \frac{\delta}{3K}$,
    \[
        \norm{\bfw_{k,r} - \bfw_{0,r}}_2 \leq \underbrace{O\paren{\frac{nK}{\lambda_0\delta\sqrt{m}}}}_{\mathcal{T}_1} + \underbrace{O\paren{\frac{K^2}{\delta}\paren{\frac{\varepsilon_2n}{\lambda_0} + \sqrt{\frac{\varepsilon_1n}{m}}}}}_{\mathcal{T}_2}
    \]
    Setting $m = \Omega\paren{\frac{n^4}{\lambda_0^4\delta^2\tau^2}}$ guarantees that $T_1 \leq \frac{C_1}{2}\cdot\frac{\tau \lambda_0}{n} = \frac{R}{2}$ and set $\varepsilon_2 \leq O\paren{\frac{\delta\lambda_0}{nK^2}}, \varepsilon_1\leq O\paren{\frac{\delta m}{K^4n}}$ gives that $T_2\leq \frac{C_1}{2}\cdot\frac{\tau \lambda_0}{n} = \frac{R}{2}$. Combining the bound on $\mathcal{T}_1$ and $\mathcal{T}_2$ and taking a union bound gives that, with probability at least $1-\frac{\delta}{3}$, it holds that
    \[
        \norm{\bfw_{k,r} - \bfw_{0,r}}_2 \leq R;\quad \forall k\in[K]
    \]
    This shows that we must have $K' > K$, which completes the proof.
\end{proof}
%%%%%%%
\begin{lemma}
    \label{lem:kernel_diff}
    Let $\bfH^\infty$ be defined in (\ref{eq:infinite_NTK_def}), and let $\bfH_k$ be defined as
    \[
        \bfH_{k,ij} = \frac{\bfx_i^\top\bfx_j}{m}\sum_{r\in S_i}\indy{\bfw_{k,r}^\top\bfx_i\geq 0;\bfw_{k,r}^\top\bfx_j\geq 0}
    \]
    Fix any $R$. Assume that $\bfw_{k,r}\in\mathcal{B}\paren{\bfw_{0,r}, R}$ for all $r\in[m]$. If $\bfw_{0,r}\sim\mathcal{N}\paren{\bm{0},\tau^2\bfI}$, and $m = \Omega\paren{}$, then we have that 
    \[
        \norm{\bfH_k - \bfH^\infty}_F \leq \frac{\lambda_0}{6} + \frac{1}{m^2}\sum_{i,j=1}^n\left|S_i^\perp\right|^2 + \frac{2nR}{\tau}
    \]
\end{lemma}
\begin{proof}
    We define $\hat{\bfH}_k$ as follows
    \[
        \hat{\bfH}_{k,ij} = \frac{\bfx_i^\top\bfx_j}{m}\sum_{r=1}^m\indy{\bfw_{k,r}^\top\bfx_i\geq 0;\bfw_{k,r}^\top\bfx_j\geq 0}
    \]
    Then we have that 
    \[
        \norm{\bfH_k - \bfH^\infty}_F\leq \norm{\bfH_k - \hat{\bfH}_k}_F + \norm{\hat{\bfH}_k - \hat{\bfH}_0}_F + \norm{\hat{\bfH}_0 - \bfH^\infty}_F
    \]
    By Lemma 3.1 in \cite{du2018gradient}, we have that with probability at least $1-\delta$, we have that $\norm{\hat{\bfH}_0 - \bfH^\infty}_F\leq \frac{\lambda_0}{6}$ when $m = \Omega\paren{\frac{n^2}{\lambda_0^2}\log\frac{n}{\delta}}$. By Lemma 3.2 in \cite{song2020quadraticsufficesoverparametrizationmatrix}, we have that with probability at least $1 - n^2\exp{-\frac{mR}{\tau}}$, it holds that $\norm{\hat{\bfH}_k -\hat{\bfH}_0} \leq \frac{2nR}{\tau}$. Lastly, for the first term, we have
    \begin{align*}
        \norm{\bfH_k - \hat{\bfH}_k}_F^2 & \leq \sum_{i,j=1}^n\paren{\bfH_{k,ij} - \hat{\bfH}_{k,ij}}^2\\
        & = \frac{1}{m^2}\sum_{i,j=1}^n\paren{\bfx_i^\top\bfx_j\sum_{r\in S_i^\perp}\indy{\bfw_{k,r}^\top\bfx_i\geq 0;\bfw_{k,r}^\top\bfx_j\geq 0}}^2\\
        & \leq \frac{1}{m^2}\sum_{i,j=1}^n\left|S_i^\perp\right|^2
    \end{align*}
    Combining the three bounds gives the desired result.
\end{proof}
%%%%%%%%%%%%%%%%%%%%%%%%%%%%%%%%%%%%%%%%%%%%%%%%%%%%%%%%%%%%%
\subsection{Proof of Theorem~\ref{thm:conv_gaussian_mask}}
\label{sec:proof_conv_form}
%%%%%%%%%%% proof of main theorem %%%%%%%%%%%%%
We view Gaussian input masking as a special case of the general stochastic
training framework in Section~\ref{sec:conv_guarantee}.  Recall that in that
framework, the randomness at iteration $k$ is denoted by $\bm\xi_k$, and the
update rule is
\[
\bfW_{k+1} = \bfW_k - \eta \nabla_{\bfW}\hat{\calL}(\bfW_k,\bm\xi_k).
\tag{\ref{eq:general_sgd}}
\]
In the Gaussian-masked setting we take
\[
\bm\xi_k \equiv \bfC_k, \qquad
\hat{\calL}(\bfW,\bm\xi_k) \equiv \calL_{\bfC_k}(\bfW),
\]
where $\bfC_k$ is the multiplicative Gaussian mask at iteration $k$ and
$\calL_{\bfC_k}$ is the masked loss.  Thus
\[
\nabla_{\bfw_r}\hat{\calL}(\bfW,\bm\xi_k)
\equiv
\nabla_{\bfw_r}\calL_{\bfC_k}(\bfW),
\]
and the update \eqref{eq:general_sgd} coincides with the masked gradient
descent rule \eqref{eq:gaussin_input_mask_gd}.  Therefore, to apply
Theorem~\ref{thm:gen_sto_conv} to training with Gaussian input masks, it
suffices to verify that Assumption~\ref{asump:randomness} holds with suitable
$\varepsilon_1,\varepsilon_2,\varepsilon_3,\gamma$, and that these parameters
satisfy the smallness conditions of Theorem~\ref{thm:gen_sto_conv} under the
constraints \eqref{eq:kappa_req}–\eqref{eq:mix_req}.

Throughout the proof we condition on the high-probability NTK event of
Theorem~\ref{thm:gen_sto_conv}, on which
\begin{itemize}[leftmargin=*]
  \item the minimum eigenvalue of the empirical NTK satisfies
  $\lambda_{\min}(\bfH_k)\ge \lambda_0/2$ for all $k\in[K]$,
  \item all first-layer weights remain in a ball of radius
  $R = C_1\tau\lambda_0/n$ around their initialization, i.e.,
  $\|\bfw_{k,r}-\bfw_{0,r}\|_2\le R$ for all $k\in[K],r\in[m]$,
  \item the data are bounded as in Assumption~\ref{asump:data}.
\end{itemize}
The probability of this event is at least
$1 - 2\delta - n^2\exp(-n^3/(\delta^2\tau^2\lambda_0^3))$, as in
Theorem~\ref{thm:gen_sto_conv}. All inequalities below hold on this event.

Comparing Corollary~\ref{cor:exp_loss_ub} with
Assumption~\ref{asump:randomness}(\ref{eq:general_exp_loss_ub}), we identify
\begin{equation}
\label{eq:eps1_exact}
    \varepsilon_1(\bfW)
    =
    2mn\kappa^2R_{\bfu}^2
     + 
    mn\big(\kappa^2R_{\bfu}^2 + \kappa R_{\bfw}\big)\phi_{\max}(\bfW)^2
     + 
    mn\kappa^2\big(R_{\bfu}^2+1\big)\psi_{\max}(\bfW)^2.
\end{equation}

On the NTK event, the weights stay close to initialization, hence their norms
are uniformly bounded; using Assumption~\ref{asump:data} and the definition of
$R_{\bfw}$ and $R_{\bfu}$, we obtain:
\begin{equation}
\label{eq:Rw_bound}
    R_{\bfw}(\bfW_k)
    :=
    \max_{r\in[m]}\|\bfw_{k,r}\|_2
     \le 
    C_w\tau \sqrt{d}
\end{equation}
\begin{equation}
\label{eq:Ru_bound}
    R_{\bfu}(\bfW_k)
    :=
    \max_{r\in[m],i\in[n]}\|\bfw_{k,r}\odot\bfx_i\|_2
     \le 
    C_u\tau\sqrt{d}.
\end{equation}
for constants $C_w,C_u>0$.  
% \subsection*{Proof of why $R_w(\mathbf{W}_k)\le C_w \tau$}
% Recall that the first–layer weights are initialized as
% \[
%     \mathbf w_{0,r} \sim \mathcal N(\mathbf 0,\tau^2 I_d),
%     \qquad r=1,\dots,m,
% \]
% and that during training, on the NTK event of Theorem~\ref{thm:gen_sto_conv},
% each row stays in a small ball around its initialization:
% \begin{equation}
%     \label{eq:NTK_ball}
%     \|\mathbf w_{k,r} - \mathbf w_{0,r}\|_2  \le  R,
%     \qquad
%     R := C_1  \tau \frac{\lambda_0}{n},
%     \qquad \forall \enspace k\le K, \enspace r\in[m].
% \end{equation}

% We define
% \[
%     \mathbf z_r := \frac{1}{\tau} \mathbf w_{0,r}.
% \]
% Then each coordinate satisfies
% \[
%     (z_r)_j = \frac{(w_{0,r})_j}{\tau},
%     \qquad j=1,\dots,d.
% \]
% Because $(w_{0,r})_j\sim\mathcal N(0,\tau^2)$, dividing by $\tau$ gives
% $(z_r)_j\sim\mathcal N(0,1)$. The coordinates remain independent, so
% $\mathbf z_r$ is a standard Gaussian vector:
% \[
%     \mathbf z_r \sim \mathcal N(\mathbf 0, I_d).
% \]
% And we can write:
% \begin{equation}
%     \label{eq:w0_equals_tau_z}
%     \mathbf w_{0,r} = \tau \mathbf z_r,
%     \qquad \forall r\in[m].
% \end{equation}

% We next analyze the norm of $\mathbf z_r$.
% Since $\mathbf z_r\sim\mathcal N(\mathbf 0,I_d)$, its squared norm satisfies
% \[
%     \|\mathbf z_r\|_2^2 = \sum_{j=1}^d z_{r,j}^2 \sim \chi^2_d,
% \]
% a chi–square random variable with $d$ degrees of freedom.
% A standard concentration inequality for $\chi^2_d$ (see, e.g.,
% Laurent--Massart) states that for all $t>0$,
% \begin{equation}
%     \label{eq:chi2_tail}
%     \Pr\big(\|\mathbf z_r\|_2^2 \ge d + 2\sqrt{dt} + 2t\big)
%      \le  e^{-t}.
% \end{equation}
% Choosing $t=d$ yields
% \[
%     d + 2\sqrt{d t} + 2t
%     = d + 2\sqrt{d^2} + 2d
%     = d + 2d + 2d
%     = 5d,
% \]
% and therefore
% \[
%     \Pr\big(\|\mathbf z_r\|_2^2 \ge 5d\big)
%      \le  e^{-d}.
% \]
% Equivalently,
% \[
%     \Pr\big(\|\mathbf z_r\|_2^2 \le 5d\big)
%      \ge  1 - e^{-d}.
% \]
% As d increases the RHS becomes 1 as $e^{-d} \approx0$ (for example, for $d=8$, $e^{-8}=0.00035$), so with high probability it holds that
% \[
% \|\mathbf z_r\|_2^2 \le 5d
% \]
% Taking the square root:
% \[
%     \|\mathbf z_r\|_2 \le \sqrt{5d} =: C_0\sqrt{d}.
% \]
% Thus there exist absolute constants $C_0>0$ and $c>0$ such that
% \begin{equation}
%     \label{eq:z_norm_concentration}
%     \Pr\big(\|\mathbf z_r\|_2 \le C_0\sqrt{d}\big)  \ge  1 - e^{-cd}.
% \end{equation}

% Combining \eqref{eq:w0_equals_tau_z} and \eqref{eq:z_norm_concentration},
% we obtain a high–probability bound on the norm at initialization:
% \begin{equation}
%     \label{eq:w0_norm_bound}
%     \|\mathbf w_{0,r}\|_2
%     = \|\tau \mathbf z_r\|_2
%     = \tau \|\mathbf z_r\|_2
%     \le \tau   C_0 \sqrt{d}.
% \end{equation}
%  Combining the triangle inequality with ~\eqref{eq:NTK_ball} and ~\eqref{eq:w0_norm_bound} we obtain
% \begin{align*}
%     \|\mathbf w_{k,r}\|_2
%     &\le
%     \|\mathbf w_{0,r}\|_2 + \|\mathbf w_{k,r}-\mathbf w_{0,r}\|_2
%     \nonumber\\
%     &\le
%     \|\mathbf w_{0,r}\|_2 + R
%     \nonumber\\
%     &\le
%     \tau   C_0 \sqrt{d} + C_1 \tau \frac{\lambda_0}{n}
%     \\
%     &=
%     (C_0 \sqrt{d} + C_1 \frac{\lambda_0}{n}) \tau
%     \\
%     &=
%     C_w\tau\sqrt{d}
% \end{align*}

% Finally,
% \begin{equation}
%     \label{eq:wkr_final_bound}
%     \|\mathbf w_{k,r}\|_2  \le  C_w \tau \sqrt{d},
%     \qquad \forall k\le K,  \forall r\in[m],
% \end{equation}
% Taking the maximum over $r$ gives the desired bound:
% \[
%     R_w(\mathbf W_k)
%     :=
%     \max_{r\in[m]}\|\mathbf w_{k,r}\|_2
%      \le 
%     C_w \tau \sqrt{d},
%     \qquad \forall k\le K,
% \]
% with high probability. Intuitively, this shows that all first–layer
% weights remain on the same scale as their initialization $\tau$
% throughout training; the NTK dynamics only move each row by a small
% amount compared to its initial magnitude.

\begin{lemma}[Bound on $R_{\bfw}$]
\label{lem:Rw_bound}
On the NTK event of Theorem~\ref{thm:gen_sto_conv}, there exists an absolute constant $C_w > 0$ such that for all iterations $k \le K$:
\begin{equation}
    R_{\bfw}(\mathbf{W}_k) := \max_{r \in [m]} \|\mathbf{w}_{k,r}\|_2 \le C_w \tau \sqrt{d}.
\end{equation}
\end{lemma}

\begin{proof}
Recall that the first--layer weights are initialized as $\mathbf{w}_{0,r} \sim \mathcal{N}(\mathbf{0}, \tau^2 \mathbf{I}_d)$ for $r=1,\dots,m$. During training, the NTK event of Theorem~\ref{thm:gen_sto_conv} ensures that each row stays in a small ball around its initialization:
\begin{equation}
    \label{eq:NTK_ball}
    \|\mathbf{w}_{k,r} - \mathbf{w}_{0,r}\|_2 \le R, \qquad R := C_1 \tau \frac{\lambda_0}{n}, \qquad \forall k \le K, r \in [m].
\end{equation}

We define $\mathbf{z}_r := \frac{1}{\tau}\mathbf{w}_{0,r}$. Each coordinate satisfies $(z_r)_j \sim \mathcal{N}(0,1)$, making $\mathbf{z}_r$ a standard Gaussian vector $\mathcal{N}(\mathbf{0}, \mathbf{I}_d)$. Its squared norm follows a chi-square distribution: $\|\mathbf{z}_r\|_2^2 \sim \chi^2_d$. Using the Laurent--Massart concentration inequality, for $t=d$:
\[
    \Pr\big( \|\mathbf{z}_r\|_2^2 \ge 5d \big) \le e^{-d}.
\]
Thus, with high probability, $\|\mathbf{z}_r\|_2 \le \sqrt{5d}$. Defining $C_0 = \sqrt{5}$, we obtain the initialization bound:
\begin{equation}
    \label{eq:w0_norm_bound}
    \|\mathbf{w}_{0,r}\|_2 = \tau \|\mathbf{z}_r\|_2 \le C_0 \tau \sqrt{d}.
\end{equation}

By a union bound over $r \in [m]$, this holds for all rows with probability at least $1 - me^{-d}$. Combining the triangle inequality with \eqref{eq:NTK_ball} and \eqref{eq:w0_norm_bound}, we find:
\begin{align*}
    \|\mathbf{w}_{k,r}\|_2 &\le \|\mathbf{w}_{0,r}\|_2 + \|\mathbf{w}_{k,r} - \mathbf{w}_{0,r}\|_2 \\
    &\le C_0 \tau \sqrt{d} + C_1 \tau \frac{\lambda_0}{n} \\
    &= \tau \sqrt{d} \left( C_0 + \frac{C_1 \lambda_0}{n \sqrt{d}} \right).
\end{align*}
Since $\lambda_0/n$ is $O(1)$ and $d \ge 1$, we define the absolute constant $C_w := C_0 + \frac{C_1 \lambda_0}{n \sqrt{d}}$. Taking the maximum over $r \in [m]$ yields:
\begin{equation}
    R_w(\mathbf{W}_k) = \max_{r \in [m]} \|\mathbf{w}_{k,r}\|_2 \le C_w \tau \sqrt{d}.
\end{equation}
Intuitively, since the movement term $\frac{C_1 \lambda_0}{n \sqrt{d}}$ vanishes as $d \to \infty$, the weights remain on the same scale as their initialization throughout training.
\end{proof}

\begin{lemma}[Bound on $R_u$]
\label{lem:Ru_bound}
Suppose Assumption~\ref{asump:data} holds, such that the input data is bounded in $\ell_\infty$-norm by $B_x := \max_{i\in[n]}\|x_i\|_\infty$. On the NTK event of Theorem~\ref{thm:gen_sto_conv}, there exists an absolute constant $C_u > 0$ such that, for all iterations $k \le K$:
\begin{equation}
    R_{u}(W_k) := \max_{r \in [m], i \in [n]} \|w_{k,r} \odot x_i\|_2 \le C_u \tau \sqrt{d}.
\end{equation}
\end{lemma}

\begin{proof}
Consider any iteration $k \le K$, neuron $r \in [m]$, and sample index $i \in [n]$. We analyze the squared $\ell_2$-norm of the Hadamard product by pulling out the maximum coordinate of the input vector:
\begin{align*}
    \|w_{k,r} \odot x_i\|_2^2 &= \sum_{j=1}^d w_{k,r,j}^2 x_{i,j}^2 \\
    &\le \left( \max_{1 \le j \le d} x_{i,j}^2 \right) \sum_{j=1}^d w_{k,r,j}^2 \\
    &= \|x_i\|_\infty^2 \|w_{k,r}\|_2^2.
\end{align*}
Taking the square root of both sides, we obtain the inequality:
\[ \|w_{k,r} \odot x_i\|_2 \le \|x_i\|_\infty \|w_{k,r}\|_2 \le B_x \|w_{k,r}\|_2. \]
Taking the maximum over all $r \in [m]$ and $i \in [n]$ yields:
\begin{equation}
    \label{eq:Ru_vs_Rw_enhanced}
    R_u(W_k) \le B_x R_w(W_k).
\end{equation}

From the weight stability bound previously established (Proof of $R_w$), we know that on the NTK event, the weights are bounded by $R_w(W_k) \le C_w \tau \sqrt{d}$, where $C_w$ is an absolute constant. Substituting this into \eqref{eq:Ru_vs_Rw_enhanced}:
\[ R_u(W_k) \le B_x (C_w \tau \sqrt{d}). \]
Defining the absolute constant $C_u := B_x C_w$ completes the proof. Note that since $B_x$ is a fixed property of the dataset and $C_w$ is independent of $k$, $C_u$ is a valid absolute constant for the problem.
\end{proof}
% %%%%%%
% Substituting into the expression for
% $\varepsilon_1(\bfW_k)$, and denoting
% $\hat\phi_{\max} := \max_k\phi_{\max}(\bfW_k)$,
% $\hat\psi_{\max} := \max_k\psi_{\max}(\bfW_k)$, we obtain the uniform bound
% \begin{align}
% \varepsilon_1
% :=
% \max_{k\in[K]}\varepsilon_1(\bfW_k)
% &\le
% O\big(\kappa^2\tau^2 mn d\big)
% + O\big(\kappa\tau^2 mn d^2(\hat\phi_{\max}^2+\hat\psi_{\max}^2)\big).
% \label{eq:eps1_scaling}
% \end{align}
% %%%%%%%
We an iteration $k\in[K]$ and denote, for brevity,
\[
    R_{\bfw} := R_{\bfw}(\bfW_k),
    \quad
    R_{\bfu} := R_{\bfu}(\bfW_k),
    \quad
    \phi_k := \phi_{\max}(\bfW_k),
    \quad
    \psi_k := \psi_{\max}(\bfW_k).
\]
Then \eqref{eq:eps1_exact} becomes
\begin{equation}
\label{eq:eps1Wk_expanded}
    \varepsilon_1(\bfW_k)
    =
    2mn\kappa^2R_{\bfu}^2
     + 
    mn\big(\kappa^2R_{\bfu}^2 + \kappa R_{\bfw}\big)\phi_k^2
     + 
    mn\kappa^2\big(R_{\bfu}^2+1\big)\psi_k^2.
\end{equation}
We now bound each of the three terms on the right-hand side using
\eqref{eq:Rw_bound},\eqref{eq:Ru_bound}.

\medskip\noindent
\emph{(i) First term.}
Using $R_{\bfu}^2 \le C_u^2\tau^2 d$ from \eqref{eq:Ru_bound}, we get
\begin{align}
    2mn\kappa^2R_{\bfu}^2
    &\le
    2mn\kappa^2\cdot C_u^2\tau^2 d
    \nonumber\\
    &=
    (2C_u^2) \kappa^2\tau^2 m n d
     = 
    O\big(\kappa^2\tau^2 m n d\big).
    \label{eq:eps1_term1}
\end{align}

\medskip\noindent
\emph{(ii) Middle term.}
\[
    mn\big(\kappa^2R_{\bfu}^2 + \kappa R_{\bfw}\big)\phi_k^2
    =
    mn\kappa^2R_{\bfu}^2\phi_k^2
     + 
    mn\kappa R_{\bfw}\phi_k^2.
\]

For the $\kappa^2R_{\bfu}^2\phi_k^2$ component, we use $R_{\bfu}^2\le C_u^2\tau^2 d$ from \eqref{eq:Ru_bound}:
\begin{align}
    mn\kappa^2R_{\bfu}^2\phi_k^2
    &\le
    mn\kappa^2 (C_u^2\tau^2 d) \phi_k^2
    \nonumber\\
    &=
    C_u^2 \kappa^2\tau^2 m n d \phi_k^2
     = 
    O\big(\kappa^2\tau^2 m n d \phi_k^2\big).
    \label{eq:eps1_term2a}
\end{align}

For the $\kappa R_{\bfw}\phi_k^2$ component, we use $R_{\bfw}\le C_w\tau \sqrt{d}$ from \eqref{eq:Rw_bound}:
\begin{align}
    mn\kappa R_{\bfw}\phi_k^2
    &\le
    mn\kappa (C_w\tau \sqrt{d}) \phi_k^2
    \nonumber\\
    &=
    C_w \kappa\tau m n \sqrt{d} \phi_k^2
     = 
    O\big(\kappa\tau m n \sqrt{d} \phi_k^2\big).
    \label{eq:eps1_term2b}
\end{align}

\medskip\noindent
\emph{(iii) Last term.}
For the last term, using \eqref{eq:Ru_bound}, it is true that
$R_{\bfu}^2+1 \le C_u^2\tau^2 d + 1$, so there exists a constant $C_u'>0$
such that $R_{\bfu}^2 + 1  \le  C_u'\tau^2 d$ for big enough $d$.\\
Hence,
\begin{align}
    mn\kappa^2\big(R_{\bfu}^2+1\big)\psi_k^2
    &\le
    mn\kappa^2(C_u'\tau^2 d) \psi_k^2
    \nonumber\\
    &=
    C_u' \kappa^2\tau^2 m n d \psi_k^2
     = 
    O\big(\kappa^2\tau^2 m n d \psi_k^2\big).
    \label{eq:eps1_term3}
\end{align}

Combining \eqref{eq:eps1Wk_expanded} with
\eqref{eq:eps1_term1}, \eqref{eq:eps1_term2a}, \eqref{eq:eps1_term2b},
and \eqref{eq:eps1_term3}, we obtain
\begin{align}
    \varepsilon_1(\bfW_k)
    &\le
    O\big(\kappa^2\tau^2 m n d\big)
     + 
    O\big(\kappa^2\tau^2 m n d \phi_k^2\big)
     + 
    O\big(\kappa\tau m n \sqrt{d} \phi_k^2\big)
     + 
    O\big(\kappa^2\tau^2 m n d \psi_k^2\big)
    \nonumber\\
    &=
    O\big(\kappa^2\tau^2 m n d\big)
    +
    O\big(\kappa^2\tau^2 m n d(\phi_k^2 + \psi_k^2)\big)
    +
    O\big(\kappa\tau m n \sqrt{d} \phi_k^2\big).
    \label{eq:eps1Wk_combined}
\end{align}
We denote
\[
    \hat\phi_{\max} := \max_{k\in[K]}\phi_{\max}(\bfW_k),
    \qquad
    \hat\psi_{\max} := \max_{k\in[K]}\psi_{\max}(\bfW_k),
\]
so that for each $k$,
\[
    \phi_k^2 \le \hat\phi_{\max}^2,
    \qquad
    \psi_k^2 \le \hat\psi_{\max}^2.
\]
Substituting these into \eqref{eq:eps1Wk_combined} yields
\begin{align}
    \varepsilon_1(\bfW_k)
    &\le
    O\big(\kappa^2\tau^2 m n d\big)
    +
    O\big(\kappa^2\tau^2 m n d(\hat\phi_{\max}^2 + \hat\psi_{\max}^2)\big)
    +
    O\big(\kappa\tau m n \sqrt{d}  \hat\phi_{\max}^2\big),
    \qquad \forall k\in[K].
    \label{eq:eps1Wk_uniform}
\end{align}
Taking the maximum over $k\in[K]$ does not change the right-hand side,
so
\begin{equation}
\label{eq:eps1_before_simplify}
    \varepsilon_1
    :=
    \max_{k\in[K]}\varepsilon_1(\bfW_k)
     \le 
    O\big(\kappa^2\tau^2 m n d\big)
    +
    O\big(\kappa^2\tau^2 m n d(\hat\phi_{\max}^2 + \hat\psi_{\max}^2)\big)
    +
    O\big(\kappa\tau m n \sqrt{d}  \hat\phi_{\max}^2\big).
\end{equation}
or equivalently: 
\begin{equation}
\label{eq:eps1_final}
    \boxed{\varepsilon_1
     \le 
    O\big(\kappa^2\tau^2 m n d(\hat\phi_{\max}^2 + \hat\psi_{\max}^2 +1)\big)
    +
    O\big(\kappa\tau m n \sqrt{d}  \hat\phi_{\max}^2\big)}
\end{equation}

Matching Corollary~\ref{cor:exp_grad_err_ub} with Assumption~\ref{asump:randomness}(\ref{eq:general_exp_grad_err}), we read off
\[
\varepsilon_3(\bfW)
= O\Big(\frac{\sigma_{\max}(\bfX)\phi_{\max}(\bfW)}{\sqrt{m}}\Big)
\]
and
\begin{align}
\varepsilon_2(\bfW)
&=
O\big((n\kappa^2B_{\bfx}^2R_{\bfw} + n\kappa R_{\bfu}\sqrt{d})\phi_{\max}(\bfW)\big)
\nonumber\\
&\quad
+ O\big(n\kappa R_{\bfu}\psi_{\max}(\bfW)+\kappa^2\sqrt{m}B_{\bfx}^2R_{\bfw}\big).
\label{eq:eps2_raw}
\end{align}

Using $\|\bfx_i\|_2\le 1$ from Assumption~\ref{asump:data}, we have
$B_{\bfx}\le 1$.  On the NTK event, for all $k\in[K]$ we have the uniform
bounds
\[
R_{\bfw}(\bfW_k) \le C_w\tau\sqrt{d},
\qquad
R_{\bfu}(\bfW_k) \le C_u\tau\sqrt{d},
\]
for some absolute constants $C_w,C_u>0$ (cf.\ the bounds proved earlier for
$R_{\bfw}$ and $R_{\bfu}$).
Substituting these into \eqref{eq:eps2_raw} yields:
\begin{align}
\varepsilon_2(\bfW_k)
&\le
O\Big(
   \big(n\kappa^2 C_w\tau\sqrt{d}
      + n\kappa C_u\tau\sqrt{d} \sqrt{d}
   \big)\phi_{\max}(\bfW_k)
\Big)
\nonumber\\
&\quad
+ O\Big(
   n\kappa C_u\tau\sqrt{d} \psi_{\max}(\bfW_k)
   + \kappa^2\sqrt{m} C_w\tau\sqrt{d}
\Big)
\nonumber\\[0.3em]
&=
O\big(n\kappa^2\tau\sqrt{d} \phi_{\max}(\bfW_k)\big)
+ O\big(n\kappa\tau d \phi_{\max}(\bfW_k)\big)
\nonumber\\
&\quad
+ O\big(n\kappa\tau\sqrt{d} \psi_{\max}(\bfW_k)\big)
+ O\big(\kappa^2\tau\sqrt{m d}\big).
\label{eq:eps2_subst_sqrt_d}
\end{align}

Since $d\ge 1$, we have $\sqrt{d}\le d$, and hence each term containing
$\sqrt{d}$ can be upper bounded by the corresponding expression with $d$ in
place of $\sqrt{d}$.  Therefore,
\begin{align}
\varepsilon_2(\bfW_k)
&\le
O\big(n\kappa^2\tau d \phi_{\max}(\bfW_k)\big)
+ O\big(n\kappa\tau d \phi_{\max}(\bfW_k)\big)
\nonumber\\[-0.2em]
&\quad
+ O\big(n\kappa\tau d \psi_{\max}(\bfW_k)\big)
+ O\big(\kappa^2\tau\sqrt{m}d\big)
\nonumber\\[0.2em]
&\le
O\big(\kappa\tau n d\big(
      \phi_{\max}(\bfW_k)+\psi_{\max}(\bfW_k)
   \big)\big)
\nonumber\\
&\quad
+ O\big(\kappa^2\tau d\big(
      n \phi_{\max}(\bfW_k)+\sqrt{m}
   \big)\big).
\label{eq:eps2_scaling_pointwise_correct}
\end{align}

To obtain a uniform bound over the whole training trajectory, define
\[
\hat\phi_{\max} := \max_{k\in[K]}\phi_{\max}(\bfW_k),
\qquad
\hat\psi_{\max} := \max_{k\in[K]}\psi_{\max}(\bfW_k).
\]
Taking the maximum over $k$ in \eqref{eq:eps2_scaling_pointwise_correct} yields
\begin{equation*}
\label{eq:eps2_scaling_correct}
\varepsilon_2
:= \max_{k\in[K]}\varepsilon_2(\bfW_k)
 \le 
O\big(\kappa\tau n d(\hat\phi_{\max}+\hat\psi_{\max})\big)
+ O\big(\kappa^2\tau d\big(n\hat\phi_{\max}+\sqrt{m}\big)\big).
\end{equation*}
% In particular, the leading linear-in-$\kappa$ term scales like
% $\kappa\tau n d(\hat\phi_{\max}+\hat\psi_{\max})$, while the remaining terms
% are higher order in $\kappa$.
We can simplify $\varepsilon_2$ further as:
\begin{align*}
    O\big(\kappa\tau n d(\hat\phi_{\max}+\hat\psi_{\max})\big)
+ O\big(\kappa^2\tau d\big(n\hat\phi_{\max}+\sqrt{m}\big)\big) &= O\big(\kappa\tau n d(\hat\phi_{\max}+\hat\psi_{\max})+\kappa^2\tau d\big(n\hat\phi_{\max}+\sqrt{m}\big)\big) \\
&= O\big(\kappa\tau n d(\hat\phi_{\max}+\hat\psi_{\max})+\kappa^2\tau dn\hat\phi_{\max}+\kappa^2\tau d\sqrt{m}\big) \\
&\leq O\big(\kappa\tau n d(\hat\phi_{\max}+\hat\psi_{\max})+\kappa\tau dn(\hat\phi_{\max}+\hat\psi_{\max})+\kappa^2\tau d\sqrt{m}\big) 
%\text{for $\kappa \leq 1 \Rightarrow \kappa^2 \leq \kappa$}
\\
&= O\big(2\kappa\tau n d(\hat\phi_{\max}+\hat\psi_{\max})+\kappa^2\tau d\sqrt{m}\big) 
\end{align*}
Thus, 
\begin{equation}
\label{eq:eps2_scaling_correct}
\boxed{\varepsilon_2 \leq
O\big(\kappa\tau n d(\hat\phi_{\max}+\hat\psi_{\max})\big)+ O\big(\kappa^2\tau d\sqrt{m}\big)}
\end{equation}
%%%%%%%%%%%%%% PAPER MATERIAL %%%%%%%%%%%%%%
Theorem~\ref{thm:gen_sto_conv} requires
$\varepsilon_2$ to satisfy, for some absolute constant $c_0>0$,
\begin{equation}
\label{eq:eps2_requirement}
    \varepsilon_2
    \;\le\;
    c_0 \frac{\delta\lambda_0}{nK^2}
\end{equation}
Combining \eqref{eq:eps2_scaling_correct} and
\eqref{eq:eps2_requirement}, we get that the following inequality must hold: 
\begin{equation}
\label{eq:eps2_ineq_raw}
    C_1 \kappa\tau n d\bigl(\hat\phi_{\max}+\hat\psi_{\max}\bigr)
    +
    C_2 \kappa^2\tau d\sqrt{m}
    \;\le\;
    c_0 \frac{\delta\lambda_0}{nK^2}.
\end{equation}

Let us denote
\[
    a := C_1 \kappa\tau n d\bigl(\hat\phi_{\max}+\hat\psi_{\max}\bigr),
    \qquad
    b := C_2 \kappa^2\tau d\sqrt{m},
    \qquad
    R := c_0 \frac{\delta\lambda_0}{nK^2}.
\]
Then \eqref{eq:eps2_ineq_raw} can be written as $a+b\le R$.

A sufficient way to enforce this inequality is to
require that each term $a$ and $b$ is at most $R/2$:
\begin{equation}
\label{eq:half_trick}
    a \;\le\; \frac{R}{2},
    \qquad
    b \;\le\; \frac{R}{2}
\end{equation}
Indeed, if \eqref{eq:half_trick} holds, then
\[
    a+b \;\le\; \frac{R}{2} + \frac{R}{2} = R,
\]
so \eqref{eq:eps2_ineq_raw} is automatically satisfied.

Imposing $a\le R/2$ yields:
\[
    C_1 \kappa\tau n d\bigl(\hat\phi_{\max}+\hat\psi_{\max}\bigr)
    \;\le\;
    \frac{c_0}{2} \frac{\delta\lambda_0}{nK^2}. 
\]
Thus,
\begin{equation}
\label{eq:kappa_lin}
    \kappa
    \;\le\;
    \kappa_{\mathrm{lin}_2}
    :=
    \frac{c_0}{2C_1} 
    \frac{\delta\lambda_0}{\tau n^2 d K^2\bigl(\hat\phi_{\max}+\hat\psi_{\max}\bigr)}.
\end{equation}

Similarly, imposing $b\le R/2$ gives:
\begin{align*}
    C_2 \kappa^2\tau d\sqrt{m}
    \;\le\;
    \frac{c_0}{2} \frac{\delta\lambda_0}{nK^2} \enspace \Rightarrow \\
    \kappa^2
    \;\le\;
    \frac{c_0}{2C_2} 
    \frac{\delta\lambda_0}{\tau d\sqrt{m} nK^2}
\end{align*}
and therefore
\begin{equation}
\label{eq:kappa_quad}
    \kappa
    \;\le\;
    \kappa_{\mathrm{quad}_2}
    :=
    \sqrt{\frac{c_0}{2C_2}} 
    \sqrt{\frac{\delta\lambda_0}{\tau d\sqrt{m} nK^2}}.
\end{equation}
To ensure \eqref{eq:eps2_requirement} holds, it is sufficient that
\eqref{eq:kappa_lin} and \eqref{eq:kappa_quad} both hold. 
Equivalently,
\[
    \kappa
    \;\le\;
    \min\{\kappa_{\mathrm{lin}_2}, \kappa_{\mathrm{quad}_2}\}.
\]
In big-$O$ notation we may write this as
\begin{equation}
\label{eq:kappa_min_form}
    \kappa
    =
    O\!\left(
        \frac{\delta\lambda_0}{\tau n^2 d K^2(\hat\phi_{\max}+\hat\psi_{\max})}
    \right)
    \quad\text{and}\quad
    \kappa
    =
    O\!\left(
        \sqrt{\frac{\delta\lambda_0}{\tau d\sqrt{m} nK^2}}
    \right).
\end{equation}

Similarly, for $\varepsilon_1$ the general stochastic convergence theorem requires that
\[
\varepsilon_1
\;\le\;
c_1 \frac{\delta m}{nK^4},
\tag{4}
\label{eq:eps1_requirement}
\]
for some absolute constant $c_1>0$.

Combining \eqref{eq:eps1_final} and \eqref{eq:eps1_requirement} we get that the following must hold:
\begin{align*}
\label{eq:eps1_req_expanded}
O\big(\kappa^2\tau^2 m n d(\hat\phi_{\max}^2 + \hat\psi_{\max}^2 +1)\big)+ O\big(\kappa\tau m n \sqrt{d}  \hat\phi_{\max}^2\big) \;\le\;
c_1 \frac{\delta m}{nK^4} \enspace \Rightarrow \\
C_a \kappa^2 \tau^2 n d\bigl(\hat\phi_{\max}^2 + \hat\psi_{\max}^2+1 \bigr)
+
C_b \kappa \tau n \sqrt{d} \hat\phi_{\max}^2
\;\le\;
c_1 \frac{\delta}{nK^4}.
\end{align*}

Similarly, we impose
\begin{align*}
C_a \kappa^2 \tau^2 n d\bigl(1 + \hat\phi_{\max}^2 + \hat\psi_{\max}^2\bigr)
\;\le\;
\frac{c_1}{2} \frac{\delta}{nK^4} \enspace \Rightarrow \\
\kappa^2
\;\le\;
\frac{c_1}{2C_a} 
\frac{\delta}
     {\tau^2 n^2 d K^4\bigl(\hat\phi_{\max}^2 + \hat\psi_{\max}^2 +1 \bigr)} \Rightarrow
\end{align*}
\begin{equation}
\label{eq:kappa_from_eps1_quad}
\kappa
\;\le\;
\kappa_{\mathrm{quad}_1}
:=
\sqrt{\frac{c_1}{2C_a}}\cdot
\frac{\sqrt{\delta}}
     {\tau n \sqrt{d} K^2 \sqrt{\hat\phi_{\max}^2 + \hat\psi_{\max}^2+1}}.
\end{equation}
and
\begin{align*}
C_b \kappa \tau n \sqrt{d} \hat\phi_{\max}^2
\;\le\;
\frac{c_1}{2} \frac{\delta}{nK^4} \enspace \Rightarrow
\end{align*}
\begin{equation}
\label{eq:kappa_from_eps1_lin}
\kappa
\;\le\;
\kappa_{\mathrm{lin_1}}
:=
\frac{c_1}{2C_b} 
\frac{\delta}
     {\tau n^2 \sqrt{d} K^4 \hat\phi_{\max}^2}.
\end{equation}

Thus, the $\varepsilon_1$ requirement
\eqref{eq:eps1_requirement} is guaranteed whenever
\[
\kappa
\;\le\;
\min\bigl\{\kappa_{\mathrm{quad}_1},\;\kappa_{\mathrm{lin}_1}\bigr\}.
\]

Combining all constraints from $\varepsilon_1$
and $\varepsilon_2$, we see that a sufficient set of conditions is
\[
\kappa
\;\le\;
\min\bigl\{
\kappa_{\mathrm{lin}_1},
\kappa_{\mathrm{quad}_1},
\kappa_{\mathrm{lin}_2},
\kappa_{\mathrm{quad}_2}
\bigr\}.
\]
Equivalently, in big-$O$ notation,
\begin{equation}
\label{eq:kappa_min_bigO}
\kappa
=
O\!\left(
\min\left\{
\frac{\delta\lambda_0}{\tau n^2 d K^2(\hat\phi_{\max}+\hat\psi_{\max})},
\;
\frac{\delta}{\tau n^2 \sqrt{d} K^4 \hat\phi_{\max}^2},
\;
\sqrt{\frac{\delta\lambda_0}{\tau d\sqrt{m} nK^2}},
\;
\sqrt{\frac{\delta}{\tau^2 n^2 d K^4(\hat\phi_{\max}^2+\hat\psi_{\max}^2+1)}}
\right\}
\right).
\end{equation}
Instead of carrying this minimum in the theorem statement, we define a slightly more restrictive but cleaner condition that implies all of the above bounds: 
\begin{equation}
\label{eq:kappa_clean}
    \kappa
    \;=\;
    O\!\left(
        \frac{\sqrt{\delta\lambda_0}}
             {\tau^2 K^2\bigl(m^{1/4}\sqrt{d} + n d\bigr) \big(\hat\phi_{\max} + \hat\psi_{\max}\big)}
    \right).
\end{equation}

%%%%%%%%%%%%%%%%%%%%%%%
For $\varepsilon_3$, we combine \eqref{eq:general_exp_grad_err} and \eqref{cor:exp_grad_err_ub} which yields that:
\[
\varepsilon_3
=
O\Big(\frac{\sigma_{\max}(\bfX)\phi_{\max}}{\sqrt{m}}\Big)
\]
We assume that 
\[
\sigma_{\max}(\bfX)\hat\phi_{\max}\le C \lambda_0/\sqrt{n}
\]
for some constant $C>0$ (assumption~\eqref{eq:mix_req}), and therefore
\begin{equation}
\varepsilon_3
 \le 
O\Big(\frac{\lambda_0}{\sqrt{mn}}\Big),
\label{eq:eps3_scaling}
\end{equation}
which is exactly the form required in Theorem~\ref{thm:gen_sto_conv}.
%%%%%%%%%%%%%%%%%%%%%%%%%%
%Smoothness condition: identification of $\gamma$
% Big O-s in main theorem
% \begin{equation}
% \label{eq:eps1_before_simplify}
%     \varepsilon_1
%     :=
%     \max_{k\in[K]}\varepsilon_1(\bfW_k)
%      \le 
%     O\big(\kappa^2\tau^2 m n d\big)
%     +
%     O\big(\kappa^2\tau^2 m n d(\hat\phi_{\max}^2 + \hat\psi_{\max}^2)\big)
%     +
%     O\big(\kappa\tau m n \sqrt{d}  \hat\phi_{\max}^2\big).
% \end{equation}
% \begin{equation}
% \label{eq:eps2_scaling_correct}
% \boxed{\varepsilon_2 \leq
% O\big(\kappa\tau n d(\hat\phi_{\max}+\hat\psi_{\max})\big)+ O\big(\kappa^2\tau d\sqrt{m}\big)}
% \end{equation}

Theorem \ref{thm:gen_sto_conv} includes the factor
\begin{align*}
    O\paren{\frac{mn}{\lambda_0^2}\cdot\varepsilon_2^2 + \varepsilon_1} \stackrel{\eqref{eq:eps2_scaling_correct}, \eqref{eq:eps1_before_simplify}}{=} 
\end{align*}
So we study the quantity $\frac{mn}{\lambda_0^2}\cdot\varepsilon_2^2$.
\begin{align}
    \frac{mn}{\lambda_0^2}\cdot\varepsilon_2^2 &=
    \frac{mn}{\lambda_0^2}\cdot \Big(\kappa^2\tau^2 n^2 d^2(\hat\phi_{\max}+\hat\psi_{\max})^2+ \kappa^4\tau^2 d^2 m\Big) \notag\\
    &\leq \frac{2}{\lambda_0^2} \kappa^2\tau^2 n^3 d^2m(\hat\phi_{\max}^2+\hat\psi_{\max}^2)
    + \frac{1}{\lambda_0^2}\kappa^4\tau^2 d^2 m^2n \notag\\
    &= \textcolor{purple}{O\paren{\kappa \tau^2 m n^3 d^2\big(\hat\phi_{\max}^2+\hat\psi_{\max}^2\big)}}
    + \textcolor{teal}{O\paren{\kappa^2\tau^2m^2nd^2}}
    \label{eq:eps2_bigO}
\end{align}

because $(a+b)^2\leq2(a^2+b^2)$ and $\kappa^4\leq\kappa^2\leq\kappa$ for $\kappa\leq1$

Furthermore,
\begin{align}
\varepsilon_1 &= O\big(\kappa^2\tau^2 m n d\big)
    +
    O\big(\kappa^2\tau^2 m n d(\hat\phi_{\max}^2 + \hat\psi_{\max}^2)\big)
    +
    O\big(\kappa\tau m n \sqrt{d}  \hat\phi_{\max}^2\big) \notag\\
    &\leq 
    \textcolor{teal}{O\big(\kappa^2\tau^2 m^2 n d^2\big)} +
    \textcolor{purple}{O\big(\kappa \tau^2 m n^3 d^2(\hat\phi_{\max}^2 + \hat\psi_{\max}^2)\big)} + \textcolor{orange}{O\big(\kappa\tau m n \sqrt{d}  \hat\phi_{\max}^2\big)}
    \label{eq:eps1_bigO}
\end{align}
%%%%%%%%%%%%%%%%%%%%%%%%%%%
Comment:

Tried bounding $O\big(\kappa\tau m n \sqrt{d}  \hat\phi_{\max}^2\big)$ in the term $O\big(\kappa\tau^2 m n^3 d^2  (\hat\phi_{\max}^2 + \hat\psi_{\max}^2) \big)$ as  $\hat\phi_{\max}^2 \leq \hat\phi_{\max}^2 + \hat\psi_{\max}^2$ and $\sqrt{d} \leq d^2$ (since in our case it definitely holds that $d>1$) but it does not necessarily hold that $\tau \leq \tau^2$
%%%%%%%%%%%%%%%%%%%%%%%%%%%
Therefore, 
\begin{align*}
    O\paren{\frac{mn}{\lambda_0^2}\cdot\varepsilon_2^2 + \varepsilon_1} \stackrel{\eqref{eq:eps1_bigO}, \eqref{eq:eps2_bigO}}{=} \textcolor{teal}{O\big(\kappa^2\tau^2 m^2 n d^2\big)} +
    \textcolor{purple}{O\big(\kappa \tau^2 m n^3 d^2(\hat\phi_{\max}^2 + \hat\psi_{\max}^2)\big)} + \textcolor{orange}{O\big(\kappa\tau m n \sqrt{d}  \hat\phi_{\max}^2\big)}
\end{align*}
Thus, the expected loss is bounded by:
\[
\mathbb{E}[\mathcal{L}(\mathbf{W}_K)] \le \left(1 - \frac{\eta \lambda_0}{2}\right)^K \mathcal{L}(\mathbf{W}_0) + \textcolor{teal}{O\big(\kappa^2\tau^2 m^2 n d^2\big)} +
    \textcolor{purple}{O\big(\kappa \tau^2 m n^3 d^2(\hat\phi_{\max}^2 + \hat\psi_{\max}^2)\big)} + \textcolor{orange}{O\big(\kappa\tau m n \sqrt{d}  \hat\phi_{\max}^2\big)}
\]
%%%%%%%%%%%%%%%%%%%%%%%%%%%%%%%%%%%%%%%%%%%%%%%%%%%%%%%%%%%%%
\subsubsection{Proof of Corollary~\ref{cor:exp_loss_ub}}
\begin{proof}
    By Theorem~\ref{thm:exp_loss_approx}, we have that 
    \begin{align*}
        \left|\E_{\bfC}\left[\calL_{\bfC}\paren{\bm{\theta}}\right] - \calL\paren{\bm{\theta}}\right| & \leq \underbrace{\left|\frac{1}{2}\sum_{i=1}^n\paren{\paren{\hat{f}\paren{\bm{\theta},\bfx_i} - y_i}^2 - \paren{{f}\paren{\bfW,\bfx_i} - y_i}^2}\right|}_{\mathcal{T}_1}\\
        &\quad\quad\quad + \underbrace{\frac{\kappa^2}{2m}\sum_{i=1}^n\norm{\sum_{r=1}^ma_r\paren{\bfw_r\odot\bfx_i}\bm{\Phi}_1\paren{\frac{\bfw_r^\top\bfx_i}{\kappa\norm{\bfw_r\odot\bfx_i}_2}}}_2^2}_{\mathcal{T}_2}\\
        & \quad\quad\quad + mn\paren{\kappa^2R_{\bfu}^2\psi_{\max}^2 + \paren{\kappa^2R_{\bfu}^2 + \kappa R_{\bfw}}\phi_{\max}^2}
    \end{align*}
    Now, we can bound $\mathcal{T}_1$ and $\mathcal{T}_2$ separately. For $\mathcal{T}_1$, we have
    \begin{align*}
        & \left|\frac{1}{2}\sum_{i=1}^n\paren{\paren{\hat{f}\paren{\bm{\theta},\bfx_i} - y_i}^2 - \paren{{f}\paren{\bfW,\bfx_i} - y_i}^2}\right| \\
        & \quad\quad\quad \leq \frac{1}{2}\sum_{i=1}^n\paren{\paren{\hat{f}\paren{\bm{\theta},\bfx_i} - {f}\paren{\bfW,\bfx_i}}^2 + 2\left|\paren{\hat{f}\paren{\bm{\theta},\bfx_i} - {f}\paren{\bfW,\bfx_i}}\paren{\hat{f}\paren{\bm{\theta},\bfx_i} - {f}\paren{\bfW,\bfx_i}}\right|}\\
        & \quad\quad\quad \leq\sum_{i=1}^n\paren{\hat{f}\paren{\bm{\theta},\bfx_i} - {f}\paren{\bfW,\bfx_i}}^2 + \calL\paren{\bm{\theta}}
    \end{align*}
    Here, we can bound $\hat{f}\paren{\bfW,\bfx_i} - {f}\paren{\bfW,\bfx_i}$ as
    \begin{align*}
        \left|\hat{f}\paren{\bfW,\bfx_i} - {f}\paren{\bfW,\bfx_i}\right| & \leq \frac{1}{\sqrt{m}}\sum_{r=1}^m\left|\hat{\sigma}\paren{\bfw_r,\bfx_i} - \sigma\paren{\bfw_r^\top\bfx_i}\right|\\
        & \leq \frac{1}{\sqrt{m}}\sum_{r=1}^m\left|\bfw_r^\top\bfx_i\right|\paren{\indy{\bfw_r^\top\bfx_i\geq 0} - \bm{\Phi}_1\paren{\frac{\bfw_r^\top\bfx_i}{\kappa\norm{\bfw_r\odot\bfx_i}_2}}}\\
        & \leq \frac{1}{\sqrt{m}}\sum_{r=1}^m\left|\bfw_r^\top\bfx_i\right|\exp\paren{-\frac{\paren{\bfw_r^\top\bfx_i}^2}{2\kappa^2\norm{\bfw_r\odot\bfx_i}_2}}\\
        & \leq \frac{1}{\sqrt{m}}\sum_{r=1}^m\psi\paren{\frac{\bfw_r^\top\bfx_i}{2\kappa\norm{\bfw_r^\top\bfx_i}_2}}\\
        & \leq \sqrt{m}\psi_{\max}
    \end{align*}
    where in the third inequality we used Lemma~\ref{lem:gaus_cdf_approx}. Therefore, $\mathcal{T}_1$ can be bounded as
    \[
        \mathcal{T}_1 \leq \calL\paren{\bm{\theta}} +mn\psi_{\max}
    \]
    For $\mathcal{T}_2$, we an bound it as
    \begin{align*}
        \mathcal{T}_2 \leq 2\kappa^2\sum_{i=1}^n\sum_{r=1}^m\norm{\bfw_r\odot\bfx_i} \leq 2\kappa^2mnR_{\bfu}^2
    \end{align*}
    Plugging in $\mathcal{T}_1$ and $\mathcal{T}_2$ gives the desired result.
\end{proof}

\subsubsection{Proof of Corollary~\ref{cor:exp_grad_err_ub}}
\begin{proof}
    By Theorem~\ref{thm:exp_grad_form}, we have that
    \begin{align*}
        \norm{\E_{\bfC}\left[\nabla_{\bfw_r}\calL\paren{\bm{\theta}}\right] - \nabla_{\bfw_r}\calL\paren{\bm{\theta}}}_2 & = \norm{\bfg_r + \frac{3\kappa^2}{\sqrt{m}}\sum_{r'=1}^ma_{r'}\text{Diag}\paren{\bfx_i}^2\bfw_{r'}\mathbb{I}\left\{\bfw_{r}^\top\bfx_i\geq 0;\bfw_{r'}^\top\bfx_i\geq 0\right\}}_2\\
        & \leq \norm{\bfg_r}_2 + \frac{3\kappa^2}{\sqrt{m}}\sum_{r'=1}^m\norm{a_{r'}\text{Diag}\paren{\bfx_i}^2\bfw_{r'}\mathbb{I}\left\{\bfw_{r}^\top\bfx_i\geq 0;\bfw_{r'}^\top\bfx_i\geq 0\right\}}_2\\
        & \leq \norm{\bfg_r}_2 + \frac{3\kappa^2}{\sqrt{m}}\sum_{r'=1}^m\norm{\bfx_i}_{\infty}^2\norm{\bfw_r}_2\\
        & \leq \paren{\frac{\sigma_{\max}\paren{\bfX}}{\sqrt{m}}\calL\paren{\bm{\theta}}^{\frac{1}{2}} + 6n\kappa^2B_{\bfx}^2R_{\bfw} + 5n\kappa R_{\bfu}\sqrt{d}}\phi_{\max}\\
        & \quad\quad\quad + 6n\kappa R_{\bfu}\psi_{\max}+3\kappa^2\sqrt{m}B_{\bfx}^2R_{\bfw}
    \end{align*}
    Theorem~\ref{thm:exp_grad_form} states that the norm $\norm{\bfg_r}_2$ satisfies
    \[
        \norm{\bfg_r}_2 \leq \paren{6n\kappa^2B_{\bfx}^2R_{\bfw} + 5n\kappa R_{\bfu}\sqrt{d}}\phi_{\max} + \frac{\sigma_{\max}\paren{\bfX}}{\sqrt{m}}\phi_{\max}\calL\paren{\bm{\theta}}^{\frac{1}{2}}+ 6n\kappa R_{\bfu}\psi_{\max}
    \]
    Moreover, by Assumption~\ref{asump:data}, we have that $\norm{\bfx}_{\infty}\leq B_{\bfx}$. Therefore, we obtain that
    \begin{align*}
        \norm{\E_{\bfC}\left[\nabla_{\bfw_r}\calL\paren{\bm{\theta}}\right] - \nabla_{\bfw_r}\calL\paren{\bm{\theta}}}_2 & \leq \paren{\frac{\sigma_{\max}\paren{\bfX}}{\sqrt{m}}\calL\paren{\bm{\theta}}^{\frac{1}{2}} + 6n\kappa^2B_{\bfx}^2R_{\bfw} + 5n\kappa R_{\bfu}\sqrt{d}}\phi_{\max}\\
        & \quad\quad\quad + 6n\kappa R_{\bfu}\psi_{\max}+3\kappa^2\sqrt{m}B_{\bfx}^2R_{\bfw}
    \end{align*}
\end{proof}

\subsubsection{Proof of Lemma~\ref{lem:general_smoothness}}
\begin{proof}
    By the form of $\nabla_{\bfw_r}\calL_{\bfC}\paren{\mathbf{W}}$ in (\ref{eq:mask_grad}), we have:
    \begin{align*}
        \norm{\nabla_{\bfw_r}\calL_{\bfC}\paren{\mathbf{W}}}_2
        & =\norm{\frac{a_r}{\sqrt{m}}\sum_{i=1}^n\paren{f\paren{\mathbf{W},\bfx_i\odot\bfc_i} - y_i} \paren{\bfx_i\odot\bfc_i}\underbrace{\indy{\bfw_r^\top\paren{\bfx_i\odot\bfc_i}\geq 0}}_{\le 1}}_2 \\
        & \leq \frac{\left|a_r\right|}{\sqrt{m}}\sum_{i=1}^n\left|f\paren{\mathbf{W},\bfx_i\odot\bfc_i}-y_i\right|\cdot\norm{\bfx_i\odot\bfc_i}_2\\
        & \leq \frac{\sqrt{n}}{\sqrt{m}}\left(\sum_{i=1}^n (f\paren{\mathbf{W},\bfx_i\odot\bfc_i}-y_i)^2\right)^{1/2} \cdot \|\bfc_i\|_\infty \|\bfx_i\|_2 \enspace \text{using \ref{lem:CS}} \\
        & \leq C \frac{\sqrt{n}}{\sqrt{m}} \left(2\calL_{\bfC}\paren{\mathbf{W}}\right)^{1/2} \enspace \text{using} \enspace \|\bfc_i\|_\infty \le C \enspace \text{and} \enspace \norm{\bfx_i}_2\leq 1\\
        & \leq C\sqrt{2}\frac{\sqrt{n}}{\sqrt{m}}\calL_{\bfC}\paren{\mathbf{W}}^{1/2}
    \end{align*}
    where, in the first inequality, we use the fact that the indicator function is upper-bounded by 1 and in the second inequality, we use the fact that $a_r = \pm 1$.

    and so,
    \begin{align*}
        \norm{\nabla_{\bfw_r}\calL_{\bfC}\paren{\mathbf{W}}}^2_2\leq 2C^2\cdot \frac{n}{m}\calL_{\bfC}\paren{\mathbf{W}}
    \end{align*}
\end{proof}
%%%%%%
\section{Auxiliary Results}
\subsection{Gaussian Random Variables}
\subsubsection{Conditional Expectation and Covariance}
\begin{lemma}
    \label{lem:conditional_exp}
    Let $\bfc\sim\mathcal{N}\paren{\bm{\mu},\kappa^2\bfI}$, and let $z = \bfc^\top\bfu, z' = \bfc^\top\bfv$. Then we have that
    \[
        \E_{\bfc}\left[\bfc\mid z\right] = \bm{\mu} + \frac{\bfu}{\norm{\bfu}_2^2}\paren{z - \bm{\mu}^\top\bfu}; \quad\E_{\bfc}\left[\bfc\mid z, z'\right] = \bm{\mu} + \bfs_1\paren{z - \bm{\mu}^\top\bfu} + \bfs_2\paren{z' - \bm{\mu}^\top\bfv}
    \]
    where the vectos $\bfs_1,\bfs_2$ are defined as
    \[
        \bfs_1 = \frac{\norm{\bfv}_2^2\bfu - \bfu^\top\bfv\cdot\bfv}{\norm{\bfu}_2^2\norm{\bfv}_2^2 - \paren{\bfu^\top\bfv}^2};\quad \bfs_2 = \frac{\norm{\bfu}_2^2\bfv - \bfu^\top\bfv\cdot\bfu}{\norm{\bfu}_2^2\norm{\bfv}_2^2 - \paren{\bfu^\top\bfv}^2}
    \]
\end{lemma}
\begin{proof}
    By the formula of conditional expectation, we have
    \[
        \E_{\bfc}\left[\bfc\mid z\right] = \E[\bfc] + \text{Cov}\paren{\bfc,z}\text{Var}\paren{z}^{-1}\paren{z - \E[z]}
    \]
    By the definition of $\bfc$, we have $\E[\bfc] = \bm{\mu}$. Moreover, since $z = \bfc^\top\bfu$, by Lemma~\ref{lem:inner_gaussian_rv}, we have $\E[z] = \bm{\mu}^\top\bfu$ and $\text{Var}\paren{z} = \kappa^2\norm{\bfu}_2^2$. Therefore, the covariance between $\bfc$ and $z$ is given by
    \[
        \text{Cov}\paren{\bfc,z} = \E[\paren{\bfc - \bm{\mu}}\paren{z - \inner{\bm{\mu}}{\bfv}}] = \E\left[z\bfc\right] - \bm{\mu}\inner{\bm{\mu}}{\bfv}
    \]
    where
    \[
        \E[z\bfc]_j = \sum_{j'=1}^d\E[c_{j}c_{j'}u_{j'}] = \sum_{j'=1}^d\paren{\mu_j\mu_{j'} + \mathbb{I}\{j = j'\}\kappa^2}u_{j'} = \mu_j\bm{\mu}^\top\bfu + \kappa^2u_j
    \]
    This gives
    \[
        \E[z\bfc] = \bm{\mu}^\top\bfu\cdot \bm{\mu} + \kappa^2\bfu
    \]
    Therefore
    \[
       \E_{\bfc}\left[\bfc\mid z\right] = \bm{\mu} + \frac{\bfu}{\norm{\bfu}_2^2}\paren{z - \bm{\mu}^\top\bfu}
    \]
    Similarly, since $z' = \bfc^\top\bfv$. Then
    \[
        \E_{\bfc}\left[\bfc\mid z,z'\right] = \E[\bfc] + \text{Cov}\paren{\bfc,[z, z']}\text{Cov}\paren{z, z'}^{-1}\paren{[z, z']^\top - \E[[z, z']]^\top}
    \]
    where
    \[
        \text{Cov}\paren{\bfc,[z, z']} = \kappa^2\begin{bmatrix}
            \bfu\\
            \bfv
        \end{bmatrix}; \quad \text{Cov}\paren{z, z'} = \kappa^2\begin{bmatrix}
            \norm{\bfu}_2^2 & \bfu^\top\bfv\\
            \bfu^\top\bfv& \norm{\bfv}_2^2 
        \end{bmatrix}
    \]
    This gives
    \begin{align*}
        \E_{\bfc}\left[\bfc\mid z,z'\right] & = \bm{\mu} + \begin{bmatrix}
            \bfu\\
            \bfv
        \end{bmatrix}\begin{bmatrix}
            \norm{\bfu}_2^2 & \bfu^\top\bfv\\
            \bfu^\top\bfv& \norm{\bfv}_2^2 
        \end{bmatrix}^{-1}\begin{bmatrix}
            z - \bm{\mu}^\top\bfu\\
            z'-\bm{\mu}^\top\bfv
        \end{bmatrix}\\
        & = \bm{\mu} + \frac{1}{\norm{\bfu}_2^2\norm{\bfv}_2^2 - \paren{\bfu^\top\bfv}^2}\begin{bmatrix}
            \bfu\\
            \bfv
        \end{bmatrix}\begin{bmatrix}
            \norm{\bfv}_2^2 & -\bfu^\top\bfv\\
            -\bfu^\top\bfv& \norm{\bfu}_2^2 
        \end{bmatrix}^{-1}\begin{bmatrix}
            z - \bm{\mu}^\top\bfu\\
            z'-\bm{\mu}^\top\bfv
        \end{bmatrix}\\
        & = \bm{\mu} + \frac{\paren{\norm{\bfv}_2^2\bfu - \bfu^\top\bfv\cdot \bfv}\paren{z - \bm{\mu}^\top\bfu} + \paren{\norm{\bfu}_2^2\bfv - \bfu^\top\bfv\cdot \bfu}\paren{z' -\bm{\mu}^\top\bfv}}{\norm{\bfu}_2^2\norm{\bfv}_2^2 - \paren{\bfu^\top\bfv}^2}\\
        & = \bm{\mu} + \bfs_1\paren{z - \bm{\mu}^\top\bfu} + \bfs_2\paren{z'- \bm{\mu}^\top\bfv}
    \end{align*}where the vectos $\bfs_1,\bfs_2$ are defined as
    \[
        \bfs_1 = \frac{\norm{\bfv}_2^2\bfu - \bfu^\top\bfv\cdot\bfv}{\norm{\bfu}_2^2\norm{\bfv}_2^2 - \paren{\bfu^\top\bfv}^2};\quad \bfs_2 = \frac{\norm{\bfu}_2^2\bfv - \bfu^\top\bfv\cdot\bfu}{\norm{\bfu}_2^2\norm{\bfv}_2^2 - \paren{\bfu^\top\bfv}^2}
    \]
\end{proof}

\begin{lemma}
    \label{lem:conditional_cov}
    Let $\bfc\sim\mathcal{N}\paren{\bm{\mu},\kappa^2\bfI}$, and let $z = \bfc^\top\bfu, z' = \bfc^\top\bfv$. Then we have
    \[
        \text{Cov}\paren{\bfc\mid z, z'} = \kappa^2\bfI - \frac{\kappa^2\paren{\bfu\bfv^\top - \bfv\bfu^\top}^2}{\norm{\bfu}_2^2\norm{\bfv}_2^2 - \paren{\bfu^\top\bfv}^2}
    \]
\end{lemma}
\begin{proof}
    The conditional covariance of Gaussian random variables is given by
    \[
        \text{Cov}\paren{\bfc\mid z, z'} = \text{Cov}\paren{\bfc} -  \text{Cov}\paren{\bfc,[z, z']}\text{Cov}\paren{z, z'}^{-1}\text{Cov}\paren{\bfc,[z, z']}^\top
    \]
    Recall that in the previos lemma we have computed
    \begin{align*}
        \text{Cov}\paren{\bfc,[z, z']} = \kappa^2\begin{bmatrix}
                \bfu\\
                \bfv
            \end{bmatrix}; \quad \text{Cov}\paren{z, z'} = \kappa^2\begin{bmatrix}
                \norm{\bfu}_2^2 & \bfu^\top\bfv\\
                \bfu^\top\bfv& \norm{\bfv}_2^2 
            \end{bmatrix}
    \end{align*}
    Thus, we have
    \begin{align*}
        \text{Cov}\paren{\bfc\mid z, z'} & = \kappa^2\bfI - \kappa^2\begin{bmatrix}
                \bfu\\
                \bfv
            \end{bmatrix}\begin{bmatrix}
                \norm{\bfu}_2^2 & \bfu^\top\bfv\\
                \bfu^\top\bfv& \norm{\bfv}_2^2 
            \end{bmatrix}^{-1}\begin{bmatrix}
                \bfu & 
                \bfv
            \end{bmatrix}\\
            & = \kappa^2\bfI - \frac{\kappa^2}{\norm{\bfu}_2^2\norm{\bfv}_2^2- \paren{\bfu^\top\bfv}^2}\begin{bmatrix}
                \bfu\\
                \bfv
            \end{bmatrix}\begin{bmatrix}
                \norm{\bfv}_2^2 & -\bfu^\top\bfv\\
                -\bfu^\top\bfv& \norm{\bfu}_2^2 
            \end{bmatrix}\begin{bmatrix}
                \bfu & 
                \bfv
            \end{bmatrix}\\
            & = \kappa^2\bfI - \frac{\kappa^2\paren{\bfu\bfv^\top - \bfv\bfu^\top}^2}{\norm{\bfu}_2^2\norm{\bfv}_2^2 - \paren{\bfu^\top\bfv}^2}
    \end{align*}
\end{proof}
%%%%%%%
\subsubsection{Approximation of CDF}
In approximation of the Gaussian CDF, we use the following property
\begin{lemma}[\cite{chu1955Bounds}]
    Let $x \geq 0$ be given, then the following inequality holds
    \[
        \frac{1}{2}\paren{1 - \exp\paren{-\frac{x^2}{2}}}^{\frac{1}{2}}\leq \frac{1}{\sqrt{2\pi}}\int_0^x\exp\paren{-\frac{t^2}{2}}dt \leq \frac{1}{2}\paren{1 - \exp\paren{-\frac{2x^2}{\pi}}}^{\frac{1}{2}}
    \]
\end{lemma}
To start, we analyze the CDF of a single variable Gaussian random variable
\begin{lemma}
    \label{lem:gaus_cdf_approx}
    Let $\bm{\Phi}_1$ be the CDF of a standard Gaussian random variable. Then we have
    \[
        \left|\bm{\Phi}_1\paren{\alpha} - \mathbb{I}\left\{\alpha\geq 0\right\}\right|\leq \exp\paren{-\frac{\alpha^2}{2}}
    \]
\end{lemma}
\begin{proof}
    Let $\nu(x) = \frac{1}{\sqrt{2\pi}}\int_0^x\exp\paren{-\frac{t^2}{2}}dt$ for $x \geq 0$. We study the case of $\alpha\geq 0$ and $\alpha < 0$ separately. For $\alpha \geq 0$, we have
    \begin{align*}
        \bm{\Phi}_1\paren{\alpha} & = \frac{1}{\sqrt{2\pi}}\int_{-\infty}^\alpha\exp\paren{-\frac{t^2}{2}}dt\\
        & = \frac{1}{2} + \nu(\alpha)\\
        & \geq \frac{1}{2} + \frac{1}{2}\paren{1 - \exp\paren{-\frac{\alpha^2}{2}}}^{\frac{1}{2}}\\
        & \geq 1 - \exp\paren{-\frac{\alpha^2}{2}}\\
        & = \mathbb{I}\left\{\alpha \geq 0\right\} - \exp\paren{-\frac{\alpha^2}{2}}
    \end{align*}
    Since $\bm{\Phi}_1\leq 1 = \mathbb{I}\left\{\alpha \geq 0\right\}$ when $\alpha \geq 0$, we must have that, for $\alpha \geq 0$, $\left|\bm{\Phi}_1\paren{\alpha} - \mathbb{I}\left\{\alpha \geq 0\right\}\right|\leq \exp\paren{-\frac{\alpha^2}{2}}$. Similarly, for $\alpha < 0$, we have
    \begin{align*}
        \bm{\Phi}_1\paren{\alpha} & = \frac{1}{\sqrt{2\pi}}\int_{-\infty}^\alpha\exp\paren{-\frac{t^2}{2}}dt\\
        & = \frac{1}{2} - \nu(\alpha)\\
        & \leq \frac{1}{2} - \frac{1}{2}\paren{1 - \exp\paren{-\frac{\alpha^2}{2}}}^{\frac{1}{2}}\\
        & \leq \exp\paren{-\frac{\alpha^2}{2}}\\
        & = \mathbb{I}\left\{\alpha \geq 0\right\} + \exp\paren{-\frac{\alpha^2}{2}}
    \end{align*}
    Since $\bm{\Phi}_1 \geq 0 = \mathbb{I}\left\{\alpha \geq 0\right\}$ when $\alpha < 0$, we must have that, for $\alpha < 0$, $\left|\bm{\Phi}_1\paren{\alpha} - \mathbb{I}\left\{\alpha \geq 0\right\}\right|\leq \exp\paren{-\frac{\alpha^2}{2}}$.
\end{proof}
\begin{lemma}
    \label{lem:gaus_cdf2_approx}
    Let $\bm{\Phi}_2\paren{\alpha_1, \alpha_2, \rho}$ be the joint CDF of two standard Gaussian random variables with covariance $\rho$ at $\alpha_1, \alpha_2$. Then we have
    \[
        \left|\bm{\Phi}_2\paren{\alpha_1,\alpha_2,\rho} - \mathbb{I}\left\{\alpha_1\geq 0;\alpha_2\geq 0\right\}\right| \leq 2\exp\paren{-\frac{\min\left\{\alpha_1,\alpha_2\right\}^2}{2}}
    \]
\end{lemma}
\begin{proof}
    Since $z_1, z_2$ are standard Gaussian random variables with covariance $\rho$, we have $z_1\mid z_2 = \zeta\sim\mathcal{N}\paren{\rho\zeta, 1 - \rho^2}$. According to the definition of CDF,
    \begin{align*}
        \bm{\Phi}_2\paren{\alpha_1,\alpha_2,\rho} & = \int_{-\infty}^{\alpha_2}\int_{-\infty}^{\alpha_1}f_{z_1,z_2}\paren{\zeta_1,\zeta_2,\rho}d\zeta_1d\zeta_2\\
        & = \int_{-\infty}^{\alpha_2}\int_{-\infty}^{\alpha_1}f_{z_1\mid z_2=\zeta_2}\paren{\zeta_1}d\zeta_1f_{z_2}\paren{\zeta_2}d\zeta_2
    \end{align*}
    Focusing on the inner integral, we substitute $\zeta' = \frac{\zeta_1 - \rho\zeta_2}{\sqrt{1 - \rho^2}}$. Then we have $d\zeta_1 = \sqrt{1 - \rho^2}d\zeta'$. Thus
    \begin{align*}
        \int_{-\infty}^{\alpha_1}f_{z_1\mid z_2=\zeta_2}\paren{\zeta_1}d\zeta_1 &= \frac{1}{\sqrt{2\pi(1 - \rho^2)}}\int_{-\infty}^{\alpha_1}\exp\paren{-\frac{\paren{\zeta_1 - \rho\zeta_2}^2}{2(1-\rho^2)}}d\zeta_1\\
        & = \frac{1}{\sqrt{2\pi}}\int_{-\infty}^{\frac{\alpha_1 - \rho\zeta_2}{\sqrt{1 - \rho^2}}}\exp\paren{-\frac{\zeta'^2}{2}}d\zeta'\\
        & = \Phi_1\paren{\frac{\alpha_1 - \rho\zeta_2}{\sqrt{1 - \rho^2}}}\\
        & = \mathbb{I}\left\{\frac{\alpha_1 - \rho\zeta_2}{\sqrt{1 - \rho^2}} \geq 0\right\} + \varepsilon\paren{\frac{\alpha_1 - \rho\zeta_2}{\sqrt{1 - \rho^2}}}\\
        & = \mathbb{I}\left\{\alpha_1 \geq \rho\zeta_2\right\} + \varepsilon\paren{\frac{\alpha_1 - \rho\zeta_2}{\sqrt{1 - \rho^2}}}
    \end{align*}
    Where $\left|\varepsilon\paren{\frac{\alpha_1 - \rho\zeta_2}{\sqrt{1 - \rho^2}}}\right|\leq \exp\paren{-\frac{\paren{\alpha_1 - \rho\zeta_2}^2}{2(1 - \rho^2)}}$. Therefore
    \begin{align*}
        \bm{\Phi}_2\paren{\alpha_1,\alpha_2,\rho} & = \int_{-\infty}^{\alpha_2}\paren{\mathbb{I}\left\{\alpha_1 \geq \rho\zeta_2\right\} + \varepsilon\paren{\frac{\alpha_1 - \rho\zeta_2}{\sqrt{1 - \rho^2}}}}f_{z_2}\paren{\zeta_2}d\zeta_2\\
        & = \underbrace{\int_{-\infty}^{\alpha_2}\mathbb{I}\left\{\alpha_1 \geq \rho\zeta_2\right\}f_{z_2}\paren{\zeta_2}d\zeta_2}_{\mathcal{I}_1} +  \underbrace{\int_{-\infty}^{\alpha_2}\varepsilon\paren{\frac{\alpha_1 - \rho\zeta_2}{\sqrt{1 - \rho^2}}}f_{z_2}\paren{\zeta_2}d\zeta_2}_{\mathcal{I}_2}
    \end{align*}
    For $\mathcal{I}_1$, we have
    \begin{align*}
        \mathcal{I}_1 & = \int_{-\infty}^{\alpha_2}\mathbb{I}\left\{\alpha_1 \geq \rho\zeta_2\right\}f_{z_2}\paren{\zeta_2}d\zeta_2\\
        & = \int_{-\infty}^{\min\left\{\frac{\alpha_1}{\rho},\alpha_2\right\}}f_{z_2}\paren{\zeta_2}d\zeta_2\\
        & = \Phi_1\paren{\min\left\{\frac{\alpha_1}{\rho},\alpha_2\right\}}\\
        & = \mathbb{I}\left\{\min\left\{\frac{\alpha_1}{\rho},\alpha_2\right\} \geq 0\right\} + \varepsilon\paren{\min\left\{\frac{\alpha_1}{\rho},\alpha_2\right\}}
    \end{align*}
    Notice that, when $\alpha_1 \geq 0$ and $\alpha_2\geq 0$, we must have $\min\left\{\frac{\alpha_1}{\rho},\alpha_2\right\}\geq 0$. Conversely, when $\min\left\{\frac{\alpha_1}{\rho},\alpha_2\right\} \geq 0$, we must have that $\alpha_1\geq 0$ and $\alpha_2\geq 0$. Therefore $\mathbb{I}\left\{\min\left\{\frac{\alpha_1}{\rho},\alpha_2\right\} \geq 0\right\} = \mathbb{I}\left\{\alpha_1\geq 0;\alpha_2\geq 0\right\}$. Thus, we have
    \[
        \bm{\Phi}_2\paren{\alpha_1,\alpha_2,\rho} = \mathbb{I}\left\{\alpha_1\geq 0;\alpha_2\geq 0\right\} + \varepsilon\paren{\min\left\{\frac{\alpha_1}{\rho},\alpha_2\right\}} + \mathcal{I}_2
    \]
    For $\mathcal{I}_2$, we have
    \begin{align*}
        \left|\mathcal{I}_2\right| & \leq \int_{-\infty}^{\alpha_2}\left|\varepsilon\paren{\frac{\alpha_1 - \rho\zeta_2}{\sqrt{1 - \rho^2}}}\right|f_{z_2}\paren{\zeta_2}d\zeta_2\\
        & \leq \int_{-\infty}^{\alpha_2}\exp\paren{-\frac{\paren{\alpha_1 - \rho\zeta_2}^2}{2(1 - \rho^2)}}f_{z_2}\paren{\zeta_2}d\zeta_2\\
        & = \frac{1}{\sqrt{2\pi}}\int_{-\infty}^{\alpha_2}\exp\paren{-\frac{\paren{\alpha_1 - \rho\zeta_2}^2}{2(1 - \rho^2)} - \frac{\zeta_2^2}{2}}d\zeta_2\\
        & = \frac{1}{\sqrt{2\pi}}\exp\paren{-\frac{\alpha_1^2}{2}}\int_{-\infty}^{\alpha_2}\exp\paren{-\frac{\paren{\zeta_2 - \rho\alpha_1}^2}{2\paren{1-\rho^2}}}d\zeta_2\\
        & \leq \exp\paren{-\frac{\alpha_1^2}{2}}
    \end{align*}
    Moreover, since $\left|\varepsilon\paren{\min\left\{\frac{\alpha_1}{\rho},\alpha_2\right\}}\right|\leq \exp\paren{-\frac{1}{2}\min\left\{\frac{\alpha_1}{\rho},\alpha_2\right\}^2}$, we have that
    \begin{align*}
        \left|\bm{\Phi}_2\paren{\alpha_1,\alpha_2,\rho} - \mathbb{I}\left\{\alpha_1\geq 0;\alpha_2\geq 0\right\}\right| & \leq \exp\paren{-\frac{1}{2}\min\left\{\frac{\alpha_1}{\rho},\alpha_2\right\}^2} + \exp\paren{-\frac{\alpha_1^2}{2}}\\
        & \leq \exp\paren{-\frac{\min\left\{\alpha_1,\alpha_2\right\}^2}{2}} + \exp\paren{-\frac{\alpha_1^2}{2}}
    \end{align*}
    Exchanging $\alpha_1$ and $\alpha_2$, we have
    \begin{align*}
        \left|\bm{\Phi}_2\paren{\alpha_1,\alpha_2,\rho} - \mathbb{I}\left\{\alpha_1\geq 0;\alpha_2\geq 0\right\}\right| & \leq \exp\paren{-\frac{1}{2}\min\left\{\frac{\alpha_1}{\rho},\alpha_2\right\}^2} + \exp\paren{-\frac{\alpha_1^2}{2}}\\
        & \leq \exp\paren{-\frac{\min\left\{\alpha_1,\alpha_2\right\}^2}{2}} + \exp\paren{-\frac{\alpha_2^2}{2}}
    \end{align*}
    Thus, we have
    \[
        \left|\bm{\Phi}_2\paren{\alpha_1,\alpha_2,\rho} - \mathbb{I}\left\{\alpha_1\geq 0;\alpha_2\geq 0\right\}\right| \leq 2\exp\paren{-\frac{\min\left\{\alpha_1,\alpha_2\right\}^2}{2}}
    \]
\end{proof}
%%%%%%%%
\subsubsection{Uni-variate Coupled Expectation}
\begin{lemma}
    \label{lem:first_ord_term_exp}
    let $z_1, z_2\sim\mathcal{N}\paren{0, 1}$ with $\text{Cov}\paren{z_1, z_2} = \rho$. Let $a, b \in\R$. Define
    \[
        T_1 = \exp\paren{-\frac{a^2}{2}}\bm{\Phi}_1\paren{\frac{\rho a - b}{\sqrt{1-\rho^2}}};\quad T_2 = \exp\paren{-\frac{b^2}{2}}\bm{\Phi}_1\paren{\frac{\rho b - a}{\sqrt{1-\rho^2}}}
    \]
    Then we have that
    \[
        \E\left[z_1\indy{z_1\geq a; z_2\geq b}\right] =\frac{1}{\sqrt{2\pi}}\paren{T_1 + \rho T_2};\quad \E\left[z_2\indy{z_1\geq a; z_2\geq b}\right] = \frac{1}{\sqrt{2\pi}}\paren{T_2 + \rho T_1}
    \]
\end{lemma}
\begin{proof}
    Writing the expectation in integral form, we have that
    \[
        \E\left[z_1\indy{z_1\geq a; z_2\geq b}\right] = \int_b^\infty\int_a^\infty z_1 f\paren{z_1, z_2}dz_1dz_2 = \int_b^\infty\paren{\int_a^\infty z_1f\paren{z_1\mid z_2}dz_1}f\paren{z_2}dz_2
    \]
    Since $z_1,z_2\sim\mathcal{N}\paren{0, 1}$ with covariance $\rho$, we have that $z_1\mid z_2\sim\mathcal{N}\paren{\rho z_2,1-\rho^2}$. Therefore, let $\rho' = 1-\rho^2$, we have
    \[
        f\paren{z_1\mid z_2} = \frac{1}{\sqrt{2\pi\rho'}}\exp\paren{-\frac{\paren{z_1-\rho z_2}^2}{2\rho'}}
    \]
    Thus, by Lemma~\ref{lem:first_order_gaus_int}, we have that
    \begin{align*}
        \int_a^\infty z_1f\paren{z_1\mid z_2}dz_1 & = \frac{1}{\sqrt{2\pi\rho'}}\int_a^\infty z_1\exp\paren{-\frac{\paren{z_1-\rho z_2}^2}{2\rho'}}dz_1\\
        & = \frac{1}{\sqrt{2\pi\rho'}}\paren{\rho'\exp\paren{-\frac{\paren{a - \rho z_2}^2}{2\rho'}}+\rho z_2\sqrt{2\pi \rho'}\bm{\Phi}_1\paren{\frac{\rho z_2-a}{\sqrt{\rho'}}}}\\
        & = \sqrt{\frac{\rho'}{2\pi}}\exp\paren{-\frac{\paren{a - \rho z_2}^2}{2\rho'}} + \rho z_2\bm{\Phi}_1\paren{\frac{\rho z_2 - a}{\sqrt{\rho'}}}
    \end{align*}
    where we set $\kappa=\sqrt{\rho'}$ and $\mu = \rho z_2$ in Lemma~\ref{lem:first_order_gaus_int}. Plugging into the original integral gives
    \begin{align*}
        \E\left[z_1\indy{z_1\geq a; z_2\geq b}\right] & = \int_b^\infty\paren{\sqrt{\frac{\rho'}{2\pi}}\exp\paren{-\frac{\paren{a - \rho z_2}^2}{2\rho'}} + \rho z_2\bm{\Phi}_1\paren{\frac{\rho z_2 - a}{\sqrt{\rho'}}}}f\paren{z_2}dz_2\\
        & = \frac{\sqrt{\rho'}}{2\pi}\int_b^\infty\exp\paren{-\frac{\paren{a - \rho z_2}^2}{2\rho'} - \frac{z_2^2}{2}}dz_2 + \rho\int_b^\infty z_2\bm{\Phi}_1\paren{\frac{\rho z_2 - a}{\sqrt{\rho'}}}f\paren{z_2}dz_2
    \end{align*}
    Notice that
    \[
        \exp\paren{-\frac{\paren{a - \rho z_2}^2}{2\rho'} - \frac{z_2^2}{2}} = \exp\paren{-\frac{a^2 - 2\rho az_2 + z_2^2}{2\rho'}} = \exp\paren{-\frac{\paren{z_2 - \rho a}^2}{2\rho'}}\cdot\exp\paren{-\frac{a^2}{2}}
    \]
    From a previous result, we have an identity for the conditional probability, which expresses the CDF term as an integral:
\begin{equation}
    \label{eq:cond_prob}
    \bm{\Phi}_1\left(\frac{\rho z_2 - a}{\sqrt{1-\rho^2}}\right) = \int_a^\infty f(z_1 \mid z_2) dz_1
\end{equation}
and so:
    \begin{align*}
    \rho \int_b^\infty z_2 \bm{\Phi}_1\left(\frac{\rho z_2 - a}{\sqrt{1-\rho^2}}\right) f(z_2) dz_2= 
    \rho \int_b^\infty z_2 \int_a^\infty f(z_1 \mid z_2) dz_1 f(z_2) dz_2
    \end{align*}
    Moreover, by applying Lemma~\ref{lem:cdf_forms}, we have
    \begin{align*}
        \E\left[z_1\indy{z_1\geq a; z_2\geq b}\right] & = \frac{\sqrt{\rho'}}{2\pi}\exp\paren{-\frac{a^2}{2}}\int_b^\infty\exp\paren{-\frac{\paren{z_2 - \rho a}^2}{2\rho'}}dz_2 + \rho \int_b^\infty z_2 \int_a^\infty f(z_1 \mid z_2) dz_1 f(z_2) dz_2\\
        & = \frac{\rho'}{\sqrt{2\pi}}\exp\paren{-\frac{a^2}{2}}\int_b^\infty f\paren{z_2\mid z_1 = a}dz_2 + \rho \int_b^\infty \int_a^\infty z_2 \underbrace{f(z_1 \mid z_2) f(z_2)}_{f(z_1,z_2)} dz_1  dz_2 \\ 
        & = \frac{\rho'}{\sqrt{2\pi}}\exp\paren{-\frac{a^2}{2}}\bm{\Phi}_1\paren{\frac{\rho a - b}{\sqrt{1-\rho^2}}} + \rho\E\left[z_2\indy{z_1\geq a; z_2\geq b}\right]
    \end{align*}
    Therefore, we can conclude that
    \begin{equation}
        \label{eq:first_order_exp_t1}
        \E\left[z_1\indy{z_1\geq a; z_2\geq b}\right] - \rho\E\left[z_2\indy{z_1\geq a; z_2\geq b}\right] = \frac{\rho'}{\sqrt{2\pi}}\exp\paren{-\frac{a^2}{2}}\bm{\Phi}_1\paren{\frac{\rho a - b}{\sqrt{1-\rho^2}}} = \frac{\rho'}{\sqrt{2\pi}}T_1
    \end{equation}
    Switching $z_1, z_2$ and $a, b$ gives
    \begin{equation}
        \label{eq:first_order_exp_t2}
        \E\left[z_2\indy{z_1\geq a; z_2\geq b}\right] - \rho\E\left[z_1\indy{z_1\geq a; z_2\geq b}\right] = \frac{\rho'}{\sqrt{2\pi}}\exp\paren{-\frac{b^2}{2}}\bm{\Phi}_1\paren{\frac{\rho b - a}{\sqrt{1-\rho^2}}} = \frac{\rho'}{\sqrt{2\pi}}T_2
    \end{equation}
    Solving for $\E\left[z_1\indy{z_1\geq a; z_2\geq b}\right]$ and $\E\left[z_2\indy{z_1\geq a; z_2\geq b}\right]$ from (\ref{eq:first_order_exp_t1}) and (\ref{eq:first_order_exp_t2}) gives
    \[
        \E\left[z_1\indy{z_1\geq a; z_2\geq b}\right] =\frac{1}{\sqrt{2\pi}}\paren{T_1 + \rho T_2};\quad 
        \E\left[z_2\indy{z_1\geq a; z_2\geq b}\right] = \frac{1}{\sqrt{2\pi}}\paren{T_2 + \rho T_1}
    \]
\end{proof}

\begin{lemma}
    \label{lem:sq_term_exp}
    Let $z_1, z_2\sim\mathcal{N}\paren{0, 1}$ with $\text{Cov}\paren{z_1, z_2} = \rho$. Let $a, b \in\R$. Define
    \[
        T_1 = \exp\paren{-\frac{a^2}{2}}\bm{\Phi}_1\paren{\frac{\rho a - b}{\sqrt{1-\rho^2}}};\quad T_2 = \exp\paren{-\frac{b^2}{2}}\bm{\Phi}_1\paren{\frac{\rho b - a}{\sqrt{1-\rho^2}}}
    \]
    Then we have that
    \begin{align*}
        \E\left[z_1^2\indy{z_1\geq a; z_2\geq b}\right] & = \frac{\rho\sqrt{1-\rho^2}}{2\pi}\exp\paren{-\frac{a^2 - 2\rho ab + b^2}{2\paren{1-\rho^2}}} + \bm{\Phi}_2\paren{-a,-b,\rho} + \frac{1}{\sqrt{2\pi}}\paren{aT_1 + \rho^2 bT_2}\\
        \E\left[z_2^2\indy{z_1\geq a; z_2\geq b}\right] & = \frac{\rho\sqrt{1-\rho^2}}{2\pi}\exp\paren{-\frac{a^2 - 2\rho ab + b^2}{2\paren{1-\rho^2}}} + \bm{\Phi}_2\paren{-a,-b,\rho} + \frac{1}{\sqrt{2\pi}}\paren{bT_2 + \rho^2 aT_1}
    \end{align*}
\end{lemma}
\begin{proof}
    Again, we write the expectation in the integral form to get that
    \[
        \E\left[z_1^2\indy{z_1\geq a; z_2\geq b}\right] = \int_b^\infty\int_a^\infty z_1^2f\paren{z_1, z_2}dz_1dz_2 = \int_b^\infty\paren{\int_a^\infty z_1^2f\paren{z_1\mid z_2}dz_1}f\paren{z_2}dz_2
    \]
    Since $z_1, z_2\sim\mathcal{N}\paren{0, 1}$ with $\text{Cov}\paren{z_1, z_2} = \rho$, we have that $z_1\mid z_2 \sim\mathcal{N}\paren{\rho z_2, 1-\rho^2}$. Define $\rho' = 1-\rho^2$. By Lemma~\ref{lem:second_order_gaus_int}, we have that
    \begin{align*}
        \int_a^\infty z_1^2f\paren{z_1\mid z_2}dz_1 & = \frac{1}{\sqrt{2\pi\rho'}}\int_a^\infty z_1^2\exp\paren{-\frac{\paren{z_1-\rho z_2}^2}{2\rho'}}dz_1\\
        & = \frac{1}{\sqrt{2\pi\rho'}}\paren{\rho'\paren{a+\rho z_2}\exp\paren{-\frac{\paren{a-\rho z_2}^2}{2\rho'}} + \sqrt{2\pi\rho'}\paren{\rho'+\rho^2z_2^2}\bm{\Phi}_1\paren{\frac{\rho z_2-a}{\sqrt{\rho'}}}}\\
        & = \rho^2z_2^2\bm{\Phi}_1\paren{\frac{\rho z_2-a}{\sqrt{\rho'}}} + \rho z_2\sqrt{\frac{\rho'}{2\pi}}\exp\paren{-\frac{\paren{a-\rho z_2}^2}{2\rho'}}\\
        & \quad\quad\quad + a\sqrt{\frac{\rho'}{2\pi}}\exp\paren{-\frac{\paren{a-\rho z_2}^2}{2\rho'}} + \rho'\bm{\Phi}_1\paren{\frac{\rho z_2-a}{\sqrt{\rho'}}}
    \end{align*}
    where we set $\kappa=\sqrt{\rho'}$ and $\mu = \rho z_2$ in Lemma~\ref{lem:second_order_gaus_int}. Therefore, we have
    \begin{align*}
        \E\left[z_1^2\indy{z_1\geq a; z_2\geq b}\right] & = \rho^2\int_b^\infty z_2^2\bm{\Phi}_1\paren{\frac{\rho z_2-a}{\sqrt{\rho'}}}f\paren{z_2}dz_2\\
        &\quad\quad\quad + \rho\sqrt{\frac{\rho'}{2\pi}}\int_b^\infty z_2\exp\paren{-\frac{\paren{a-\rho z_2}^2}{2\rho'}}f\paren{z_2}dz_2\\
        & \quad\quad\quad +a\sqrt{\frac{\rho'}{2\pi}}\int_b^\infty\exp\paren{-\frac{\paren{a-\rho z_2}^2}{2\rho'}}f\paren{z_2}dz_2\\
        &\quad\quad\quad + \rho'\int_b^\infty\bm{\Phi}_1\paren{\frac{\rho z_2-a}{\sqrt{\rho'}}}f\paren{z_2}dz_2
    \end{align*}
    To start, by Lemma~\ref{lem:cdf_forms}, we have
    \[
        \bm{\Phi}_1\paren{\frac{\rho z_2-a}{\sqrt{\rho'}}}f\paren{z_2} = \int_a^\infty f\paren{z_1\mid z_2}f\paren{z_2}dz_1 = \int_a^\infty f\paren{z_1,z_2}dz_1
    \]
    Therefore, for the first term, we have
    \[
        \int_b^\infty z_2^2\bm{\Phi}_1\paren{\frac{\rho z_2-a}{\sqrt{\rho'}}}f\paren{z_2}dz_2 = \int_b^\infty\int_a^\infty z_2^2f\paren{z_1,z_2}dz_2 = \E\left[z_2^2\indy{z_1\geq a;z_2\geq b}\right]
    \]
    For the last term, we have
    \[
        \int_b^\infty\bm{\Phi}_1\paren{\frac{\rho z_2-a}{\sqrt{\rho'}}}f\paren{z_2}dz_2 = \int_b^\infty\int_a^\infty f\paren{z_1,z_2}dz_1dz_2 = \bm{\Phi}_2\paren{-a, -b, \rho}
    \]
    Next, we notice that
    \[
        \exp\paren{-\frac{\paren{a-\rho z_2}^2}{2\rho'}}f\paren{z_2} = \frac{1}{\sqrt{2\pi}}\exp\paren{-\frac{a^2 - 2\rho az_2 + z_2^2}{2\rho'}} = \frac{1}{\sqrt{2\pi}}\exp\paren{-\frac{a^2}{2}}\exp\paren{-\frac{\paren{z_2 - \rho a}^2}{2\rho'}}
    \]
    Therefore, for the second term, we apply Lemma~\ref{lem:first_order_gaus_int} to get that
    \begin{align*}
        \int_b^\infty z_2\exp\paren{-\frac{\paren{a-\rho z_2}^2}{2\rho'}}f\paren{z_2}dz_2 & = \frac{1}{\sqrt{2\pi}}\exp\paren{-\frac{a^2}{2}}\int_b^\infty z_2\exp\paren{-\frac{\paren{z_2 - \rho a}^2}{2\rho'}}dz_2\\
        & = \frac{1}{\sqrt{2\pi}}\exp\paren{-\frac{a^2}{2}}\paren{\rho'\exp\paren{-\frac{\paren{\rho a - b}^2}{2\rho'}} + \sqrt{2\pi\rho'}\rho a\bm{\Phi}_1\paren{\frac{\rho a- b}{\sqrt{\rho'}}}}\\
        & = \frac{\rho'}{\sqrt{2\pi}}\exp\paren{-\frac{a^2 - 2\rho ab + b^2}{2\rho'}}+ \rho a\sqrt{\rho'}\exp\paren{-\frac{a^2}{2}}\bm{\Phi}_1\paren{\frac{\rho a- b}{\sqrt{\rho'}}}
    \end{align*}
    where we set $\kappa=\sqrt{\rho'}$ and $\mu =\rho a$ in Lemma~\ref{lem:first_order_gaus_int}. For the third term, we have
    \begin{align*}
        \int_b^\infty\exp\paren{-\frac{\paren{a-\rho z_2}^2}{2\rho'}}f\paren{z_2}dz_2 & = \frac{1}{\sqrt{2\pi}}\exp\paren{-\frac{a^2}{2}}\int_b^\infty\exp\paren{-\frac{\paren{z_2 - \rho a}^2}{2\rho'}}dz_2\\
        & = \sqrt{\rho'}\exp\paren{-\frac{a^2}{2}}\int_b^\infty f\paren{z_2\mid z_1 = a}dz_2\\
        & = \sqrt{\rho'}\exp\paren{-\frac{a^2}{2}}\bm{\Phi}_1\paren{\frac{\rho a - b}{\sqrt{\rho'}}}
    \end{align*}
    Combining all four terms gives
    \begin{align*}
        \E\left[z_1^2\indy{z_1\geq a; z_2\geq b}\right] & = \rho\sqrt{\frac{\rho'}{2\pi}}\paren{\frac{\rho'}{\sqrt{2\pi}}\exp\paren{-\frac{a^2 - 2\rho ab + b^2}{2\rho'}}+ \rho a\sqrt{\rho'}\exp\paren{-\frac{a^2}{2}}\bm{\Phi}_1\paren{\frac{\rho a- b}{\sqrt{\rho'}}}}\\
        &\quad\quad\quad + a\sqrt{\frac{\rho'}{2\pi}}\cdot \sqrt{\rho'}\exp\paren{-\frac{a^2}{2}}\bm{\Phi}_1\paren{\frac{\rho a - b}{\sqrt{\rho'}}}\\
        & \quad\quad\quad + \rho^2\E\left[z_2^2\indy{z_1\geq a; z_2\geq b}\right] + \rho'\bm{\Phi}_2\paren{-a,-b,\rho}\\
        & = \frac{\rho'^{\frac{3}{2}}\rho}{2\pi}\exp\paren{-\frac{a^2 - 2\rho ab + b^2}{2\rho'}} + \frac{\rho'a}{\sqrt{2\pi}}\paren{\rho^2+1}\exp\paren{-\frac{a^2}{2}}\bm{\Phi}_1\paren{\frac{\rho a- b}{\sqrt{\rho'}}}\\
        &\quad\quad\quad + \rho^2\E\left[z_2^2\indy{z_1\geq a; z_2\geq b}\right] + \rho'\bm{\Phi}_2\paren{-a,-b,\rho}
    \end{align*}
    Therefore, we can conclude that
    \begin{align*}
        & \E\left[z_1^2\indy{z_1\geq a; z_2\geq b}\right] - \rho^2\E\left[z_2^2\indy{z_1\geq a; z_2\geq b}\right]\\
        &\quad\quad\quad  = \frac{\rho'^{\frac{3}{2}}\rho}{2\pi}\exp\paren{-\frac{a^2 - 2\rho ab + b^2}{2\rho'}} + \frac{\rho' a}{\sqrt{2\pi}}\paren{\rho^2+1}T_1 + \rho'\bm{\Phi}_2\paren{-a,-b,\rho}\\
        & \E\left[z_2^2\indy{z_1\geq a; z_2\geq b}\right] - \rho^2\E\left[z_1^2\indy{z_1\geq a; z_2\geq b}\right]\\
        &\quad\quad\quad  = \frac{\rho'^{\frac{3}{2}}\rho}{2\pi}\exp\paren{-\frac{a^2 - 2\rho ab + b^2}{2\rho'}} + \frac{\rho' b}{\sqrt{2\pi}}\paren{\rho^2+1}T_2 + \rho'\bm{\Phi}_2\paren{-a,-b,\rho}
    \end{align*}
    Solving for $\E\left[z_1^2\indy{z_1\geq a; z_2\geq b}\right]$ and $\E\left[z_2^2\indy{z_1\geq a; z_2\geq b}\right]$ gives
    \begin{align*}
        \E\left[z_1^2\indy{z_1\geq a; z_2\geq b}\right] & = \frac{\rho\sqrt{\rho'}}{2\pi}\exp\paren{-\frac{a^2 - 2\rho ab + b^2}{2\rho'}} + \bm{\Phi}_2\paren{-a,-b,\rho} + \frac{1}{\sqrt{2\pi}}\paren{aT_1 + \rho^2 bT_2}\\
        \E\left[z_2^2\indy{z_1\geq a; z_2\geq b}\right] & = \frac{\rho\sqrt{\rho'}}{2\pi}\exp\paren{-\frac{a^2 - 2\rho ab + b^2}{2\rho'}} + \bm{\Phi}_2\paren{-a,-b,\rho} + \frac{1}{\sqrt{2\pi}}\paren{bT_2 + \rho^2 aT_1}
    \end{align*}
    Write $\rho' = 1-\rho^2$ gives the desired result.
\end{proof}
%%%%%%%%%%%%%%%%%
\begin{lemma}
    \label{lem:prod_term_exp}
    Let $z_1, z_2\sim\mathcal{N}\paren{0, 1}$ with $\text{Cov}\paren{z_1, z_2} = \rho$. Let $a, b \in\R$. Define
    \[
        T_1 = \exp\paren{-\frac{a^2}{2}}\bm{\Phi}_1\paren{\frac{\rho a - b}{\sqrt{1-\rho^2}}};\quad T_2 = \exp\paren{-\frac{b^2}{2}}\bm{\Phi}_1\paren{\frac{\rho b - a}{\sqrt{1-\rho^2}}}
    \]
    Then we have that
    \[
        \E\left[z_1z_2\indy{z_1\geq a;z_2\geq b}\right] =  \frac{\sqrt{1-\rho^2}}{2\pi}\exp\paren{-\frac{a^2 - 2\rho ab + b^2}{2\paren{1-\rho^2}}} + \rho\bm{\Phi}_2\paren{-a,-b,\rho} + \frac{\rho}{\sqrt{2\pi}}\paren{aT_1 + bT_2}
    \]
\end{lemma}
\begin{proof}
    To start, we write out the integral form of the expectation as
    \[
        \E\left[z_1z_2\indy{z_1\geq a;z_2\geq b}\right] = \int_b^\infty\int_a^\infty z_1z_2f\paren{z_1,z_2}dz_1dz_2 = \int_b^\infty\paren{\int_a^\infty z_1f\paren{z_1\mid z_2}dz_1}z_2f\paren{z_2}dz_2
    \]
    Since $z_1,z_2\sim\mathcal{N}\paren{0, 1}$ with covariance $\rho$, we have that $z_1\mid z_2\sim\mathcal{N}\paren{\rho z_2,1-\rho^2}$. Therefore, let $\rho' = 1-\rho^2$, we have
    \[
        f\paren{z_1\mid z_2} = \frac{1}{\sqrt{2\pi\rho'}}\exp\paren{-\frac{\paren{z_1-\rho z_2}^2}{2\rho'}}
    \]
    Thus, by Lemma~\ref{lem:first_order_gaus_int}, we have that
    \begin{align*}
        \int_a^\infty z_1f\paren{z_1\mid z_2}dz_1 & = \frac{1}{\sqrt{2\pi\rho'}}\int_a^\infty z_1\exp\paren{-\frac{\paren{z_1-\rho z_2}^2}{2\rho'}}dz_1\\
        & = \frac{1}{\sqrt{2\pi\rho'}}\paren{\rho'\exp\paren{-\frac{\paren{a - \rho z_2}^2}{2\rho'}}+\rho z_2\sqrt{2\pi \rho'}\bm{\Phi}_1\paren{\frac{\rho z_2-a}{\sqrt{\rho'}}}}\\
        & = \sqrt{\frac{\rho'}{2\pi}}\exp\paren{-\frac{\paren{a - \rho z_2}^2}{2\rho'}} + \rho z_2\bm{\Phi}_1\paren{\frac{\rho z_2 - a}{\sqrt{\rho'}}}
    \end{align*}
    where we set $\kappa = \sqrt{\rho'}$ and $\mu = \rho z_2$. Therefore
    \begin{align*}
        \E\left[z_1z_2\indy{z_1\geq a;z_2\geq b}\right] & = \int_b^\infty\paren{\sqrt{\frac{\rho'}{2\pi}}\exp\paren{-\frac{\paren{a - \rho z_2}^2}{2\rho'}} + \rho z_2\bm{\Phi}_1\paren{\frac{\rho z_2 - a}{\sqrt{\rho'}}}}z_2f\paren{z_2}dz_2\\
        & = \sqrt{\frac{\rho'}{2\pi}}\int_b^\infty z_2\exp\paren{-\frac{\paren{a - \rho z_2}^2}{2\rho'}}f\paren{z_2}dz_2 + \rho\int_b^\infty z_2^2\bm{\Phi}_1\paren{\frac{\rho z_2 - a}{\sqrt{\rho'}}}f\paren{z_2}dz_2
    \end{align*}
    To start, by Lemma~\ref{lem:cdf_forms}, we have
    \[
        \bm{\Phi}_1\paren{\frac{\rho z_2 - a}{\sqrt{\rho'}}}f\paren{z_2} = \int_a^\infty f\paren{z_1\mid z_2}dz_1f\paren{z_2} = \int_a^\infty f\paren{z_1, z_2}dz_1
    \]
    Therefore, for the second term, we have
    \[
        \int_b^\infty z_2^2\bm{\Phi}_1\paren{\frac{\rho z_2 - a}{\sqrt{\rho'}}}f\paren{z_2}dz_2 = \int_b^\infty \int_a^\infty z_2^2f\paren{z_1,z_2}dz_1dz_2 = \E\left[z_2^2\indy{z_1\geq a;z_2\geq b}\right]
    \]
    For the first term, we have
    \[
        \exp\paren{-\frac{\paren{a - \rho z_2}^2}{2\rho'}}f\paren{z_2} = \frac{1}{\sqrt{2\pi}}\exp\paren{-\frac{z_2^2 - 2\rho az_2 + a^2}{2\rho'}} = \frac{1}{\sqrt{2\pi}}\exp\paren{-\frac{a^2}{2}}\exp\paren{-\frac{\paren{z_2-\rho a}^2}{2\rho'}}
    \]
    Therefore, by Lemma~\ref{lem:first_order_gaus_int}, the first term can be written as
    \begin{align*}
        \int_b^\infty z_2\exp\paren{-\frac{\paren{a - \rho z_2}^2}{2\rho'}}f\paren{z_2}dz_2 & = \frac{1}{\sqrt{2\pi}}\exp\paren{-\frac{a^2}{2}}\int_b^\infty z_2\exp\paren{-\frac{\paren{z_2-\rho a}^2}{2\rho'}}dz_2\\
        & = \frac{1}{\sqrt{2\pi}}\exp\paren{-\frac{a^2}{2}}\paren{\rho'\exp\paren{-\frac{\paren{\rho a- b}^2}{2\rho'}} + \sqrt{2\pi\rho'}\rho a\bm{\Phi}_1\paren{\frac{\rho a - b}{\sqrt{\rho'}}}}\\
        & = \frac{\rho'}{\sqrt{2\pi}}\exp\paren{-\frac{a^2-2\rho ab + b^2}{2\rho'}} + \rho a\sqrt{\rho'}\exp\paren{-\frac{a^2}{2}}\bm{\Phi}_1\paren{\frac{\rho a - b}{\sqrt{\rho'}}}
    \end{align*}
    where we set $\kappa=\sqrt{\rho'}$ and $\mu = \rho a$ in Lemma~\ref{lem:first_order_gaus_int}. Combining the two terms, we have
    \begin{align*}
        \E\left[z_1z_2\indy{z_1\geq a;z_2\geq b}\right] & = \sqrt{\frac{\rho'}{2\pi}}\paren{\frac{\rho'}{\sqrt{2\pi}}\exp\paren{-\frac{a^2-2\rho ab + b^2}{2\rho'}} + \rho a\sqrt{\rho'}\exp\paren{-\frac{a^2}{2}}\bm{\Phi}_1\paren{\frac{\rho a - b}{\sqrt{\rho'}}}}\\
        &\quad\quad\quad + \rho \E\left[z_2^2\indy{z_1\geq a;z_2\geq b}\right]\\
        & = \frac{\rho'^{\frac{3}{2}}}{2\pi}\exp\paren{-\frac{a^2-2\rho ab + b^2}{2\rho'}} + \frac{a\rho\rho'}{\sqrt{2\pi}}T_1 +  \rho \E\left[z_2^2\indy{z_1\geq a;z_2\geq b}\right]
    \end{align*}
    From Lemma~\ref{lem:sq_term_exp}, we have that
    \[
        \E\left[z_2^2\indy{z_1\geq a; z_2\geq b}\right] = \frac{\rho\sqrt{\rho'}}{2\pi}\exp\paren{-\frac{a^2 - 2\rho ab + b^2}{2\rho'}} + \bm{\Phi}_2\paren{-a,-b,\rho} + \frac{1}{\sqrt{2\pi}}\paren{bT_2 + \rho^2 aT_1}
    \]
    Therefore
    \[
        \E\left[z_1z_2\indy{z_1\geq a;z_2\geq b}\right] = \frac{\sqrt{\rho'}}{2\pi}\exp\paren{-\frac{a^2 - 2\rho ab + b^2}{2\rho'}} + \rho\bm{\Phi}_2\paren{-a,-b,\rho} + \frac{\rho}{\sqrt{2\pi}}\paren{aT_1 + bT_2}
    \]
    Write $\rho' = 1-\rho^2$ gives the desired result.
\end{proof}

\begin{lemma}
    \label{lem:exp_prod_relu}
    Let $z_1\sim\mathcal{N}\paren{\mu_1,\kappa_1^2}$ and $z_2\sim\mathcal{N}\paren{\mu_2,\kappa_2^2}$, with $\text{Cov}\paren{z_1,z_2} = \kappa_1\kappa_2\rho$. Let $a,b\in\R$. Then we have
    \begin{align*}
        \E\left[z_1z_2\indy{z_1\geq a;z_2\geq b}\right] & = \paren{\mu_1\mu_2+\kappa_1\kappa_2\rho}\bm{\Phi}\paren{\frac{\mu_1-a}{\kappa_1},\frac{\mu_2-b}{\kappa_2},\rho}\\
        & + \frac{\kappa_1\kappa_2}{2\pi}\exp\paren{-\frac{1}{2\paren{1-\rho^2}}\paren{\frac{\paren{\mu_1-a}^2}{\kappa_1^2} - \frac{2\rho}{\kappa_1\kappa_2}\paren{\mu_1-a}\paren{\mu_2-b} + \frac{\paren{\mu_2-b}^2}{\kappa_2^2}}}\\
        &\quad\quad\quad + \frac{1}{\sqrt{2\pi}}\paren{\paren{\kappa_2\rho a + \kappa_1\mu_2}T_1 + \paren{\kappa_1\rho b + \kappa_2\mu_1}T_2}
    \end{align*}
    Here $T_1, T_2$ are defined as
    \begin{align*}
        T_1 = \exp\paren{-\frac{\paren{a- \mu_1}^2}{2\kappa_1^2}}\bm{\Phi}_1\paren{\frac{1}{\sqrt{1-\rho^2}}\paren{\frac{\rho\paren{a-\mu_1}}{\kappa_1} - \frac{b-\mu_2}{\kappa_2}}}\\ T_2 = \exp\paren{-\frac{\paren{b- \mu_2}^2}{2\kappa_2^2}}\bm{\Phi}_1\paren{\frac{1}{\sqrt{1-\rho^2}}\paren{\frac{\rho\paren{b-\mu_2}}{\kappa_2} - \frac{a-\mu_1}{\kappa_1}}}
    \end{align*}
\end{lemma}
\begin{proof}
    Let $\hat{z}_1 = \frac{z_1-\mu_1}{\kappa_1}$ and $\hat{z}_1 = \frac{z_2 - \mu_2}{\kappa_2}$. Then we have $\hat{z}_1,\hat{z}_2\sim\mathcal{N}\paren{0, 1}$. Moreover,
    \[
        \text{Cov}\paren{\hat{z}_1,\hat{z}_2} = \E\left[\hat{z}_1\hat{z}_2\right] = \frac{1}{\kappa_1\kappa_2}\E\left[\paren{z_1-\mu_1}\paren{z_2 - \mu_2}\right] = \rho
    \]
    Since $z_1 = \kappa_1\hat{z}_1 + \mu_1$ and $z_2 = \kappa_2\hat{z}_2 + \mu_2$, we have
    \begin{align*}
        \E\left[z_1z_2\indy{z_1\geq a;z_2\geq b}\right] & = \E\left[\paren{\kappa_1\hat{z}_1 + \mu_1}\paren{\kappa_2\hat{z}_2 + \mu_2}\indy{\hat{z}_1\geq \frac{a-\mu_1}{\kappa_1};\hat{z}_2\geq \frac{b - \mu_2}{\kappa_2}}\right]\\
        & = \kappa_1\kappa_2\E\left[\hat{z}_1\hat{z}_2\indy{\hat{z}_1\geq\hat{a};\hat{z}_2\geq\hat{b}}\right] + \mu_1\mu_2\E\left[\indy{\hat{z}_1\geq\hat{a};\hat{z}_2\geq\hat{b}}\right]\\
        & \quad\quad\quad + \kappa_1\mu_2\E\left[\hat{z}_1\indy{\hat{z}_1\geq\hat{a};\hat{z}_2\geq\hat{b}}\right] +\kappa_2\mu_1\E\left[\hat{z}_2\indy{\hat{z}_1\geq\hat{a};\hat{z}_2\geq\hat{b}}\right]
    \end{align*}
    where we re-defined $\hat{a} = \frac{a- \mu_1}{\kappa_1}$ and $\hat{b} = \frac{b - \mu_2}{\kappa_2}$. By Lemma~\ref{lem:cdf_forms}, Lemma~\ref{lem:first_ord_term_exp}, and Lemma~\ref{lem:prod_term_exp}, we have
    \begin{gather*}
        \E\left[\hat{z}_1\indy{\hat{z}_1\geq \hat{a}; \hat{z}_2\geq \hat{b}}\right] =\frac{1}{\sqrt{2\pi}}\paren{T_1 + \rho T_2};\quad \E\left[\hat{z}_2\indy{\hat{z}_1\geq \hat{a}; \hat{z}_2\geq \hat{b}}\right] = \frac{1}{\sqrt{2\pi}}\paren{T_2 + \rho T_1}\\
        \E\left[\hat{z}_1\hat{z}_2\indy{\hat{z}_1\geq \hat{a};\hat{z}_2\geq \hat{b}}\right] =  \frac{\sqrt{1-\rho^2}}{2\pi}\exp\paren{-\frac{\hat{a}^2 - 2\rho \hat{a}\hat{b} + \hat{b}^2}{2\paren{1-\rho^2}}} + \rho\bm{\Phi}_2\paren{-\hat{a},-\hat{b},\rho} + \frac{\rho}{\sqrt{2\pi}}\paren{\hat{a}T_1 + \hat{b}T_2}
    \end{gather*}
    and $\E\left[\indy{\hat{z}_1\geq\hat{a};\hat{z}_2\geq\hat{b}}\right] = \bm{\Phi}_2\paren{-\hat{a}, -\hat{b}, \rho}$. Here, $T_1, T_2$ are defined as
    \[
        T_1 = \exp\paren{-\frac{\hat{a}^2}{2}}\bm{\Phi}\paren{\frac{\rho\hat{a}-\hat{b}}{\sqrt{1-\rho^2}}};\quad T_2 = \exp\paren{-\frac{\hat{b}^2}{2}}\bm{\Phi}\paren{\frac{\rho\hat{b}-\hat{a}}{\sqrt{1-\rho^2}}}
    \]
    Plugging in the value of $\hat{a}$ and $\hat{b}$ gives the desired result.
\end{proof}
%%%%%%%%%%%%%%%%%%%
\subsubsection{Multi-variate Coupled Expectation}
\begin{lemma}
    \label{lem:first_order_indy_exp}
    Let $\bfc\sim\mathcal{N}\paren{\bm{\mu},\kappa^2\bfI}$, and let $\bfu\in\R^d$. Define $z = \bfc^\top\bfu$. Then we have
    \[
        \E\left[\bfc\indy{z\geq 0}\right] = \bm{\mu}\Phi_1\paren{-\frac{\bm{\mu}^\top\bfu}{\kappa\norm{\bfu}_2}} + \frac{\kappa}{\sqrt{2\pi}}\exp\paren{-\frac{\paren{\bm{\mu}^\top\bfu}^2}{2\kappa^2\norm{\bfu}_2^2}}\cdot\frac{\bfu}{\norm{\bfu}_2}
    \]
\end{lemma}
\begin{proof}
    According to the law of total expectation, 
    \[
        \E\left[\bfc\indy{\bfc^\top\bfu\geq 0}\right] = \E_{z}\left[\E_{\bfc}\left[\bfc\indy{z\geq 0}\mid z\right]\right] = \E_{z}\left[\E_{\bfc}\left[\bfc\mid z\right]\indy{z\geq 0}\right]
    \]
    By Lemma~\ref{lem:conditional_exp}, we have that $\E_{\bfc}\left[\bfc\mid z\right] = \bm{\mu} + \frac{\bfu}{\norm{\bfu}_2^2}\paren{z - \bm{\mu}^\top\bfu}$. Therefore
    \[
        \E\left[\bfc\indy{\bfc^\top\bfu\geq 0}\right] = \frac{\bfu}{\norm{\bfu}_2^2}\E_{z}\left[z\indy{z\geq 0}\right] + \paren{\bm{\mu} - \frac{\bm{\mu}^\top\bfu\cdot\bfu}{\norm{\bfu}_2^2}}\E_{z}\left[\indy{z\geq 0}\right]
    \]
    By definition, $\E_{z}\left[\indy{z\geq 0}\right] = \Pr\paren{z\geq 0} = 1 - \Pr\paren{z\leq 0}$. Since, by Lemma~\ref{lem:inner_gaussian_rv}, $z\sim\mathcal{N}\paren{\bm{\mu}^\top\bfu,\kappa^2\norm{\bfu}_2^2}$, we have that
    \[
        \Pr\paren{z\leq 0} = \Pr\paren{\frac{z-\bm{\mu}^\top\bfu}{\kappa\norm{\bfu}_2}\leq -\frac{\bm{\mu}^\top\bfu}{\norm{\bfu}_2}} = \bm{\Phi}_1\paren{-\frac{\bm{\mu}^\top\bfu}{\kappa\norm{\bfu}_2}}
    \]
    Moreover, let $\hat{z} = \frac{z-\bm{\mu}^\top\bfu}{\kappa\norm{\bfu}_2}$, then we have
    \begin{align*}
        \E_z\left[z\indy{z\geq 0}\right] & =\E_{\hat{z}}\left[\paren{ \kappa\norm{\bfu}_2\hat{z} + \bm{\mu}^\top\bfu}\indy{\hat{z}\geq -\frac{\bm{\mu}^\top\bfu}{\kappa\norm{\bfu}_2}}\right]\\
        & = \kappa\norm{\bfu}_2\E_{\hat{z}}\left[\hat{z}\indy{\hat{z}\geq -\frac{\bm{\mu}^\top\bfu}{\kappa\norm{\bfu}_2}}\right] + \bm{\mu}^\top\bfu\E_z\left[z\geq 0\right]
    \end{align*}
    By the PDF of $\hat{z}$, we have
    \[
        \E_{\hat{z}}\left[\hat{z}\indy{z\geq 0}\right] = \frac{1}{\sqrt{2\pi}}\int_a^\infty z\exp\paren{-\frac{z^2}{2}}dz = -\frac{1}{\sqrt{2\pi}}\exp\paren{-\frac{z^2}{2}}\vert_a^\infty = \frac{1}{\sqrt{2\pi}}\exp\paren{-\frac{a^2}{2}}
    \]
    Therefore
    \[
        \E_z\left[z\indy{z\geq 0}\right] = \frac{\kappa\norm{\bfu}_2}{\sqrt{2\pi}}\exp\paren{-\frac{\paren{\bm{\mu}^\top\bfu}^2}{2\kappa^2\norm{\bfu}_2^2}} + \bm{\mu}^\top\bfu\E_z\left[z\geq 0\right]
    \]
    Plugging in gives
    \[
        \E\left[\bfc\indy{\bfc^\top\bfu\geq 0}\right] = \bm{\mu}\Phi_1\paren{-\frac{\bm{\mu}^\top\bfu}{\kappa\norm{\bfu}_2}} + \frac{\kappa}{\sqrt{2\pi}}\exp\paren{-\frac{\paren{\bm{\mu}^\top\bfu}^2}{2\kappa^2\norm{\bfu}_2^2}}\cdot\frac{\bfu}{\norm{\bfu}_2}
    \]
\end{proof}

\begin{lemma}
    \label{lem:exp_grad_2nd}
    Let $\bm{c}\sim\mathcal{N}\paren{\bm{\mu},\kappa^2\bfI}$ with $\kappa \leq 1$. Let $\bfu,\bfv$ be given, and let $z_1 = \bfc^\top\bfu, z_2 = \bfc^\top\bfv$. Then we have that
    \[
        \norm{\E_{\bfc}\left[\bfc\mathbb{I}\left\{z_1\geq 0\right\}\right] - \bm{\mu}\bm{\Phi}_1\paren{\frac{\bm{\mu}^\top\bfu}{\kappa\norm{\bfu}_2}}}_{\infty}\leq \kappa\norm{\bfu}_{2}\exp\paren{-\frac{\paren{\bm{\mu}^\top\bfu}^2}{2\kappa^2\norm{\bfu}_2}}
    \]
\end{lemma}
\begin{proof}
Given the form of the conditional expectation, we have
\[
    \E_{\bfc}\left[\bfc\mathbb{I}\left\{z_1\geq 0\right\}\right] = \int_0^\infty\E_{\bfc}\left[\bfc\mid z_1\right]f_1(z_1)dz_1 = \int_0^\infty\paren{\bm{\mu} + \frac{\bfu}{\norm{\bfu}_2^2}\paren{z_1 - \bm{\mu}^\top\bfu}}f_1(z_1)dz_1
\]
Since $z_1 = \bfc^\top\bfu$, we must have that $z_1\sim\mathcal{N}\paren{\bm{\mu}^\top\bfu,\kappa^2\norm{\bfu}_2^2}$. Define $z' = \frac{z_1 - \bm{\mu}^\top\bfu}{\kappa\norm{\bfu}_2}$, then we have that $z_1 = \kappa\norm{\bfu}_2z' + \bm{\mu}^\top\bfu$
\begin{align*}
    \E_{\bfc}\left[\bfc\mathbb{I}\left\{z_1\geq 0\right\}\right] & = \int_{\frac{- \bm{\mu}^\top\bfu}{\kappa\norm{\bfu}_2}}^{\infty}\paren{\mu+\kappa\bfu\cdot z'}f\paren{z'}dz'\\
    & = \mu\cdot \int_{-\infty}^{\frac{\bm{\mu}^\top\bfu}{\kappa\norm{\bfu}_2}}f\paren{z'}dz' + \kappa\bfu\cdot\int_{-\infty}^{\frac{\bm{\mu}^\top\bfu}{\kappa\norm{\bfu}_2}}z'f\paren{z'}dz'\\
    & = \bm{\mu}\bm{\Phi}_1\paren{\frac{\bm{\mu}^\top\bfu}{\kappa\norm{\bfu}_2}} + \kappa\bfu\exp\paren{-\frac{\paren{\bm{\mu}^\top\bfu}^2}{2\kappa^2\norm{\bfu}_2}}
\end{align*}
where in the last equality we applied Lemma~\ref{lem:first_order_gaus_int} with $a = 0$ and $\kappa = 1$ in the lemma. Therefore, we have that
\[
    \norm{\E_{\bfc}\left[\bfc\mathbb{I}\left\{z_1\geq 0\right\}\right] - \bm{\mu}\bm{\Phi}_1\paren{\frac{\bm{\mu}^\top\bfu}{\kappa\norm{\bfu}_2}}}_{\infty}\leq \kappa\norm{\bfu}_{\infty}\exp\paren{-\frac{\paren{\bm{\mu}^\top\bfu}^2}{2\kappa^2\norm{\bfu}_2}}
\]
\end{proof}
Now, we shall dive into $\E_{\bfc}\left[\bfc\bfc^\top\mathbb{I}\left\{z_1\geq 0;z_2\geq 0\right\}\right]$.
\begin{lemma}
    \label{lem:exp_grad_1st_prod}
    Let $\bm{c}\sim\mathcal{N}\paren{\bm{\mu},\kappa^2\bfI}$ with $\norm{\bm{\mu}}_2\geq 2$ and $\kappa \leq 1$. Let $\bfu,\bfv$ be given, and let $z_1 = \bm{\mu}^\top\bfu, z_2 = \bm{\mu}^\top\bfv$. Then we have that
    \[
        \norm{\E_{\bfc}\left[\bfc\bfc^\top\mathbb{I}\left\{z_1\geq 0;z_2\geq 0\right\}\right]\bfv - \paren{\bm{\mu}\bm{\mu}^\top\bfv + 3\kappa^2\bfv}\bm{\Phi}_1\paren{\frac{\bm{\mu^\top\bfu}}{2\kappa\norm{\bfv}_2}}\bm{\Phi}_1\paren{\frac{\bm{\mu^\top\bfv}}{2\kappa\norm{\bfv}_2}}} \leq \Delta
    \]
    where $\Delta$ is given by
    \[
        \Delta = 2\kappa\norm{\bfv}_2\paren{\norm{\bm{\mu}}_2\paren{\phi\paren{\frac{\bm{\mu^\top\bfu}}{2\kappa\norm{\bfu}_2}} + \phi\paren{\frac{\bm{\mu^\top\bfu}}{2\kappa\norm{\bfv}_2}}} +\norm{\bm{\mu}}_{\infty}\paren{\psi\paren{\frac{\bm{\mu^\top\bfu}}{2\kappa\norm{\bfu}_2}} + \psi\paren{\frac{\bm{\mu^\top\bfu}}{2\kappa\norm{\bfu}_2}}}}
    \]
\end{lemma}
%%%%%%%%%%
\begin{proof}
    We have
    \[
        \E_{\bfc}\left[\bfc\bfc^\top\mathbb{I}\left\{z_1\geq 0;z_2\geq 0\right\}\right] = \int_{z_1,z_2\geq 0}\E_{\bfc}\left[\bfc\bfc^\top\mid z_1, z_2\right]f\paren{z_1, z_2}dz_1dz_2
    \]
    Notice that
    \[
        \E_{\bfc}\left[\bfc\bfc^\top\mid z_1, z_2\right] = \text{Cov}\paren{\bfc\mid z_1,z_2} + \E[\bfc\mid z_1,z_2]\E[\bfc\mid z_1,z_2]^\top
    \]
    By the form of $\text{Cov}\paren{\bfc\mid z_1,z_2}$, we have
    \[
        \text{Cov}\paren{\bfc\mid z_1,z_2} = \kappa^2\bfI - \frac{\kappa^2\paren{\bfu\bfv^\top - \bfv\bfu^\top}^2}{\norm{\bfv}_2^2\norm{\bfu}_2^2 - \inner{\bfv}{\bfu}^2}  := \bfM
    \]
    By the form of $\E[\bfc\mid z_1,z_2]$ we have
    \[
        \E[\bfc\mid z_1,z_2] = z_1\cdot\bfs_1 + z_2\cdot\bfs_2 + \bfs_3
    \]
    where 
    \begin{gather*}
        \bfs_1 = \frac{\norm{\bfv}_2^2\bfu - \inner{\bfu}{\bfv}\bfv}{\norm{\bfv}_2^2\norm{\bfu}_2^2 - \inner{\bfv}{\bfu}^2};  \bfs_2 = \frac{\norm{\bfu}_2^2\bfv - \inner{\bfu}{\bfv}\bfu}{\norm{\bfv}_2^2\norm{\bfu}_2^2 - \inner{\bfv}{\bfu}^2}\\
        \bfs_3 = \bm{\mu} - \frac{\paren{\norm{\bfv}_2^2\bfu - \inner{\bfu}{\bfv}\bfv}\inner{\bm{\mu}}{\bfu} + \paren{\norm{\bfu}_2^2\bfv - \inner{\bfu}{\bfv}\bfu}\inner{\bm{\mu}}{\bfv}}{\norm{\bfv}_2^2\norm{\bfu}_2^2 - \inner{\bfv}{\bfu}^2}
    \end{gather*}
    Therefore
    \[
        \E\left[\bfc\mid z_1,z_2\right]\E\left[\bfc\mid z_1,z_2\right]^\top = \paren{z_1\cdot\bfs_1 + z_2\cdot\bfs_2 + \bfs_3}\paren{z_1\cdot\bfs_1 + z_2\cdot\bfs_2 + \bfs_3}^\top
    \]
    Recall that we are interested in
    \[
        \E_{\bfc}\left[\bfc\bfc^\top\mathbb{I}\left\{z_1\geq 0;z_2\geq 0\right\}\right] = \int_{z_1,z_2\geq 0}\paren{\text{Cov}\paren{\bfc\mid z_1,z_2} + \E[\bfc\mid z_1,z_2]\E[\bfc\mid z_1,z_2]^\top}f(z_1,z_2)dz_1dz_2
    \]
    Let
    \[
        \hat{z}_1 = \frac{z_1 - \bm{\mu}^\top\bfu}{\kappa\norm{\bfu}_2};\quad \hat{z}_2 = \frac{z_1 - \bm{\mu}^\top\bfv}{\kappa\norm{\bfv}_2}
    \]
    Then we have
    \[
        \hat{z}_1,\hat{z}_2\sim\mathcal{N}(0,1);\quad \text{Cov}\paren{\hat{z}_1,\hat{z}_2} = \frac{\bfu^\top\bfv}{\norm{\bfu}_2\norm{\bfv}_2}: = \rho
    \]
    Thus,
    \[
        \hat{z}_1\mid \hat{z}_2 = \gamma_2 \sim\mathcal{N}\paren{\rho\gamma_2, 1- \rho^2};\quad \hat{z}_2\mid \hat{z}_1 = \gamma_1 \sim\mathcal{N}\paren{\rho\gamma_1, 1- \rho^2}
    \]
    Moreover,
    \begin{align*}
        \E[\bfc\mid z_1,z_2] & = z_1\cdot\bfs_1 + z_2\cdot\bfs_2 + \bfs_3\\
        & = \paren{\kappa\norm{\bfu}_2\hat{z}_1 + \bm{\mu}^\top\bfu}\bfs_1 + \paren{\kappa\norm{\bfv}_2\hat{z}_2 + \bm{\mu}^\top\bfv}\bfs_2 + \bfs_3
    \end{align*}
    Redefine
    \[
        \hat{\bfs}_1 = \kappa\norm{\bfu}_2\bfs_1;\quad\hat{\bfs}_2 = \kappa\norm{\bfv}_2\bfs_2;\quad\hat{\bfs}_3 = \bm{\mu}^\top\bfu\cdot\bfs_1 + \bm{\mu}^\top\bfv\cdot\bfs_2 + \bfs_3 = \bm{\mu}
    \]
    Then we have
    \[
        \E[\bfc\mid z_1,z_2] = \hat{\bfs}_1\hat{z}_1+\hat{\bfs}_2\hat{z}_2 + \hat{\bfs}_3
    \]
    In this case, let $\hat{f}$ be the joint PDF of $\hat{z}_1$ and $\hat{z}_2$, then we have
    \begin{align*}
        \hat{f}\paren{\hat{z}_1,\hat{z}_2} & = \frac{1}{2\pi\sqrt{1 - \rho^2}}\exp\paren{-\frac{1}{2(1-\rho^2)}\paren{z_1^2 - 2\rho z_1z_2 + z_2^2}}\\
        & = \frac{\kappa^2\norm{\bfu}_2\norm{\bfv}_2}{2\pi\kappa^2\norm{\bfu}_2\norm{\bfv}_2\sqrt{1-\rho^2}}\exp\paren{- 
     \frac{1}{2(1-\rho^2)}\paren{z_1^2 - 2\rho z_1z_2 + z_2^2}}\\
        & = \kappa^2\norm{\bfu}_2\norm{\bfv}_2f\paren{\hat{z}_1,\hat{z}_2}
    \end{align*}
    Therefore, since $dz_1 = \kappa\norm{\bfu}_2d\hat{z}_1$, and $dz_2 = \kappa\norm{\bfv}_2d\hat{z}_2$, we have
    \begin{align*}
        \E_{\bfc}\left[\bfc\bfc^\top\mathbb{I}\left\{z_1\geq 0;z_2\geq 0\right\}\right] & = \int_{z_1,z_2\geq 0}\paren{\text{Cov}\paren{\bfc\mid z_1,z_2} + \E[\bfc\mid z_1,z_2]\E[\bfc\mid z_1,z_2]^\top}f(z_1,z_2)dz_1dz_2\\
        & = \int_{-\frac{\bm{\mu}^\top\bfv}{\kappa\norm{\bfv}_2}}^\infty\int_{-\frac{\bm{\mu}^\top\bfu}{\kappa\norm{\bfu}_2}}^\infty\paren{\text{Cov}\paren{\bfc\mid z_1,z_2} + \E[\bfc\mid z_1,z_2]\E[\bfc\mid z_1,z_2]^\top}\hat{f}(\hat{z}_1,\hat{z}_2)d\hat{z}_1d\hat{z}_2\\
        & = \hat{\bfs}_1\hat{\bfs}_1^\top\underbrace{\int_{-\frac{\bm{\mu}^\top\bfv}{\kappa\norm{\bfv}_2}}^\infty\int_{-\frac{\bm{\mu}^\top\bfu}{\kappa\norm{\bfu}_2}}^\infty\hat{z}_1^2\hat{f}(\hat{z}_1,\hat{z}_2)d\hat{z}_1d\hat{z}_2}_{\mathcal{I}_1}\\
        & \quad\quad\quad + \hat{\bfs}_2\hat{\bfs}_2^\top\underbrace{\int_{-\frac{\bm{\mu}^\top\bfv}{\kappa\norm{\bfv}_2}}^\infty\int_{-\frac{\bm{\mu}^\top\bfu}{\kappa\norm{\bfu}_2}}^\infty\hat{z}_2^2\hat{f}(\hat{z}_1,\hat{z}_2)d\hat{z}_1d\hat{z}_2}_{\mathcal{I}_2}\\
        & \quad\quad\quad + \paren{\hat{\bfs}_1\hat{\bfs}_2^\top + \hat{\bfs}_2\hat{\bfs}_1^\top}\underbrace{\int_{-\frac{\bm{\mu}^\top\bfv}{\kappa\norm{\bfv}_2}}^\infty\int_{-\frac{\bm{\mu}^\top\bfu}{\kappa\norm{\bfu}_2}}^\infty\hat{z}_1\hat{z}_2\hat{f}(\hat{z}_1,\hat{z}_2)d\hat{z}_1d\hat{z}_2}_{\mathcal{I}_3}\\
        & \quad\quad\quad + \paren{\hat{\bfs}_1\hat{\bfs}_3 + \hat{\bfs}_3\hat{\bfs}_1^\top}\underbrace{\int_{-\frac{\bm{\mu}^\top\bfv}{\kappa\norm{\bfv}_2}}^\infty\int_{-\frac{\bm{\mu}^\top\bfu}{\kappa\norm{\bfu}_2}}^\infty\hat{z}_1\hat{f}(\hat{z}_1,\hat{z}_2)d\hat{z}_1d\hat{z}_2}_{\mathcal{I}_4}\\
        & \quad\quad\quad + \paren{\hat{\bfs}_2\hat{\bfs}_3^\top + \hat{\bfs}_3\hat{\bfs}_2^\top}\underbrace{\int_{-\frac{\bm{\mu}^\top\bfv}{\kappa\norm{\bfv}_2}}^\infty\int_{-\frac{\bm{\mu}^\top\bfu}{\kappa\norm{\bfu}_2}}^\infty\hat{z}_2\hat{f}(\hat{z}_1,\hat{z}_2)d\hat{z}_1d\hat{z}_2}_{\mathcal{I}_5}\\
        & \quad\quad\quad + \paren{\hat{\bfs}_3\hat{\bfs}_3^\top + \bfM}\underbrace{\int_{-\frac{\bm{\mu}^\top\bfv}{\kappa\norm{\bfv}_2}}^\infty\int_{-\frac{\bm{\mu}^\top\bfu}{\kappa\norm{\bfu}_2}}^\infty \hat{f}(\hat{z}_1,\hat{z}_2)d\hat{z}_1d\hat{z}_2}_{\mathcal{I}_6}
    \end{align*}
    Since our goal is to study the term $\E_{\bfc}\left[\bfc\bfc^\top\indy{z_1\geq 0;z_2\geq 0}\right]\bfv$, we need to understand the terms $\mathcal{I}_1$ to $\mathcal{I}_6$, as well as understanding the matrix-vector product in front of these terms. To start, under some standard computation, we have
    \begin{align*}
        \hat{\bfs}_1^\top\bfu &  = \kappa\norm{\bfu}_2\cdot \frac{\norm{\bfv}_2^2\bfu^\top\bfu - \bfv^\top\bfu\cdot\bfv^\top\bfu}{\norm{\bfv}_2^2\norm{\bfu}_2^2 - \paren{\bfv^\top\bfu}^2} = \kappa\norm{\bfu}_2\\
        \hat{\bfs}_1^\top\bfv &  = \kappa\norm{\bfu}_2\cdot \frac{\norm{\bfv}_2^2\bfu^\top\bfv - \bfv^\top\bfu\cdot\bfv^\top\bfv}{\norm{\bfv}_2^2\norm{\bfu}_2^2 - \paren{\bfv^\top\bfu}^2} = 0\\
        \hat{\bfs}_2^\top\bfu &  = \kappa\norm{\bfv}_2\cdot \frac{\norm{\bfu}_2^2\bfv^\top\bfu - \bfv^\top\bfu\cdot\bfu^\top\bfu}{\norm{\bfv}_2^2\norm{\bfu}_2^2 - \paren{\bfv^\top\bfu}^2} = 0\\
        \hat{\bfs}_2^\top\bfv &  = \kappa\norm{\bfv}_2\cdot \frac{\norm{\bfu}_2^2\bfv^\top\bfv - \bfv^\top\bfu\cdot\bfu^\top\bfv}{\norm{\bfv}_2^2\norm{\bfu}_2^2 - \paren{\bfv^\top\bfu}^2} = \kappa\norm{\bfv}_2\\
    \end{align*}
    Therefore, the following must holds
    \begin{gather*}
        \hat{\bfs}_1\hat{\bfs}_1^\top\bfv = \bm{0};\quad \hat{\bfs}_2\hat{\bfs}_2^\top\bfv = \kappa\norm{\bfv}_2\hat{\bfs}_2;\quad \paren{\hat{\bfs}_1\hat{\bfs}_2^\top + \hat{\bfs}_2\hat{\bfs}_1^\top}\bfv = \kappa\norm{\bfv}_2\hat{\bfs}_1; \\
        \paren{\hat{\bfs}_1\hat{\bfs}_3^\top + \hat{\bfs}_3\hat{\bfs}_1^\top}\bfv = \bm{\mu}^\top\bfv\cdot \hat{\bfs}_1;\quad \paren{\hat{\bfs}_2\hat{\bfs}_3^\top + \hat{\bfs}_3\hat{\bfs}_2^\top}\bfv = \bm{\mu}^\top\bfv\cdot \hat{\bfs}_2 + \kappa\norm{\bfv}_2\bm{\mu}
    \end{gather*}
    Lastly, we have
    \begin{align*}
        \paren{\hat{\bfs}_3\hat{\bfs}_3^\top + \bfM}\bfv & = \bm{\mu}^\top\bfv \cdot \bm{\mu} + \kappa^2\bfv - \frac{\kappa^2}{\norm{\bfv}_2^2\norm{\bfu}_2^2 - \paren{\bfv^\top\bfu}^2}\paren{\bfu\bfv^\top-\bfv\bfu^\top}^2\bfv\\
        & = \bm{\mu}^\top\bfv \cdot \bm{\mu} + \kappa^2\bfv - \frac{\kappa^2}{\norm{\bfv}_2^2\norm{\bfu}_2^2 - \paren{\bfv^\top\bfu}^2}\paren{\bfu\bfv^\top -\bfv\bfu^\top}\paren{\bfu^\top\bfv\cdot\bfv - \norm{\bfv}_2^2\bfu}\\
        & = \bm{\mu}^\top\bfv \cdot \bm{\mu} + \kappa^2\bfv\\
        &\quad\quad\quad - \frac{\kappa^2}{\norm{\bfv}_2^2\norm{\bfu}_2^2 - \paren{\bfv^\top\bfu}^2}\paren{\bfu^\top\bfv\cdot \norm{\bfv}_2^2\bfu - \paren{\bfu^\top\bfv}^2\bfv - \bfu^\top\bfv\cdot\norm{\bfv}_2^2\bfu + \norm{\bfv}_2^2\norm{\bfu}_2^2\bfv}\\
        &=\bm{\mu}^\top\bfv \cdot \bm{\mu} + \kappa^2\bfv + \kappa^2\bfv\\
        & = \bm{\mu}^\top\bfv \cdot \bm{\mu} + 2\kappa^2\bfv
    \end{align*}
    Therefore, we can write $\E_{\bfc}\left[\bfc\bfc^\top\indy{z_1\geq 0;z_2\geq 0}\right]\bfv$ as
    \begin{equation}
        \label{eq:mat_exp_vec_prod_form1}
        \begin{aligned}
            \E_{\bfc}\left[\bfc\bfc^\top\indy{z_1\geq 0;z_2\geq 0}\right]\bfv & = \kappa\norm{\bfv}_2\hat{\bfs}_2\cdot\mathcal{I}_2 + \kappa\norm{\bfv}_2\hat{\bfs}_1\cdot\mathcal{I}_3 + \bm{\mu}^\top\bfv \cdot\hat{\bfs}_1\cdot \mathcal{I}_4\\
            &\quad\quad\quad + \paren{\bm{\mu}^\top\bfv\cdot \hat{\bfs}_2 + \kappa\norm{\bfv}_2\bm{\mu}}\mathcal{I}_5 + \paren{\bm{\mu}^\top\bfv \cdot \bm{\mu} + 2\kappa^2\bfv}\mathcal{I}_6\\
            & = \kappa\norm{\bfv}_2\paren{\mathcal{I}_2\cdot\hat{\bfs}_2 + \mathcal{I}_3\cdot\hat{\bfs}_1} + \bm{\mu}^\top\bfv\paren{\mathcal{I}_4\cdot\hat{\bfs}_1+\mathcal{I}_5\cdot\hat{\bfs}_2}\\
            & \quad\quad\quad + \paren{\kappa\norm{\bfv}_2\cdot\mathcal{I}_5 + \bm{\mu}^\top\bfv\cdot\mathcal{I}_6}\bm{\mu} + 2\kappa^2\mathcal{I}_6\bfv
        \end{aligned}
    \end{equation}
    By the definition of $\mathcal{I}_2$ to $\mathcal{I}_6$, we first notice that
    \begin{align*}
        \mathcal{I}_6 & = \int_{-\frac{\bm{\mu}^\top\bfu}{\kappa\norm{\bfu}_2}}^\infty\int_{-\frac{\bm{\mu}^\top\bfu}{\kappa\norm{\bfu}_2}}^\infty f(\hat{z}_1,\hat{z}_2)d\hat{z}_1d\hat{z}_2\\
        & = \mathbb{P}\paren{\hat{z}_1 \geq -\frac{\bm{\mu}^\top\bfu}{\kappa\norm{\bfu}_2};\hat{z}_2\geq -\frac{\bm{\mu}^\top\bfv}{\kappa\norm{\bfv}_2}}\\
        & = \mathbb{P}\paren{\hat{z}_1\leq\frac{\bm{\mu}^\top\bfu}{\kappa\norm{\bfu}_2};\hat{z}_2\leq \frac{\bm{\mu}^\top\bfv}{\kappa\norm{\bfv}_2}}\\
        & = \Phi_2\paren{\frac{\bm{\mu}^\top\bfu}{\kappa\norm{\bfu}_2},\frac{\bm{\mu}^\top\bfv}{\kappa\norm{\bfv}_2},\rho}
    \end{align*}
    Moreover, we can invoke Lemma~\ref{lem:first_ord_term_exp}, Lemma~\ref{lem:sq_term_exp}, and Lemma~\ref{lem:prod_term_exp} to get that
    \begin{gather*}
        \mathcal{I}_4 =\frac{1}{\sqrt{2\pi}}\paren{T_1 + \rho T_2};\quad \mathcal{I}_5 = \frac{1}{\sqrt{2\pi}}\paren{T_2 + \rho T_1}\\
        \mathcal{I}_2 = \frac{\rho\sqrt{1-\rho^2}}{2\pi}\exp\paren{-\frac{a^2 - 2\rho ab + b^2}{2\paren{1-\rho^2}}} + \bm{\Phi}_2\paren{-a,-b,\rho} + \frac{1}{\sqrt{2\pi}}\paren{bT_2 + \rho^2 aT_1}\\
        \mathcal{I}_3 =  \frac{\sqrt{1-\rho^2}}{2\pi}\exp\paren{-\frac{a^2 - 2\rho ab + b^2}{2\paren{1-\rho^2}}} + \rho\bm{\Phi}_2\paren{-a,-b,\rho} + \frac{\rho}{\sqrt{2\pi}}\paren{aT_1 + bT_2}
    \end{gather*}
    where $T_1, T_2$ and $a, b$ are defined as
    \begin{gather*}
        T_1 = \exp\paren{-\frac{a^2}{2}}\bm{\Phi}_1\paren{\frac{\rho a - b}{\sqrt{1 - \rho^2}}};\quad T_2 = \exp\paren{-\frac{b^2}{2}}\bm{\Phi}_1\paren{\frac{\rho b - a}{\sqrt{1 - \rho^2}}}\\
        a = -\frac{\bm{\mu}^\top\bfu}{\kappa\norm{\bfu}_2};\quad b = -\frac{\bm{\mu}^\top\bfv}{\kappa\norm{\bfv}_2};\quad \rho = \frac{\bfu^\top\bfv}{\norm{\bfu}_2\norm{\bfv}_2}
    \end{gather*}
    To ease our computation, we define 
    \[
        E = \frac{1-\rho^2}{2\pi}\exp\paren{-\frac{a^2 - 2\rho ab + b^2}{2\paren{1-\rho^2}}}; F = \bm{\Phi}_2\paren{-a,-b,\rho}
    \]
    Then the terms $\mathcal{I}_2$ to $\mathcal{I}_6$ can be written as
    \begin{equation}
        \label{eq:I_terms_simp}
        \begin{gathered}
            \mathcal{I}_2 = \rho E + F + \frac{1}{\sqrt{2\pi}}\paren{ b T_2 + \rho^2a T_1};\quad \mathcal{I}_3 = E + \rho F + \frac{1}{\sqrt{2\pi}}\paren{ b T_2 + a T_1}\\
            \mathcal{I}_4 =\frac{1}{\sqrt{2\pi}}\paren{T_1 + \rho T_2};\quad \mathcal{I}_5 = \frac{1}{\sqrt{2\pi}}\paren{T_2 + \rho T_1};\quad\mathcal{I}_6 = F
            \end{gathered}
    \end{equation}
    Now, the trick of evaluating (\ref{eq:mat_exp_vec_prod_form1}) is to re-write $\hat{\bfs}_1$ and $\hat{\bfs}_2$ as below
    \begin{equation}
        \label{eq:s1_s2_simp}
        \begin{aligned}
            \hat{\bfs}_1 & = \frac{\kappa\norm{\bfu}_2}{ \norm{\bfu}_2^2\norm{\bfv}_2^2 - \paren{\bfu^\top\bfv}^2}\cdot\paren{\norm{\bfv}_2^2\bfu - \bfu^\top\bfv\cdot\bfv}\\
            & = \frac{\kappa\paren{\norm{\bfv}_2^2\bfu - \bfu^\top\bfv\cdot\bfv}}{\norm{\bfu}_2\norm{\bfv}_2^2\paren{1 - \rho^2}}\\
            & = \frac{\kappa}{1-\rho^2}\cdot\frac{\bfu}{\norm{\bfu}}_2 - \frac{\kappa\rho}{1 - \rho^2}\cdot\frac{\bfv}{\norm{\bfv}_2}\\
            & = \frac{\kappa}{1-\rho^2}\paren{\frac{\bfu}{\norm{\bfu}}_2 - \rho\cdot \frac{\bfv}{\norm{\bfv}_2}}\\
            \hat{\bfs}_2 & = \frac{\kappa\norm{\bfv}_2}{ \norm{\bfu}_2^2\norm{\bfv}_2^2 - \paren{\bfu^\top\bfv}^2}\cdot\paren{\norm{\bfu}_2^2\bfv - \bfu^\top\bfv\cdot\bfu}\\
            & = \frac{\kappa\paren{\norm{\bfu}_2^2\bfv - \bfu^\top\bfv\cdot\bfu}}{\norm{\bfu}_2\norm{\bfv}_2^2\paren{1 - \rho^2}}\\
            & = \frac{\kappa}{1-\rho^2}\cdot\frac{\bfv}{\norm{\bfv}}_2 - \frac{\kappa\rho}{1 - \rho^2}\cdot\frac{\bfu}{\norm{\bfu}_2}\\
            & = \frac{\kappa}{1-\rho^2}\paren{\frac{\bfv}{\norm{\bfv}_2} - \rho\cdot\frac{\bfu}{\norm{\bfu}_2}}
        \end{aligned}
    \end{equation}
    Now, we can simplify (\ref{eq:mat_exp_vec_prod_form1}) with (\ref{eq:I_terms_simp}) and (\ref{eq:s1_s2_simp}). To start, for the terms $\mathcal{I}_2\cdot\hat{\bfs}_2 + \mathcal{I}_3\cdot\hat{\bfs}_1$ we have
    \begin{align*}
        \mathcal{I}_2\cdot\hat{\bfs}_2 + \mathcal{I}_3\cdot\hat{\bfs}_1 & = \frac{\kappa}{1-\rho^2}\paren{\rho E + F + \frac{1}{\sqrt{2\pi}}\paren{bT_2 + \rho^2aT_1}}\paren{\frac{\bfv}{\norm{\bfv}_2} - \rho\cdot\frac{\bfu}{\norm{\bfu}_2}}\\
        &\quad\quad\quad + \frac{\kappa}{1-\rho^2}\paren{E + \rho F + \frac{\rho}{\sqrt{2\pi}}\paren{bT_2 + aT_1}}\paren{\frac{\bfu}{\norm{\bfu}_2} - \rho\cdot \frac{\bfv}{\norm{\bfv}_2}}\\
        & = \frac{\kappa E}{1-\rho^2}\paren{\rho\paren{\frac{\bfv}{\norm{\bfv}_2} - \rho\cdot\frac{\bfu}{\norm{\bfu}_2}} + \paren{\frac{\bfu}{\norm{\bfu}_2} - \rho\cdot \frac{\bfv}{\norm{\bfv}_2}}}\\
        &\quad\quad\quad + \frac{\kappa F}{1-\rho^2}\paren{\paren{\frac{\bfv}{\norm{\bfv}_2} - \rho\cdot\frac{\bfu}{\norm{\bfu}_2}} + \rho\paren{\frac{\bfu}{\norm{\bfu}_2} - \rho\cdot \frac{\bfv}{\norm{\bfv}_2}}}\\
        & \quad\quad\quad + \frac{\kappa\rho aT_1}{\sqrt{2\pi}\paren{1-\rho^2}}\paren{\rho\paren{\frac{\bfv}{\norm{\bfv}_2} - \rho\cdot\frac{\bfu}{\norm{\bfu}_2}} + \paren{\frac{\bfu}{\norm{\bfu}_2} - \rho\cdot \frac{\bfv}{\norm{\bfv}_2}}}\\
        &\quad\quad\quad + \frac{\kappa b T_2}{\sqrt{2\pi}\paren{1-\rho^2}}\paren{\paren{\frac{\bfv}{\norm{\bfv}_2} - \rho\cdot\frac{\bfu}{\norm{\bfu}_2}} + \rho\paren{\frac{\bfu}{\norm{\bfu}_2} - \rho\cdot \frac{\bfv}{\norm{\bfv}_2}}}\\
        & = \kappa E\cdot\frac{\bfu}{\norm{\bfu}_2} + \kappa F\cdot\frac{\bfv}{\norm{\bfv}_2} + \frac{\kappa\rho a T_1}{\sqrt{2\pi}}\cdot\frac{\bfu}{\norm{\bfu}_2} + \frac{\kappa b T_2}{\sqrt{2\pi}}\cdot\frac{\bfv}{\norm{\bfv}_2}\\
        & = \kappa\paren{\paren{E + \frac{\rho a T_1}{\sqrt{2\pi}}}\frac{\bfu}{\norm{\bfu}_2} + \paren{F + \frac{b T_2}{\sqrt{2\pi}}}\frac{\bfv}{\norm{\bfv}_2}}
    \end{align*}
    Similarly, for the term $\mathcal{I}_4\hat{\bfs}_1 + \mathcal{I}_4\hat{\bfs}_2$, we have
    \begin{align*}
        \mathcal{I}_4\hat{\bfs}_1 + \mathcal{I}_4\hat{\bfs}_2 & = \frac{\kappa}{\sqrt{2\pi}\paren{1-\rho^2}}\paren{\paren{T_1 + \rho T_2}\paren{\frac{\bfv}{\norm{\bfv}_2} - \rho\cdot\frac{\bfu}{\norm{\bfu}_2}} + \paren{T_2 + \rho T_1}\paren{\frac{\bfu}{\norm{\bfu}_2} - \rho\cdot\frac{\bfv}{\norm{\bfv}_2}}}\\
        & = \frac{\kappa T_1}{\sqrt{2\pi}\paren{1-\rho^2}}\paren{\paren{\frac{\bfv}{\norm{\bfv}_2} - \rho\cdot\frac{\bfu}{\norm{\bfu}_2}} + \rho\paren{\frac{\bfu}{\norm{\bfu}_2} - \rho\cdot\frac{\bfv}{\norm{\bfv}_2}}}\\
        &\quad\quad\quad +\frac{\kappa T_2}{\sqrt{2\pi}\paren{1-\rho^2}}\paren{\rho\paren{\frac{\bfv}{\norm{\bfv}_2} - \rho\cdot\frac{\bfu}{\norm{\bfu}_2}} + \paren{\frac{\bfu}{\norm{\bfu}_2} - \rho\cdot\frac{\bfv}{\norm{\bfv}_2}}}\\
        & = \frac{\kappa T_1}{\sqrt{2\pi}}\cdot\frac{\bfv}{\norm{\bfv}_2} + \frac{\kappa T_2}{\sqrt{2\pi}}\cdot\frac{\bfu}{\norm{\bfu}_2}\\
        & = \frac{\kappa}{\sqrt{2\pi}}\paren{T_1\cdot \frac{\bfv}{\norm{\bfv}_2} + T_2\cdot \frac{\bfu}{\norm{\bfu}_2}}
    \end{align*}
    Applying these evaluations, (\ref{eq:mat_exp_vec_prod_form1}) becomes
    \begin{equation}
        \label{eq:mat_exp_vec_prod_form2}
        \begin{aligned}
            \E_{\bfc}\left[\bfc\bfc^\top\indy{z_1\geq 0;z_2\geq 0}\right]\bfv & = \kappa^2\norm{\bfv}_2\paren{\paren{E + \frac{\rho a T_1}{\sqrt{2\pi}}}\frac{\bfu}{\norm{\bfu}_2} + \paren{F + \frac{b T_2}{\sqrt{2\pi}}}\frac{\bfv}{\norm{\bfv}_2}}\\
            &\quad\quad\quad + \frac{\kappa\bm{\mu}^\top\bfv}{\sqrt{2\pi}}\paren{T_1\cdot \frac{\bfv}{\norm{\bfv}_2} + T_2\cdot \frac{\bfu}{\norm{\bfu}_2}}\\
            &\quad\quad\quad + \frac{\kappa\norm{\bfv}_2}{\sqrt{2\pi}}\paren{T_2 + \rho T_1}\bm{\mu} + \bm{\mu}^\top\bfv\cdot F\cdot \bm{\mu} + 2\kappa^2F\bfv\\
            & = \underbrace{\kappa^2\norm{\bfv}_2\paren{E\cdot \frac{\bfu}{\norm{\bfu}_2} + \frac{\rho a T_1}{\sqrt{2\pi}}\cdot \frac{\bfu}{\norm{\bfu}_2} + \frac{b T_2}{\sqrt{2\pi}}\cdot \frac{\bfv}{\norm{\bfv}_2}}}_{\bfg_1}\\
            & \quad\quad\quad + \underbrace{\frac{\kappa}{\sqrt{2\pi}}\paren{\bm{\mu}^\top\bfv\paren{T_1\cdot \frac{\bfv}{\norm{\bfv}_2} + T_2\cdot \frac{\bfu}{\norm{\bfu}_2}} + \norm{\bfv}_2\paren{T_2 + \rho T_1}\bm{\mu}}}_{\bfg_2}\\
            & \quad\quad\quad + F\paren{\bm{\mu}\bm{\mu}^\top\bfv + 3\kappa^2\bfv}
        \end{aligned}
    \end{equation}
    Then we have that
    \begin{equation}
        \label{eq:exp_mat_vec_diff}
        \begin{aligned}
             & \E_{\bfc}\left[\bfc\bfc^\top\indy{z_1\geq 0;z_2\geq 0}\right]\bfv - \paren{\bm{\mu}\bm{\mu}^\top\bfv + 3\kappa^2\bfv}\bm{\Phi}_1\paren{\frac{\bm{\mu}^\top\bfu}{\kappa\norm{\bfu}_2}}\bm{\Phi}_1\paren{\frac{\bm{\mu}^\top\bfv}{\kappa\norm{\bfv}_2}} \\
             & \quad\quad\quad = \bfg_1 + \bfg_2 + \mathcal{C}\paren{\frac{\bm{\mu}^\top\bfu}{\kappa\norm{\bfu}_2}, \frac{\bm{\mu}^\top\bfv}{\kappa\norm{\bfv}_2}, \frac{\bfu^\top\bfv}{\norm{\bfu}_2\norm{\bfv}_2}}
        \end{aligned}
    \end{equation}
    The proof then proceed by estimating the magnitude of the three terms. To start, we need to bound $T_1$ and $T_2$. In particular, since $\bm{\Phi}_1$ is the CDF, its magnitude must be bounded by $1$. Therefore
    \[
        0\leq T_1 \leq \exp\paren{-\frac{a^2}{2}};\quad 0\leq T_2 \leq \exp\paren{-\frac{b^2}{2}}
    \]
    Therefore, the $\ell_{\infty}$ norm of $\bfg_2$ is bounded by
    \begin{equation}
        \label{eq:g2_bound}
        \begin{aligned}
            \norm{\bfg_2}_{\infty} & \leq \frac{\kappa}{\sqrt{2\pi}}\paren{\bm{\mu}^\top\bfv\paren{\exp\paren{-\frac{a^2}{2}} + \exp\paren{-\frac{b^2}{2}}}} + \norm{\bfv}_2\norm{\bm\mu}_2\paren{\exp\paren{-\frac{a^2}{2}} + \rho\exp\paren{-\frac{b^2}{2}}}\\
            & \leq \frac{2\kappa}{\sqrt{2\pi}}\norm{\bfv}_2\norm{\bm{\mu}}_2\paren{\exp\paren{-\frac{a^2}{2}} + \exp\paren{-\frac{b^2}{2}}}\\
            & \leq \kappa\norm{\bfv}_2\norm{\bm{\mu}}_2\paren{\phi\paren{\frac{a}{2}} + \phi\paren{\frac{b}{2}}}
        \end{aligned}
    \end{equation}
    Next, for $E$, we have
    \begin{align*}
        E & = \frac{1-\rho^2}{2\pi}\exp\paren{-\frac{a^2 - 2\rho ab + b^2}{2\paren{1-\rho^2}}}\\
        & =\frac{1-\rho^2}{4\pi}\paren{\exp\paren{-\frac{a^2 - 2\rho ab + \rho^2b^2}{2\paren{1-\rho^2}} - \frac{b^2}{2}} + \exp\paren{-\frac{\rho^2a^2 - 2\rho ab + b^2}{2\paren{1-\rho^2}} - \frac{a^2}{2}}}\\
        & \leq \frac{1}{4\pi}\paren{\exp\paren{-\frac{a^2}{2}} + \exp\paren{-\frac{b^2}{2}}}\\
        & \leq  \frac{1}{4\pi}\paren{\phi\paren{\frac{a}{2}} + \phi\paren{\frac{b}{2}}}
    \end{align*}
    Therefore, the magnitude of $\bfg_1$ can be bounded by
    \begin{equation}
        \label{eq:g1_bound}
        \begin{aligned}
            \norm{\bfg_1}_2 & \leq \kappa^2\norm{\bfv}_2\paren{\left|E\right| + \frac{\left|a\right|T_1}{\sqrt{2\pi}} + \frac{\left|b\right|T_2}{\sqrt{2\pi}}}\\
            & \leq \kappa^2\norm{\bfv}_2\paren{\frac{1}{4\pi}\paren{\phi\paren{\frac{a}{2}} + \phi\paren{\frac{b}{2}}} + \frac{\left|a\right|}{\sqrt{2\pi}}\exp\paren{-\frac{a^2}{2}} + \frac{\left|b\right|}{\sqrt{2\pi}}\exp\paren{-\frac{b^2}{2}}}\\
            & \leq \kappa^2\norm{\bfv}_2\paren{\frac{1}{4\pi}\paren{\phi\paren{\frac{a}{2}} + \phi\paren{\frac{b}{2}}} + \psi\paren{\frac{a}{2}} + \psi\paren{\frac{b}{2}}}
        \end{aligned}
    \end{equation}
    Moreover, by the bound of the Gaussian Copula function, we have that
    \[
        \left|\mathcal{C}\paren{a,b,\rho}\right|\leq \frac{1}{4}\exp\paren{-\frac{a^2+b^2}{4}}
    \]
    Therefore, we have that
    \[
        \mathcal{C}\paren{\frac{\bm{\mu}^\top\bfu}{\kappa\norm{\bfu}_2}, \frac{\bm{\mu}^\top\bfv}{\kappa\norm{\bfv}_2}, \frac{\bfu^\top\bfv}{\norm{\bfu}_2\norm{\bfv}_2}} \leq \frac{1}{4}\exp\paren{-\frac{a^2}{4}}\exp\paren{-\frac{b^2}{4}} = \frac{1}{4}\phi\paren{\frac{a}{2}}\phi\paren{\frac{b}{2}}
    \]
    Combining the results gives
    \begin{align*}
        & \norm{\E_{\bfc}\left[\bfc\bfc^\top\indy{z_1\geq 0;z_2\geq 0}\right]\bfv - \paren{\bm{\mu}\bm{\mu}^\top\bfv + 3\kappa^2\bfv}\bm{\Phi}_1\paren{\frac{\bm{\mu}^\top\bfu}{\kappa\norm{\bfu}_2}}\bm{\Phi}_1\paren{\frac{\bm{\mu}^\top\bfv}{\kappa\norm{\bfv}_2},}}_2\\
        & \quad\quad\quad \leq \kappa^2\norm{\bfv}_2\paren{\frac{1}{4\pi}\paren{\phi\paren{\frac{a}{2}} + \phi\paren{\frac{b}{2}}} + \psi\paren{\frac{a}{2}} + \psi\paren{\frac{b}{2}}} + \kappa\norm{\bfv}_2\norm{\bm{\mu}}_2\paren{\phi\paren{\frac{a}{2}} + \phi\paren{\frac{b}{2}}}\\
        & \quad\quad\quad\quad\quad + \frac{1}{4}\paren{\left|\bm{\mu}^\top\bfv\right|\norm{\bm{\mu}}_{\infty} + 3\kappa^2\norm{\bfv}_2}\phi\paren{\frac{a}{2}}\phi\paren{\frac{b}{2}}\\
        & \quad\quad\quad = \norm{\bfv}_2\paren{\paren{\kappa^2 + \kappa\norm{\bm{\mu}}_2}\paren{\phi\paren{\frac{a}{2}} + \phi\paren{\frac{b}{2}}} + \paren{\kappa^2+\kappa\norm{\bm{\mu}}_{\infty}}\paren{\psi\paren{\frac{a}{2}} + \psi\paren{\frac{b}{2}}}}\\
        & \quad\quad\quad \leq 2\kappa\norm{\bfv}_2\paren{\norm{\bm{\mu}}_2\paren{\phi\paren{\frac{a}{2}} + \phi\paren{\frac{b}{2}}} + \norm{\bm{\mu}}_{\infty}\paren{\psi\paren{\frac{a}{2}} + \psi\paren{\frac{b}{2}}}}
    \end{align*}
\end{proof}

\subsubsection{Other Results}
\begin{lemma}
    \label{lem:inner_gaussian_rv}
    Let $\bfc\sim\mathcal{N}\paren{\bm{\mu},\kappa^2\bfI}$, and let $\bfu\in\R^d$ be a vector. Define $z = \bfc^\top\bfu$. Then we have that $z\sim\mathcal{N}\paren{\bm{\mu}^\top\bfu,\kappa^2\norm{\bfu}_2^2}$.
\end{lemma}
\begin{proof}
    Since $z = \bfc^\top\bfu$ where $\bfc\sim\mathcal{N}\paren{\bm{\mu},\kappa^2\bfI}$. Then the moment generating function of $z$ is given by
    \begin{align*}
        M_z(t) & = \E\left[\exp\paren{zt}\right]\\
        & = \E\left[\prod_{j=1}^d\exp\paren{c_ju_jt}\right]\\
        & = \prod_{j=1}^d\E\left[\exp\paren{c_ju_jt}\right]\\
        & = \prod_{j=1}^d\exp\paren{u_j\mu_jt + \frac{1}{2}u_j^2\kappa^2t^2} \\
        & = \exp\paren{\paren{\sum_{j=1}^du_j\mu_j}t + \frac{1}{2}\paren{\sum_{j=1}^du_j^2}\kappa^2t^2}\\
        & = \exp\paren{\bm{\mu}^\top\bfu\cdot t + \frac{1}{2}\norm{\bfu}_2^2\kappa^2t^2}
    \end{align*}
    Therefore, $z\sim\mathcal{N}\paren{\bm{\mu}^\top\bfu,\kappa^2\norm{\bfu}_2^2}$.
\end{proof}

\begin{lemma}
    \label{lem:first_order_gaus_int}
    Let $\kappa, \mu, a\in\R$ be given such that $\kappa > 0$. Then we have that
    \[
        \int_{a}^\infty z\exp\paren{-\frac{\paren{z-\mu}^2}{2\kappa^2}}dz = \kappa^2\exp\paren{-\frac{\paren{\mu - a}^2}{2\kappa^2}} + \kappa\mu\sqrt{2\pi}\bm{\Phi}_1\paren{\frac{\mu-a}{\kappa}}
    \]
\end{lemma}
\begin{proof}
    We use a change of variable $z' = \frac{z - \mu}{\kappa}$. Then we have that $z = \kappa z' + \mu$, and $dz = \kappa dz'$. Therefore
    \begin{align*}
        \int_{a}^\infty z\exp\paren{-\frac{\paren{z-\mu}^2}{2\kappa^2}}dz & = \int_{\frac{a - \mu}{\kappa}}^\infty \paren{\kappa z' +\mu}\exp\paren{-\frac{z'^2}{2}}\kappa dz'\\
        & = \kappa^2\int_{\frac{a - \mu}{\kappa}}^\infty z'\exp\paren{-\frac{z'^2}{2}}dz' + \kappa\mu\int_{\frac{a - \mu}{\kappa}}^\infty\exp\paren{-\frac{z'^2}{2}}dz'\\
        & = \kappa^2\exp\paren{-\frac{z'^2}{2}}\mid _{\infty}^{\frac{a-\mu}{\kappa}} + \kappa\mu\sqrt{2\pi}\paren{1-\bm{\Phi}_1\paren{\frac{a-\mu}{\kappa}}}\\
        & = \kappa^2\exp\paren{-\frac{\paren{\mu - a}^2}{2\kappa^2}} + \kappa\mu\sqrt{2\pi}\bm{\Phi}_1\paren{\frac{\mu-a}{\kappa}}
    \end{align*}
\end{proof}

\begin{lemma}
    \label{lem:second_order_gaus_int}
    Let $\kappa, \mu, a\in\R$ be given such that $\kappa > 0$. Then we have that
    \[
        \int_a^\infty z^2\exp\paren{-\frac{\paren{z-\mu}^2}{2\kappa^2}}dz = \kappa^2\paren{a+\mu}\exp\paren{-\frac{\paren{\mu-a}^2}{2\kappa^2}} + \sqrt{2\pi}\kappa\paren{\kappa^2+\mu^2}\bm{\Phi}_1\paren{\frac{\mu-a}{\kappa}}
    \]
\end{lemma}
\begin{proof}
    To start, let $z' = \frac{z - \mu}{\kappa}$. Then we have that $z = \kappa z' + \mu$, and $dz = \kappa dz'$. Therefore
    \begin{align*}
        \int_{a}^\infty z^2\exp\paren{-\frac{\paren{z-\mu}^2}{2\kappa^2}}dz & = \kappa\int_{\frac{a-\mu}{\kappa}}^\infty\paren{\kappa z' + \mu}^2\exp\paren{-\frac{z'^2}{2}}dz'\\
        & = \kappa^3\int_{\frac{a-\mu}{\kappa}}^\infty z'^2\exp\paren{-\frac{z'^2}{2}}dz' + 2\kappa^2\mu\int_{\frac{a-\mu}{\kappa}}^\infty z'\exp\paren{-\frac{z'^2}{2}}dz'\\
        & \quad\quad\quad + \kappa\mu^2\int_{\frac{a-\mu}{\kappa}}^\infty\exp\paren{-\frac{z'^2}{2}}dz'
    \end{align*}
    Notice that for the third term, we have that
    \[
        \int_{\frac{a-\mu}{\kappa}}^\infty\exp\paren{-\frac{z'^2}{2}}dz' = \sqrt{2\pi}\paren{1 - \bm{\Phi}_1\paren{\frac{a - \mu}{\kappa}}} = \sqrt{2\pi}\bm{\Phi}_1\paren{\frac{\mu-a}{\kappa}}
    \]
    For the second term, we can directly apply Lemma~\ref{lem:first_order_gaus_int} with $\kappa = 1, \mu = 0$ to get that
    \[
        \int_{\frac{a-\mu}{\kappa}}^\infty z'\exp\paren{-\frac{z'^2}{2}}dz' = \exp\paren{-\frac{\paren{a-\mu}^2}{2\kappa^2}}
    \]
    For the first term, we apply integration by parts with $u(z') = -z'$ and $v(z') = \exp\paren{-\frac{z'^2}{2}}$. In particular, notice that $v'(z') = -z'\exp\paren{-\frac{z'^2}{2}}$ and $u'(z') = -1$. Therefore
    \begin{align*}
        \int_{\frac{a-\mu}{\kappa}}^\infty z'^2\exp\paren{-\frac{z'^2}{2}}dz' & = \int_{\frac{a-\mu}{\kappa}}^\infty  u(z')dv(z')\\
        & = u(z')v(z')\mid_{\frac{a-\mu}{\kappa}}^\infty  - \int_{\frac{a-\mu}{\kappa}}^\infty  v(z')du(z')\\
        & = -z'\exp\paren{-\frac{z'^2}{2}}\mid_{\frac{a-\mu}{\kappa}}^\infty  + \int_{\frac{a-\mu}{\kappa}}^\infty  \exp\paren{-\frac{z'^2}{2}}dz'\\
        & = \frac{a - \mu}{\kappa}\exp\paren{-\frac{\paren{\mu-a}^2}{2\kappa^2}} + \sqrt{2\pi}\paren{1-\bm{\Phi}_1\paren{\frac{a-\mu}{\kappa}}}\\
        & = \frac{a - \mu}{\kappa}\exp\paren{-\frac{\paren{\mu-a}^2}{2\kappa^2}} + \sqrt{2\pi}\bm{\Phi}_1\paren{\frac{\mu-a}{\kappa}}
    \end{align*}
    Putting things together, we have that
    \begin{align*}
        \int_{a}^\infty z^2\exp\paren{-\frac{\paren{z-\mu}^2}{2\kappa^2}}dz & = \kappa^3\paren{\frac{a - \mu}{\kappa}\exp\paren{-\frac{\paren{\mu-a}^2}{2\kappa^2}} + \sqrt{2\pi}\bm{\Phi}_1\paren{\frac{\mu-a}{\kappa}}}\\
        &\quad\quad\quad + 2\kappa^2\mu\exp\paren{-\frac{\paren{a-\mu}^2}{2\kappa^2}} + \kappa\mu^2\sqrt{2\pi}\bm{\Phi}_1\paren{\frac{\mu-a}{\kappa}}\\
        & = \paren{\kappa^2\paren{a - \mu} + 2\kappa^2\mu}\exp\paren{-\frac{\paren{\mu-a}^2}{2\kappa^2}} + \sqrt{2\pi}\paren{\kappa^3 + \kappa\mu^2}\bm{\Phi}_1\paren{\frac{\mu-a}{\kappa}}\\
        &  = \kappa^2\paren{a+\mu}\exp\paren{-\frac{\paren{\mu-a}^2}{2\kappa^2}} + \sqrt{2\pi}\kappa\paren{\kappa^2+\mu^2}\bm{\Phi}_1\paren{\frac{\mu-a}{\kappa}}
    \end{align*}
\end{proof}

\begin{lemma}
    \label{lem:cdf_forms}
    Let $z_1,z_2\sim\mathcal{N}\paren{0,1}$ with $\text{Cov}\paren{z_1, z_2} = \rho$. Then we have that
    \[
        \int_a^\infty f\paren{z_1\mid z_2}dz_1 = \bm{\Phi}_1\paren{\frac{\rho z_2 - a}{\sqrt{1-\rho^2}}}
    \]
\end{lemma}
\begin{proof}
    Since $z_1,z_2\sim\mathcal{N}\paren{0,1}$ with $\text{Cov}\paren{z_1, z_2} = \rho$, we have that $z_1\mid z_2 \sim\mathcal{N}\paren{\rho z_2, 1-\rho^2}$. Therefore, using a change of variable $z' = \frac{z_1 - \rho z_2}{\sqrt{1-\rho^2}}$, we have 
    \begin{align*}
        \int_a^\infty f\paren{z_1\mid z_2}dz_1 & = \frac{1}{\sqrt{2\pi\paren{1-\rho^2}}}\int_a^\infty\exp\paren{-\frac{\paren{z_1-\rho z_2}^2}{2\paren{1-\rho^2}}}dz_1\\
        & = \frac{1}{\sqrt{2\pi}}\int_{\frac{a - \rho z_2}{\sqrt{1-\rho^2}}}^\infty\exp\paren{-\frac{z'^2}{2}}dz'\\
        & = 1 - \bm{\Phi}_1\paren{\frac{a-\rho z_2}{\sqrt{1-\rho^2}}}\\
        & = \bm{\Phi}_1\paren{\frac{\rho z_2 - a}{\sqrt{1-\rho^2}}}
    \end{align*}
\end{proof}

\begin{lemma}
    \label{lem:joint_cdf_forms}
    Let $z_1,z_2\sim\mathcal{N}\paren{0,1}$ with $\text{Cov}\paren{z_1, z_2} = \rho$. Then we have that
    \[
        \bm{\Phi}_2\paren{-a, -b, \rho} = \int_a^\infty\int_b^\infty f\paren{z_1, z_2}dz_2dz_1 = \int_b^\infty\bm{\Phi}_1\paren{\frac{\rho z_2 - a}{\sqrt{1-\rho^2}}}f\paren{z_2}dz_2
    \]
\end{lemma}
\begin{proof}
    Let $z_1',z_2'\sim\mathcal{N}\paren{0, 1}$ with $\text{Cov}\paren{z_1, z_2} = \rho$, and define $z_1 = -z_1', z_2 = -z_2'$. Then we have that $z_1,z_2\sim\mathcal{N}\paren{0,1}$ with $\text{Cov}\paren{z_1, z_2} = \rho$. By symmetry, we have $f\paren{z_1, z_2} = f\paren{-z_1, -z_2} = f\paren{z_1', z_2'}$. Moreover, $dz_2'dz_1' = \paren{-dz_2}\paren{-dz_1} = dz_2dz_1$. Thus
    \[
        \bm{\Phi}_2\paren{-a, -b, \rho} = \int_{-\infty}^{-a}\int_{-\infty}^{-b}f\paren{z_1', z_2'}dz_2'dz_1' = \int_{a}^\infty\int_b^\infty f\paren{z_1,z_2}dz_2dz_1
    \]
    Recall that $f\paren{z_1, z_2} = f\paren{z_1\mid z_2}f\paren{z_2}$. Then we can apply Lemma~\ref{lem:cdf_forms} to get that
    \[
        \int_{a}^\infty\int_b^\infty f\paren{z_1,z_2}dz_2dz_1 =\int_b^\infty\paren{\int_a^\infty f\paren{z_1\mid z_2}dz_1}f\paren{z_2}dz_2 = \int_b^\infty\bm{\Phi}_1\paren{\frac{\rho z_2-a}{\sqrt{1-\rho^2}}}f\paren{z_2}dz_2
    \]
\end{proof}

\begin{lemma}
    \label{lem:exp_relu}
    Let $z\sim\mathcal{N}\paren{\mu, \kappa^2}$, and let $a\in\R$. Then we have
    \[
        \E\left[z\indy{z\geq a}\right] = \frac{\kappa}{\sqrt{2\pi}}\exp\paren{-\frac{\paren{\mu-a}^2}{2\kappa^2}} + \mu\bm{\Phi}_1\paren{\frac{\mu-a}{\kappa}}
    \]
\end{lemma}
\begin{proof}
    Define $\hat{z} = \frac{z - \mu}{\kappa}$. Then we have that $\hat{z} \sim\mathcal{N}\paren{0, 1}$. Since $z = \kappa\hat{z} + \mu$, we have
    \begin{align*}
        \E\left[z\indy{z\geq 0}\right] & = \E\left[\paren{\kappa\hat{z} + \mu}\indy{\hat{z}\geq \frac{a-\mu}{\kappa}}\right]\\
        & = \kappa\E\left[\hat{z}\indy{\hat{z}\geq \frac{a-\mu}{\kappa}}\right] + \mu\E\left[\indy{\hat{z}\geq \frac{a-\mu}{\kappa}}\right]
    \end{align*}
    Notice that $\E\left[\indy{\hat{z}\geq \frac{a-\mu}{\kappa}}\right] = \Pr\paren{\hat{z}\geq \frac{a-\mu}{\kappa}} = \bm{\Phi}_1\paren{\frac{\mu-a}{\kappa}}$. Moreover
    \begin{align*}
        \E\left[\hat{z}\indy{\hat{z}\geq -\frac{\mu}{\kappa}}\right] = \frac{1}{\sqrt{2\pi}}\int_{\frac{a-\mu}{\kappa}}^\infty \hat{z}\exp\paren{-\frac{z^2}{2}}dz = \frac{1}{\sqrt{2\pi}}\exp\paren{-\frac{\paren{a-\mu}^2}{2\kappa^2}}
    \end{align*}
    Therefore
    \[
        \E\left[z\indy{z\geq 0}\right] = \frac{\kappa}{\sqrt{2\pi}}\exp\paren{-\frac{\paren{a-\mu}^2}{2\kappa^2}} + \mu\bm{\Phi}_1\paren{\frac{\mu-a}{\kappa}}
    \]
\end{proof}

\begin{lemma}
    \label{lem:arcsin_bound}.
    Let $x\in[-1,1]$. Then we have that $|\arcsin x|\leq \frac{\pi}{2}\cdot |x|$.
\end{lemma}
\begin{proof}
    To start, consider the case $x > 0$. Define $f(x) = \frac{\arcsin x}{x}$. Then we have that
    \[
        f'(x) = x^{-2}\paren{\frac{x}{\sqrt{1-x^2}} - \arcsin x};\quad\quad f''(x) = x^{-3}\paren{\frac{3x^3-2x}{\paren{1-x^2}^{\frac{3}{2}}} + 2 \arcsin x}
    \]
    For all $x\in (0, 1]$, we have that $1 - x^2 \geq 0$. Notice that by the Taylor expansion of $\arcsin x$, we have
    \[
        \arcsin x \leq x + \frac{x^3}{6}
    \]
    when $x\in(0, 1]$. Therefore
    \[
        \frac{3x^3-2x}{\paren{1-x^2}^{\frac{3}{2}}} + 2 \arcsin x \geq \frac{3x^3-2x + 2x\paren{1-x^2}^{\frac{3}{2}} + \frac{x^3}{3}\paren{1-x^2}^{\frac{3}{2}}}{\paren{1-x^2}^{\frac{3}{2}}} \geq \frac{3x^3 - 2x\paren{1 - \paren{1-x^2}^{\frac{3}{2}}}}{\paren{1-x^2}^{\frac{3}{2}}}
    \]
    Since $\paren{1-x^2}^{\frac{3}{2}} \geq \paren{1-x^2}^3 \geq 1- 3x^4 + 2x^6$, we must have that
    \[
        3x^3 - 2x\paren{1 - \paren{1-x^2}^{\frac{3}{2}}} \leq 3x^3 - 6x^5 + 4x^7 = 3x^3\paren{1 - 2x^2 + x^4} + x^7 = 3x^3\paren{1-x^2}^2 + x^7 \geq 0
    \]
    Thus, we must have that $f''(x) \geq 0$. Therefore, for any $\epsilon \in (0, 0.1]$, we have that for $x\in[\epsilon, 1]$
    \begin{align*}
        f\paren{x} & \leq f\paren{\paren{(1+\epsilon)x - \epsilon}\cdot 1 + (1-x)\cdot \epsilon}\\
        & \leq (1-x)\cdot f(\epsilon) + \paren{(1+\epsilon)x - \epsilon}\cdot f(1)\\
        & = (1-x)\cdot f\paren{\epsilon} + \frac{\pi}{2}\cdot x - \frac{\pi}{2}\cdot \epsilon\paren{1-x}\\
        & \leq \frac{\pi}{2}\cdot x + 1.002\paren{1-x}\\
        & \leq \frac{\pi}{2}
    \end{align*}
    This gives that $f\paren{x}\leq \frac{\pi}{2}$ for all $x\in(0,1]$. Since $f\paren{x}$ is an even function, we have $f\paren{x}\leq \frac{\pi}{2}$ for all $x\in[-1, 0)$. Therefore, $|\arcsin x|\leq \frac{\pi}{2}|x|$ when $x\in[-1,1]\setminus\{0\}$. When $x = 0$, we have $\arcsin x = 0$. This completes the proof. 
\end{proof}
%\subsection*{some notes on $\gamma$}
\begin{lemma}
\label{lem:CS}
Let $\mathbf{u} := (|e_1|, ..., |e_n|)  \in \mathbb{R}^n$ and $\mathbf{1} := (1, ..., 1) \in \mathbb{R}^n.$ Then
\[
\sum_{i=1}^n |e_i|
= u^\top \mathbf{1}.
\]
By the Cauchy--Schwarz inequality,
\[
u^\top \mathbf{1}
\le \|u\|_2  \|\mathbf{1}\|_2.
\]
Moreover,
\[
\|u\|_2
= \left(\sum_{i=1}^n u_i^2\right)^{1/2}
= \left(\sum_{i=1}^n |e_i|^2\right)^{1/2}
= \left(\sum_{i=1}^n e_i^2\right)^{1/2},
\qquad
\|\mathbf{1}\|_2
= \left(\sum_{i=1}^n 1^2\right)^{1/2}
= \sqrt{n}.
\]
Combining the above gives
\[
\sum_{i=1}^n |e_i|
\le
\sqrt{n}\left(\sum_{i=1}^n e_i^2\right)^{1/2}.
\]
\end{lemma}
% Lemma correction!!
\begin{lemma}
    \label{lem:general_smoothness}
    Assume that Assumption~\ref{asump:data} holds. Then, we have:
    \[
\norm{\nabla_{\bfw_r}\calL_{\bfC}\paren{\mathbf{W}}}_2^2\leq \frac{\sqrt{n}}{\sqrt{m}}\calL_{\bfC}\paren{\mathbf{W}}^{\frac{1}{2}}.
    \]
\end{lemma}
\begin{proof}
    By the form of $\nabla_{\bfw_r}\calL_{\bfC}\paren{\mathbf{W}}$ in (\ref{eq:mask_grad}), we have:
    \begin{align*}
        \norm{\nabla_{\bfw_r}\calL_{\bfC}\paren{\mathbf{W}}}_2
        & =\norm{\frac{a_r}{\sqrt{m}}\sum_{i=1}^n\paren{f\paren{\mathbf{W},\bfx_i\odot\bfc_i} - y_i} \paren{\bfx_i\odot\bfc_i}\underbrace{\indy{\bfw_r^\top\paren{\bfx_i\odot\bfc_i}\geq 0}}_{\le 1}}_2 \\
        & \leq \frac{\left|a_r\right|}{\sqrt{m}}\sum_{i=1}^n\left|f\paren{\mathbf{W},\bfx_i\odot\bfc_i}-y_i\right|\cdot\norm{\bfx_i\odot\bfc_i}_2\\
        & \leq \frac{\sqrt{n}}{\sqrt{m}}\left(\sum_{i=1}^n (f\paren{\mathbf{W},\bfx_i\odot\bfc_i}-y_i)^2\right)^{1/2} \cdot \|\bfc_i\|_\infty \|\bfx_i\|_2 \enspace \text{using \ref{lem:CS}} \\
        & \leq C \frac{\sqrt{n}}{\sqrt{m}} \left(2\calL_{\bfC}\paren{\mathbf{W}}\right)^{1/2} \enspace \text{using \ref{lem:c_infty_bound} with }\|\bfc_i\|_\infty \;\le\; 1 + \kappa\sqrt{2\log\!\Big(\frac{2d}{\delta}\Big)} = C \enspace \text{and} \enspace \norm{\bfx_i}_2\leq 1\\
        & \leq C\sqrt{2}\frac{\sqrt{n}}{\sqrt{m}}\calL_{\bfC}\paren{\mathbf{W}}^{1/2}
    \end{align*}
    where, in the first inequality, we use the fact that the indicator function is upper-bounded by 1 and in the second inequality, we use the fact that $a_r = \pm 1$.

    and so,
    \begin{align*}
        \norm{\nabla_{\bfw_r}\calL_{\bfC}\paren{\mathbf{W}}}^2_2\leq 2C^2\cdot \frac{n}{m}\calL_{\bfC}\paren{\mathbf{W}}
    \end{align*}
\end{proof}

%%%%
\begin{lemma}[High-probability $\ell_\infty$ bound for Gaussian masks]
\label{lem:c_infty_bound}
Fix $\delta\in(0,1)$. Let $\bfc_i\in\mathbb R^d$ be a Gaussian mask with independent coordinates
\[
\bfc_i \sim \mathcal N(\bm 1,\kappa^2 \bfI_d),
\qquad\text{i.e.,}\qquad
c_{i,j} = 1+\kappa g_{i,j},\;\; g_{i,j}\stackrel{\text{i.i.d.}}{\sim}\mathcal N(0,1).
\]
Then, with probability at least $1-\delta$, we have
\[
\|\bfc_i\|_\infty \;\le\; 1 + \kappa\sqrt{2\log\!\Big(\frac{2d}{\delta}\Big)}.
\]
\end{lemma}

\begin{proof}
Since $\bfc_i\sim\mathcal N(\bm 1,\kappa^2\bfI_d)$ with independent coordinates, each coordinate can be written as
\[
c_{i,j} = 1+ \kappa g_{i,j},
\qquad\text{where } g_{i,j}\sim\mathcal N(0,1) \enspace i.i.d.
\]
Hence
\[
\|\bfc_i\|_\infty = \max_{j\in[d]}|c_{i,j}| = \max_{j\in[d]}|1+\kappa g_{i,j}|.
\]
We observe that 
\begin{align*}
|c_{i,j}|
&=|1+(c_{i,j}-1)| \\
&\le |1|+|c_{i,j}-1| \enspace \text{triangle inequality} \\
&= 1+|\kappa g_{i,j}| 
\end{align*}
We will first bound $\max_j |c_{i,j}-1| = \max_j |\kappa g_{i,j}|$, and then convert this into a bound on $\|\bfc_i\|_\infty$.

For brevity, we denote $g_{i,j}$ as $g\sim\mathcal N(0,1)$ and we are going to show that
\begin{equation}
\label{eq:gauss_tail}
\boxed{\Pr(|g|\ge t)\le 2e^{-t^2/2}}
\end{equation}
By symmetry of the standard normal distribution,
\[
\Pr(|g|\ge t)=\Pr(g\ge t)+\Pr(g\le -t)=2\Pr(g\ge t).
\]
So it suffices to upper bound $\Pr(g\ge t)$.

For any $\lambda>0$, since the exponential is monotone increasing we have:
\[
g\ge t \Rightarrow  \lambda g \ge \lambda t  \Rightarrow e^{\lambda g}\ge e^{\lambda t}
\]
and so 
\begin{align*}
    \Pr(g\ge t) = \Pr(e^{\lambda g}\ge e^{\lambda t})
\end{align*}
By Markov's inequality, for any nonnegative random variable $X$ and any $a>0$,
\[
\Pr(X\ge a)\le \frac{ \mathbb E[X]}{a}
\]
Applying this with $X=e^{\lambda g}$ and $a=e^{\lambda t}$ gives
\begin{equation}
\label{eq:markov_step}
\Pr(g\ge t)=\Pr\big(e^{\lambda g}\ge e^{\lambda t}\big)
\le \frac{\mathbb E[e^{\lambda g}]}{e^{\lambda t}}
= \mathbb E[e^{\lambda g}]e^{-\lambda t}
\end{equation}
Computation of the moment generating function $\mathbb E[e^{\lambda g}]$:

The standard normal density is
\[
\varphi(x)=\frac{1}{\sqrt{2\pi}}e^{-x^2/2},\qquad x\in\mathbb R.
\]
Therefore,
\begin{align}
\mathbb E[e^{\lambda g}]
&=\int_{-\infty}^{\infty} e^{\lambda x} \varphi(x) dx
=\frac{1}{\sqrt{2\pi}}\int_{-\infty}^{\infty} \exp\!\Big(\lambda x-\frac{x^2}{2}\Big) dx.
\label{eq:mgf_int}
\end{align}
We now study the exponent:
\[
\lambda x-\frac{x^2}{2}
=-\frac{1}{2}\big(x^2-2\lambda x\big)
=-\frac{1}{2}\big((x-\lambda)^2-\lambda^2\big)
=\frac{\lambda^2}{2}-\frac{(x-\lambda)^2}{2}.
\]
Plugging this into \eqref{eq:mgf_int} yields
\begin{align*}
\mathbb E[e^{\lambda g}]
&=\frac{1}{\sqrt{2\pi}}\int_{-\infty}^{\infty}
\exp\!\Big(\frac{\lambda^2}{2}-\frac{(x-\lambda)^2}{2}\Big) dx
\\
&=e^{\lambda^2/2}\cdot \frac{1}{\sqrt{2\pi}}
\int_{-\infty}^{\infty} \exp\!\Big(-\frac{(x-\lambda)^2}{2}\Big) dx.
\end{align*}
We set $u=x-\lambda$ (change of variables with $dx=du$)
\[
\int_{-\infty}^{\infty} \exp\!\Big(-\frac{(x-\lambda)^2}{2}\Big) dx
=
\int_{-\infty}^{\infty} \exp\!\Big(-\frac{u^2}{2}\Big) du
=
\sqrt{2\pi}.
\]
Hence
\begin{equation}
\label{eq:mgf_closed}
\mathbb E[e^{\lambda g}] = e^{\lambda^2/2}.
\end{equation}
Substituting \eqref{eq:mgf_closed} into \eqref{eq:markov_step} gives
\[
\Pr(g\ge t)\le \exp\!\Big(\frac{\lambda^2}{2}-\lambda t\Big),
\qquad \forall \lambda>0.
\]
The right-hand side is a valid bound for \emph{every} $\lambda>0$, so we choose $\lambda$ to make it as small as possible.
Define
\[
f(\lambda):=\frac{\lambda^2}{2}-\lambda t.
\]
Then $f'(\lambda)=\lambda-t$, so the unique minimizer is $\lambda=t$ (and $f''(\lambda)=1>0$ confirms it is a minimum).
Plugging $\lambda=t$ gives
\[
\Pr(g\ge t)\le \exp\!\Big(\frac{t^2}{2}-t^2\Big)=e^{-t^2/2}.
\]
Using $\Pr(|g|\ge t)=2\Pr(g\ge t)$ proves \eqref{eq:gauss_tail}.

Now, we focus on one of the mask coordinates by fixing a coordinate $j\in[d]$. Since $c_{i,j}-1=\kappa g_{i,j}$ with $g_{i,j}\sim\mathcal N(0,1)$, for any $u\ge 0$ we have
\[
\Pr\big(|c_{i,j}-1|\ge u\big)
=
\Pr\big(|g_{i,j}|\ge u/\kappa\big)
\le
2\exp\!\Big(-\frac{u^2}{2\kappa^2}\Big),
\]
where we applied \eqref{eq:gauss_tail} with $t=u/\kappa$.

Let's define the event
\[
\mathcal E_i(u) := \Big\{\max_{j\in[d]} |c_{i,j}-1|\le u\Big\}.
\]
Its complement is the event that at least one coordinate deviates by more than $u$:
\[
\mathcal E_i(u)^c = \Big\{\exists j\in[d]\ \text{s.t.}\ |c_{i,j}-1|>u\Big\}.
\]
By the union bound,
\begin{align*}
\Pr\big(\mathcal E_i(u)^c\big)
&=
\Pr\Big(\bigcup_{j=1}^d \{|c_{i,j}-1|>u\}\Big)
\le \sum_{j=1}^d \Pr\big(|c_{i,j}-1|>u\big)
\\
&\le \sum_{j=1}^d 2\exp\!\Big(-\frac{u^2}{2\kappa^2}\Big)
= 2d\exp\!\Big(-\frac{u^2}{2\kappa^2}\Big).
\end{align*}
Choose $u$ so that the right-hand side is at most $\delta$:
\[
2d\exp\!\Big(-\frac{u^2}{2\kappa^2}\Big)\le \delta
\quad\Longleftrightarrow\quad
-\frac{u^2}{2\kappa^2}\le \log\!\Big(\frac{\delta}{2d}\Big)
\quad\Longleftrightarrow\quad
u\ge \kappa\sqrt{2\log\!\Big(\frac{2d}{\delta}\Big)}.
\]
Set
\[
u := \kappa\sqrt{2\log\!\Big(\frac{2d}{\delta}\Big)}.
\]
Then $\Pr(\mathcal E_i(u)^c)\le \delta$, i.e.\ $\Pr(\mathcal E_i(u))\ge 1-\delta$, and on $\mathcal E_i(u)$ we have
\[
|c_{i,j}-1|\le u \quad \forall j\in[d].
\]

On $\mathcal E_i(u)$, for each coordinate $j$,
\[
|c_{i,j}|
=|1+(c_{i,j}-1)|
\le |1|+|c_{i,j}-1|
\le 1+u.
\]
Taking the maximum over $j$ yields, with probability at least $1-\delta$,
\[
\|\bfc_i\|_\infty
=\max_{j\in[d]}|c_{i,j}|
\le 1+u
= 1+\kappa\sqrt{2\log\!\Big(\frac{2d}{\delta}\Big)}.
\]
\end{proof}

\end{document} 
